%?deps @bib-commons

\documentclass[12pt,openany]{book}
\usepackage{aliascnt,amsmath,amssymb,amsthm,bm,calc,colonequals,enumitem,etoolbox,fancyhdr,imakeidx,mathtools,mleftright,parskip,refcount,setspace,textgreek,xcolor,xstring}
\usepackage[a4paper,twoside,inner=4cm,outer=3cm,top=3cm,bottom=2.5cm,headheight=2cm]{geometry}
\usepackage[UKenglish]{babel}
\usepackage[hidelinks,unicode]{hyperref}

\usepackage{bib-commons}

\makeatletter
% add unnumbered chapters to toc
\newcommand\auxitoc[1]\relax
\let\chapter@original\chapter
\def\chapter{\@ifstar\chapter@star\chapter@original}
\newcommand\chapter@star[1]{\chapter@original*{#1}\auxitoc{#1}}

\newcommand\university[1]{\gdef\@university{#1}}
\newcommand\school[1]{\gdef\@school{#1}}
\newcommand\degree[1]{\gdef\@degree{#1}}
\renewcommand\maketitle{%
  \thispagestyle{empty}
  \begin{titlepage}
    \begin{center}
      \singlespacing
      \vspace*{1cm}
      \rule\textwidth{1.5pt} \\
      {\Large \bfseries \@title \par \ }\\[-.85cm]
      \rule\textwidth{1.5pt} \\
      \vspace{1.8cm}
      {\Large \bfseries \@author \par}
      \vspace{0.5cm}
      {\Large \@university \\[.04cm]
       \large \@school \par}
      \vspace{4cm}
      {\large Submitted in accordance with the requirements \\ for the degree of \@degree \par}
      \vspace{2cm}
      {\large \bfseries \@date \par}
    \end{center}
  \end{titlepage}}
\makeatother

% default page style
\fancypagestyle{main}{%
  \fancyhf{}
  \fancyhead[C]{\thepage}}
% page style for intellectual property statement
\fancypagestyle{cover}{\fancyhf{}}
% page style for chapter starting pages
\fancypagestyle{plain}{%
  \fancyhf{}
  \renewcommand\headrulewidth{0pt}}
\pagestyle{main}

\setlocalecaption{UKenglish}{bib}{References}

\appto\frontmatter{%
  \renewcommand\auxitoc[1]{\addcontentsline{toc}{section}{\protect #1}}}
\appto\backmatter{%
  \renewcommand\auxitoc[1]{\addcontentsline{toc}{chapter}{\protect\numberline{}#1}}}

% frontmatter has chapter-like sections
\newcommand\frontsection[1]{\clearpage\section*{\LARGE #1}\auxitoc{#1}}

\counterwithout{footnote}{chapter}

% adjust amsrefs reference visual styles 
\renewcommand\eprint[1]{\texttt{#1}}
\renewcommand\PrintThesisType[1]{\StrSubstitute{#1}{Ph.D.}{PhD}}

\onehalfspacing
\allowdisplaybreaks

\begin{filecontents}{\jobname.mst}
item_0 "\n  {\\bfseries "
item_1 "\n  \\item "
item_2 "\n    \\subitem "
item_01 "\\par}\n  \\item "
item_x1 "\\par}\n  \\item "
item_12 "\n    \\subitem "
item_x2 "\n    \\subitem "

delim_0 ""
delim_1 "\\quad"
delim_2 "\\quad"
\end{filecontents}
\makeindex
\newcommand\indg[1]{\index{A@General!#1}}
\newcommand\inda[1]{\index{B@Arithmetic!#1}}
\newcommand\inds[1]{\index{C@Set theory!#1}}

\newcommand\newaliastheorem[2]{%
  \newaliascnt{#1}{theorem}
  \newtheorem{#1}[#1]{#2}
  \aliascntresetthe{#1}
  \expandafter\def\csname #1autorefname\endcsname{#2}
}

\theoremstyle{plain}
\newtheorem{theorem}{Theorem}[section]
\newaliastheorem{lemma}{Lemma}
\newaliastheorem{proposition}{Proposition}
\newaliastheorem{corollary}{Corollary}
\newaliastheorem{conjecture}{Conjecture}

\newtheorem*{theorem*}{Theorem}
\newtheorem*{lemma*}{Lemma}
\newtheorem*{refthm}{\autoref{\currthmref}}

\theoremstyle{definition}
\newaliastheorem{definition}{Definition}
\newaliastheorem{problem}{Problem}

\newcommand\uppar[1]{{\upshape(}#1{\upshape)}}
\newcommand\uplabel[1]{\textup{(#1)}}

% content commands

\newcommand\tuple[1]{\left\langle #1 \right\rangle}

\newcommand\god[1]{\left\ulcorner #1 \right\urcorner}

\newcommand\opA[1]{\mathop{\forall #1}}
\newcommand\opE[1]{\mathop{\exists #1}}
\newcommand\oplam[1]{\mathop{\lambda #1.}}

\newcommand\defeqv{\quad \ratio\Longleftrightarrow \quad}

\DeclareMathOperator\Con{Con}

\DeclareMathOperator\TJ{TJ}
\DeclareMathOperator\WO{WO}

\DeclareMathOperator\Ind{Ind}
\DeclareMathOperator\TC{TC}

% powerset is not an operator
% so that PP(x) gives no extra gaps
\newcommand\Pow{\mathcal{P}} 

% first-order semantics
\DeclareMathOperator\dd{def}
\DeclareMathOperator\Tr{Tr}

% ordinal notations
\DeclareMathOperator\field{field}
\DeclareMathOperator\otyp{otyp}

% plump ordinals
\DeclareMathOperator\relpl{relpl}

% formal theories
\newcommand\HA{\mathrm{HA}}
\newcommand\PA{\mathrm{PA}}

\newcommand\ACA{\mathrm{ACA}}
\newcommand\ATR{\mathrm{ATR}}
\newcommand\BI{\mathrm{BI}}
\newcommand\AC{\text{-}\mathrm{AC}}
\newcommand\DC{\text{-}\mathrm{DC}}
\newcommand\CA{\text{-}\mathrm{CA}}
\newcommand\TI{\text{-}\mathrm{TI}}
\newcommand\ID{\mathrm{ID}}

\newcommand\Zs{\mathrm{Z}_2}

\newcommand\KP{\mathrm{KP}}
\newcommand\KPP{\mathrm{KP}\mleft(\Pow\mright)}

\newcommand\CZF{\mathrm{CZF}}
\newcommand\CZFP{\mathrm{CZF}\mleft(\Pow\mright)}

\newcommand\Exp{\mathrm{Exp}}

\newcommand\REM{\mathrm{REM}}
\newcommand\LPO{\mathrm{LPO}}
\newcommand\RDC{\mathrm{RDC}}

\newcommand\ACg{\mathrm{AC}_\text{global}}

\newcommand\IKP{\mathrm{IKP}}
\newcommand\IKPP{\mathrm{IKP}\mleft(\Pow\mright)}
\newcommand\IZF{\mathrm{IZF}}

\newcommand\ZF{\mathrm{ZF}}
\newcommand\ZFC{\mathrm{ZFC}}

\newcommand\CM{\mathrm{CM}}
\newcommand\GWO{\mathrm{GWO}}

\newcommand\MLV{\mathrm{ML}_1\mathrm{V}}

% ordinal
\newcommand\Ord{\mathrm{Ord}}
\newcommand\PlOrd{\mathrm{PlOrd}}

% pcas
\newcommand\Kl{\mathrm{Kl}}
\newcommand\Hyp{\mathrm{Hyp}}

% distinguished operations
\newcommand\bk{\mathbf{k}}
\newcommand\bs{\mathbf{s}}
\newcommand\bp{\mathbf{p}}
\newcommand\br{\mathbf{r}}

% distinguished operations on natural numbers
\newcommand\bsN{\mathbf{s}_\mathbb{N}}
\newcommand\bpN{\mathbf{p}_\mathbb{N}}
\newcommand\bdN{\mathbf{d}_\mathbb{N}}

% distinguished operations on sets
\newcommand\bi{\mathbf{i}}
\newcommand\bx{\mathbf{x}}

% finite tuple operations
\newcommand\bl{\mathbf{l}}
\newcommand\bt{\mathbf{t}}

% realisers for set equality
\newcommand\ior{\iota_\textup{r}}
\newcommand\ios{\iota_\textup{s}}
\newcommand\iot{\iota_\textup{t}}
\newcommand\ioso{\iota_\textup{s1}}
\newcommand\iost{\iota_\textup{s2}}

% MLTT universes
\newcommand\bU{\mathbf{U}}
\newcommand\bV{\mathbf{V}}

% recursive indices
\newcommand\Rec{\mathrm{Rec}}

% Kleene-Brouwer ordering
\newcommand\KB{\mathrm{KB}}

% plump ordinals
\newcommand\pl{\mathrm{pl}}
\newcommand\pls{{\pl +}}
\newcommand\PlBSig{\mathrm{PlB}\Sigma_1}

% standard names
\newcommand\St{\mathrm{St}}

\title{Higher-Order Conceptual Mathematics and Intuitionistic Set Theory: Well-Orderings, Ordinals and G\"odel's $L$}
\author{Shuwei Wang}
\university{The University of Leeds}
\school{School of Mathematics}
\degree{Doctor of Philosophy}
\date{June 2026}

\begin{document}

\maketitle

\thispagestyle{cover}

The candidate confirms that the work submitted is his own and that appropriate credit has been given where reference has been made to the work of others. The thesis incorporates material from solely-authored arXiv manuscripts \cites{wang-sw25-cm,wang-sw25-antichain-pre,wang-sw26-plump,wang-sw26-czf-v-l-pre}, as indicated in the main text. \cite{wang-sw25-antichain-pre} has additionally been published online as an open-access FirstView article in the \emph{Journal of Symbolic Logic} at the time of submission.

This copy has been supplied on the understanding that it is copyright material and that no quotation from the thesis may be published without proper acknowledgement.

The right of Shuwei Wang to be identified as Author of this work has been asserted by Shuwei Wang in accordance with the Copyright, Designs and Patents Act 1988.

\frontmatter

\frontsection{Acknowledgements}

I would like to express my foremost gratitude to the primary supervisor of the project, Professor Michael Rathjen, who has been very knowledgeable in all areas relevant to this thesis and provided me with key guidance and insights throughout the duration of the programme. I also thank my other supervisors, Paul Shafer and Asaf Karagila, who took over the administrative duties when Michael retired a few months ago.

I am also grateful to the brilliant people at the Leeds Logic Group, especially my former and current colleagues in the Proof \& Computation reading group, including Emanuele Frittaion and Giorgio Genovesi, for all the interesting discussions we had over the years. And I thank Asaf (again!) for giving the graduate lectures on forcing and large cardinals, which greatly helped with my foundational understanding of set theory.

I appreciate the many conferences that kindly allowed me to present my work and connect with fellow researchers all over the world. My particular gratitude goes to Anton Freund and Keita Yokoyama, who have given me valuable opportunities to deliver long, special session talks as a PhD candidate at Computability in Europe and the Asian Logic Conference in 2025.

I also thank Keita for his detailed comments about my work on Weaver's semi-intuitionistic arithmetic during our multiple conversations. I thank Jem Lord, James E.\ Hanson and Paul Taylor for their in-person or online correspondence and their professional advice in the area of intuitionistic set theory.

Academics aside, I am additionally thankful to my old friends for remaining in contact and for all the fun online as I moved to Leeds on my own. I thank my parents and my extended family for being as supportive as usual during all these years where I pursued my studies abroad and couldn't visit them often.

Finally, I shall thank the Engineering and Physical Sciences Research Council [EP/W523860/1] for their financial support, which made the project possible. I also acknowledge the help of generative AI tools (ChatGPT 5.5, OpenAI) in proofreading this thesis, in accordance with the Leeds postgraduate proofreading policy.

\frontsection{Abstract}

This thesis studies the construction of realisability models from various partial combinatory algebras in formal theories of both arithmetic and set theory. The focus is on the representation of `set objects' (or higher-order objects, respectively) and the implementation of transfinite recursion along well-orderings and set-theoretic ordinals for the purpose of interpreting specific axioms.

For the arithmetic part, we study a semi-intuitionistic theory of third-order conceptual mathematics, $\CM$, proposed by Nik Weaver in \cite{weaver09-cm}. We provide an ordinal analysis through a realisability interpretation in fragments of classical second-order arithmetic and use it to gauge the effect of adding different instances of the transfinite induction scheme on a global well-ordering to the proof-theoretic strength of the system.

The rest of the thesis then implements several realisability arguments over different partial combinatory algebras for constructive Zermelo--Fraenkel set theory. We obtain multiple consistency and independence results surrounding the intuitionistic ordinals and G\"odel's constructible universe $L$ in this way. We also analyse and resolve the open problem of whether adding the axiom $V = L$ will increase the consistency strength of certain constructive set theories, giving a negative answer.

The set-theoretic contribution of this thesis is the method of incomparable codings, which utilises the existence of incomparable ordinals in intuitionistic set theory to encode information about arbitrary sets into a single ordinal. In conjunction with Paul Taylor's notion of plump ordinals in \cite{taylor96-intuitionistic-ordinals}, this leads to a very useful mechanism for verifying exotic properties of $L$ in intuitionistic models.

\tableofcontents

\mainmatter

\chapter{Introduction}

Realisability interpretation is a technique that dates back to Kleene \cite{kleene45-int-number-theory} and serves as an important method for intuitionistic model construction. Compared to other approaches to the semantics of intuitionistic logic such as Kripke frames or Heyting-valued forcing, realisability arguments utilise a partial combinatory algebra to equip a theory with a constructive, often computational interpretation, and can provide key intuition as to why certain principles, such as dependent choice or the collection scheme in set theory, carry much less proof-theoretic strength when used in an intuitionistic context.

The main interest of this thesis is to examine different ways in which the construction of `set objects', both in higher-order arithmetic and in set theories, can be carried out in a realisability model. We aim to study the degree to which various notions of computation --- Turing machines and other topics in higher recursion theory: hyperarithmetic functions, $\alpha$-recursion and set recursion --- can interact with infinitary objects in a realisability model and support transfinite constructions along well-orderings and set-theoretic ordinals.

More specifically, we will look at various extensions and fragments of a semi-intuitionistic theory of third-order arithmetic, and of constructive Zermelo--Fraenkel set theory, $\CZF$. We use a number of realisability constructions to interpret them in well-studied classical systems such as fragments of second-order arithmetic and extensions of Kripke--Platek set theory. As a consequence, we obtain many equiconsistency and independence results for different axioms in these contexts.

In \autoref{chap:pca-intro}, we begin by introducing the formal definition for standard and weak partial combinatory algebras. We will present a general formulation of the realisability conditions for logical connectives and quantifiers (which every realisability model in later chapters will inherit and modify slightly), and observe carefully that the rules of intuitionistic logic can be implemented in this system.

We shall also make an important distinction between \emph{interpreted} realisability models, where the domain of the model remains a subset of the combinatory algebra, and \emph{uninterpreted} realisability models, where entities in the model can be arbitrary constructs in the base theory and application terms only operate on partial information extracted from them. Many arguments in later chapters will exemplify and take advantage of their difference in behaviour.

In \autoref{chap:arith}, we will examine the third-order theory of conceptual mathematics, $\CM$, proposed by Nik Weaver in \cite{weaver09-cm} as an alternative foundation for mainstream mathematics. An ordinal analysis is performed through a realisability argument to determine that:
\def\currthmref{cor:cm-oa}
\begin{refthm}
  $\CM \equiv_{\Con} \Sigma^1_1\DC$, and we have $\left|\CM\right| = \left|\Sigma^1_1\DC\right| = \varphi_{\varepsilon_0}\mleft(0\mright)$.
\end{refthm}

Additionally, we give a detailed discussion of the extension of $\CM$ by a global well-ordering axiom that Weaver mentioned briefly in \cite{weaver09-cm}*{section 2.3}. We will extend our realisability interpretation to show that an unrestricted transfinite induction scheme easily raises the proof-theoretic strength of the extension to the Bachmann--Howard ordinal, $\vartheta_{\varepsilon_{\Omega + 1}}\mleft(0\mright)$. A weaker, predicatively flavoured restriction will be presented, along with some consistency strength bounds for other fragments of this theory of a global well-ordering.

In \autoref{chap:ist-intro}, we then turn towards the study of intuitionistic set theories. We will present the formal statements of the axioms of a hierarchy of set theories: $\IKP$, $\CZF$, $\CZFP$ and $\IZF$. In particular, we will investigate $\CZF$ and the standard realisability interpretation of this theory in itself by Rathjen \cite{rathjen06-czf-realizability}.

The chapter also covers the construction of G\"odel's constructible universe $L$ in an intuitionistic context, and its various exotic behaviours in the given realisability model. We will construct specific examples to conclude that:
\def\currthmref{cor:czf-real-alpha-l-alpha-mismatch}
\begin{refthm}
  $\CZF$ is consistent with $\opE{\alpha \in \Ord} \neg L_\alpha \cap \Ord = \alpha$.
\end{refthm}

\def\currthmref{cor:exo-omega-same-l}
\begin{refthm}
  $\CZF$ is consistent with $\opE{\alpha, \beta \in \Ord} \left(\alpha \neq \beta \land L_\alpha = L_\beta\right)$.
\end{refthm}

In \autoref{chap:incom-coding}, we will develop a method of \emph{incomparable codings} to make use of these intuitionistic ordinals that are incomparable under the $\in$- or $\subseteq$-relation. We show that, given the existence of sufficiently many such incomparable objects, one may encode information about the internal structure of an arbitrary set into a single ordinal, and thus implement scenarios where a highly `non-constructive' set can be found inside $L$. We cover the key result from my manuscript \cite{wang-sw25-antichain-pre} that:
\def\currthmref{thm:set-perp-v-eq-l-equiv}
\begin{refthm}[$\IKP$]
  Assume there exist weakly incomparable $\alpha, \beta \in \Ord$, then the following are equivalent:
  \begin{enumerate}[label=\uplabel{\roman*}]
    \item $\Ord \subseteq L$,
    \item $V = L$.
  \end{enumerate}
\end{refthm}

As another application of the technique, we will also work with a variant of Aczel's realisability interpretation of $\CZF$ in Martin-L\"of type theory \cite{aczel78-type-theoretic-cst} and use incomparable codings to prove $V = L$ inside the model. This allows us to justify the following equiconsistency result, where the classical, trivial proof of the similar fact in $\ZF$ fails to apply (cf.\ \autoref{sec:overview-v-eq-l} for more detail):
\def\currthmref{cor:czf-equi-v-eq-l}
\begin{refthm}
  $\CZF \equiv_{\Con} \CZF + V = L$.
\end{refthm}

In \autoref{chap:plump-power}, we give a formal treatment of Paul Taylor's notion of plump ordinals in \cite{taylor96-intuitionistic-ordinals}, and demonstrate how it leads to a semi-automatic process for constructing incomparable coding objects. We will combine this result with the weak realisability interpretation of $\CZFP$ in a classical base theory $\KPP$ in \cite{rathjen12-czf-powerset} to extend the construction above to a model satisfying the powerset axiom, and show that $\CZFP$ also remains equiconsistent with its extension $\CZFP + V = L$.

In \autoref{chap:large-pcas}, finally, we will discuss how to embed a larger classical fragment into an intuitionistic realisability universe via the use of large combinatory algebras including the $\alpha$-recursive and set recursive functions. We will explore how such realisability models may support a direct construction of large pairwise incomparable functions, unlike the hierarchy method employed in \autoref{chap:plump-power}, and use this to prove a realisability-based analogue of an existing result of mine in \cite{wang-sw26-plump}*{section 5} for Heyting-valued forcing:
\def\currthmref{thm:alpha-rec-pow-in-l}
\begin{refthm}[$\CZF^- + \REM$]
  For any ordinal $\beta$, if there exists an admissible ordinal $\alpha$ such that $\beta \subseteq \alpha$, then there exists a \uppar{weak} pca $\mathcal{A}$ such that $\beta \subseteq \mathcal{A}$ and the corresponding realisability model
  \[V\mleft(\mathcal{A}\mright) \vDash \Pow\mleft(\check{\beta}\mright) \subseteq L.\]
\end{refthm}

\chapter{Partial combinatory algebras and realisability}
\label{chap:pca-intro}

To begin this thesis, we will formally define the notions of a \emph{partial combinatory algebra} (henceforth abbreviated as pca) and realisability models. Most things covered in this chapter should be well-known folklore (cf.\ Beeson \cite{beeson85-constructive-maths}*{Chapters VI \& VII}), but we will give a detailed account here with the aim of presenting a self-contained exposition on the realisability arguments employed in the thesis.

Additionally, we observe that the notations for application terms vary a lot in existing literature, and this chapter shall serve as the notational convention to be followed in the rest of the chapters. In \autoref{sec:real-con-int-logic}, we also present a Hilbert-style system similar to Kleene's original formulation \cite{kleene52-metamathematics}, which should be understood as the precise meaning of `first-order intuitionistic logic' throughout this thesis.

\section{Partial combinatory algebras}
\label{sec:pca}

A \emph{partial combinatory algebra} (\emph{pca}\indg{pca}) is a set\footnote{Here we develop some general theory about pcas without fixing a particular formal base theory, so we use the word `set' in an informal sense, and we actually allow the domain $\mathcal{A}$ to be a proper definable class in some theories as well.} $\mathcal{A}$ together with a partial function $\mathrm{ap} : \mathcal{A}^2 \rightharpoonup \mathcal{A}$ denoting the \emph{function application} operation and two distinguished objects $\bk, \bs \in \mathcal{A}$ satisfying their respective defining equations.

Before laying out these equations, we first stipulate that, since function application is the sole core instrumentation in all our realisability models, we will reserve the juxtaposition notation for it, i.e.\ we will simply write $ab$ to denote $\mathrm{ap}\mleft(a, b\mright)$. Numerical multiplication, instead, will always be denoted with a dot as $a \cdot b$ in this thesis.

Following conventions, juxtaposition here shall associate to the left, i.e.\ the expression $abc$ should be read as $\mleft(ab\mright) c$ or $\mathrm{ap}\mleft(\mathrm{ap}\mleft(a, b\mright), c\mright)$. We also write $ab {\downarrow}$ and say that the application \emph{converges} to indicate that $ab = \mathrm{ap}\mleft(a, b\mright) \in \mathcal{A}$ is defined; we say the application \emph{diverges} otherwise.

Now, the axioms of a pca shall enforce that:
\begin{enumerate}[label=\uplabel{\roman*}]
  \item for any $a, b \in \mathcal{A}$, $\bk a {\downarrow}$, $\bk ab {\downarrow}$, and
        \[\bk ab = a;\]
  \item for any $a, b, c \in \mathcal{A}$, $\bs a {\downarrow}$, $\bs ab {\downarrow}$, and\footnote{When reasoning about partial combinatory algebras, we follow usual convention and allow one to write potentially undefined terms like $ac(bc)$ without requiring that $ac {\downarrow}$ or $bc {\downarrow}$ in the first place; we consider the term $ac(bc)$ to diverge if either $ac$ or $bc$ already diverges.

          Additionally, for simplicity we will abuse the equality symbol between two application terms in this thesis, instead of introducing a new symbol $\simeq$ as in some other literature; when we write $t = s$ where $t$ and $s$ are potentially undefined application terms, this refers to the assertion that either $t$ and $s$ both diverge, or that they both converge and are equal.}
        \[\bs abc = ac \mleft(bc\mright).\]
\end{enumerate}

In practice, we will work with objects in a pca through $\lambda$-abstraction. We say that the \emph{formal application terms} of $\mathcal{A}$ over a finite enumeration of variables $x_1, \dots, x_k$ are defined inductively as follows:
\begin{enumerate}[label=\uplabel{\arabic*}]
  \item any variable $x_i$ or any element in $\mathcal{A}$ is a formal application term;
  \item if $t, s$ are formal application terms, then $\mleft(ts\mright)$ is a formal application term;
  \item nothing else is a formal application term.
\end{enumerate}
We shall abuse notation and identify any closed formal application term $t$ (i.e.\ a formal application term in which no variables occur) with the result of its computation, i.e.\ either an element in $\mathcal{A}$ or undefined, and write that $t {\downarrow}$ if the computation converges. For any formal application term $t$ over variables $x_1, \dots, x_k$ and any $a_1, \dots, a_k \in \mathcal{A}$, the \emph{evaluation} $t\mleft[a_1/x_1, \ldots, a_k/x_k\mright]$ is then the result of substituting each $x_i$ by the corresponding $a_i$. We verify that:

\begin{lemma}[$\lambda$-abstraction\indg{lambda-abstraction@$\lambda$-abstraction}]
  \label{lem:lambda-abstraction}
  For any formal application term $t$ over variables $x, y_1, \ldots, y_k$, there is a formal application term, henceforth denoted as $\oplam{x} t$, over variables $y_1, \ldots, y_k$ only, such that
  \begin{enumerate}[label=\uplabel{\roman*}]
    \item for any $b_1, \ldots, b_k \in \mathcal{A}$, $\mleft(\oplam{x} t\mright)\mleft[b_1/y_1, \ldots, b_k/y_k\mright] {\downarrow}$;
    \item for any $a, b_1, \ldots, b_k \in \mathcal{A}$,
          \[\mleft(\oplam{x} t\mright)\mleft[b_1/y_1, \ldots, b_k/y_k\mright] a = t\mleft[a/x, b_1/y_1, \ldots, b_k/y_k\mright].\]
  \end{enumerate}
\end{lemma}

\begin{proof}
  We prove this by induction on the complexity of $t$. If $t = x$, take $\oplam{x} t = \bs \bk \bk$; if $t = y_i$ or $t = a \in \mathcal{A}$, simply take $\oplam{x} t = \bk t$.

  If $t = \left(rs\right)$ where $r, s$ are formal application terms, by the inductive hypothesis we then have formal application terms $\oplam{x} r$ and $\oplam{x} s$. Taking $\oplam{x} t = \bs \mleft(\oplam{x} r\mright) \mleft(\oplam{x} s\mright)$ would suffice.
\end{proof}

As an immediate consequence, any formal application term of the form $\oplam{x} t$ where $t$ has no variables other than $x$ (including the case where $t$ contains additional $\lambda$-terms which bind more variables) is a closed formal application term which \emph{must} converge, i.e.\ $\oplam{x} t$ is a definite element in $\mathcal{A}$.

In what follows, we will abbreviate $\oplam{x_1} \oplam{x_2} \ldots \oplam{x_k} t$ as $\oplam{x_1 x_2 \cdots x_k} t$, or simply as $\oplam{\vec{x}} t$ where $\vec{x} = x_1, \ldots, x_k$ denotes a sequence of variables; correspondingly, we will write $ts_1 \cdots s_k$ as $t \vec{s}$ where $\vec{s} = \tuple{s_1, \ldots, s_k}$ is a tuple of (formal) application terms.

\begin{corollary}[Recursion theorem]
  \label{cor:recursion-theorem}
  There must be a distinguished element $\br \in \mathcal{A}$ such that for any $f, a \in \mathcal{A}$, $\br f {\downarrow}$ and
  \[\br fa = fa \mleft(\br f\mright).\]
\end{corollary}

\begin{proof}
  We build the well-known `Y-combinator'. Let $t = \oplam{fga} fa \mleft(gg\mright)$, and observe that
  \[tf \mleft(tf\mright) a = fa \mleft(tf \mleft(tf\mright)\mright).\]
  Therefore, simply taking $\br = \oplam{f} tf \mleft(tf\mright)$ would suffice. Note that $\br f {\downarrow}$ is ensured because for any $f, g \in \mathcal{A}$, we have
  \[tfg = \mleft(\oplam{a} fa \mleft(gg\mright)\mright) {\downarrow}\]
  by \autoref{lem:lambda-abstraction}.
\end{proof}

\subsection*{Pairing functions}

We shall maintain the convention that every pca $\mathcal{A}$ is equipped with a distinguished \emph{pairing function}\indg{pairing function}, i.e.\ some $\bp \in \mathcal{A}$ together with unpairing operations $\bp_0, \bp_1 \in \mathcal{A}$ such that for any $a, b \in \mathcal{A}$,
\[\bp_0 \mleft(\bp ab\mright) = a \quad \text{and} \quad \bp_1 \mleft(\bp ab\mright) = b.\]

We do not require that the same (uniformly-defined) pairing function is used in all the pcas considered in this thesis, since in some scenarios we may want our pairing function to possess additional properties --- for example, when the domain $\mathcal{A} = \mathbb{N}$ (cf.\ the next section for pcas over natural numbers), we may want the pairing function to be a bijection $\mathbb{N}^2 \rightarrow \mathbb{N}$, so that $\bp_0, \bp_1$ will also be total operations. However, do note that some default construction will work for an arbitrary pca, such as the following one mentioned in \cite{frittaion-rathjen20-realizability-ist} (which is a typical Church encoding):

\begin{lemma}[\cite{frittaion-rathjen20-realizability-ist}*{footnote 5}]
  Let $\bp = \oplam{abc} cab$, $\bp_0 = \oplam{p} p \mleft(\oplam{ab} a\mright)$ and $\bp_1 = \oplam{p} p \mleft(\oplam{ab} b\mright)$. Then $\bp, \bp_0, \bp_1$ give a pairing function.
\end{lemma}

\begin{proof}
  Trivial.
\end{proof}

For convenience, we will also adopt the notation $\bp_{i_0 i_1 \cdots i_n}$ where each $i_k \in \left\{0, 1\right\}$ as the sequential composition of the unpairing operations, i.e.
\[\bp_{i_0 i_1 \cdots i_n} = \oplam{x} \bp_{i_n} \mleft(\bp_{i_{n - 1}}  \cdots \mleft(\bp_{i_0} x\mright) \cdots \mright).\]
Take note of the reversal of indices as I find this notation more intuitive to work with.

\section{Natural numbers in pcas}

We shall follow the terminology of \cite{frittaion-rathjen20-realizability-ist} and use `pca over natural numbers' to refer to a pca with the infrastructure to represent natural numbers and recursive functions on them. This thesis will not concern itself with a full formal theory of such structures, but rather views them as mere machinery within an established base theory to facilitate realisability constructions. Therefore, we directly say that a \emph{pca over natural numbers}\indg{pca!over natural numbers} is a pca $\mathcal{A}$ with additional distinguished objects $\bsN, \bpN, \bdN \in \mathcal{A}$ satisfying that:
\begin{enumerate}[label=\uplabel{\roman*}]
  \item the natural numbers (in the base theory) $\mathbb{N} \subseteq \mathcal{A}$;
  \item for any $n \in \mathbb{N}$, $\bsN n = n + 1$ and\footnote{Some sources would additionally require that $\bpN 0 {\downarrow} = 0$. We do not demand this here, since one can always use $\bdN$ in the next clause to test for zero first before applying $\bpN$ safely.} $\bpN \mleft(\bsN n\mright) = n$;
  \item for any $n, m \in \mathbb{N}$, $a, b \in \mathcal{A}$,
        \[\bdN nmab = \left\{\begin{aligned}
             & a &  & \text{if $n = m$},    \\
             & b &  & \text{if $n \neq m$}.
          \end{aligned}\right.\]
\end{enumerate}

There is a subtlety to note here that, given our definition of term application on a pca, a term would only converge when all of its subterms converge, even when it is the case operation $\bdN nmab$ that should always disregard one of its last two arguments. In order to have a `lazy' alternative that does not require the convergence of all its subterms, we will enforce the abbreviation that
\[\bdN^* nmab = \bdN nm \mleft(\oplam{z} a\mright) \mleft(\oplam{z} b\mright) 0\]
throughout this thesis, where $z$ is a distinct new variable and $a, b$ can be complex (formal) application terms so that $\oplam{z} a$, $\oplam{z} b$ are implemented according to \autoref{lem:lambda-abstraction}, and $\bdN nm \mleft(\oplam{z} a\mright) \mleft(\oplam{z} b\mright)$ is guaranteed to converge. Note that $\bdN^*$ itself cannot be an actual element in $\mathcal{A}$, due to precisely the evaluation rule mentioned above.

Now, it is easy to verify that this pca $\mathcal{A}$ must contain a name for every partial recursive function on natural numbers, in the following sense:

\begin{lemma}
  \label{lem:rec-in-pca}
  For any partial recursive function $f : \mathbb{N}^k \rightharpoonup \mathbb{N}$, there is some $e \in \mathcal{A}$ such that for any tuple $\vec{n} = \tuple{n_1, \ldots, n_k} \in \mathbb{N}^k$,
  \[f\mleft(\vec{n}\mright) {\downarrow} \rightarrow e \vec{n} {\downarrow} = f\mleft(\vec{n}\mright).\]
\end{lemma}

\begin{proof}
  We use the standard characterisation that the collection of partial recursive functions is the closure of the successor, constant and projection functions under composition, primitive recursion and minimisation (cf.\ \cite{soare87-re-book}*{Chapter I}). Now, given any $g, h \in \mathcal{A}$, primitive recursion
  \[fx\vec{y} = \left\{\begin{aligned}
       & g\vec{y}                                                                       &  & \text{if $x = 0$}, \\
       & h \mleft(\bpN x\mright) \mleft(f \mleft(\bpN x\mright) \vec{y}\mright) \vec{y} &  & \text{if $x > 0$}
    \end{aligned}\right.\]
  is given simply by $f = \br \mleft(\oplam{xt \vec{y}} \bdN^* x0 \mleft(g \vec{y}\mright) \mleft(h \mleft(\bpN x\mright) \mleft(t \mleft(\bpN x\mright) \vec{y}\mright) \vec{y}\mright)\mright)$; minimisation
  \[f\vec{y} = \min \left\{x \in \mathbb{N} : gx\vec{y} {\downarrow} = 0 \land \opA{z < x} gz\vec{y} {\downarrow}\right\}\]
  (if such $x$ exists) is then $f = \oplam{\vec{y}} \br \mleft(\oplam{xt} \bdN^* \mleft(gx \vec{y}\mright) 0 x \mleft(t \mleft(\bsN x\mright)\mright)\mright) 0$, where we recurse upwards (as long as $gx\vec{y} {\downarrow}$) until the first $x$ such that $gx\vec{y} = 0$.
\end{proof}

Indeed, consider the application function $\mathrm{ap} : \mathbb{N}^2 \rightharpoonup \mathbb{N}$ given by Turing machine computation, i.e.\ $ea = b$ if the Turing machine with index $e$ halts on input $a$ and produces output $b$, then it is easy to observe that the distinguished operations $\bk, \bs, \bsN, \bpN, \bdN$ are all themselves recursive (for $\bs$ by the existence of the universal machine). This gives the standard pca on domain $\mathbb{N}$ known as \emph{Kleene's first algebra}\indg{Kleene's first algebra, $\Kl$}, which we denote as $\Kl$.

\section{Realisability of intuitionistic logic}
\label{sec:real-con-int-logic}

Given a first-order language $\mathcal{L}$ and a fixed domain $U$, in this thesis we will work with terms/formulae together with some parameters from $U$ as formal objects in the base theory.

A \emph{term with parameters} is a term $t \in \mathcal{L}$ together with a finite tuple of variable assignments $a_1/x_1, \ldots, a_k/x_k$ where each $a_i \in U$. This will be presented using the substitution notation $t\mleft[a_1/x_1, \ldots, a_k/x_k\mright]$ and, whenever there is no ambiguity, we may abuse notation and write, for example, simply $f\mleft(a\mright)$ instead of $f\mleft(x\mright)\mleft[a/x\mright]$. Likewise, a \emph{formula with parameters} is a formula $\varphi \in \mathcal{L}$ together with a finite tuple of variable assignments $a_1/x_1, \ldots, a_k/x_k$ where each $a_i \in U$, denoted as $\varphi\mleft[a_1/x_1, \ldots, a_k/x_k\mright]$; and we will sometimes abuse notation and write $\varphi\mleft(a\mright)$ instead.

We say that a term/formula with parameters is \emph{closed} if all free variables of the expression appear in the accompanying variable assignments; otherwise it is \emph{open} and its \emph{free variables} are those that lack assignments.

Given a first-order language $\mathcal{L}$, an \emph{interpretation} of $\mathcal{L}$-terms (with parameters) in a pca $\mathcal{A}$ consists of a domain $U \subseteq \mathcal{A}$ together with, for each $k$-ary function symbol $f$ in $\mathcal{L}$, some $f^\mathcal{A} \in \mathcal{A}$ such that for any $\vec{a} \in U^k$, $f^\mathcal{A} \vec{a} {\downarrow} \in U$. The interpretation clearly extends to all $\mathcal{L}$-terms, where for each closed term $t$ we correlate some $t^\mathcal{A} \in U$; and for each term $t$ with (an explicit, complete enumeration of) free variables $x_1, \ldots, x_k$, we have a function $t^\mathcal{A} \in \mathcal{A}$ such that for any $\vec{a} = \tuple{a_1, \ldots, a_k} \in U^k$,
\[t^\mathcal{A} \vec{a} {\downarrow} = t\mleft[a_1/x_1, \ldots, a_k/x_k\mright]^\mathcal{A}.\]

Given some \emph{interpretation} of $\mathcal{L}$-terms in a pca $\mathcal{A}$ over natural numbers\footnote{Note that the sole need for a pca over natural numbers here is for the disjunction clause, where one builds a disjoint union using $0$ and $1$. This can be avoided through a Church encoding of disjoint unions instead, for example, by requiring that $d \Vdash \varphi \lor \psi$ if and only if either $\bp_0 d {\downarrow} \Vdash \varphi$ and for any $f, g \in \mathcal{A}$, $\bp_1 dfg = f \mleft(\bp_0 d\mright)$, or $\bp_0 d {\downarrow} \Vdash \psi$ and for any $f, g \in \mathcal{A}$, $\bp_1 dfg = g \mleft(\bp_0 d\mright)$ (so that the realiser $d$ could be constructed as either $\bp e \mleft(\oplam{fg} fe\mright)$ for $e \Vdash \varphi$ or $\bp e \mleft(\oplam{fg} ge\mright)$ for $e \Vdash \psi$). But we opt for the simpler, conventional approach using natural numbers in this thesis.}, a collection\footnote{There is a caveat that sometimes, this collection will not be internally definable inside the base theory. See the discussion on page \pageref{loc:real-con-def-caveat} for more detail.} of \emph{realisability conditions}\indg{realisability conditions} then consists of pairs $\tuple{d, \varphi}$ where $d \in \mathcal{A}$ and $\varphi$ is a closed $\mathcal{L}$-formula (with parameters), where we say that `$d$ realises $\varphi$' and write $d \Vdash \varphi$ if $\tuple{d, \varphi}$ is in the collection.

How atomic $\varphi$'s are realised is usually specified on a case-by-case basis, depending on whether the truth conditions of the atomic formulae in our desired model are computable by the pca; however, most realisability models would extend its atomic conditions to the collection of all formulae through the following inductive clauses:
\begin{align*}
  d \Vdash \varphi \land \psi       & \defeqv \bp_0 d {\downarrow} \Vdash \varphi \land \bp_1 d {\downarrow} \Vdash \psi,                                                              \\
  d \Vdash \varphi \lor \psi        & \defeqv \left(\bp_0 d = 0 \land \bp_1 d {\downarrow} \Vdash \varphi\right) \lor \left(\bp_0 d = 1 \land \bp_1 d {\downarrow} \Vdash \psi\right), \\
  d \Vdash \neg \varphi             & \defeqv \opA{e \in \mathcal{A}} \neg e \Vdash \varphi,                                                                                           \\
  d \Vdash \varphi \rightarrow \psi & \defeqv \opA{e \in \mathcal{A}} \left(e \Vdash \varphi \rightarrow de {\downarrow} \Vdash \psi\right),                                           \\
  d \Vdash \opA{x} \varphi          & \defeqv \opA{a \in U} da {\downarrow} \Vdash \varphi\mleft[a/x\mright],                                                                          \\
  d \Vdash \opE{x} \varphi          & \defeqv \bp_0 d {\downarrow} \in U \land \bp_1 d {\downarrow} \Vdash \varphi\mleft[\bp_0 d/x\mright].
\end{align*}

To align better with the usual notion of proofs in first-order logic (where having open formulae as rows in a proof is crucial to many systems), we shall further extend the notion to open formulae: for any formula (with parameters) $\varphi \in \mathcal{L}$ with its free variables enumerated as $x_1, \ldots, x_k$, we say that $d \in \mathcal{A}$ realises $\varphi$ if and only if it realises the universal closure $\opA{x_1, \ldots, x_k} \varphi$, i.e.\ if and only if for any $\vec{a} = \tuple{a_1, \ldots, a_k} \in U^k$,
\[d \vec{a} \Vdash \varphi\mleft[a_1/x_1, \ldots, a_k/x_k\mright].\]
Obviously the order of the variable enumeration does not affect the existence of a realiser, nor does the inclusion of additional free variables that do not actually occur in $\varphi$ in the enumeration.

We can now fix a formal, intuitionistic proof system of first-order logic, for example the following Hilbert-style system that was (roughly) originally used by Kleene \cite{kleene52-metamathematics}*{Chapter IV} (cf.\ \cite{moschovakis-joan24-intuitionistic-logic} and \cite{troelstra-dalen88-constructivism-1}*{Chapter 2}) with the axioms
\begin{enumerate}[label=\uplabel{\arabic*}]
  \item $\varphi \rightarrow \left(\psi \rightarrow \varphi\right)$,
  \item $\left(\varphi \rightarrow \psi\right) \rightarrow \left(\left(\varphi \rightarrow \left(\psi \rightarrow \sigma\right)\right) \rightarrow \left(\varphi \rightarrow \sigma\right)\right)$,
  \item $\varphi \rightarrow \left(\psi \rightarrow \varphi \land \psi\right)$,
  \item $\varphi \land \psi \rightarrow \varphi$,
  \item $\varphi \land \psi \rightarrow \psi$,
  \item $\varphi \rightarrow \varphi \lor \psi$,
  \item $\psi \rightarrow \varphi \lor \psi$,
  \item \label{itm:int-logic-disj-elim} $\left(\varphi \rightarrow \sigma\right) \rightarrow \left(\left(\psi \rightarrow \sigma\right) \rightarrow \left(\varphi \lor \psi \rightarrow \sigma\right)\right)$,
  \item $\left(\varphi \rightarrow \psi\right) \rightarrow \left(\left(\varphi \rightarrow \neg \psi\right) \rightarrow \neg \varphi\right)$,
  \item $\neg \varphi \rightarrow \left(\varphi \rightarrow \psi\right)$,
  \item $\opA{x} \varphi \rightarrow \varphi\mleft[t/x\mright]$ where $t$ is any term which is free\footnote{In other words, no free variables in $t$ will become bound when substituted for $x$ in $\varphi$.} for $x$ in $\varphi$,
  \item $\varphi\mleft[t/x\mright] \rightarrow \opE{x} \varphi$ where $t$ is any term which is free for $x$ in $\varphi$,
\end{enumerate}
for any formulae $\varphi$, $\psi$ and $\sigma$, together with the following inference rules:
\begin{enumerate}[label=\uplabel{\roman*}]
  \item \emph{modus ponens}: from $\varphi$ and $\varphi \rightarrow \psi$, deduce $\psi$;
  \item \emph{$\forall$-introduction}: from $\varphi \rightarrow \psi$ where $x$ does not occur free in $\varphi$, deduce $\varphi \rightarrow \opA{x} \psi$;
  \item \emph{$\exists$-elimination}: from $\varphi \rightarrow \psi$ where $x$ does not occur free in $\psi$, deduce $\opE{x} \varphi \rightarrow \psi$.
\end{enumerate}

\begin{proposition}[Soundness of intuitionistic logic in realisability models]
  \label{prop:real-con-int-logic}
  Any logical axiom of the intuitionistic system is realised in any realisability model \uppar{of the form above}. Moreover, for any inference rule of the system, if all antecedents are realised, then so is the consequent.
\end{proposition}

\begin{proof}
  Given any logical axiom (which is potentially an open formula), let $\vec{x}$ enumerate the free variables in it. Then it is easy to observe that:
  \begin{enumerate}[label=\uplabel{\arabic*}]
    \item $\oplam{\vec{x} de} d$ realises $\varphi \rightarrow \left(\psi \rightarrow \varphi\right)$;
    \item $\oplam{\vec{x} fgd} gd \mleft(fd\mright)$ realises $\left(\varphi \rightarrow \psi\right) \rightarrow \left(\left(\varphi \rightarrow \left(\psi \rightarrow \sigma\right)\right) \rightarrow \left(\varphi \rightarrow \sigma\right)\right)$;
    \item $\oplam{\vec{x}} \bp$ realises $\varphi \rightarrow \left(\psi \rightarrow \varphi \land \psi\right)$;
    \item $\oplam{\vec{x} d} \bp_0 d$ realises $\varphi \land \psi \rightarrow \varphi$;
    \item $\oplam{\vec{x} d} \bp_1 d$ realises $\varphi \land \psi \rightarrow \psi$;
    \item $\oplam{\vec{x} d} \bp 0d$ realises $\varphi \rightarrow \varphi \lor \psi$;
    \item $\oplam{\vec{x} d} \bp 1d$ realises $\psi \rightarrow \varphi \lor \psi$;
    \item $\oplam{\vec{x} fgd} \bdN \mleft(\bp_0 d\mright) 0fg \mleft(\bp_1 d\mright)$, whose inner term evaluates to $f \mleft(\bp_1 d\mright)$ when the first disjunct of $\varphi \lor \psi$ is realised and to $g \mleft(\bp_1 d\mright)$ otherwise, realises $\left(\varphi \rightarrow \sigma\right) \rightarrow \left(\left(\psi \rightarrow \sigma\right) \rightarrow \left(\varphi \lor \psi \rightarrow \sigma\right)\right)$;
    \item $\oplam{\vec{x} fg} 0$ realises $\left(\varphi \rightarrow \psi\right) \rightarrow \left(\left(\varphi \rightarrow \neg \psi\right) \rightarrow \neg \varphi\right)$, since suppose otherwise\footnote{This appears to be a proof-by-contradiction argument, but it applies in an intuitionistic base theory as well, since the statement to be proved, $\opA{e \in \mathcal{A}} \neg e \Vdash \varphi$, is itself a negated statement, thus it suffices to simply derive a contradiction.}, i.e.\ $f \Vdash \varphi \rightarrow \psi$, $g \Vdash \varphi \rightarrow \neg \psi$, yet some $d \Vdash \varphi$, then $fd \Vdash \psi$ and $gd \Vdash \neg \psi$, contradicting the stipulation that $\neg \psi$ is realised only when no realisers of $\psi$ exist;
    \item $\oplam{\vec{x} f} 0$ realises $\neg \varphi \rightarrow \left(\varphi \rightarrow \psi\right)$, since the existence of $f \Vdash \neg \varphi$ and $g \Vdash \varphi$ again implies a contradiction. Thus $f \Vdash \neg \varphi$, $g \Vdash \varphi$ entails $0g {\downarrow} \Vdash \psi$ as \emph{ex falso quodlibet}.
          \xdef\intlogicproofquantcount{\number\value{enumi}+1}
  \end{enumerate}

  For the quantifier axioms, assume that $t^\mathcal{A} \in \mathcal{A}$ takes the same enumeration of arguments $\vec{x}$ above as input, then
  \begin{enumerate}[label=\uplabel{\arabic*},resume]
    \item $\oplam{\vec{x} f} f \mleft(t^\mathcal{A} \vec{x}\mright)$ realises $\opA{x} \varphi \rightarrow \varphi\mleft[t/x\mright]$;
    \item $\oplam{\vec{x} d} \bp \mleft(t^\mathcal{A} \vec{x}\mright) d$ realises $\varphi\mleft[t/x\mright] \rightarrow \opE{x} \varphi$.
  \end{enumerate}

  For the inference rules, check that:
  \begin{enumerate}[label=\uplabel{\roman*}]
    \item Assume that $d \Vdash \varphi$ and $f \Vdash \varphi \rightarrow \psi$, which (without loss of generality) both involve the same enumeration of free variables $\vec{x}$, then $\oplam{\vec{x}} f \vec{x} \mleft(d \vec{x}\mright)$ realises $\psi$;
          \xdef\intlogicrulesquantcount{\number\value{enumi}+1}
    \item Assume $f \Vdash \varphi \rightarrow \psi$ where the free variables are enumerated as $x, \vec{y}$ and $x$ does not occur free in $\varphi$, then $\oplam{\vec{y} dx} fx \vec{y} d$ realises $\varphi \rightarrow \opA{x} \psi$;
    \item Assume $f \Vdash \varphi \rightarrow \psi$ where the free variables are enumerated as $x, \vec{y}$ and $x$ does not occur free in $\psi$, then $\oplam{\vec{y} d} f \mleft(\bp_0 d\mright) \vec{y} \mleft(\bp_1 d\mright)$ realises $\opE{x} \varphi \rightarrow \psi$. \qedhere
  \end{enumerate}
\end{proof}

It follows immediately, by induction on the length of a proof, that given any intuitionistic theory $T$ and a realisability model that realises all axioms of $T$, any theorem of $T$ is realised as well.

Now, there is an \textbf{important} caveat\label{loc:real-con-def-caveat} concerning the formalisation of this argument: depending on the choice of a base theory and the pca $\mathcal{A}$ --- especially when working with a `tight' interpretation between two theories of the same proof-theoretic strength --- the inductively-formulated collection of all realisability conditions may not even form a definable class in the base theory. So we actually have two distinct kinds of formal realisability arguments:
\begin{enumerate}[label=\uplabel{\alph*}]
  \item When the collection $\left\{\tuple{d, \varphi} : d \Vdash \varphi\right\}$ of all realisability conditions is definable inside some base theory $S$, then the aforementioned induction on the length of a proof can also be carried out inside $S$, and thus the realisability argument simply yields a proof of $S \vdash \Con T$ since no $d \in \mathcal{A}$ can realise $\bot$ by stipulation.
  \item Alternatively, when the said collection cannot be defined internally, we could instead formulate a metatheoretic family of collections
        \[r_\varphi = \left\{\tuple{d, \overline{\varphi}\mleft[\vec{a}/\vec{x}\mright]} : \vec{a} \in U^k, d \Vdash \overline{\varphi}\mleft[\vec{a}/\vec{x}\mright]\right\},\]
        each of which alone is definable inside the theory $S$, where $\overline{\varphi}$ is the $S$-term denoting the `standard' formula $\varphi$ in the metatheory, for example a numerical constant representing the G\"odel numbering of $\varphi$. It then follows that $r_\bot = \varnothing$, but the induction on proof length must be performed in the metatheory, resulting in an argument for $\vdash \Con S \rightarrow \Con T$ there.
\end{enumerate}

The latter option is equivalent to the traditional interpretation-by-\emph{translation} approach (cf.\ \cite{beeson85-constructive-maths}*{Chapter VII}) which maintains that there is a metatheoretic function mapping each formula $\varphi$ in the language of theory $T$ to a formula in the language of $S$ with one additional free variable $d$ that captures the denotation of `$d$ realises $\varphi$'. One ensures that proofs in intuitionistic logic are preserved and $\bot$ is unchanged under the translation, thus verifying that the consistency of $S$ implies that of $T$.

In this thesis, however, I prefer the current presentation as I want to put more emphasis on the objects in the chosen domain $U$ inside the base theory and the fact that we can still operate on these intuitionistic constructs (and speak about which formulae in the language of $T$ `hold' on them) in a very model-theoretic way, even in the second scenario. This intuition shall become very relevant in the later chapters on intuitionistic set theory.

\subsection*{Equality axioms}
\indg{equality axioms}

Across the variety of formal theories considered in this thesis, we allow some freedom in treating the equality symbol as either a primitive, logical relation symbol or a defined abbreviation (usually through extensional equality for sets) --- or, in the case of higher-order languages, as a primitive symbol for first-order variables, but as abbreviations for variables of higher orders, cf.\ \autoref{sec:cm-axioms}.

In the case where equality is primitive, we will need the following additional logical axioms:
\begin{enumerate}[label=\uplabel{E\arabic*}]
  \item \label{itm:eq-refl-axiom} reflexivity, $x = x$;
  \item symmetry, $x = y \rightarrow y = x$;
  \item transitivity, $x = y \rightarrow \left(y = z \rightarrow x = z\right)$;
  \item \label{itm:eq-sub-axiom} the substitution scheme, $x = y \rightarrow \left(\varphi\mleft(x\mright) \rightarrow \varphi\mleft(y\mright)\right)$ for any atomic formula $\varphi\mleft(x\mright)$.
\end{enumerate}

Since each realisability model will provide its own interpretation of a primitive equality symbol, the equality axioms need to be verified on a case-by-case basis. However, in several realisability interpretations we see a common pattern of enforcing
\[d \Vdash x = y \defeqv x = y\]
for any $d \in \mathcal{A}$ and $x, y \in U$, in which case we say that equality is \emph{realised trivially}. We then have the following and do not have to worry about the equality axioms further:
\begin{lemma}
  \label{lem:eq-axiom-trivial}
  If a realisability model realises equality trivially, then all the equality axioms are immediately realised.
\end{lemma}

\begin{proof}
  We simply check that:
  \begin{enumerate}[label=\uplabel{E\arabic*}]
    \item $\oplam{x} 0$ realises $x = x$;
    \item $\oplam{xyd} 0$ realises $x = y \rightarrow y = x$;
    \item $\oplam{xyzde} 0$ realises $x = y \rightarrow \left(y = z \rightarrow x = z\right)$;
    \item $\oplam{xy\vec{z}de} e$ realises $x = y \rightarrow \left(\varphi\mleft(x\mright) \rightarrow \varphi\mleft(y\mright)\right)$, where $\vec{z}$ enumerates the remaining free variables in $\varphi\mleft(x\mright)$, since for any $a, b \in U$, $\vec{c} \in U^k$, if $a = b$, then $\varphi\mleft(x\mright)\mleft[a/x, \vec{c}/\vec{z}\mright]$ and $\varphi\mleft(y\mright)\mleft[b/y, \vec{c}/\vec{z}\mright]$ are simply the same formula (with same parameters). \qedhere
  \end{enumerate}
\end{proof}

\section{Uninterpreted realisability models}
\label{sec:uninterp-real}

In other variants of realisability arguments, we often wish to choose a domain $U$ that is not restricted by the size of the involved pca $\mathcal{A}$, for example, when we want to build an uncountable set-theoretic universe but still use a countable pca such as $\Kl$ for the realisers. In this thesis we call such constructions \emph{uninterpreted realisability models}\indg{realisability conditions!uninterpreted}, where we make the following modifications to the standard argument in \autoref{sec:real-con-int-logic}:
\begin{enumerate}[label=\uplabel{\roman*}]
  \item The domain $U$ is a collection (i.e.\ set or definable class) in the base theory but not necessarily a subset of $\mathcal{A}$. For each term with parameters $t \in \mathcal{L}$, its interpretation $t^\mathcal{A}$ should be a function that maps assignments of its free variables to an element in $U$, but it need not be computable by $\mathcal{A}$.
  \item The realisability conditions for quantifiers are instead:
        \begin{align*}
          d \Vdash \opA{x} \varphi & \defeqv \opA{a \in U} d \Vdash \varphi\mleft[a/x\mright], \\
          d \Vdash \opE{x} \varphi & \defeqv \opE{a \in U} d \Vdash \varphi\mleft[a/x\mright].
        \end{align*}
  \item For any formula (with parameters) $\varphi \in \mathcal{L}$ with its free variables enumerated as $x_1, \ldots, x_k$, we still say that $d \in \mathcal{A}$ realises $\varphi$ if and only if it realises the universal closure $\opA{x_1, \ldots, x_k} \varphi$. But now this means that a uniform realiser $d$ must satisfy
        \[d \Vdash \varphi\mleft[a_1/x_1, \ldots, a_k/x_k\mright]\]
        for any $\vec{a} = \tuple{a_1, \ldots, a_k} \in U^k$.
\end{enumerate}

\begin{proposition}
  Soundness of intuitionistic logic as in \autoref{prop:real-con-int-logic} still holds for any uninterpreted realisability model.
\end{proposition}

\begin{proof}
  It suffices to check the axioms and inference rules involving quantifiers. Observe that:
  \begin{enumerate}[label=\uplabel{\arabic*},start=\intlogicproofquantcount]
    \item $\oplam{d} d$ simply realises $\opA{x} \varphi \rightarrow \varphi\mleft[t/x\mright]$;
    \item $\oplam{d} d$ also realises $\varphi\mleft[t/x\mright] \rightarrow \opE{x} \varphi$, since let $\vec{y}$ enumerate the free variables in this axiom, then for any corresponding tuple $\vec{a} \in U^k$, if $d$ realises $\varphi\mleft[t^\mathcal{A}\mleft(\vec{a}\mright)/x, \vec{a}/\vec{y}\mright]$, then it also realises $\opE{x} \varphi\mleft[\vec{a}/\vec{y}\mright]$ by definition.
  \end{enumerate}

  For the inference rules:
  \begin{enumerate}[label=\uplabel{\roman*},start=\intlogicrulesquantcount]
    \item Assume that $f \Vdash \varphi \rightarrow \psi$ where the free variables are enumerated as $x, \vec{y}$ and $x$ does not occur free in $\varphi$, then for any corresponding $\vec{b} \in U^k$, if $d \Vdash \varphi\mleft[\vec{b}/\vec{y}\mright]$, then $fd \Vdash \psi\mleft[a/x, \vec{b}/\vec{y}\mright]$ for any $a \in U$. Thus $\oplam{d} fd$ realises $\varphi \rightarrow \opA{x} \psi$;
    \item Assume that $f \Vdash \varphi \rightarrow \psi$ where the free variables are enumerated as $x, \vec{y}$ and $x$ does not occur free in $\psi$, then for any corresponding $\vec{b} \in U^k$, if $d \Vdash \opE{x} \varphi\mleft[\vec{b}/\vec{y}\mright]$, i.e.\ $d \Vdash \varphi\mleft[a/x, \vec{b}/\vec{y}\mright]$ for some $a \in U$, then $fd \Vdash \psi\mleft[a/x, \vec{b}/\vec{y}\mright]$ which is the same formula as $\psi\mleft[\vec{b}/\vec{y}\mright]$. Thus $\oplam{d} fd$ realises $\opE{x} \varphi \rightarrow \psi$. \qedhere
  \end{enumerate}
\end{proof}

\section{Weak pcas}

Finally, it is necessary to bring into this overview a notion of \emph{weak pcas}\indg{weak pca}, which consist of the same data as a pca, but with a relaxed requirement for the behaviour of the operation $\bs$:
\begin{enumerate}[label=\uplabel{\roman*'},start=2]
  \item for any $a, b, c \in \mathcal{A}$, $\bs a {\downarrow}$, $\bs ab {\downarrow}$, and
        \[ac \mleft(bc\mright) {\downarrow} \rightarrow \bs abc {\downarrow} = ac \mleft(bc\mright).\]
\end{enumerate}
Thus, $\bs abc$ may still converge when $ac \mleft(bc\mright)$ diverges. This is a notion examined in \cite{faber14-recursion-realizability}*{section 5.5}, and can be seen as an intermediate variant between a pca and an `ordered partial combinatory algebra' as considered by Hofstra and van Oosten \cite{hofstra-van-oosten03-opca}. It is pointed out in \cite{faber-van-oosten16-type-2-pca}*{section 5.4} that every weak pca can be made into a pca as defined in \autoref{sec:pca} by manually introducing a correct $\bs$-combinator. However, as we will soon see later in \autoref{sec:oa-cm}, we sometimes want to make use of non-conventional constructions that have an explicit definition for the application function, where it can be hard to require the element $\bs$ to satisfy the strict rules of a pca. Thus, one may prefer to work with a weak variant directly, and we shall observe that many general observations in previous sections can be easily modified to accommodate a weak pca:
\begin{lemma}[$\lambda$-abstraction\indg{lambda-abstraction@$\lambda$-abstraction}]
  Let $\mathcal{A}$ be a weak pca. For any formal application term $t$ over variables $x, \vec{y}$, there is a formal application term \uppar{denoted as $\oplam{x} t$} over variables $\vec{y}$ only, such that
  \begin{enumerate}[label=\uplabel{\roman*}]
    \item for any $\vec{b} \in \mathcal{A}^k$, $\mleft(\oplam{x} t\mright)\mleft[\vec{b}/\vec{y}\mright] {\downarrow}$;
    \item for any $a \in \mathcal{A}$, $\vec{b} \in \mathcal{A}^k$,
          \[t\mleft[a/x, \vec{b}/\vec{y}\mright] {\downarrow} \rightarrow \mleft(\oplam{x} t\mright)\mleft[\vec{b}/\vec{y}\mright] a {\downarrow} = t\mleft[a/x, \vec{b}/\vec{y}\mright].\]
  \end{enumerate}
\end{lemma}

\begin{proof}
  Same as \autoref{lem:lambda-abstraction}.
\end{proof}

\begin{corollary}[Recursion theorem]
  Let $\mathcal{A}$ be a weak pca. There must be a distinguished element $\br \in \mathcal{A}$ such that for any $f, a \in \mathcal{A}$, $\br f {\downarrow}$ and
  \[fa \mleft(\br f\mright) {\downarrow} \rightarrow \br fa {\downarrow} = fa \mleft(\br f\mright).\]
\end{corollary}

\begin{proof}
  Same as \autoref{cor:recursion-theorem}.
\end{proof}

Over the natural numbers, we mostly work with total functions, so nothing needs to be modified in order to define a \emph{weak pca over natural numbers}. It is also easy to observe that realisability models can be built and \autoref{prop:real-con-int-logic} continues to hold for a weak pca over natural numbers.

\chapter{A semi-intuitionistic higher-order arithmetic}
\label{chap:arith}

In this chapter, we work in the base theory of classical second-order arithmetic as developed in Simpson \cite{simpson09-soa}. We build a realisability model for a theory of semi-intuitionistic third-order arithmetic proposed by Nik Weaver and determine the weakest fragment of second-order arithmetic required to formulate this interpretation, in order to perform an ordinal analysis of Weaver's system. From \autoref{sec:oa-cm-gwo} onwards, we expand this classical mechanism to additionally analyse extensions of the system concerning a global well-ordering of the second-order objects.

The majority of content in this chapter is extracted from my arXiv manuscript \cite{wang-sw25-cm} on Weaver's system $\CM$.

\section{Weaver's conceptual mathematics}
\label{sec:cm-axioms}

Nik Weaver, who was a mathematician in the field of functional analysis, has published actively on arXiv about his ideas of a weak, essentially predicative foundation for mainstream mathematics since the 2000s. In this chapter, we focus on the formal fragment of third-order arithmetic he introduced in \cite{weaver09-cm}, as an example of his philosophical vision known as \emph{mathematical conceptualism} \cite{weaver05-conceptualism}. In that exposition, he rejects the full Cantorian universe of sets as a foundational ontological commitment and instead advocates for a restricted language involving only countably and predicatively constructed mathematical objects; as he wrote in \cite{weaver05-conceptualism}*{section 1},

\begin{quote}
  Classical set theory presents us with a picture of an incredibly vast universe. [\textellipsis] Yet virtually all important objects in mainstream mathematics are either countable or separable. This is not because the uncountable/nonseparable case has not yet been sufficiently studied but rather because on mainstream questions it tends to be either pathological or undecidable.
\end{quote}

On the other hand, unlike the traditional Weylian predicativists who take the natural numbers as given and allow essentially only finitary constructions on top of them\footnote{This is most evident from the well-known identification of the `limit of predicativity' as the proof-theoretic ordinal $\Gamma_0$ by Feferman \cite{feferman64-predicative-analysis} and Sch\"utte \cite{schutte65-predicative-well-orderings}, versus concerns raised by Howard \cite{howard96-proof-theory} and Weaver \cites{weaver09-predicativity,weaver18-predicative-well-ordering}. See page \pageref{loc:lim-pred-discuss} for a brief discussion.}, Weaver makes a hard distinction between the countable and uncountable collections. He believes that one is capable of mentally conceiving the completion of any countable process, for example through the well-known mechanism of supertasks in philosophy (see \cite{weaver09-set}*{section 3.1}); whereas uncountable processes, such as a transfinite enumeration of all real numbers, are inconceivable and thus do not deserve to be `complete' objects in the mathematical ontology.

For formal conceptual mathematics, Weaver adopts a constructive notion of truth\footnote{Weaver has a flexible attitude regarding the underlying logic of mathematical reasoning and sees both classical formulations (such as the mini-fragment of set theory he worked with in \cite{weaver05-j2}) and constructive ones (such as the system of conceptual mathematics in \cite{weaver09-cm}, which we study in this chapter) as justified under conceptualism. More detail is given in \cite{weaver09-set}*{section 2.4}.} that aligns with his stance on ontology: he views a typical instance of classical reasoning
\[\opA{x \in X} \varphi \lor \opE{x \in X} \neg \varphi\]
as the statement that either the property $\varphi$ holds for all $x \in X$, or that there is a \emph{process} for finding one counterexample. Therefore, under a conceptualist philosophy, this is always possible when the domain $X$ is countable, but very likely not the case when $X$ is uncountable. This justifies the use of semi-intuitionistic logic in conceptual mathematics, i.e.\ intuitionistic logic together with the law of excluded middle $\varphi \lor \neg \varphi$ for formulae that only quantify over natural number variables, since it only requires enumerating over countably many objects to determine the truth value of such statements.

Now, we shall not be concerned with more philosophical justifications in this thesis, but grant that Weaver's philosophical ideals explained so far motivate the following axioms in \cite{weaver09-cm}:
\begin{enumerate}[label=\uplabel{\arabic*},series=cm-axioms]
  \item \label{itm:cm-num-axiom} the number axioms, essentially the axioms of $\HA$ excluding induction;
  \item \label{itm:cm-ind-axiom} for any third-order formula $\varphi$, the full induction scheme:
        \[\varphi\mleft(0\mright) \land \opA{n} \left(\varphi\mleft(n\mright) \rightarrow \varphi\mleft(n + 1\mright)\right) \rightarrow \opA{n} \varphi\mleft(n\mright);\]
  \item \label{itm:cm-rec-axiom} and the full recursion scheme (also known as \emph{dependent choice}):
        \[\opA{n} \opA{X} \opE{Y} \varphi\mleft(n, X, Y\mright) \rightarrow \opA{X} \opE{Z} \left(\left(Z\right)_0 = X \land \opA{n} \varphi\mleft(n, \left(Z\right)_n, \left(Z\right)_{n + 1}\mright)\right)\]
        (where $Z$ is a conventional encoding of a family of second-order objects $\left(Z\right)_n$ as a single object, to be explained in more detail on page \pageref{loc:arith-set-slice}), because both induction and recursion on first-order objects concern mere iterations of countable length;
  \item \label{itm:cm-dec-comp-axiom} for any third-order formula $\varphi$ (which does not contain the new variable $X$ or $\mathbf{X}$ used for comprehension), the \emph{decidable comprehension} axioms:
        \begin{align*}
          \opA{n} \left(\varphi\mleft(n\mright) \lor \neg \varphi\mleft(n\mright)\right) & \rightarrow \opE{X} \opA{n} \left(n \in X \leftrightarrow \varphi\mleft(n\mright)\right),                   \\
          \opA{X} \left(\varphi\mleft(X\mright) \lor \neg \varphi\mleft(X\mright)\right) & \rightarrow \opE{\mathbf{X}} \opA{X} \left(X \in \mathbf{X} \leftrightarrow \varphi\mleft(X\mright)\right),
        \end{align*}
        where $X$ denotes a second-order variable and boldface $\mathbf{X}$ denotes a third-order variable;
  \item \label{itm:cm-arith-dec-axiom} and lastly, the law of excluded middle $\varphi \lor \neg \varphi$ for every arithmetic formula $\varphi$, i.e.\ formulae that do not contain second- or third-order quantifiers.
        \xdef\cmarithdecaxiomlabel{\arabic{enumi}}
\end{enumerate}

As Weaver demonstrated in \cite{weaver09-cm}*{section 3}, this formal theory is very expressive and suffices for the development of much mainstream mathematics such as real analysis, topology and functional analysis, by expressing arbitrary sets (e.g.\ $\mathbb{R}$, arbitrary topological spaces and arbitrary metric spaces) as third-order objects and their elements as second-order objects. For example, he defined and proved the following (separable) version of the Hahn--Banach theorem:

\begin{theorem}[\cite{weaver09-cm}*{Theorem 3.72}, $\CM$]
  \label{thm:cm-hahn-banach}
  Let $\mathbf{E}$ be a separable Banach space, let $\mathbf{E}_0$ be a separable closed
  subspace, and let $\mathbf{f}_0 : \mathbf{E}_0 \to \mathbb{R}$ be a bounded linear functional on $\mathbf{E}_0$. Then $\mathbf{f}_0$
  extends to a bounded linear functional $\mathbf{f}$ on $\mathbf{E}$ with $\left\|\mathbf{f}\right\| = \left\|\mathbf{f}_0\right\|$.
\end{theorem}

The separability assumptions are needed because the restricted law of excluded middle and the decidable comprehension axioms in Weaver's system together yield a deep distinction between quantifiers over countable and uncountable collections. Weaver was able to prove in his theory strong, classical results about algebraic structures that admit `countable descriptions' such as a countable dense set, which fail to apply to other structures in general. This is not a disheartening observation, as Weaver argued in \cite{weaver09-set-theory-indispensable}, since the majority of mainstream functional analysis, for example, is precisely about separable spaces.

Thus, it is a very interesting question to ask what the precise proof-theoretic strength of Weaver's system is. In \cite{rathjen23-cst}*{section 2.7}, Rathjen observed that this system is easily interpretable in the weak intuitionistic set theory $\CZF + \LPO + \RDC$, which is shown in \cite{rathjen14-czf-lpo} to have the same proof-theoretic ordinal as $\CZF$ itself, namely the Bachmann--Howard ordinal. In this chapter, we will improve this bound and show that Weaver's theory is actually equiconsistent with the classical second-order theory of $\Sigma^1_1\DC$, which is much weaker and has proof-theoretic ordinal $\varphi_{\varepsilon_0}\mleft(0\mright)$.

\subsection*{The full system $\CM$}

Before our analysis, we will first make a few clarifications and minor modifications to Weaver's original presentation of conceptual mathematics above, which do not affect the core ideas of his system.

We phrase \emph{conceptual mathematics}, $\CM$\inda{conceptual mathematics, $\CM$}, as a theory in the usual three-sorted implementation of third-order logic, where first-order variables are denoted by lowercase letters $n, m, \ldots$, second-order variables are denoted by italic uppercase letters $X, Y, \ldots$, and third-order variables are denoted by bold uppercase letters $\mathbf{X}, \mathbf{Y}, \ldots$. The vocabulary is that of first-order arithmetic, i.e.\ with the constant symbol $0$ and function symbols for successor, $+$ and $\cdot$, as well as binary membership relations between first- and second-order and between second- and third-order objects, which will be distinguished and rendered\footnote{This distinction is not made in Weaver's original formulation in \cite{weaver09-cm}. However, we need this treatment so that, when we work with a realisability model and have the same kind of realisers (i.e.\ objects in the pca) assigned to all three sorts of variables, we can still identify the order of the variable from the presentation of the formula directly.} as $n \in_1 X$ and $X \in_2 \mathbf{X}$.

Later in \autoref{sec:cm-gwo-axioms}, we will consider an extension of this theory envisioned by Weaver in \cite{weaver09-cm}*{section 2.3}, where a global well-ordering is hypothesised on the collection of all second-order objects. So we will preemptively introduce an additional binary relation symbol $\prec$ between two second-order variables. The only axiom in $\CM$ concerning it will be to assert that this is a decidable relation (as well as that second-order objects are extensional with respect to all relations, see axiom \ref{itm:cm-gwo-ext-axiom} later):
\begin{enumerate}[label=\uplabel{\arabic*},resume=cm-axioms]
  \item \label{itm:cm-gwo-dec-axiom} $\opA{X} \opA{Y} \left(X \prec Y \lor \neg X \prec Y\right)$,
\end{enumerate}
which is clearly conservative over the system without this symbol, for example, by trivially interpreting $X \prec Y$ as the decidable arithmetic formula $\opA{n} \left(n \in_1 X \leftrightarrow n \in_1 Y\right)$.

From \cite{weaver09-cm}*{section 2.2}, one can infer that Weaver would like to treat first-order equality as a primitive symbol with the standard equality axioms, while second- and third-order equalities are treated as extensional abbreviations:
\begin{align*}
  X_1 = X_2                   & \defeqv \opA{n} \left(n \in_1 X_1 \leftrightarrow n \in_1 X_2\right),                   \\
  \mathbf{X}_1 = \mathbf{X}_2 & \defeqv \opA{X} \left(X \in_2 \mathbf{X}_1 \leftrightarrow X \in_2 \mathbf{X}_2\right).
\end{align*}

However, he does not include an axiom of extensionality for these inferred equalities in his formulation, which can bring quite some problems when we reason with third-order objects. For example, to understand expressions involving set-building terms like $\left\{n : \varphi\mleft(n\mright)\right\} \in_2 \mathbf{X}$ as unambiguous abbreviations for formal formulae in the language --- where the formula above is usually read as
\[\opE{X} \left(\opA{n} \left(n \in_1 X \leftrightarrow \varphi\mleft(n\mright)\right) \land X \in_2 \mathbf{X}\right)\]
--- it is important that there are no multiple copies of the `same' set of natural numbers with different behaviours in the theory. Therefore, we would add to $\CM$ the following axioms:
\begin{enumerate}[label=\uplabel{\arabic*},resume=cm-axioms]
  \item \label{itm:cm-mem-ext-axiom} $X = Y \rightarrow \left(X \in_2 \mathbf{X} \leftrightarrow Y \in_2 \mathbf{X}\right)$,
  \item \label{itm:cm-gwo-ext-axiom} $X = Y \land X' = Y' \rightarrow \left(X \prec Y \leftrightarrow X' \prec Y'\right)$.
\end{enumerate}
With these necessary ingredients, we can show that:
\begin{lemma}[$\CM$]
  Given axioms \ref{itm:cm-mem-ext-axiom} and \ref{itm:cm-gwo-ext-axiom}, the Leibniz law holds, i.e.
  \[t = s \rightarrow \left(\varphi\mleft(t\mright) \leftrightarrow \varphi\mleft(s\mright)\right),\]
  for any \uppar{third-order} formula $\varphi$ and any terms $t, s$ of the same, correct order.
\end{lemma}

\begin{proof}
  This follows from simple induction on the complexity of $\varphi$. All atomic cases hold trivially by definition of the equality symbol, except when $\varphi\mleft(t\mright)$ is $t \in_2 \mathbf{X}$, $t \prec X$ or $X \prec t$, where the extensionality assertions above are used.
\end{proof}

Additionally, the arithmetic decidability scheme \ref{itm:cm-arith-dec-axiom} can be complicated to work with, and we observe that it can be broken down into two separate components:
\begin{enumerate}[label=\uplabel{\noexpand\enumdynlabel}]
  \def\enumdynlabel{\cmarithdecaxiomlabel$'$}
  \item \label{itm:cm-atm-dec-axiom} atomic decidability axioms:
        \[n = m \lor \neg n = m, \qquad n \in_1 X \lor \neg n \in_1 X, \qquad X \in_2 \mathbf{X} \lor \neg X \in_2 \mathbf{X};\]
        \def\enumdynlabel{\cmarithdecaxiomlabel$''$}
  \item \label{itm:cm-lpo-axiom} for any third-order formulae $\varphi\mleft(n\mright)$ and $\psi\mleft(n\mright)$, the numerical omniscience scheme, also known as the \emph{limited principle of omniscience}:
        \[\opA{n} \left(\varphi\mleft(n\mright) \lor \psi\mleft(n\mright)\right) \rightarrow \opA{n} \varphi\mleft(n\mright) \lor \opE{n} \psi\mleft(n\mright).\]
\end{enumerate}

\begin{lemma}
  \label{lem:cm-lpo}
  Over the rest of the axioms of $\mathrm{CM}$, the axiom scheme \ref{itm:cm-arith-dec-axiom} is equivalent to $\text{\ref{itm:cm-atm-dec-axiom}} + \text{\ref{itm:cm-lpo-axiom}}$.
\end{lemma}

\begin{proof}
  The forward direction is essentially \cite{weaver09-cm}*{Theorem 2.2}. We only need to check numerical omniscience \ref{itm:cm-lpo-axiom}: for any $n$, assume that $\varphi\mleft(n\mright) \lor \psi\mleft(n\mright)$, then clearly there exists some $X$ such that
  \[\left(\varphi\mleft(n\mright) \land X = \varnothing\right) \lor \left(\psi\mleft(n\mright) \land X = \left\{n\right\}\right).\]
  Now, by the recursion axiom, there exists some $Z$ such that $\left(Z\right)_n$ satisfies the condition for $X$ above for each specific $n$ (see the paragraphs after this proof for the formal definition of $\left(Z\right)_n$ using a recursive pairing function), and we can take $Y = \left\{n : n \in_1 \left(Z\right)_n\right\}$ by decidable comprehension, so that
  \[\opA{n} \left(\left(\varphi\mleft(n\mright) \land \neg n \in_1 Y\right) \lor \left(\psi\mleft(n\mright) \land n \in_1 Y\right)\right).\]
  By arithmetic decidability \ref{itm:cm-arith-dec-axiom}, we have $\neg\!\opE{n} n \in_1 Y \lor \opE{n} n \in_1 Y$ which immediately implies $\opA{n} \neg n \in_1 Y \lor \opE{n} n \in_1 Y$. Thus also $\opA{n} \varphi\mleft(n\mright) \lor \opE{n} \psi\mleft(n\mright)$ as desired.

  For the backward direction, observe that \ref{itm:cm-lpo-axiom} implies
  \begin{align*}
    \opA{n} \left(\varphi\mleft(n\mright) \lor \neg \varphi\mleft(n\mright)\right) & \rightarrow \opA{n} \varphi\mleft(n\mright) \lor \neg\!\opA{n} \varphi\mleft(n\mright), \\
    \opA{n} \left(\neg \psi\mleft(n\mright) \lor \psi\mleft(n\mright)\right)       & \rightarrow \neg\!\opE{n} \psi\mleft(n\mright) \lor \opE{n} \psi\mleft(n\mright),
  \end{align*}
  i.e.\ the collection of formulae satisfying excluded middle is closed under first-order quantification. Since this collection is trivially also closed under logical connectives, arithmetic decidability \ref{itm:cm-arith-dec-axiom} follows from atomic decidability \ref{itm:cm-atm-dec-axiom} through a straightforward induction on formula complexity.
\end{proof}

In this thesis, we will take $\mathrm{CM}$ to formally denote the collection of axioms $\text{\ref{itm:cm-num-axiom}} + \text{\ref{itm:cm-ind-axiom}} + \text{\ref{itm:cm-rec-axiom}} + \text{\ref{itm:cm-dec-comp-axiom}} + \text{\ref{itm:cm-atm-dec-axiom}} + \text{\ref{itm:cm-lpo-axiom}} + \text{\ref{itm:cm-gwo-dec-axiom}} + \text{\ref{itm:cm-mem-ext-axiom}} + \text{\ref{itm:cm-gwo-ext-axiom}}$.

Lastly, we shall use the notation $\tuple{{-}, {-}} : \mathbb{N}^2 \rightarrow \mathbb{N}$ in $\CM$ as an abbreviation to denote the recursive bijection
\[\tuple{n, m} = \left(n + m\right) \cdot \left(n + m + 1\right) / 2 + n,\]
with recursive inverses\footnote{In this section, for any $a \in \mathbb{N}$, we use $a_0, a_1$ to denote the two components when $a$ is treated as a pair of natural numbers. This is admittedly an abuse of notation, as we have previously defined $\bp_0, \bp_1$ to denote the unpairing functions in a pca. We shall clarify the ambiguity here by promising that we will never use $\bp_0, \bp_1$ to denote `the components of the natural number $\bp$' in this thesis.} $\left(-\right)_0, \left(-\right)_1 : \mathbb{N} \rightarrow \mathbb{N}$. We then follow \cite{weaver09-cm}*{Definition 3.2} and specify the encoding for families of second-order objects used in the recursion scheme \ref{itm:cm-rec-axiom} as
\label{loc:arith-set-slice}
\[\left(Z\right)_n = \left\{m : \tuple{n, m} \in Z\right\}.\]
The fact that every natural number is a code for some pair will be convenient for constructions later. In \cite{weaver09-cm}, Weaver used a different, non-surjective map $\tuple{n, m} = 2^n \cdot 3^m$. This difference does affect the precise meaning of the recursion scheme \ref{itm:cm-rec-axiom}; however, since $\CM$ contains the classical subtheory $\PA$, it can be easily shown that the different interpretations are ultimately equivalent.

\section{A realisability model for \texorpdfstring{$\CM$}{CM}}
\label{sec:oa-cm}

We begin the construction of our realisability model for $\CM$ in the classical base theory $\ACA_0$, which is defined and studied in detail in \cite{simpson09-soa}*{Chapter III}. Following \cite{simpson09-soa}*{Definition VIII.3.18}, let $\pi\mleft(e, m, a, b, X\mright)$ denote a universal $\Pi^0_1$-formula with all free variables indicated, and correspondingly a universal $\Sigma^1_1$-formula is given by $\opE{X} \pi\mleft(e, m, a, b, X\mright)$, i.e.\ for any $\Sigma^1_1$-formula $\varphi\mleft(m, a, b\mright)$ with all free variables indicated, $\ACA_0$ proves (cf.\ \cite{simpson09-soa}*{Lemma V.1.4}) that there exists $e \in \mathbb{N}$ such that
\[\opA{m} \opA{a} \opA{b} \left(\varphi\mleft(m, a, b\mright) \leftrightarrow \opE{X} \pi\mleft(e, m, a, b, X\mright)\right).\]

We now let\footnote{We pick the name $\Hyp$ from the fact that the total functions $\mathbb{N} \rightarrow \mathbb{N}$ representable in the weak pca are precisely the hyperarithmetic functions. However, this cannot be verified in a weak base theory like $\ACA_0$; we will later see as \autoref{thm:conv-s2-to-hyp} that this is provable in $\ATR_0$ (cf.\ \cite{simpson09-soa}*{Theorem VIII.3.19}).} $\Hyp$\inda{hyperarithmetic algebra, $\Hyp$} denote the domain $\mathbb{N}$ with a partial application function $\mathrm{ap} : \mathbb{N}^2 \rightharpoonup \mathbb{N}$ where we say $ea {\downarrow} = b$ if and only if $b \in \mathbb{N}$ is the \emph{unique} natural number such that
\[\opE{X} \pi\mleft(e_0, e_1, a, b, X\mright)\]
holds, where $e = \tuple{e_0, e_1}$. We will check below that it is a weak pca. To begin, we observe that:
\begin{lemma}
  \label{lem:sigma11-smn}
  Given any $\Sigma^1_1$-formula $\varphi\mleft(m, a, b\mright)$ with all free variables indicated, then there exists $e \in \mathbb{N}$ such that for any $m, a \in \mathbb{N}$, $em {\downarrow}$ always holds, and $ema {\downarrow} = b$ if and only if $b \in \mathbb{N}$ is the unique natural number such that $\varphi\mleft(m, a, b\mright)$ holds.
\end{lemma}

\begin{proof}
  Simply pick $f, n \in \mathbb{N}$ by universality such that
  \begin{gather*}
    \opA{m} \opA{a} \opA{b} \left(b = \tuple{m, a} \leftrightarrow \opE{X} \pi\mleft(f, m, a, b, X\mright)\right), \\
    \opA{m} \opA{a} \opA{b} \left(\varphi\mleft(m, a, b\mright) \leftrightarrow \opE{X} \pi\mleft(n, m, a, b, X\mright)\right).
  \end{gather*}
  Let $e = \tuple{f, n}$, then $em {\downarrow} = \tuple{n, m}$ for any $m \in \mathbb{N}$ and thus we have $ema {\downarrow} = b$ precisely when $b \in \mathbb{N}$ is the unique natural number such that $\varphi\mleft(m, a, b\mright)$ holds.
\end{proof}

Using this we can verify that:
\begin{corollary}
  \label{cor:sigma11-pca}
  $\Hyp$ is a weak pca over natural numbers.
\end{corollary}

\begin{proof}
  We use \autoref{lem:sigma11-smn} to identify the distinguished elements in the algebra. For operations that involve one parameter only, we simply ignore the parameter $m$ in the proposition:
  \begin{enumerate}[label=\uplabel{\roman*}]
    \item let $s \in \mathbb{N}$ be the index such that for any $m, a \in \mathbb{N}$, $sma {\downarrow} = b$ if and only if $b = a + 1$, and set $\bsN = s0$;
    \item let $p \in \mathbb{N}$ be the index such that for any $m, a \in \mathbb{N}$, $pma {\downarrow} = b$ if and only if $a = b + 1$, and set $\bpN = p0$;
  \end{enumerate}

  For $\bk$ which involves two parameters, we use precisely the parameters in \autoref{lem:sigma11-smn}:
  \begin{enumerate}[label=\uplabel{\roman*},resume]
    \item let $\bk \in \mathbb{N}$ simply be the index such that for any $a, b \in \mathbb{N}$, $\bk ab {\downarrow} = c$ if and only if $c = a$.
  \end{enumerate}

  For $\bs$ which involves three parameters, we invoke the proposition twice:
  \begin{enumerate}[label=\uplabel{\roman*},resume]
    \item let $s' \in \mathbb{N}$ be the index such that for any $a, b \in \mathbb{N}$, $s' ab {\downarrow} = c$ if and only if $c \in \mathbb{N}$ is the unique natural number satisfying\footnote{Note that the existence of a unique $c$ here does not also guarantee the uniqueness of the intermediate values $d$ and $e$. Thus the existence of a unique $c$ does not imply $a_0 b \mleft(a_1 b\mright) {\downarrow}$ and $\Hyp$ as constructed here is only a weak pca.}
          \[\opE{X, Y, Z} \opE{d, e} \left(\pi\mleft(a_{00}, a_{01}, b, d, X\mright) \land \pi\mleft(a_{10}, a_{11}, b, e, Y\mright) \land \pi\mleft(d_0, d_1, e, c, Z\mright)\right).\]
          Thus, if $a_0 b {\downarrow} = d$, $a_1 b {\downarrow} = e$ and $de {\downarrow} = c$, then also $s' ab {\downarrow} = c$. Now let $\bs \in \mathbb{N}$ be such that $\bs ab {\downarrow} = s' \tuple{a, b}$, i.e.\ $\bs ab {\downarrow} = c$ if and only if
          \[\opE{X} \pi\mleft(s'_0, s'_1, \tuple{a, b}, c, X\mright).\]
  \end{enumerate}

  Finally, similarly for $\bdN$ which involves four parameters:
  \begin{enumerate}[label=\uplabel{\roman*},resume]
    \item let $d' \in \mathbb{N}$ be the index such that for any $a, b \in \mathbb{N}$,
          \[d' ab {\downarrow} = \left\{\begin{aligned}
               & a_1 &  & \text{if $a_{00} = a_{01}$},    \\
               & b   &  & \text{if $a_{00} \neq a_{01}$}.
            \end{aligned}\right.\]
          Then let $d'' \in \mathbb{N}$ be such that $d'' ab {\downarrow} = d' \tuple{a, b}$ and finally $\bdN \in \mathbb{N}$ be such that $\bdN ab {\downarrow} = d'' \tuple{a, b}$. \qedhere
  \end{enumerate}
\end{proof}

We will stipulate that the standard pairing function denoted by $\bp$ in $\Hyp$ is
\[\bp ab = \tuple{a, b},\]
which can be computed in exactly the same way as in the proof above. In pca contexts we will just write $\bp n m$ for $\tuple{n, m}$, $\bp_0 n$ for $n_0$ and $\bp_1 n$ for $n_1$.

Now, consider a third-order realisability model over the weak pca $\Hyp$, where variables of all three sorts are interpreted using domains within the weak pca. A first-order variable shall range over $\mathbb{N}$, a second-order variable shall range over the definable class
\[S_2 = \left\{d \in \Hyp : \opA{n} dn {\downarrow} \in \left\{0, 1\right\}\right\},\]
and third-order variables range over
\[S_3 = \left\{d \in \Hyp : \opA{f \in S_2} df {\downarrow} \in \left\{0, 1\right\} \land \opA{e, f \in S_2} \left(e \equiv f \rightarrow de = df\right)\right\},\]
where the binary relation $\equiv$ of extensional equivalence on $S_2$ is defined as
\[e \equiv f \defeqv \opA{n} en = fn.\]
We fix an arbitrary interpretation $q \in \Hyp$ for the relation $\prec$ in the language of $\CM$ satisfying $qe e' {\downarrow} \in \left\{0, 1\right\}$ for any $e, e' \in S_2$, and
\[e \equiv f \land e' \equiv f' \rightarrow qe e' = qf f'\]
for any $e, e', f, f' \in S_2$. The precise choice of $q$ does not matter for interpreting the base theory $\CM$, so we can just take, for example, $q = \lambda nm. 0$.

The realisability conditions are then given as in \autoref{sec:real-con-int-logic}, with the atomic clauses given by
\begin{align*}
  d \Vdash t = s     & \defeqv t = s \text{ for any arithmetic terms $t$ and $s$},                          \\
  d \Vdash t \in_1 e & \defeqv \opE{n} \left(t = n \land en = 1\right) \text{ for any arithmetic term $t$}, \\
  d \Vdash e \in_2 f & \defeqv fe = 1,                                                                      \\
  d \Vdash e \prec f & \defeqv qef = 1.
\end{align*}

We now move into the extension of $\ACA_0$ known as $\Sigma^1_1\AC$ (with unrestricted induction), which adds the axiom of second-order choice
\[\opA{n} \opE{Y} \varphi\mleft(n, Y\mright) \rightarrow \opE{Z} \opA{n} \varphi\mleft(n, \left(Z\right)_n\mright)\]
for any $\Sigma^1_1$-formula $\varphi$, and induction on the natural numbers for any second-order formulae. Here we show that the axioms of $\CM$ are all realised.

Firstly, the equality axioms for the primitive, first-order equality symbol are realised trivially here. The number axioms \ref{itm:cm-num-axiom} are just true, quantifier-free formulae in the base theory (with free variables), which can be easily realised by realisers that are total on the appropriate inputs.

For schemes \ref{itm:cm-ind-axiom} and \ref{itm:cm-rec-axiom}, we use the unrestricted induction axiom in $\Sigma^1_1\AC$ with recursion in the weak pca:
\begin{lemma}
  The induction scheme \ref{itm:cm-ind-axiom} is realised:
  \[\varphi\mleft(0\mright) \land \opA{n} \left(\varphi\mleft(n\mright) \rightarrow \varphi\mleft(n + 1\mright)\right) \rightarrow \opA{n} \varphi\mleft(n\mright).\]
\end{lemma}

\begin{proof}
  Let $n, \vec{x}$ enumerate the free variables in $\varphi$, where $\vec{x}$ can include variables of any order. If $d \Vdash \varphi\mleft[0/n, \vec{a}/\vec{x}\mright]$ and $f \Vdash \opA{n} \left(\varphi\mleft(n\mright) \rightarrow \varphi\mleft(n + 1\mright)\right)\mleft[\vec{a}/\vec{x}\mright]$, then
  \[\br \mleft(\oplam{yt} \bdN^* y0d \mleft(f \mleft(\bpN y\mright) \mleft(t \mleft(\bpN y\mright)\mright)\mright)\mright) m \Vdash \varphi\mleft[m/n, \vec{a}/\vec{x}\mright]\]
  by induction\footnote{\label{ftn:cm-real-complex}Even though we never use an arbitrary amount of second-order quantifiers, the use of unrestricted induction is necessary here: suppose that $\varphi$ contains an alternating sequence of quantifiers, then the statement $d \Vdash \varphi\mleft(p\mright)$ expands to a form like
  \[\opA{n_0} \opE{d_1} \left(d n_0 {\downarrow} = d_1 \land \opA{n_1} \opE{d_2} \left(\mleft(\bp_0 d_1\mright) n_1 {\downarrow} = d_2 \land \ldots \opA{n_i} \opE{d_i} \ldots \right.\right.\]
  Here formulae of the form $dn {\downarrow} = e$ are $\Delta^1_2$ but lie within a sequence of first-order quantifiers. Since $\Sigma^1_1\AC$ does not show that the class of $\Delta^1_2$-formulae is closed under adding first-order quantifiers, it cannot prove that this statement has bounded complexity.} on $m \in \mathbb{N}$. Thus $\oplam{\vec{x} df} \br \mleft(\oplam{yt} \bdN^* y0d \mleft(f \mleft(\bpN y\mright) \mleft(t \mleft(\bpN y\mright)\mright)\mright)\mright)$ realises the induction axiom as desired.
\end{proof}

\begin{lemma}
  \label{lem:cm-rec-real}
  The recursion scheme \ref{itm:cm-rec-axiom} is realised:
  \[\opA{n} \opA{X} \opE{Y} \varphi\mleft(n, X, Y\mright) \rightarrow \opA{X} \opE{Z} \left(\left(Z\right)_0 = X \land \opA{n} \varphi\mleft(n, \left(Z\right)_n, \left(Z\right)_{n + 1}\mright)\right)\]
\end{lemma}

\begin{proof}
  Let $n, X, Y, \vec{x}$ enumerate the free variables in $\varphi$, where $\vec{x}$ can include variables of any order. Given $f \Vdash \opA{n} \opA{X} \opE{Y} \varphi\mleft(n, X, Y\mright)\mleft[\vec{a}/\vec{x}\mright]$ and $s \in S_2$, we compute
  \[g = \br \mleft(\oplam{yt} \bdN^* y0s \mleft(\bp_0 \mleft(f \mleft(\bpN y\mright) \mleft(t \mleft(\bpN y\mright)\mright)\mright)\mright)\mright),\]
  then $g0 = s$ and $g \mleft(\bsN n\mright)$ is the `next' set $Y$ after some $X$ represented by $gn$, as given by $f$. Thus, the set $Z$ can be computed as
  \[z = \oplam{m} g \mleft(\bp_0 m\mright) \mleft(\bp_1 m\mright) \in S_2.\]
  Now $\left(Z\right)_0 = X$ can be interpreted as
  \[\opA{n} \opE{y} \left(y = \tuple{0, n} \land \left(y \in_1 Z \leftrightarrow n \in_1 X\right)\right)\]
  and realised simply by $\oplam{n} \bp \mleft(\bp 0n\mright) \mleft(\bp 0 \mleft(\bp \mleft(\oplam{d} 0\mright) \mleft(\oplam{d} 0\mright)\mright)\mright)$, whereas $\bp_1 \mleft(fn \mleft(gn\mright)\mright)$ realises\footnote{To be completely formal, we will need to understand $\varphi\mleft(n, \left(Z\right)_n, \left(Z\right)_{n + 1}\mright)$ as
    \[\opE{X} \opE{Y} \left(\left(Z\right)_n = X \land \left(Z\right)_{n + 1} = Y \land \varphi\mleft(n, X, Y\mright)\right),\]
    so the full realiser would be $\bp \mleft(gn\mright) \mleft(\bp \mleft(g \mleft(\bsN n\mright)\mright) \mleft(\bp d \mleft(\bp_1 \mleft(fn \mleft(gn\mright)\mright)\mright)\mright)\mright)$ where $d$ is the uninteresting realiser for $\left(Z\right)_n = X \land \left(Z\right)_{n + 1} = Y$.} $\varphi\mleft(n, \left(Z\right)_n, \left(Z\right)_{n + 1}\mright)$ by induction on $n \in \mathbb{N}$.
\end{proof}

The decidable comprehension axioms \ref{itm:cm-dec-comp-axiom} are easy cases here:
\begin{lemma}
  The following comprehension axioms are realised:
  \begin{align*}
    \opA{n} \left(\varphi\mleft(n\mright) \lor \neg \varphi\mleft(n\mright)\right) & \rightarrow \opE{X} \opA{n} \left(n \in X \leftrightarrow \varphi\mleft(n\mright)\right),                   \\
    \opA{X} \left(\varphi\mleft(X\mright) \lor \neg \varphi\mleft(X\mright)\right) & \rightarrow \opE{\mathbf{X}} \opA{X} \left(X \in \mathbf{X} \leftrightarrow \varphi\mleft(X\mright)\right),
  \end{align*}
\end{lemma}

\begin{proof}
  Let $n, \vec{x}$ enumerate the free variables in $\varphi\mleft(n\mright)$. Given any realiser $f \Vdash {\opA{n} \left(\varphi\mleft(n\mright) \lor \neg \varphi\mleft(n\mright)\right)}\mleft[\vec{a}/\vec{x}\mright]$, the set $X$ we need is simply
  \[s = \oplam{n} 1 - \bp_0 \mleft(fn\mright) \in S_2,\]
  and we have $\oplam{d} \bp_1 \mleft(fm\mright) \Vdash \mleft(n \in X \rightarrow \varphi\mleft(n\mright)\mright)\mleft[m/n, s/X, \vec{a}/\vec{x}\mright]$. Thus we know that $\oplam{\vec{x} f} \bp \mleft(\oplam{n} 1 - \bp_0 \mleft(fn\mright)\mright) \mleft(\oplam{n} \bp \mleft(\oplam{d} \bp_1 \mleft(fn\mright)\mright) \mleft(\oplam{d} 0\mright)\mright)$ realises the (decidable) second-order comprehension axiom.

  The third-order case is entirely analogous.
\end{proof}

For the atomic decidability axioms \ref{itm:cm-atm-dec-axiom}, simply observe that:
\begin{itemize}
  \item $t = s \lor \neg t = s$ is realised by
        \[\oplam{\vec{x}} \left\{\begin{aligned}
             & \bp 00 &  & \text{if $t = s$},    \\
             & \bp 10 &  & \text{if $t \neq s$},
          \end{aligned}\right.\]
        where $\vec{x}$ enumerate the free variables in the terms $t$ and $s$ and are used in the computation to distinguish the cases;
  \item $t \in_1 e \lor \neg t \in_1 e$ is realised by $\oplam{\vec{x} e} \bp \mleft(1 - et\mright) 0$ where $\vec{x}$ enumerate the free variables and are used to compute $t$;
  \item $e \in_2 f \lor \neg e \in_2 f$ is realised by $\oplam{ef} \bp \mleft(1 - fe\mright) 0$.
\end{itemize}
And likewise the decidability of the binary relation \ref{itm:cm-gwo-dec-axiom}, $e \prec f \lor \neg e \prec f$, is realised by $\oplam{ef} \bp \mleft(1 - qef\mright) 0$.

\begin{lemma}
  \label{lem:cm-num-omni-real}
  The numerical omniscience scheme \ref{itm:cm-lpo-axiom} is realised:
  \[\opA{n} \left(\varphi\mleft(n\mright) \lor \psi\mleft(n\mright)\right) \rightarrow \opA{n} \varphi\mleft(n\mright) \lor \opE{n} \psi\mleft(n\mright).\]
\end{lemma}

\begin{proof}
  Let $\vec{x}$ enumerate the free variables in $\opA{n} \left(\varphi\mleft(n\mright) \lor \psi\mleft(n\mright)\right)$ and suppose that $f \Vdash {\opA{n} \left(\varphi\mleft(n\mright) \lor \psi\mleft(n\mright)\right)}\mleft[\vec{a}/\vec{x}\mright]$. By $\Sigma^1_1$-choice, it is then the case that
  \[\opE{Z} \opA{n} \opE{!b} \pi\mleft(f_0, f_1, n, b, \left(Z\right)_n\mright)\]
  where the number $b$ corresponding to the same $n \in \mathbb{N}$ is necessarily unique. By \autoref{lem:sigma11-smn}, there is $e \in \Hyp$ such that for any $f \in \mathbb{N}$, $ef {\downarrow} = d$ if and only if
  \begin{enumerate}[label=\uplabel{\alph*}]
    \item either $\opE{Z} \opA{n} \opE{b} \left(\pi\mleft(f_0, f_1, n, b, \left(Z\right)_n\mright) \land b_0 = 0\right)$ and $d = \tuple{0, 0}$,
    \item or $\opE{Z} \opE{n, b} \left(\pi\mleft(f_0, f_1, n, b, \left(Z\right)_n\mright) \land b_0 = 1 \land d = \tuple{1, \tuple{n, b_1}}\right)$, where $n$ is the least such number, i.e.\ also $\opA{k < n} \opA{b} \left(\pi\mleft(f_0, f_1, k, b, \left(Z\right)_k\mright) \rightarrow b_0 = 0\right)$.
  \end{enumerate}
  It is then easy to verify that $f \Vdash {\opA{n} \left(\varphi\mleft(n\mright) \lor \psi\mleft(n\mright)\right)}\mleft[\vec{a}/\vec{x}\mright]$ implies $ef {\downarrow}$, where either $\opA{n} \bp_0 \mleft(fn\mright) = 0$ and $ef = \tuple{0, 0}$, or $ef = \tuple{1, \tuple{m, d}}$ for some $m, d \in \mathbb{N}$ such that $fm = \tuple{1, d}$, i.e.\ $d \Vdash \psi\mleft[m/n, \vec{a}/\vec{x}\mright]$.

  Hence, $\oplam{\vec{x} f} \bdN \mleft(\bp_0 \mleft(ef\mright)\mright) 0 \mleft(\bp 0 \mleft(\oplam{y} \bp_1 \mleft(fy\mright)\mright)\mright) \mleft(\bp 1 \mleft(\bp_1 \mleft(ef\mright)\mright)\mright)$ realises the axiom as desired.
\end{proof}

Finally, observe that in the set-up we have ensured for any $d \in S_3$,
\[\opA{e, f \in S_2} \left(e \equiv f \rightarrow de = df\right),\]
and likewise for the fixed interpretation $q$ of the relation $\prec$,
\[\opA{e, e', f, f' \in S_2} \left(e \equiv f \land e' \equiv f' \rightarrow qe e' = qf f'\right).\]
Thus the extensionality assertions \ref{itm:cm-mem-ext-axiom} and \ref{itm:cm-gwo-ext-axiom} will both be realised by trivial realisers of the form $\oplam{\vec{x} d} \bp \mleft(\oplam{e} 0\mright) \mleft(\oplam{e} 0\mright)$.

Combining the lemmata above and \autoref{prop:real-con-int-logic}, we can conclude that:
\begin{theorem}
  Let $\varphi$ be any theorem of $\CM$, then $\Sigma^1_1\AC$ proves there exists $d \in \Hyp$ such that $d \Vdash \varphi$.
\end{theorem}

\subsection*{An ordinal analysis of $\CM$}

Let $\Sigma^1_1\DC$ denote the extension of $\ACA_0$ by unrestricted induction and the axiom of $\Sigma^1_1$-dependent choice (cf.\ \cite{simpson09-soa}*{section VII.6}), $\Delta^1_1\CA$ denote the extension of $\ACA_0$ by unrestricted induction and the axiom of $\Delta^1_1$-comprehension (cf.\ \cite{simpson09-soa}*{section VII.5}), and $\Sigma^1_1\AC^i$ denote the intuitionistic fragment of $\Sigma^1_1\AC$. It is well-known that:
\begin{theorem}
  The following theories have the same proof-theoretic strength:
  \[\Sigma^1_1\AC^i \equiv_{\Con} \Sigma^1_1\AC \equiv_{\Con} \Sigma^1_1\DC.\]
\end{theorem}

\begin{proof}
  The proof-theoretic equivalence of $\Sigma^1_1\AC^i$ and $\Sigma^1_1\AC$ is due to Peter Aczel \cite{aczel77-martin-lof}. Harvey Friedman showed in \cite{friedman67-thesis} that one can interpret both $\Sigma^1_1\AC$ and $\Sigma^1_1\DC$ in $\Delta^1_1\CA$, thus the three theories are equivalent in strength since $\Delta^1_1$-comprehension is implied by $\Sigma^1_1$-choice.
\end{proof}

Since $\Sigma^1_1\AC^i$ is an obvious subtheory of $\CM$ and our analysis above shows that $\CM$ can be interpreted in $\Sigma^1_1\AC$, it is an immediate corollary that:
\begin{corollary}
  \label{cor:cm-oa}
  $\CM \equiv_{\Con} \Sigma^1_1\DC$. In particular, since the exact proof-theoretic strength of $\Sigma^1_1\DC$ is analysed in \cite{friedman70-sigma11-12-ac}, we have $\left|\CM\right| = \left|\Sigma^1_1\DC\right| = \varphi_{\varepsilon_0}\mleft(0\mright)$.
\end{corollary}

\section{Axiomatising a global well-ordering}
\label{sec:cm-gwo-axioms}

In \cite{weaver09-cm}*{section 2.3}, Weaver suggested that the axioms of $\CM$ do not fully represent the conception that `second-order objects [appear] in well-ordered stages' and briefly considered extending the system to include the axiomatisation of a global well-ordering on the collection of second-order objects. Formally, this extension would include the following axioms:
\begin{enumerate}[label=\uplabel{W\arabic*}]
  \item \label{itm:cm-gwo-ord-axiom} the binary relation $\prec$ between second-order objects is a linear ordering;
  \item \label{itm:cm-gwo-ind-axiom} the transfinite induction scheme
        \[\opA{X} \left(\opA{Y} \left(Y \prec X \rightarrow \varphi\mleft(Y\mright)\right) \rightarrow \varphi\mleft(X\mright)\right) \rightarrow \opA{X} \varphi\mleft(X\mright),\]
        for a certain class of formulae;
  \item \label{itm:cm-gwo-init-seg-axiom} the axiom of countable initial segments
        \[\opA{X} \opE{Z} \opA{Y} \left(Y \prec X \rightarrow \opE{n} Y = \left(Z\right)_n\right).\]
\end{enumerate}
For the axiom scheme \ref{itm:cm-gwo-ind-axiom} of transfinite induction, Weaver made explicit that an extension justified in his conceptualist philosophy cannot assert this for every formula $\varphi$ --- especially those that contain the relation $\prec$ itself --- since the collection of all second-order objects is uncountable and from a predicativist perspective, Weaver does not want to assert that they can be enumerated in a way that supports induction using properties of the entire well-ordering (cf.\ \cite{weaver09-predicativity}*{section 2.5}).

However, in his brief discussion of the extension, Weaver never managed to suggest a plausible class of formulae to which the transfinite induction scheme can be restricted. We will give a detailed analysis of his proposal in \cite{weaver09-cm}*{section 2.3} and other possible alternatives of a `predicative' global well-ordering later in \autoref{sec:cm-gwo-pred}. For now, we shall let $\GWO$\inda{conceptual mathematics, $\CM$!with GWO@with global well-ordering axioms, $\GWO$} denote the entire collection of axioms $\text{\ref{itm:cm-gwo-ord-axiom}} + \text{\ref{itm:cm-gwo-ind-axiom}} + \text{\ref{itm:cm-gwo-init-seg-axiom}}$ where the transfinite induction scheme \ref{itm:cm-gwo-ind-axiom} applies to any third-order formula $\varphi$ in the language. We will examine the expressive power of the full extension $\CM + \GWO$ in this section and provide an ordinal analysis for it in \autoref{sec:oa-cm-gwo}, which acts as an upper bound for the strength of all its (predicative) fragments.

\subsection*{Zorn's lemma and applications}

In \cite{weaver09-cm}*{section 3}, Weaver gave a very detailed exposition of how ordinary mathematics, including theorems of real analysis, topology and functional analysis, can be treated inside $\CM$, in order to justify its viability as a foundational theory of mathematics. However, he never explored how the global well-ordering axioms add to this development --- whether they fulfil the role of a `genuine choice axiom' for mathematics as he suggested --- apart from a comment after his proof of the separable Hahn--Banach theorem in $\CM$ (previously mentioned in this thesis as \autoref{thm:cm-hahn-banach}) that:
\begin{quote}
  The same result holds, with the same proof, for extensions from separable subspaces of nonseparable spaces, but this requires the well-ordering $\prec$.
\end{quote}

In this section, we will prove inside $\CM + \GWO$ a general result analogous to Zorn's lemma and demonstrate that the global well-ordering axioms indeed serve their purpose as a choice principle and constitute a useful extension if one wants to use $\mathrm{CM}$ as a foundational theory. This shall serve as a typical sample of how much transfinite induction is needed by mainstream mathematics and facilitate our later discussion of the (usable) predicative fragments of $\GWO$ in \autoref{sec:cm-gwo-pred}.

To begin, we will first observe that omniscience holds for a bounded second-order quantifier:

\begin{lemma}[$\CM + \GWO$]
  \label{lem:cm-gwo-prec-bounded-so-omni}
  For any formula $\varphi\mleft(X\mright)$ and $\psi\mleft(X\mright)$, we have
  \[\opA{X} \left(\opA{Y \prec X} \left(\varphi\mleft(Y\mright) \lor \psi\mleft(Y\mright)\right) \rightarrow \left(\opA{Y \prec X} \varphi\mleft(Y\mright) \lor \opE{Y \prec X} \psi\mleft(Y\mright)\right)\right).\]
  Obviously the same also holds when every $\prec$ is replaced by $\preceq$ in the proposition.
\end{lemma}

\begin{proof}
  The axiom of countable initial segments asserts that for any second-order set $X$, there exists $Z$ satisfying
  \[\opA{Y \prec X} \opE{n} Y = \left(Z\right)_n.\]
  Assume that $\opA{Y \prec X} \left(\varphi\mleft(Y\mright) \lor \psi\mleft(Y\mright)\right)$. For each $n \in \mathbb{N}$, we know by atomic decidability that $\left(Z\right)_n \prec X \lor \neg \left(Z\right)_n \prec X$, and in the former case it follows that $\varphi\mleft(\left(Z\right)_n\mright) \lor \psi\mleft(\left(Z\right)_n\mright)$. By numerical omniscience, we have
  \[\opA{n} \left(\left(Z\right)_n \prec X \rightarrow \varphi\mleft(\left(Z\right)_n\mright)\right) \lor \opE{n} \left(\left(Z\right)_n \prec X \land \psi\mleft(\left(Z\right)_n\mright)\right).\]
  The former disjunct implies $\opA{Y \prec X} \varphi\mleft(Y\mright)$, while the latter disjunct implies $\opE{Y \prec X} \psi\mleft(Y\mright)$.
\end{proof}

We also need a construction to encode countable collections of second-order objects as one single second-order object in $\CM$. Instead of the obvious coding using slices $\left(Z\right)_n$, we will use a slightly modified construction so that we can accommodate the empty or finite collections more conveniently: for any second-order set $X$, let $\left[X\right]_n = \left\{m : \tuple{n, 0} \in_1 X \land \tuple{n, m + 1} \in_1 X\right\}$, and define a new binary relation $\in^*$ between two second-order objects as
\[Y \in^* X \defeqv \opE{n} \left(\tuple{n, 0} \in_1 X \land Y = \left[X\right]_n\right).\]

It follows that we can denote the following simple abbreviations:
\begin{align*}
  X \subseteq^* Y          & \defeqv \opA{Z \in^* X} Z \in^* Y,             \\
  X \subseteq^* \mathbf{X} & \defeqv \opA{Z \in^* X} Z \in_2 \mathbf{X},    \\
  X =^* Y                  & \defeqv X \subseteq^* Y \land Y \subseteq^* X,
\end{align*}
and immediately observe that omniscience for countable collections,
\[\opA{Y \in^* X} \left(\varphi\mleft(Y\mright) \lor \psi\mleft(Y\mright)\right) \rightarrow \left(\opA{Y \in^* X} \varphi\mleft(Y\mright) \lor \opE{Y \in^* X} \psi\mleft(Y\mright)\right),\]
follows directly from numerical omniscience, thus the relations defined here are all decidable, i.e.\ the law of excluded middle holds for them.

In the same way, we can also denote certain set-forming expressions as second-order terms, such as:
\begin{align*}
  \left\{X_1, \ldots, X_k\right\}                    & = \left\{\tuple{n, m} : 1 \leq n \leq k \land \left(m = 0 \lor \opE{i \in_1 X_n} m = i + 1\right)\right\},                    \\
  \bigcup X                                          & = \left\{\tuple{\tuple{n, m}, i} : \tuple{n, 0} \in_1 X \land \tuple{n, \tuple{m, i} + 1} \in_1 X\right\},                    \\
  \left\{Y \in^* X : \varphi\mleft(Y\mright)\right\} & = \left\{\tuple{n, m} : \tuple{n, 0} \in_1 X \land \tuple{n, m} \in_1 X \land \varphi\mleft(\left[X\right]_n\mright)\right\},
\end{align*}
where in the last definition, the formula $\varphi\mleft(Y\mright)$ must be decidable. It is easy to verify that these collections `contain' the same second-order sets as their third-order counterparts.

Now, an analogue of Zorn's lemma\inds{Zorn's lemma} can be stated as follows\footnote{\label{ftn:zorn-lemm-proof-sigma}The proof of Zorn's lemma presented here is a slightly modified version of the proof in \cite{wang-sw25-cm}*{Proposition 3.2}. The readers should observe that the presentation in this thesis only invokes the transfinite induction axiom for formulae with a single unbounded second-order existential quantifier (on top of existing quantifiers in $\varphi$), which will be important for discussions in \autoref{sec:cm-gwo-pred}.}:
\begin{proposition}[$\CM + \GWO$]
  \label{prop:cm-gwo-zorn}
  Let $\varphi\mleft(X\mright)$ be a decidable formula \uppar{possibly with other parameters} where $X$ denotes a countable second-order collection encoded as above, i.e.\ we have $\opA{X} \left(\varphi\mleft(X\mright) \lor \neg \varphi\mleft(X\mright)\right)$. If $\varphi$ is downward closed, i.e.
  \[\opA{X, Y} \left(\left(\varphi\mleft(X\mright) \land Y \subseteq^* X\right) \rightarrow \varphi\mleft(Y\mright)\right),\]
  and also closed under unions of countable chains, i.e.
  \[\opA{X} \left(\left(\opA{Y, Z \in^* X} \left(Y \subseteq^* Z \lor Z \subseteq^* Y\right) \land \opA{Y \in^* X} \varphi\mleft(Y\mright)\right) \rightarrow \varphi\mleft(\bigcup X\mright)\right),\]
  then there exists some $\mathbf{X}$ satisfying $\opA{X \subseteq^* \mathbf{X}} \varphi\mleft(X\mright)$, which is maximal among such subsets in the sense that
  \[\opA{Y} \left(Y \not\in \mathbf{X} \rightarrow \opE{Z \subseteq^* \mathbf{X}} \neg \varphi\mleft(Z \cup \left\{Y\right\}\mright)\right).\]
\end{proposition}

\begin{lemma}[$\CM + \GWO$]
  \label{lem:cm-gwo-countable-family-bound}
  For any $X$, there exists $Y$ such that
  \[\opA{n} \left(X\right)_n \prec Y.\]
  The same immediately holds for the countable collection encoding we defined above, i.e.\ for any $X$, there exists $Y$ such that $\opA{Z \in^* X} Z \prec Y$.
\end{lemma}

\begin{proof}
  By the axiom of countable initial segments,
  \[\opA{n} \opE{W} \opA{Y \preceq \left(X\right)_n} \opE{m} Y = \left(W\right)_m,\]
  and by the recursion axiom, we can find some $U$ satisfying that
  \[\opA{n} \opA{Y \preceq \left(X\right)_n} \opE{m} Y = \left(\left(U\right)_n\right)_m.\]
  We shall diagonalise and construct
  \[Y_0 = \left\{\tuple{n, m} : \tuple{n, m} \not\in_1 \left(\left(U\right)_n\right)_m\right\},\]
  so that for any $n$, $\neg Y_0 \preceq \left(X\right)_n$. By linearity of the well-ordering, it must follow that $\opA{n} \left(X\right)_n \prec Y_0$.

  For a countable family $X$ encoded through $\in^*$ instead, simply apply the same process to
  \[X' = \left\{\tuple{n, m} : \tuple{n, m + 1} \in X\right\}.\]
  Clearly $Z = \left[X\right]_n$ implies $\opE{n} Z = \left(X'\right)_n$, and thus $\opA{Z \in^* X} Z \prec Y_0$ for $Y_0$ constructed as above.
\end{proof}

\begin{proof}[Proof of \autoref{prop:cm-gwo-zorn}]
  Define formulae
  \begin{align*}
    \psi\mleft(X, Z\mright)   & \defeqv \varphi\mleft(Z\mright) \land \left(\varphi\mleft(Z \cup \left\{X\right\}\mright) \rightarrow X \in^* Z\right), \\
    \theta\mleft(X, Z\mright) & \defeqv \opA{Y \in^* Z} Y \preceq X \land \opA{Y \preceq X} \psi\mleft(Y, Z^{\preceq Y}\mright).
  \end{align*}
  where $Z^{\preceq Y} = \left\{W \in^* Z : W \preceq Y\right\}$. Fix any $Z_1, Z_2$, and we first show by transfinite induction on $X$ that
  \[\opA{X} \left(\theta\mleft(X, Z_1^{\preceq X}\mright) \land \theta\mleft(X, Z_2^{\preceq X}\mright) \rightarrow Z_1^{\preceq X} =^* Z_2^{\preceq X}\right).\]
  This is because for any $Y \prec X$, $\theta\mleft(X, Z^{\preceq X}\mright)$ clearly implies $\theta\mleft(Y, Z^{\preceq Y}\mright)$. Thus from the inductive hypothesis, we know that $\opA{Y \prec X} Z_1^{\preceq Y} =^* Z_2^{\preceq Y}$, and clearly $Z_1^{\prec X} =^* \bigcup_{Y \prec X} Z_1^{\preceq Y} =^* \bigcup_{Y \prec X} Z_2^{\preceq Y} =^* Z_2^{\prec X}$. Now assume that $X \in^* Z_1$, then
  \[Z_1^{\preceq X} =^* Z_1^{\prec X} \cup \left\{X\right\} =^* Z_2^{\prec X} \cup \left\{X\right\} =^* Z_2^{\preceq X} \cup \left\{X\right\},\]
  so $\varphi\mleft(Z_2^{\preceq X} \cup \left\{X\right\}\mright)$ holds and thus $X \in^* Z_2$. By symmetry, the converse can be shown assuming $X \in^* Z_2$, so we can conclude that $Z_1^{\preceq X} =^* Z_2^{\preceq X}$.

  For any $X, Z$, observe that $\theta\mleft(X, Z\mright)$ implies $Z =^* Z^{\preceq X}$. Thus the induction above implies that for any $X$, the second-order object $Z$ satisfying $\theta\mleft(X, Z\mright)$ is unique up to $=^*$, i.e.\
  \[\opA{X, Z_1, Z_2} \left(\theta\mleft(X, Z_1\mright) \land \theta\mleft(X, Z_2\mright) \rightarrow Z_1 =^* Z_2\right).\]

  We then show by another transfinite induction that $\opA{X} \opE{Z} \theta\mleft(X, Z\mright)$. Fix $X$ and pick $W$ by the axiom of countable initial segments such that $\opA{Y \prec X} \opE{n} Y = \left(W\right)_n$, and the inductive hypothesis would imply
  \[\opA{n} \opE{Z} \left(\left(W\right)_n \prec X \rightarrow \theta\mleft(\left(W\right)_n, Z\mright)\right).\]
  By the recursion axiom in $\CM$, we find some second-order $U$ satisfying that $\opA{n} \left(\left(W\right)_n \prec X \rightarrow \theta\mleft(\left(W\right)_n, \left(U\right)_n\mright)\right)$, and now define
  \[Z' = \left\{\tuple{n, m} : \left(W\right)_n \prec X \land \left(m = 0 \lor \opE{i \in_1 \left(U\right)_n} m = i + 1\right)\right\},\]
  essentially the countable collection $\left\{\left(U\right)_n : \left(W\right)_n \prec X\right\}$. For each $Y \in^* Z'$, we then have $\varphi\mleft(Y\mright)$. $Z'$ is additionally a chain because for any $\left[Z'\right]_m, \left[Z'\right]_n \in^* Z'$, we can assume without loss of generality that $\left(W\right)_m \preceq \left(W\right)_n$, and thus
  \[\left[Z'\right]_m =^* \left[Z'\right]_n^{\preceq \left(W\right)_m} \subseteq^* \left[Z'\right]_n\]
  by uniqueness above. Hence $\varphi\mleft(\bigcup Z'\mright)$. It also follows that $\left[Z'\right]_n = \left(\bigcup Z'\right)^{\preceq \left(W\right)_n}$ for any $n \in \mathbb{N}$ such that $\left(W\right)_n \prec X$.

  Therefore, it suffices to set $Z = \left(\bigcup Z'\right) \cup \left\{X\right\}$ if $\varphi\mleft(\left(\bigcup Z'\right) \cup \left\{X\right\}\mright)$ holds and $Z = \bigcup Z'$ otherwise (which is possible because $\varphi$ is decidable). It is easy to verify that $\theta\mleft(X, Z\mright)$ holds indeed.

  Finally, by the unique existence of some $Z$ satisfying $\theta\mleft(Y, Z\mright)$ up to $=^*$, we can decidably comprehend the diagonal set
  \[\mathbf{X} = \left\{Y : \opE{Z} \left(\theta\mleft(Y, Z\mright) \land Y \in^* Z\right)\right\}.\]
  For any countable subset $X \subseteq^* \mathbf{X}$, by \autoref{lem:cm-gwo-countable-family-bound}, there exists $Y_0$ such that $\opA{Y \in^* X} Y \prec Y_0$. Let $Z_0$ be a countable collection satisfying $\theta\mleft(Y_0, Z_0\mright)$, then our definition of $\mathbf{X}$ and the uniqueness of $Z_0$ ensure that for any $Y \preceq Y_0$,
  \[Y \in_2 \mathbf{X} \leftrightarrow Y \in^* Z_0.\]
  In other words, $X \subseteq^* Z_0$, and $\varphi\mleft(X\mright)$ follows immediately from $\varphi\mleft(Z_0\mright)$ by downward closedness.

  On the other hand, suppose that some $Y \not\in_2 \mathbf{X}$. There must exist $Z$ satisfying $\theta\mleft(Y, Z\mright)$ and $Y \not\in^* Z$, which, by our stipulation of $\psi$ at the beginning of this proof, implies $\neg \varphi\mleft(Z \cup \left\{Y\right\}\mright)$. The uniqueness of $Z$ ensures that also $Z \subseteq^* \mathbf{X}$. Therefore, we have checked that $\mathbf{X}$ has all the desired properties.
\end{proof}

This proposition can be applied just like the classical Zorn's lemma to show that one can find maximal subsets satisfying certain properties in uncountable collections, as long as the corresponding decidability criterion can be verified through intuitionistic reasoning. For example, consider any field $\mathbf{F}$ and an $\mathbf{F}$-vector space $\mathbf{V}$, with their respective unary and binary operations also given as third-order objects, following the conventions of \cite{weaver09-cm}*{section 3}, then it follows that:
\begin{corollary}[$\CM + \GWO$, \cite{wang-sw25-cm}*{Corollary 3.3}]
  Suppose that the formula
  \[\opE{Y} \left(\opA{n \leq k} \left(Y\right)_n \in \mathbf{F} \land \opE{n \leq k} \left(Y\right)_n \neq 0_\mathbf{F} \land 0_\mathbf{V} = \sum_{n = 0}^k \left(Y\right)_n \cdot \left(X\right)_n\right),\]
  i.e.\ informally the statement `$\left(X\right)_0, \ldots, \left(X\right)_k$ are linearly dependent as vectors in $\mathbf{V}$', is decidable \uppar{whenever $\left(X\right)_0, \ldots, \left(X\right)_k \in_2 \mathbf{V}$}, then $\mathbf{V}$ has a basis $\mathbf{X} \subseteq \mathbf{V}$ such that every vector $V \in \mathbf{V}$ is spanned by finitely many elements in $\mathbf{X}$.
\end{corollary}

This applies to show the existence of a basis in the following scenarios:
\begin{enumerate}[label=\uplabel{\alph*}]
  \item when the field $\mathbf{F}$ is \emph{countable} --- for example, when showing that the Hamel basis for the reals $\mathbb{R}$ as a vector space over the rationals $\mathbb{Q}$ exists;
  \item when the vector space $\mathbf{V}$ has an \emph{inner product} $\mathbf{V}^2 \rightarrow \mathbf{F}$;
  \item when $\mathbf{F}$ is a complete $\mathbb{R}$-normed field with a countable dense set $C \subseteq^* \mathbf{F}$ with respect to the induced norm, and $\mathbf{V}$ has a compatible norm $\left\|-\right\| : \mathbf{V} \rightarrow \mathbf{R}$ --- for example, when $\mathbf{V}$ is any real or complex Banach space.
\end{enumerate}
The proofs of these results are sketched in \cite{wang-sw25-cm}, but are omitted in this thesis since they are not related to the topic of realisability models. At any rate, I think they demonstrate the relevance of the global well-ordering axioms to mathematics founded on the system $\CM$ and justify our endeavour to analyse the proof-theoretic commitments in them.

\section{Realising the axioms of \texorpdfstring{$\GWO$}{GWO}}
\label{sec:oa-cm-gwo}

Now, we look back at the same realisability model for $\CM$ as in \autoref{sec:oa-cm}, but in order to realise the axioms of $\GWO$, we will need a stronger base theory in which we can actually construct a well-ordering on our domain for second-order variables,
\[S_2 = \left\{d \in \Hyp : \opA{n} dn {\downarrow} \in \left\{0, 1\right\}\right\}.\]
We shall do this by developing the theory of hyperarithmetic sets in the theory $\ATR_0$, consisting of $\ACA_0$ and the additional axiom of arithmetic transfinite recursion (cf.\ \cite{simpson09-soa}*{section V.2}), and show that the (total) functions denoted by indices in $S_2$ are precisely the hyperarithmetic functions.

\subsection*{Hyperarithmetic theory in $\ATR_0$}

The constructions we will need here roughly follow \cite{simpson09-soa}*{section VIII.3}: let $\tilde{\pi}\mleft(e, m, X\mright)$ again\footnote{We use $\tilde{\pi}$ since, strictly speaking, this formula has a different number of free variables than $\pi$ considered in \autoref{sec:oa-cm}.} denote a universal $\Pi^0_1$-formula with all free variables indicated, such that for any $\Pi^0_1$-formula $\varphi\mleft(e', m, X\mright)$ with the same number of free variables, $\ACA_0$ proves that
\[\opA{e'} \opE{e} \opA{m} \opA{X} \left(\varphi\mleft(e', m, X\mright) \leftrightarrow \tilde{\pi}\mleft(e, m, X\mright)\right).\]
We can then formalise relative Turing computability in second-order arithmetic using $\tilde{\pi}$ by saying that a set $Y \subseteq \mathbb{N}$ is \emph{$X$-recursive} and write $Y \leq_T X$ if there exists $e = \tuple{e_0, e_1} \in \mathbb{N}$ such that
\[Y = \left\{m : \tilde{\pi}\mleft(e_0, m, X\mright)\right\} = \left\{m : \neg \tilde{\pi}\mleft(e_1, m, X\mright)\right\}.\]
We say that $e$ is the \emph{$X$-recursive index} of $Y$, and we write
\[\Rec^X = \left\{e \in \mathbb{N} : \opA{m} \left(\tilde{\pi}\mleft(e_0, m, X\mright) \leftrightarrow \neg \tilde{\pi}\mleft(e_1, m, X\mright)\right)\right\}\]
to denote the set of all $X$-recursive indices\inda{recursive indices, $\Rec^X$}. For any $e \in \Rec^X$, let $R^X_e$ denote the set defined by $e$. The superscript $X$ will be omitted when $X = \varnothing$.

\begin{definition}[Ordinal notations]
  Let $\mathcal{O}_+$ denote the set of indices $\alpha \in \mathbb{N}$ where $\alpha_0 \in \Rec$, $R_{\alpha_0}$ is a weak linear ordering on the set
  \[\field\mleft(R_{\alpha_0}\mright) = \left\{n \in \mathbb{N} : \opE{m} \left(\tuple{n, m} \in R_{\alpha_0} \lor \tuple{m, n} \in R_{\alpha_0}\right)\right\}\]
  and $\alpha_1 \in \field\mleft(R_{\alpha_0}\mright)$. The canonical partial ordering on $\mathcal{O}_+$ is given by
  \[\alpha \leq_\mathcal{O} \beta \defeqv \alpha_0 = \beta_0 \land \tuple{\alpha_1, \beta_1} \in R_{\alpha_0}.\]

  We then let $\mathcal{O}$\inda{ordinal notations, $\mathcal{O}$} denote the class of indices $\alpha$ such that $<_\mathcal{O}$ is a well-ordering on the initial segment $\left\{\beta \in \mathcal{O}_+ : \beta <_\mathcal{O} \alpha\right\}$, i.e.
  \[\mathcal{O} = \left\{\alpha \in \mathcal{O}_+ : \neg\!\opE{f : \mathbb{N} \rightarrow \mathcal{O}_+} \left(f\mleft(0\mright) = \alpha \land \opA{n} f\mleft(n + 1\mright) <_\mathcal{O} f\mleft(n\mright)\right)\right\}\]
  Given any $\alpha \in \mathcal{O}$, we say its \emph{order type} is precisely the order type\footnote{We will not fix a concrete object representing an `order type' here in second-order arithmetic. Rather, given two well-orderings $\tuple{X, \leq_X}$ and $\tuple{Y, \leq_Y}$, we simply treat statements $\otyp\mleft(X, \leq_X\mright) = \otyp\mleft(Y, \leq_Y\mright)$ and $\otyp\mleft(X, \leq_X\mright) \leq \otyp\mleft(Y, \leq_Y\mright)$ as abbreviations for the respective assertions that corresponding comparison maps exist.} of this initial segment
  \[\otyp\mleft(\alpha\mright) = \otyp\mleft(\left\{\beta \in \mathcal{O}_+ : \beta <_\mathcal{O} \alpha\right\}, \leq_\mathcal{O}\mright).\]
\end{definition}

\begin{definition}[$H$-sets]
  The \emph{Turing jump}\inda{Turing jump} of a set $X \subseteq \mathbb{N}$ is given as
  \[\TJ\mleft(X\mright) = \left\{\tuple{e, m} : \tilde{\pi}\mleft(e, m, X\mright)\right\}.\]
  For an ordinal notation $\alpha$, the \emph{$\alpha$-iterated Turing jumps}\inda{Turing jump!$\alpha$-iterated}, or as Simpson \cite{simpson09-soa} calls it, an \emph{$H$-set} for the ordinal $\alpha$, is some set $Y$ satisfying the following formula:
  \[H\mleft(\alpha, Y\mright) \defeqv \alpha \in \mathcal{O}_+ \land Y = \left(Y\right)_{<_\mathcal{O} \alpha} \land \opA{\beta <_\mathcal{O} \alpha} \left(Y\right)_\beta = \TJ\mleft(\left(Y\right)_{<_\mathcal{O} \beta}\mright),\]
  where the truncation of $Y$ before $\alpha$ is
  \[\left(Y\right)_{<_\mathcal{O} \alpha} = \left\{\tuple{\beta, m} : \beta <_\mathcal{O} \alpha \land m \in \left(Y\right)_{\beta}\right\}.\]
\end{definition}

As in \cite{simpson09-soa}*{Theorem VIII.3.15}, $\ATR_0$ suffices to prove $\opA{\alpha \in \mathcal{O}} \opE{Y} H\mleft(\alpha, Y\mright)$. We additionally know that:
\begin{proposition}[$\ACA_0$, \cite{simpson09-soa}*{Lemma VIII.3.14}]
  \label{prop:ord-cmp-red-to-hyp}
  Given $\alpha, \beta \in \mathcal{O}$ and some $Y$ such that $H\mleft(\alpha, Y\mright)$, then the order types $\otyp\mleft(\alpha\mright)$ and $\otyp\mleft(\beta\mright)$ are comparable, and the comparison map is $\leq_T Y$.
\end{proposition}

Thus, any $\alpha, \beta \in \mathcal{O}$ are comparable in $\ATR_0$, and we can verify that the comparison function is effectively computed by our weak pca $\Hyp$:
\begin{corollary}[$\ATR_0$]
  \label{cor:hyp-well-order-cmp}
  There exists $C \in \Hyp$ such that for any $\alpha \in \mathcal{O}$, $\beta \in \mathbb{N}$, $C\alpha\beta {\downarrow} \in \left\{0, 1\right\}$ and
  \[C\alpha\beta = 1 \leftrightarrow \beta \in \mathcal{O} \land \otyp\mleft(\beta\mright) \leq \otyp\mleft(\alpha\mright).\]
\end{corollary}

\begin{proof}
  By \autoref{lem:sigma11-smn}, there is $C \in \Hyp$ such that for any $\alpha, \beta \in \mathbb{N}$, $C\alpha\beta {\downarrow} = d$ if and only if
  \begin{enumerate}[label=\uplabel{\alph*}]
    \item \label{itm:hyp-well-order-cmp-con1} either $d = 1$ and there exists $Y$ such that $H\mleft(\alpha, Y\mright)$, $\beta \in \mathcal{O}_+$ and there exists $e \in \Rec^Y$ such that $R^Y_e$ is an order-isomorphism between the set $\left\{\gamma \in \mathcal{O}_+ : \gamma <_\mathcal{O} \beta\right\}$ and an initial segment of $\left\{\gamma \in \mathcal{O}_+ : \gamma <_\mathcal{O} \alpha\right\}$;
    \item or $d = 0$ and there exists $Y$ such that $H\mleft(\alpha, Y\mright)$ and the remainder of condition \ref{itm:hyp-well-order-cmp-con1} does not hold.
  \end{enumerate}
  It is easy to see that this definition is $\Sigma^1_1$ and precisely one of the two conditions must apply for any $\alpha \in \mathcal{O}$ since the corresponding $H$-set $Y$ is necessarily unique by arithmetic transfinite induction, thus indeed $C\alpha\beta {\downarrow} \in \left\{0, 1\right\}$ for any $\alpha \in \mathcal{O}$ and $\beta \in \mathbb{N}$.

  Now suppose that $C\alpha\beta = 1$, then $\beta \in \mathcal{O}$ because any strictly decreasing sequence below $\beta$ will correspond to a strictly decreasing sequence below $\alpha$ through $R^Y_e$. Conversely, suppose that $\beta \in \mathcal{O}$ and $\otyp\mleft(\beta\mright) \leq \otyp\mleft(\alpha\mright)$, then \autoref{prop:ord-cmp-red-to-hyp} ensures that the comparison function is $Y$-recursive, so $C\alpha\beta$ must evaluate to $1$ indeed.
\end{proof}

Next, we consider a representation of all hyperarithmetic sets as the collection of indices
\[\mathcal{H} = \left\{a : a_0 \in \mathcal{O} \land \opE{Y} \left(H\mleft(a_0, Y\mright) \land a_1 \in \Rec^Y\right)\right\}.\]
We say that a set $X \subseteq \mathbb{N}$ is \emph{hyperarithmetic}\inda{hyperarithmetic set} if there exists $a \in \mathcal{H}$ and $Y \subseteq \mathbb{N}$ such that $H\mleft(a_0, Y\mright)$ and $X = R^Y_{a_1}$. We say that $a = \tuple{a_0, a_1}$ is a \emph{hyperarithmetic index} of $X$.

We shall effectivise the Kleene--Souslin theorem \cite{simpson09-soa}*{Theorem VIII.3.19}, and state that the weak pca $\Hyp$ computes a function from $S_2$ to $\mathcal{H}$ that maps characteristic functions $\mathbb{N} \rightarrow \left\{0, 1\right\}$ to their corresponding hyperarithmetic sets, i.e.
\begin{theorem}[$\ATR_0$]
  \label{thm:conv-s2-to-hyp}
  There exists $H \in \Hyp$ such that
  \begin{enumerate}[label=\uplabel{\roman*}]
    \item for each $e \in S_2$, $He {\downarrow} \in \mathcal{H}$ is a hyperarithmetic index of $\left\{n : en = 1\right\}$;
    \item for $d, e \in S_2$, $Hd = He$ if and only if $d \equiv e$.
  \end{enumerate}
\end{theorem}

\begin{proof}
  We will replicate some key constructions in Simpson's proof, and then show that they can be effectively specified by a $\Sigma^1_1$-formula. For any $e \in S_2$, it is easy to see that $en = 1$ if and only if $\opE{X} \pi\mleft(e_0, e_1, n, 1, X\mright)$, which is $\Sigma^1_1$. Observe that $en = 1$ also if and only if $en \neq 0$, if and only if
  \[\opA{X} \neg \pi\mleft(e_0, e_1, n, 0, X\mright),\]
  which is $\Pi^1_1$; so $en = 1$ is a $\Delta^1_1$-assertion.

  By Kleene normal form theorem for $\Pi^1_1$-formulae (\cite{simpson09-soa}*{Theorem V.1.4}), there exists a $\Delta^0_0$-formula $\theta\mleft(e, n, \tau\mright)$ such that $\mathrm{ACA}_0$ proves
  \[\opA{e \in S_2} \opA{n} \left(en = 1 \leftrightarrow \opA{f} \opE{k} \neg \theta\mleft(e, n, f \upharpoonright k\mright)\right)\]
  where $f$ ranges over second-order objects representing functions $\mathbb{N} \rightarrow \mathbb{N}$, and $f \upharpoonright k = \tuple{f\mleft(0\mright), \ldots, f\mleft(k - 1\mright)}$ is a finite tuple encoded as a first-order object. We then consider recursive trees
  \[T^{(e)}_n = \left\{\tau \in \mathbb{N}^{< \mathbb{N}} : \opA{k \leq \left|\tau\right|} \theta\mleft(e, n, \tau \upharpoonright k\mright)\right\}\]
  and let $\KB^{(e)}_n = \KB\mleft(T^{(e)}_n\mright)$ be the Kleene--Brouwer ordering of $T^{(e)}_n$, i.e. $\sigma \leq_\KB \tau$ if either $\sigma \sqsupseteq \tau$ or that $\sigma$ has a smaller component at the first index where the two differ (cf.\ \cite{simpson09-soa}*{Section V.1}). It follows that
  \begin{align*}
    en = 1 \quad & \text{iff} \quad \opA{f} \opE{k} \neg \theta\mleft(e, n, f \upharpoonright k\mright), \\
                 & \text{iff} \quad \text{$T^{(e)}_n$ has no path},                                      \\
                 & \text{iff} \quad \text{$\KB^{(e)}_n$ is a recursive well-ordering}.
  \end{align*}

  Now, we work in $\ATR_0$ where $\opA{\alpha \in \mathcal{O}} \opE{Y} H\mleft(\alpha, Y\mright)$. Then Simpson's argument applies to justify that the collection of well-orderings $\KB^{(e)}_n$ (for all $n \in \mathbb{N}$, given a fixed $e$) has bounded order type in $\mathcal{O}$ --- otherwise, $\mathcal{O}$ can be defined precisely as the $\Sigma^1_1$-class of $\beta \in \mathcal{O}_+$ such that $C\alpha\beta = 1$ for some $\alpha = \KB^{(e)}_n$ where $n \in \mathbb{N}$ and $en = 1$, contradicting \cite{simpson09-soa}*{Lemma VIII.3.4}.

  Let $\alpha \in \mathcal{O}$ be such a bound and $\Omega \subseteq \mathbb{N}$ be such that $H\mleft(\alpha, \Omega\mright)$. By \autoref{prop:ord-cmp-red-to-hyp}, the order types of any initial segment of $\KB^{(e)}_n$ and any $\beta \leq_\mathcal{O} \alpha$ will all be comparable via comparison maps $\leq_T \Omega$, which means for any such well-orderings $A, B$, a comparison statement $\otyp\mleft(A\mright) \leq \otyp\mleft(B\mright)$ will be an arithmetic formula `there exists $a \in \mathbb{N}$ such that $R^\Omega_a$ is an order isomorphism between $A$ and an initial segment of $B$'.

  Therefore, by $\Delta^1_1$-comprehension (cf.\ \cite{simpson09-soa}*{Theorem V.5.1}),
  \[\Gamma = \left\{\beta \leq_\mathcal{O} \alpha : \opA{n} \left(en = 1 \rightarrow \otyp\mleft(\KB^{(e)}_n\mright) \leq \otyp\mleft(\beta\mright)\right)\right\}\]
  is a set and we can find $\lambda = \min \Gamma$. For any $\gamma <_\mathcal{O} \lambda$, there must exist $n \in \mathbb{N}$ such that $en = 1$ and $\otyp\mleft(\gamma\mright) < \otyp\mleft(\KB^{(e)}_n\mright)$. In fact, these two characteristics determine $\lambda$ uniquely up to order isomorphism, i.e.\ we can prove the following lemma that:
  \begin{lemma*}
    Suppose that $\rho \in \mathcal{O}_+$ satisfies both
    \begin{enumerate}[label=\uplabel{\roman*}]
      \item $\otyp\mleft(\KB^{(e)}_n\mright) \leq \otyp\mleft(\rho\mright)$ for any $n$ such that $en = 1$,
      \item and for any $\gamma <_\mathcal{O} \rho$, there exists $n \in \mathbb{N}$ such that $en = 1$ and $\otyp\mleft(\gamma\mright) < \otyp\mleft(\KB^{(e)}_n\mright)$,
    \end{enumerate}
    then $\rho \in \mathcal{O}$ and $\otyp\mleft(\rho\mright) = \otyp\mleft(\lambda\mright)$.
  \end{lemma*}

  \begin{proof}
    Firstly, assume that $\rho \not\in \mathcal{O}$, then there exists a strictly decreasing sequence $f : \mathbb{N} \rightarrow \mathcal{O}_+$ such that $f\mleft(0\mright) = \rho$. We have $f\mleft(1\mright) <_\mathcal{O} \rho$, hence we can consider some order-preserving map $g : \left\{\gamma : \gamma <_\mathcal{O} f\mleft(1\mright)\right\} \rightarrow \KB^{(e)}_n$. Now $m \mapsto g\mleft(f\mleft(m + 2\mright)\mright)$ is a descending sequence in $\KB^{(e)}_n$, a contradiction.

    Therefore, we must have $\rho \in \mathcal{O}$. We then know in $\ATR_0$ that the order types of $\rho$ and $\lambda$ are comparable, but it is easy to verify that we can have neither $\otyp\mleft(\rho\mright) > \otyp\mleft(\lambda\mright)$ nor $\otyp\mleft(\rho\mright) < \otyp\mleft(\lambda\mright)$.
  \end{proof}

  Finally, given any $\lambda \in \mathcal{O}$, we can remove the ambiguity around the order-isomorphic ordinal notations by requiring that $\lambda$ is minimal among all order-isomorphic elements in $\mathcal{O}_+$, with respect to the natural ordering on $\mathbb{N}$. Such a minimum exists by $\Delta^0_0$-induction because, using some $Y \subseteq \mathbb{N}$ satisfying $H\mleft(\lambda, Y\mright)$ as a parameter, the set of elements in $\mathcal{O}_+$ order-isomorphic to $\lambda$ is an arithmetic set. The rest of Simpson's argument in \cite{simpson09-soa}*{Theorem VIII.3.19} applies to show that if $H\mleft(\lambda + \omega, Y\mright)$, then $\left\{n : en = 1\right\} \leq_T Y$.

  Therefore, by \autoref{lem:sigma11-smn}, we can stipulate $He {\downarrow} = a$ when there exists $\alpha \in \mathcal{O}_+$ and $Y \subseteq \mathbb{N}$ satisfying all of the following:
  \begin{enumerate}[label=\uplabel{\roman*}]
    \item $H\mleft(a_0, Y\mright)$ and\footnote{Here we fix $\alpha + \omega$ to mean the following specific recursive well-ordering in $\mathcal{O}$ which can be computed uniformly from $\alpha$: order the set $\left\{2\beta + 2 : \beta <_\mathcal{O} \alpha\right\}$ by
            \[2\beta_1 + 2 < 2\beta_2 + 2 \defeqv \beta_1 <_\mathcal{O} \beta_2,\]
            and append the sequence $1, 3, 5, 7, \ldots, 0$ in this order. Then the initial segment before $0$ has order type $\alpha + \omega$.} $\otyp\mleft(a_0\mright) = \otyp\mleft(\alpha + \omega\mright)$;
    \item for any $n \in \mathbb{N}$, $en = 1 \rightarrow \otyp\mleft(\KB^{(e)}_n\mright) \leq \otyp\mleft(\alpha\mright)$;
    \item \label{itm:conv-s2-to-hyp-ord-uniq} for any $\beta <_\mathcal{O} \alpha$, there exists $n \in \mathbb{N}$ such that $en = 1$ and $\otyp\mleft(\beta\mright) < \otyp\mleft(\KB^{(e)}_n\mright)$;
    \item for any $\gamma \in \mathcal{O}_+$, if there is a $Y$-recursive order-isomorphism between $\gamma$ and $a_0$, then $\gamma \geq a_0$;
    \item $a_1 \in \Rec^Y$ and for any $n \in \mathbb{N}$, $en = 1 \leftrightarrow n \in R^Y_{a_1}$;
    \item \label{itm:conv-s2-to-hyp-rec-ind-uniq} for any $b \in \Rec^Y$, if $n \in R^Y_{a_1} \leftrightarrow n \in R^Y_b$ holds for all $n \in \mathbb{N}$, then $b \geq a_1$.
  \end{enumerate}
  Using the fact that $en = 1$ is $\Delta^1_1$ and the comparison maps for the concerned order types are all $Y$-recursive, this definition can be written as a $\Sigma^1_1$-formula, and conditions \ref{itm:conv-s2-to-hyp-ord-uniq}, \ref{itm:conv-s2-to-hyp-rec-ind-uniq} ensure that $a = \tuple{a_0, a_1}$ satisfying them is necessarily unique. Thus for any $e \in S_2$, we see that $He {\downarrow}$.

  To verify that $Hd = He$ if and only if $d \equiv e$, simply observe that the value of $He$ is uniquely determined by the set $\left\{n : en = 1\right\}$.
\end{proof}

The (rough) inverse is much easier and can be constructed in $\ACA_0$:
\begin{proposition}[$\mathrm{ACA}_0$]
  There exists\footnote{Formally speaking, $H \in \Hyp$ is only a natural number. But I feel it is understandable here to abuse notation and let $H^{-1}$ denote another natural number, which computes roughly the inverse of what $H$ computes.} $H^{-1} \in \Hyp$ such that, for any $a \in \mathcal{H}$ and any $Y \subseteq \mathbb{N}$ satisfying $H\mleft(a_0, Y\mright)$, we have $H^{-1} a {\downarrow} \in S_2$ and
  \[\opA{n} \left(H^{-1} an = 1 \leftrightarrow n \in R^Y_{a_1}\right).\]
\end{proposition}

\begin{proof}
  By \autoref{lem:sigma11-smn}, pick $H^{-1} \in \Hyp$ such that $H^{-1} an {\downarrow} = i$ if and only if
  \[\opE{Y} \left(H\mleft(a_0, Y\mright) \land \left(\left(\varphi\mleft(a, n, Y\mright) \land i = 1\right) \lor \left(\neg \varphi\mleft(a, n, Y\mright) \land i = 0\right)\right)\right)\]
  where $\varphi\mleft(a, n, Y\mright) = a_1 \in \Rec^Y \land n \in R^Y_{a_1}$ is arithmetic.
\end{proof}

Do note that $H$ and $H^{-1}$ are inverses of each other only up to extensional equivalences, i.e.\ for any $e \in S_2$, we have $H^{-1} \mleft(He\mright) \equiv e$, and for any $a \in \mathcal{H}$, $a$ and $H \mleft(H^{-1} a\mright)$ are hyperarithmetic indices for the same set.

\subsection*{Realising the axioms of $\GWO$}

By \autoref{lem:sigma11-smn}, let $R \in \Hyp$ be such that $Ra {\downarrow} = d$ if and only if there exists $Y \subseteq \mathbb{N}$ such that $H\mleft(a_0, Y\mright)$ and either $a_1 \in \Rec^Y$ and $d = 1$ or $a_1 \not\in \Rec^Y$ and $d = 0$, so that for any $a \in \mathbb{N}$, if $a_0 \in \mathcal{O}$ then $Ra {\downarrow} \in \left\{0, 1\right\}$ and
\[Ra = 1 \leftrightarrow a \in \mathcal{H}.\]
Let $C^* \in \Hyp$ be the function
\[\oplam{ab} \left\{\begin{aligned}
     & 1 &  & \text{if $C \mleft(\bp_0 a\mright) \mleft(\bp_0 b\mright) = 1 \land \left(C \mleft(\bp_0 b\mright) \mleft(\bp_0 a\mright) = 0 \lor b < a\right)$} \land Rb = 1, \\
     & 0 &  & \text{otherwise},
  \end{aligned}\right.\]
defined using $C$ from \autoref{cor:hyp-well-order-cmp}. To be pedantic, observe that when $\bp_0 a = a_0 \in \mathcal{O}$, $C \mleft(\bp_0 a\mright) \mleft(\bp_0 b\mright) {\downarrow}$ for any $b \in \mathbb{N}$; and $C \mleft(\bp_0 a\mright) \mleft(\bp_0 b\mright) = 1$ implies that $\bp_0 b = b_0 \in \mathcal{O}$, in which case $C \mleft(\bp_0 b\mright) \mleft(\bp_0 a\mright)$ and $Rb$ will also converge, hence they will be safe to evaluate when guarded behind the lazy $\bdN^* \mleft(C \mleft(\bp_0 a\mright) \mleft(\bp_0 b\mright)\mright) 1$. Thus, as long as $a \in \mathcal{H}$, we have $C^* ab {\downarrow}$ for any $b \in \mathbb{N}$, and it is easy to observe that $C^* ab = 1$ implies $b \in \mathcal{H}$.

\begin{proposition}[$\mathrm{ATR}_0$]
  \label{prop:hyp-ind-cmp-wo}
  The relation $\prec$ on $\mathcal{H}$ given by
  \[b \prec a \defeqv C^* ab = 1\]
  is a well-ordering.
\end{proposition}

\begin{proof}
  $\prec$ is a linear ordering simply because any two elements in $\mathcal{O}$ are comparable, and $b \in \mathcal{H}$ ensures that $Rb = 1$.

  Suppose $\prec$ is not a well-ordering, i.e.\ we have a strictly decreasing function $f : \mathbb{N} \rightarrow \mathcal{H}$. Then $f_0 : \mathbb{N} \rightarrow \mathcal{O}$ given by $f_0\mleft(n\mright) = f\mleft(n\mright)_0$ satisfies
  \[\otyp\mleft(f_0\mleft(0\mright)\mright) \geq \otyp\mleft(f_0\mleft(1\mright)\mright) \geq \cdots.\]
  We have $\otyp\mleft(f_0\mleft(0\mright)\mright) \geq \otyp\mleft(f_0\mleft(n\mright)\mright)$ for any $n \in \mathbb{N}$, so there must be a unique $\gamma_n \leq_\mathcal{O} f_0\mleft(0\mright)$ such that $\otyp\mleft(f_0\mleft(n\mright)\mright) = \otyp\mleft(\gamma_n\mright)$. Now, fix $\Omega \subseteq \mathbb{N}$ such that $H\mleft(f_0\mleft(0\mright), \Omega\mright)$, then all order type statements concerning $f_0\mleft(n\mright)$ will be arithmetic relative to $\Omega$. Thus, by arithmetic comprehension, there exists a function $f' : n \mapsto \gamma_n$ and we immediately have $f_0\mleft(0\mright) = f'\mleft(0\mright) \geq_\mathcal{O} f'\mleft(1\mright) \geq_\mathcal{O} \cdots$. Since $f_0\mleft(0\mright) \in \mathcal{O}$, $f'$ must stabilise after some $N \in \mathbb{N}$, implying that $\otyp\mleft(f_0\mleft(m\mright)\mright) = \otyp\mleft(f_0\mleft(n\mright)\mright)$ for any $m, n > N$.

  Then, for any $m > N$, we will have $f\mleft(m\mright) \succ f\mleft(m + 1\mright)$ yet $\otyp\mleft(f\mleft(m\mright)_0\mright) = \otyp\mleft(f\mleft(m + 1\mright)_0\mright)$, thus we must have $f\mleft(m\mright) > f\mleft(m + 1\mright)$ by the definition of $C^*$. Thus, $f$ restricted to $\left\{m : m > N\right\}$ is a strictly decreasing sequence of natural numbers. This contradicts $\Delta^0_0$-induction. Therefore, such $f$ cannot exist and $\prec$ is indeed a well-ordering.
\end{proof}

Now, move to the theory $\ATR$ (with unrestricted induction on the natural numbers, which is only needed for the induction and recursion schemes of $\CM$) and note that this is an extension of $\Sigma^1_1\AC$ (cf.\ \cite{simpson09-soa}*{Theorem V.8.3}). We reuse the realisability conditions in \autoref{sec:oa-cm}, but require that the interpretation $q$ for $\prec$ is chosen as
\[q = \oplam{ef} C^* \mleft(Hf\mright) \mleft(He\mright).\]
\autoref{thm:conv-s2-to-hyp} guarantees that $Hd = He \leftrightarrow d \equiv e$, and ensures that $q$ satisfies
\[\opA{e, e', f, f' \in S_2} \left(e \equiv f \land e' \equiv f' \rightarrow qe e' = qf f'\right).\]

The axioms \ref{itm:cm-gwo-ord-axiom} stating that $\prec$ is a linear ordering are simply propositional compositions of atomic formulae of the form $X \prec Y$ and $X = Y$. Since their truths are computable\footnote{To decide whether $e = f$ should be realised, i.e.\ whether $e \equiv f$ for $e, f \in S_2$, we can either observe that
  \[\opA{X} \opA{Y} \left(X = Y \lor \neg X = Y\right)\] follows from numerical omniscience and use the mechanisms in \autoref{lem:cm-num-omni-real}, or just conveniently check whether $He = Hf$.} in $\Hyp$ for any assignments from $S_2$ and $q$ indeed computes a well-ordering on $S_2$ (up to extensional equivalence) by \autoref{prop:hyp-ind-cmp-wo}, these axioms are realised by any realiser that is total on the appropriate inputs. \ref{itm:cm-gwo-init-seg-axiom} is another easy case:
\begin{lemma}
  The axiom of countable initial segments \ref{itm:cm-gwo-init-seg-axiom} is realised:
  \[\opA{X} \opE{Z} \opA{Y} \left(Y \prec X \rightarrow \opE{n} Y = \left(Z\right)_n\right).\]
\end{lemma}

\begin{proof}
  Consider any $x \in S_2$, then for any $y \in \mathbb{N}$, $C^* \mleft(Hx\mright) y {\downarrow}$ and it evaluates to $1$ if and only if $y \in \mathcal{H}$ and $y \prec Hx$ as defined in \autoref{prop:hyp-ind-cmp-wo}. Thus, we simply need some $Z$ containing the second-order objects $H^{-1} y \in S_2$ for all $y$ such that $C^* \mleft(Hx\mright) y = 1$ (up to extensional equivalence).

  Formally, we shall define
  \[z = \oplam{n} \bdN^* \mleft(C^* \mleft(Hx\mright) \mleft(\bp_0 n\mright)\mright) 1 \mleft(H^{-1} \mleft(\bp_0 n\mright) \mleft(\bp_1 n\mright)\mright) 0.\]
  For any $y \in S_2$ such that $qyx = 1$, we see that
  \[z \mleft(\bp \mleft(Hy\mright) n\mright) = H^{-1} \mleft(Hy\mright) n = yn\]
  by extensional equivalence between $y$ and $H^{-1} \mleft(Hy\mright)$, thus it is easy to check that $\bp \mleft(Hy\mright) \mleft(\oplam{m} \mleft(\bp \mleft(\bp \mleft(Hy\mright) m\mright) d\mright)\mright)$ realises $\opE{n} y = \left(z\right)_{n}$ interpreted as
  \[\opE{n} \opA{m} \opE{k} \left(k = \tuple{n, m} \land \left(m \in_1 y \leftrightarrow k \in_1 z\right)\right),\]
  where $d$ is the uninteresting realiser for the entire quantifier-free subformula inside.
\end{proof}

Now, $\ATR$ does not suffice to realise the transfinite induction scheme \ref{itm:cm-gwo-ind-axiom} for formula $\varphi$ of arbitrary complexity, since the base theory cannot support transfinite induction on well-orderings for any second-order formulae (cf.\ \autoref{ftn:cm-real-complex} on page \pageref{ftn:cm-real-complex} for the quantifier complexity of the realisability conditions). For that purpose, we require the much stronger system $\BI$\inda{bar induction, $\BI$} --- also known as $\Pi^1_\infty\TI_0$ in Simpson \cite{simpson09-soa} --- which adds to $\ACA_0$ the axiom scheme of \emph{bar induction},
\[\opA{X} \opA{<_X} \left(\WO\mleft(X, {<_X}\mright) \rightarrow \opA{x \in X} \left(\opA{y <_X x} \ \varphi\mleft(y\mright) \rightarrow \varphi\mleft(x\mright)\right) \rightarrow \opA{x \in X} \varphi\mleft(x\mright)\right)\]
for any formula $\varphi\mleft(x\mright)$, where $\WO\mleft(X, {<_X}\mright)$ denotes the $\Pi^1_1$-formula asserting that $<_X$ is a linear ordering on $X$ and there exists no strictly decreasing function $f : \mathbb{N} \rightarrow X$.

It does not matter how much induction on the natural numbers we have inside $\BI$, since $\Delta^0_0$-induction suffices to verify $\WO\mleft(\mathbb{N}, <\mright)$, and thus bar induction already implies unrestricted induction on $\mathbb{N}$. Additionally, by \cite{simpson09-soa}*{Corollary VII.2.19}, $\BI$ implies $\ATR$, so the existing axioms of $\CM$, \ref{itm:cm-gwo-ord-axiom} and \ref{itm:cm-gwo-init-seg-axiom} are already realised in $\BI$.

In order to realise \ref{itm:cm-gwo-ind-axiom}, we need the following lemma:
\begin{lemma}[$\Sigma^1_1\AC$]
  \label{lem:cm-same-real-equiv-form}
  For any $d \in \Hyp$ and $e, f \in S_2$ such that $e \equiv f$, if $d \Vdash \varphi\mleft[e/X\mright]$ for some formula $\varphi\mleft(X\mright)$ (possibly with other parameters), then also $d \Vdash \varphi\mleft[f/X\mright]$.

  Specifically, assuming $\ATR$, $d \Vdash \varphi\mleft[e/X\mright]$ if and only if $d \Vdash \varphi\mleft[H^{-1} \mleft(He\mright)/X\mright]$.
\end{lemma}

\begin{proof}
  This simply amounts to the fact that when $\varphi$ is atomic, the realiser $d$ is irrelevant by our stipulations and the statement is equivalent to saying that $\varphi\mleft[e/X\mright]$ is realised if and only if $\varphi\mleft[f/X\mright]$ is. Thus, this holds for any formula $\varphi\mleft(X\mright)$ by induction on the complexity of $\varphi$.
\end{proof}

\begin{proposition}[$\BI$]
  \label{prop:cm-gwo-ind-real}
  The axiom of unrestricted transfinite induction \ref{itm:cm-gwo-ind-axiom} is realised for any formula $\varphi\mleft(X\mright)$:
  \[\opA{X} \left(\opA{Y} \left(Y \prec X \rightarrow \varphi\mleft(Y\mright)\right) \rightarrow \varphi\mleft(X\mright)\right) \rightarrow \opA{X} \varphi\mleft(X\mright).\]
\end{proposition}

\begin{proof}
  We shall define
  \[h = \oplam{f} \br \mleft(\oplam{ag} f \mleft(H^{-1} a\mright) \mleft(\oplam{zd} g \mleft(Hz\mright)\mright)\mright),\]
  and we claim that $\oplam{\vec{x} fy} hf \mleft(Hy\mright)$ realises the axiom, where $\vec{x}$ enumerate the free variables in it.

  Fix some $f \Vdash \opA{X} \left(\opA{Y} \left(Y \prec X \rightarrow \varphi\mleft(Y\mright)\right) \rightarrow \varphi\mleft(X\mright)\right)\mleft[\vec{u}/\vec{x}\mright]$ and any $y \in S_2$. We consider the condition
  \[\psi\mleft(a\mright) \defeqv a \in \mathcal{H} \land a \preceq Hy \land a = H \mleft(H^{-1} a\mright).\]
  Note that $a \prec Hy$ if and only if $C^* \mleft(Hy\mright) a = 1$, which already implies $a \in \mathcal{H}$ and that $H \mleft(H^{-1} a\mright) {\downarrow}$, thus $\psi\mleft(a\mright)$ is equivalent to a $\Sigma^1_1$-formula. On the other hand, given $H \mleft(H^{-1} a\mright) {\downarrow}$, we have $a \neq H \mleft(H^{-1} a\mright)$ if and only if $\opE{b \neq a} H \mleft(H^{-1} a\mright) = b$. So $\neg \psi\mleft(a\mright)$ is equivalent to
  \[\left(a \neq Hy \land C^* \mleft(Hy\mright) a = 0\right) \lor \left(C^* \mleft(Hy\mright) a = 1 \land \opE{b \neq a} H \mleft(H^{-1} a\mright) = b\right),\]
  which is also $\Sigma^1_1$. Therefore, $\psi\mleft(a\mright)$ is $\Delta^1_1$ and by $\Delta^1_1$-comprehension (which is a consequence of $\ATR_0$, cf.\ \cite{simpson09-soa}*{Theorem V.5.1}, and thus holds in $\BI$), we can define the set
  \[W = \left\{a : \psi\mleft(a\mright)\right\} = \left\{a \in \mathcal{H} : a \preceq Hy, a = H \mleft(H^{-1} a\mright)\right\}\]
  Observe additionally that for any $z \in S_2$, $H \mleft(H^{-1} \mleft(Hz\mright)\mright) = Hz$ by \autoref{thm:conv-s2-to-hyp} since $H^{-1} \mleft(Hz\mright) \equiv z$, thus $Hz \in W$ as long as $Hz \preceq Hy$.

  By \autoref{prop:hyp-ind-cmp-wo}, the set $W$ is well-ordered by $\prec$ and we shall prove by bar induction that for any $a \in W$,
  \[hfa {\downarrow} \Vdash \varphi\mleft[H^{-1} a/X, \vec{u}/\vec{x}\mright].\]
  Consider any $z \in S_2$ such that $qz \mleft(H^{-1} a\mright) = 1$, i.e.\ $Hz \prec H \mleft(H^{-1} a\mright) = a$. By the inductive hypothesis, we have
  \[hf \mleft(Hz\mright) {\downarrow} \Vdash \varphi\mleft[H^{-1} \mleft(Hz\mright)/X, \vec{u}/\vec{x}\mright],\]
  which also realises $\varphi\mleft[z/X, \vec{u}/\vec{x}\mright]$ by \autoref{lem:cm-same-real-equiv-form}. This means that
  \[\oplam{zd} hf \mleft(Hz\mright) \Vdash \opA{Y} \left(Y \prec X \rightarrow \varphi\mleft(Y\mright)\right)\mleft[H^{-1} a/X, \vec{u}/\vec{x}\mright].\]
  It follows that
  \[f \mleft(H^{-1} a\mright) \mleft(\oplam{zd} hf \mleft(Hz\mright)\mright) {\downarrow} \Vdash \varphi\mleft[H^{-1} a/X, \vec{u}/\vec{x}\mright]\]
  and thus so does $hfa {\downarrow}$ by the recursion theorem. Therefore, by bar induction, it follows that
  \[hf \mleft(Hy\mright) {\downarrow} \Vdash \varphi\mleft[H^{-1} \mleft(Hy\mright)/X, \vec{u}/\vec{x}\mright],\]
  thus the same realiser also realises $\varphi\mleft[y/X, \vec{u}/\vec{x}\mright]$ by \autoref{lem:cm-same-real-equiv-form}.
\end{proof}

It follows that all the axioms of $\CM + \GWO$ are interpreted in the given realisability model in $\BI$, hence immediately:
\begin{theorem}
  Let $\varphi$ be any theorem of $\CM + \GWO$, then $\BI$ proves there exists $d \in \Hyp$ such that $d \Vdash \varphi$.
\end{theorem}

\subsection*{An ordinal analysis of $\CM + \GWO$}

To give a lower bound for the proof-theoretic strength of $\CM + \GWO$, we say that a formula $\varphi\mleft(X, x\mright)$ in the language of $\mathrm{CM}$ is a \emph{strictly positive first-order formula}, possibly with other free (first-order) variables, if it is built up according to the following rules:
\begin{enumerate}[label=\uplabel{\arabic*}]
  \item $\varphi$ can contain any atomic formulae in the language of $\PA$, as well as the atomic formula $t \in_1 X$ where $t$ is an arbitrary first-order term;
  \item $\varphi$ can contain the propositional constant `bottom', $\bot$, which may be interpreted as $\neg 0 = 1$;
  \item $\varphi$ can contain connectives $\land$, $\lor$, $\rightarrow$, with the restriction that $X$ cannot occur on the left of the connective $\rightarrow$;
  \item $\varphi$ can contain quantifiers $\opA{y}$ and $\opE{y}$ where $y$ is any first-order variable.
\end{enumerate}

Let $\theta\mleft(x\mright)$ be any formula, and we let $\varphi\mleft(\theta, x\mright)$ denote the result of substituting every occurrence of $t \in_1 X$ inside $\varphi\mleft(X, x\mright)$ by $\theta\mleft(t\mright)$, with appropriate renaming so that any other free variable in $\theta\mleft(x\mright)$ is not bound by any quantifiers in $\varphi$.

We shall prove the following:
\begin{theorem}[Inductive definitions]
  \label{thm:cm-ind-def}
  For any strictly positive first-order formula $\varphi\mleft(X, x\mright)$, there exists $\theta_\varphi\mleft(x\mright)$ such that
  \[\CM + \GWO \vdash \opA{x} \left(\varphi\mleft(\theta_\varphi, x\mright) \rightarrow \theta_\varphi\mleft(x\mright)\right),\]
  and also for any formula $\eta\mleft(x\mright)$
  \[\CM + \GWO \vdash \opA{x} \left(\varphi\mleft(\eta, x\mright) \rightarrow \eta\mleft(x\mright)\right) \rightarrow \opA{x} \left(\theta_\varphi\mleft(x\mright) \rightarrow \eta\mleft(x\mright)\right).\]
\end{theorem}

This will involve a few lemmata:
\begin{lemma}[Positivity]
  \label{lem:cm-ind-pos}
  For any formulae $\theta\mleft(x\mright), \eta\mleft(x\mright)$,
  \[\CM \vdash \opA{x} \left(\theta\mleft(x\mright) \rightarrow \eta\mleft(x\mright)\right) \rightarrow \opA{x} \left(\varphi\mleft(\theta, x\mright) \rightarrow \varphi\mleft(\eta, x\mright)\right).\]
  This trivially implies $X \subseteq Y \rightarrow \opA{x} \left(\varphi\mleft(X, x\mright) \rightarrow \varphi\mleft(Y, x\mright)\right)$ as well.
\end{lemma}

\begin{proof}
  This follows from an easy induction on the complexity of $\varphi\mleft(X, x\mright)$.
\end{proof}

\begin{lemma}
  \label{lem:cm-ind-set-witnesses}
  For any formula $\theta\mleft(x\mright)$,
  \[\CM \vdash \opA{x} \left(\varphi\mleft(\theta, x\mright) \rightarrow \opE{X} \left(\opA{z \in_1 X} \theta\mleft(z\mright) \land \varphi\mleft(X, x\mright)\right)\right)\]
\end{lemma}

\begin{proof}
  We show this by induction on the complexity of $\varphi\mleft(X, x\mright)$:

  Firstly, if $\varphi\mleft(X, x\mright)$ is an atomic formula not of the form $t \in_1 X$, then $\varphi\mleft(\theta, x\mright) \rightarrow \varphi\mleft(\varnothing, x\mright)$ trivially; if $\varphi\mleft(X, x\mright)$ is $t \in_1 X$, then $\varphi\mleft(\theta, x\mright) \rightarrow \opA{z \in_1 \left\{t\right\}} \theta\mleft(z\mright) \land \varphi\mleft(\left\{t\right\}, x\mright)$, where $t$ is any first-order term containing the variable $x$.

  The case where $\varphi$ is $\bot$ is vacuously true.

  Now, suppose that $\varphi$ is of the form $\psi \land \chi$ and $\varphi\mleft(\theta, x\mright)$ holds, then by the inductive hypothesis we have $X$ and $Y$ such that
  \[\opA{z \in_1 X \cup Y} \theta\mleft(z\mright) \land \psi\mleft(X, x\mright) \land \chi\mleft(Y, x\mright).\]
  It follows from \autoref{lem:cm-ind-pos} that we also have $\psi\mleft(X \cup Y, x\mright) \land \chi\mleft(X \cup Y, x\mright)$.

  Suppose $\varphi$ is of the form $\psi \rightarrow \chi$ and $\varphi\mleft(\theta, x\mright)$ holds, we know from stipulation that $\psi\mleft(x\mright)$ must not contain $X$. By the inductive hypothesis, we have
  \[\psi\mleft(x\mright) \rightarrow \opE{X} \left(\opA{z \in_1 X} \theta\mleft(z\mright) \land \chi\mleft(X, x\mright)\right).\] Since $\psi\mleft(x\mright)$ is arithmetic, hence decidable, we can take $Y$ to be $X$ above if $\psi\mleft(x\mright)$ holds and $\varnothing$ otherwise. Then we always have $\opA{z \in_1 Y} \theta\mleft(z\mright)$, and we also have $\psi\mleft(x\mright) \rightarrow \chi\mleft(Y, x\mright)$.

  Suppose $\varphi$ is of the form $\opA{y} \psi$ and $\varphi\mleft(\theta, x\mright)$ holds, then by the inductive hypothesis
  \[\opA{y} \opE{X} \left(\opA{z \in_1 X} \theta\mleft(z\mright) \land \psi\mleft(X, x, y\mright)\right).\]
  By the recursion axiom of $\CM$, we can find $Z$ such that
  \[\opA{y} \left(\opA{z \in_1 \left(Z\right)_y} \theta\mleft(z\mright) \land \psi\mleft(\left(Z\right)_y, x, y\mright)\right).\]
  Take $Y = \bigcup_y \left(Z\right)_y$, then $\opA{z \in_1 Y} \theta\mleft(z\mright)$ and also $\opA{y} \psi\mleft(Y, x, y\mright)$ by \autoref{lem:cm-ind-pos}.

  Finally, the cases where $\varphi$ is of the form $\psi \lor \chi$ or $\opE{y} \psi$ are trivial.
\end{proof}

\begin{proof}[Proof of \autoref{thm:cm-ind-def}]
  Let a formula $\psi\mleft(X, Z, S\mright)$ assert that both of the following hold:
  \begin{enumerate}[label=\uplabel{\roman*}]
    \item $\opA{Y \preceq X} \opE{n} Y = \left(Z\right)_n$;
    \item \label{itm:cm-ind-def-iter} $\opA{n} \left(\left(Z\right)_n \preceq X \rightarrow \opA{x} \left(x \in_1 \left(S\right)_n \leftrightarrow \varphi\mleft(\bigcup \left\{\left(S\right)_m : \left(Z\right)_m \prec \left(Z\right)_n\right\}, x\mright)\right)\right)$.
  \end{enumerate}
  Observe that if $\left(S\right)_m$ is given for all $m$ such that $\left(Z\right)_m \prec \left(Z\right)_n$, then the unique $\left(S\right)_n$ can be constructed by decidable comprehension as $\varphi$ is arithmetic. Thus, by transfinite induction on all second-order objects, it is easy to show that for each $X$, a pair $Z, S$ satisfying $\psi\mleft(X, Z, S\mright)$ exists and is unique up to matched permutations, i.e.\ for any $n, m$,
  \[\psi\mleft(X, Z, S\mright) \land \psi\mleft(X, Z', S'\mright) \land \left(Z\right)_n \preceq X \land \left(Z\right)_n = \left(Z'\right)_m \rightarrow \left(S\right)_n = \left(S'\right)_m.\]

  Now, we denote
  \[\theta_\varphi\mleft(x\mright) \defeqv \opE{X, Z, S} \left(\psi\mleft(X, Z, S\mright) \land \opE{n} \left(X = \left(Z\right)_n \land x \in_1 \left(S\right)_n\right)\right).\]
  We will verify that $\mathrm{CM} + \mathrm{GWO}$ proves the desired properties of $\theta_\varphi\mleft(x\mright)$: firstly, suppose that $\varphi\mleft(\theta_\varphi, x\mright)$ holds for some $x$. By \autoref{lem:cm-ind-set-witnesses}, we can find $A$ satisfying
  \[\opA{x \in_1 A} \theta_\varphi\mleft(x\mright) \land \varphi\mleft(A, x\mright).\]
  Using the recursion axiom, the former conjunct will imply the existence of families $X, Z, S$ such that
  \[\opA{x \in_1 A} \left(\psi\mleft(\left(X\right)_x, \left(Z\right)_x, \left(S\right)_x\mright) \land \opE{n} \left(\left(X\right)_x = \left(\left(Z\right)_x\right)_n \land x \in_1 \left(\left(S\right)_x\right)_n\right)\right).\]
  By \autoref{lem:cm-gwo-countable-family-bound}, we can find some $Y$ such that $\opA{x \in_1 A} \left(X\right)_x \prec Y$. Let $Z^*, S^*$ be such that $\psi\mleft(Y, Z^*, S^*\mright)$ holds, then it follows that
  \[A \subseteq \bigcup \left\{\left(S^*\right)_m : \left(Z^*\right)_m \prec Y\right\}.\]
  Now $\varphi\mleft(A, x\mright)$ and \autoref{lem:cm-ind-pos} will imply that $x \in_1 \left(S^*\right)_n$ where $\left(Z^*\right)_n = Y$, and thus $\theta_\varphi\mleft(x\mright)$.

  Secondly, suppose that for some formula $\eta\mleft(x\mright)$, $\opA{x} \left(\varphi\mleft(\eta, x\mright) \rightarrow \eta\mleft(x\mright)\right)$ holds. In order to have $\opA{x} \left(\theta_\varphi\mleft(x\mright) \rightarrow \eta\mleft(x\mright)\right)$, it suffices to fix $X, Z, S$ satisfying $\psi\mleft(X, Z, S\mright)$ and show by transfinite induction on $Y$ that
  \[\opA{Y} \left(Y \preceq X \rightarrow \opA{n} \left(Y = \left(Z\right)_n \rightarrow \opA{x \in_1 \left(S\right)_n} \eta\mleft(x\mright)\right)\right).\]
  Here the inductive hypothesis implies
  \[\opA{x \in_1 \bigcup \left\{\left(S\right)_m : \left(Z\right)_m \prec Y\right\}} \eta\mleft(x\mright),\]
  thus for any $x \in_1 \left(S\right)_n$, i.e.\ any $x$ satisfying $\varphi\mleft(\bigcup \left\{\left(S\right)_m : \left(Z\right)_m \prec Y\right\}, x\mright)$, \autoref{lem:cm-ind-pos} would imply $\varphi\mleft(\eta, x\mright)$ and thus $\eta\mleft(x\mright)$ as required. It follows by transfinite induction that $X = \left(Z\right)_n$ implies $\opA{x \in_1 \left(S\right)_n} \eta\mleft(x\mright)$, and thus $\theta_\varphi\mleft(x\mright)$ implies $\eta\mleft(x\mright)$.
\end{proof}

What we have shown is the existence of first-order inductive definitions\inda{inductive definition} in $\CM + \GWO$. Since Heyting arithmetic $\HA$ is trivially a subtheory of $\CM$, \autoref{thm:cm-ind-def} implies that $\CM + \GWO$ interprets the intuitionistic\footnote{Even though $\CM \vdash \PA$, we can only interpret the intuitionistic part here because we are translating the first-order predicate symbol $I_\varphi$ for the least fixed point of $\varphi$ into a non-decidable formula $\theta_\varphi\mleft(x\mright)$.} first-order theory of inductive definitions $\ID_1^i$.

The proof-theoretic ordinal of $\mathrm{ID}_1^i$ is computed in \cite{buchholz-pohlers78-iterated-id} as $\left|\ID_1^i\right| = \vartheta_{\varepsilon_{\Omega + 1}}\mleft(0\mright) = \left|\BI\right|$, the Bachmann--Howard ordinal\inda{Bachmann--Howard ordinal}. Here $\Omega$ is treated as a formal name for the first uncountable ordinal $\omega_1$, and $\vartheta_\alpha$ refers to any of the variants of ordinal collapse functions\footnote{These notations have minor differences, but they catch up at key positions like the Bachmann--Howard ordinal. Their differences are too intricate to be relevant throughout this thesis. For more detail, see discussions in \cite{pohlers93-oa}*{section 7} or the closing remarks in \cite{rathjen93-kp-fragments}.} in Buchholz \cite{buchholz75-normalfunktionen}, Sch\"utte \cite{schutte77-proof-theory} or Pohlers \cite{pohlers93-oa}. Since $\CM + \GWO$ can be realised in $\BI$, it is an immediate corollary that:
\begin{corollary}
  $\CM + \GWO \equiv_{\Con} \BI$. In particular, $\left|\CM + \GWO\right| = \left|\BI\right| = \vartheta_{\varepsilon_{\Omega + 1}}\mleft(0\mright)$.
\end{corollary}

\section{Indispensability of countable initial segments}

In this section, we use the realisability model above to showcase one key distinction between the global well-ordering axioms in the semi-intuitionistic system $\CM$ and the more traditional cases of choice principles in classical systems: that the axiom of countable initial segments \ref{itm:cm-gwo-init-seg-axiom} is indispensable for the proof-theoretic strength of the system.

The readers may contrast this with, for example, classical Kripke--Platek set theory\footnote{In this thesis, we shall use $\KP$ to denote the version of Kripke--Platek set theory that contains the axiom of infinity, which is sometimes referred to as $\KP\omega$ in literature. The hereditarily finite fragment satisfying $\KP$ without infinity is too weak, in comparison to other theories assessed in this thesis, to be relevant here.} $\KP$, where G\"odel's constructible universe $L$ is also a model of $\KP$ satisfying `there exists a global well-ordering'. Let $\omega_1^L$ denote what $L$ believes to be the first uncountable ordinal --- if it exists, or the class of all (necessarily countable) ordinals otherwise --- then it is easy to verify
\[\KP \vdash \left(\Delta_0\text{-collection}\right)^{L_{\omega_1^L}}.\]
Additionally, by G\"odel's condensation lemma, for any $\alpha \in \omega_1^L$, any bijection $f : \alpha \rightarrow \omega$ in $L$ must also lie in $L_{\omega_1^L}$, as $f$ (as a set) is hereditarily countable. Thus $L_{\omega_1^L}$ is a model of $\KP$ satisfying the refined claim `there exists a global well-ordering whose every initial segment is countable'. This shows that adding the requirement of countable initial segments does not increase the proof-theoretic strength of the well-ordering assertions over $\KP$.

This truncation cannot be done in $\CM$ however, since with a lack of the law of excluded middle, we can no longer classify the initial segments of the global well-ordering as either a countable one or an extension of all countable initial segments. Moreover, the fact that bounded second-order quantifiers, i.e.\ $\opA{Y \prec X}$ and $\opE{Y \prec X}$, are decidable, which is crucial for the applications of the well-ordering axioms we considered in \autoref{sec:cm-gwo-axioms}, relies critically on the countable initial segment assumption.

Here, we shall support these informal observations by demonstrating that the full system $\CM + \GWO$ is indeed strictly stronger than the fragment $\CM + \text{\ref{itm:cm-gwo-ord-axiom}} + \text{\ref{itm:cm-gwo-ind-axiom}}$ without the axiom of countable initial segments. Let us work in $\ATR$ and again consider the realisability conditions given in \autoref{sec:oa-cm}, where the interpretation of the well-ordering $\prec$ is now chosen as
\[q' = \oplam{ef} \left\{\begin{aligned}
     & 1 &  & \text{if $He < Hf$}, \\
     & 0 &  & \text{otherwise}.
  \end{aligned}\right.\]
Recall from \autoref{thm:conv-s2-to-hyp} that for $e, f \in S_2$, $He = Hf$ if and only if $e \equiv f$, thus $q'$ satisfies the extensionality requirement that
\[\opA{e, e', f, f' \in S_2} \left(e \equiv f \land e' \equiv f' \rightarrow q' e e' = q' f f'\right).\]

Once again, we see that the axioms \ref{itm:cm-gwo-ord-axiom} stating that $\prec$ is a linear ordering are interpreted as propositional compositions of numerical inequalities in the base theory, thus realised by any realiser that is total on the appropriate inputs. We also have:

\begin{proposition}[$\ATR$]
  Using the interpretation $q'$, the axiom of unrestricted transfinite induction \ref{itm:cm-gwo-ind-axiom} is realised:
  \[\opA{X} \left(\opA{Y} \left(Y \prec X \rightarrow \varphi\mleft(Y\mright)\right) \rightarrow \varphi\mleft(X\mright)\right) \rightarrow \opA{X} \varphi\mleft(X\mright).\]
\end{proposition}

\begin{proof}
  Let $\vec{x}$ enumerate the free variables in the axiom and consider some $f \Vdash \opA{X} \left(\opA{Y} \left(Y \prec X \rightarrow \varphi\mleft(Y\mright)\right) \rightarrow \varphi\mleft(X\mright)\right)\mleft[\vec{u}/\vec{x}\mright]$. By recursion, let
  \[h = \oplam{f} \br \mleft(\oplam{ag} f \mleft(H^{-1} a\mright) \mleft(\oplam{ye} g \mleft(Hy\mright)\mright)\mright).\]
  We show by induction on $\mathbb{N}$ that for any $a \in \mathcal{H}$ satisfying $a = H \mleft(H^{-1} a\mright)$,
  \[hfa {\downarrow} \Vdash \varphi\mleft(X\mright)\mleft[H^{-1} a/X, \vec{u}/\vec{x}\mright].\]
  This is because the inductive hypothesis and \autoref{lem:cm-same-real-equiv-form} together imply that for any $y \in S_2$ such that $Hy < H \mleft(H^{-1} a\mright) = a$ (where we always have $H \mleft(H^{-1} \mleft(Hy\mright)\mright) = Hy$), $hf \mleft(Hy\mright) {\downarrow} \Vdash \varphi\mleft(X\mright)\mleft[y/X, \vec{u}/\vec{x}\mright]$, thus the (recursive) invocation of the realiser $f$ would yield the desired result.

  Therefore, $\oplam{\vec{x} fy} hf \mleft(Hy\mright)$ realises the axiom, by \autoref{lem:cm-same-real-equiv-form} again.
\end{proof}

It follows that every theorem of $\CM + \text{\ref{itm:cm-gwo-ord-axiom}} + \text{\ref{itm:cm-gwo-ind-axiom}}$ is realised in $\ATR$ in this set-up. As argued in \cite{jager-strahm00-fixed-point}, the proof-theoretic strength of $\ATR$ is $\Gamma_{\varepsilon_0}$, much weaker than the full theory $\CM + \GWO$ at the Bachmann--Howard ordinal.

As a closing remark, we shall comment that, in fact, the only condition necessitating the use of hyperarithmetic indices and thus the developments in $\ATR_0$ here is the fact that second-order equality in the language of $\CM$ is extensional, thus the ordering $\prec$ must respect extensional equivalences, i.e.\ satisfy axiom \ref{itm:cm-gwo-ext-axiom} from our definition of $\CM$ in \autoref{sec:cm-axioms}, in order to be a linear ordering.

Observe that, without extensionality requirements, the following much simpler ordering would suffice:
\[q^\dagger = \oplam{ef} \left\{\begin{aligned}
     & 1 &  & \text{if $e < f$}, \\
     & 0 &  & \text{otherwise},
  \end{aligned}\right.\]
and the entire analysis could be completed in the base theory $\Sigma^1_1\AC$.

We should point out that a lower bound for the proof-theoretic strength of $\CM + \text{\ref{itm:cm-gwo-ord-axiom}} + \text{\ref{itm:cm-gwo-ind-axiom}}$ still remains an open problem, and I suspect that the addition of extensional assertions alone does not cause a concrete increase in strength from $\Sigma^1_1\AC$ to $\ATR$. Rather, the fact that $\ATR$ is needed here is because only in the base theory $\ATR_0$ would we have sufficiently robust infrastructure to work with the weak pca $\Hyp$. We will see more evidence and discussion for this need of a strong base theory in the next section, where we consider other weak fragments of $\CM + \GWO$.

\section{Predicative fragments of \texorpdfstring{$\GWO$}{GWO}}
\label{sec:cm-gwo-pred}

Now, as briefly mentioned in \autoref{sec:cm-gwo-axioms}, Weaver believed that a \emph{predicative} extension should never admit the full transfinite induction scheme \ref{itm:cm-gwo-ind-axiom} for arbitrary third-order formulae. In his own words in \cite{weaver09-cm},

\begin{quote}
  [transfinite induction] should only be asserted for formulas that do not contain $\prec$, for reasons having to do with the circularity involved in making sense of a relation that is well-ordered with respect to properties that are defined in terms of that relation.
\end{quote}

Indeed, our ordinal analysis of full $\CM + \GWO$ suggests that Weaver's remark aligns well with his philosophical stances, since he is known to criticise many other theories with proof-theoretic strength at the Bachmann--Howard ordinal as impredicative, such as $\BI$, $\ID_1$ or the set theories $\KP$ and $\CZF$ --- the latter two for their unrestricted set induction schemes, precisely (cf.\ \cite{weaver18-predicative-well-ordering}).

However, while the rejection of impredicative constructions is understandable from his philosophical viewpoint, Weaver's suggestion quoted here fails to resolve the problem. If we read it in the usual sense, i.e.\ we forbid the primitive relation symbol $\prec$ from occurring in $\varphi$ but still allow the occurrences of any other free parameters, then given any formula $\varphi\mleft(X\mright)$, one can simply apply the transfinite induction scheme to the formula $\tilde{\varphi}\mleft(X\mright)$ instead, where every occurrence of $X \prec Y$ has been replaced by $\tuple{X, Y} \in \mathbf{W}$. Here $\mathbf{W}$ is an unused, free third-order parameter, and $\tuple{X, Y}$ is a fixed encoding of an ordered pair of second-order objects as a single second-order object, for example, as
\[\tuple{X, Y} = \left\{\tuple{0, n} : n \in_1 X\right\} \cup \left\{\tuple{1, m} : m \in_1 Y\right\}.\]
By decidable comprehension, we can simply take $\mathbf{W} = \left\{\tuple{X, Y} : X \prec Y\right\}$ later to derive transfinite induction for the original formula $\varphi$. Therefore, such superficial constraint on $\varphi$ does not actually restrict the theory in any way.

If we are to read this proposal in a harsher way, by either forbidding the occurrence of any free (third-order) parameter in $\varphi$ at all, or maintaining that the axiom scheme of decidable comprehension is postulated for the restricted language of $\mathrm{CM}$ only and cannot be applied to formulae containing the new symbol $\prec$, then the utility of the extension will be severely impaired.

For example, as we have explored in \autoref{sec:cm-gwo-axioms}, the proof of traditional mathematical results like the existence of bases using the global well-ordering axioms will heavily involve recursive constructions along the well-ordering and then applications of the induction axiom to prove properties of such constructions. It is unclear how any of those results, including the general version of the Hahn--Banach theorem that Weaver himself claimed to have proved, can be established without accessing the initial segment, i.e.\ second-order sets $Y \prec X$ for a given $X$, in a meaningful way in the induction axiom.

Moving away from Weaver's original proposal, we would still like to find interesting weak fragments lying between $\CM$ and $\CM + \GWO$ such that both
\begin{enumerate}[label=\uplabel{\arabic*}]
  \item it has significantly weaker strength than full $\CM + \GWO$; and
  \item it can still prove most of the results surrounding Zorn's lemma that we examined in \autoref{sec:cm-gwo-axioms}.
\end{enumerate}
We shall look at some candidates along this direction.

Observe that for any decidable formula $\varphi\mleft(X\mright)$, there is some third-order object $\mathbf{X}$ such that
\[\opA{X} \left(\varphi\mleft(X\mright) \leftrightarrow X \in_2 \mathbf{X}\right)\]
by decidable comprehension, i.e.\ every decidable formula is equivalent to an atomic formula of this specific form. The class of decidable formulae is very robust since they are closed under arithmetic quantifiers $\opA{n}$, $\opE{n}$ by numerical omniscience (cf.\ \autoref{lem:cm-lpo}), and bounded second-order quantifiers $\opA{Y \prec X}$, $\opE{Y \prec X}$ as well, by \autoref{lem:cm-gwo-prec-bounded-so-omni} which only requires the axiom of countable initial segments \ref{itm:cm-gwo-init-seg-axiom}.

Likewise, we can consider the class of \emph{$\Sigma$-formulae} in the language of $\CM$ as the class of formulae $\varphi\mleft(X\mright)$ equivalent to some formula of the form $\opE{Y} \tuple{X, Y} \in_2 \mathbf{X}$. This class contains all formulae with no third-order quantifiers or second-order universal quantifiers, and only one second-order existential quantifier occurring in strictly positive positions, i.e.\ not behind a negation or in the antecedent of an implication; the class of $\Sigma$-formulae is also closed under first-order and bounded second-order quantifiers, thanks to the recursion scheme \ref{itm:cm-rec-axiom} allowing us to swap first- and second-order quantifiers.

Therefore, we may consider two single axioms
% manual margin tweak by ~4pt to accommodate long second item
\begin{enumerate}[label=\uplabel{\noexpand\enumdynlabel},leftmargin=25.5pt]
  \def\enumdynlabel{W2$_0$}
  \item \label{itm:cm-gwo-res-arith} $\opA{X} \left(\opA{Y \prec X} Y \in_2 \mathbf{X} \rightarrow X \in_2 \mathbf{X}\right) \rightarrow \opA{X} X \in_2 \mathbf{X}$,
        \def\enumdynlabel{W2$_\Sigma$}
  \item \label{itm:cm-gwo-res-sigma} $\opA{X} \left(\opA{Y \prec X} \opE{Z} \tuple{Y, Z} \in_2 \mathbf{X} \rightarrow \opE{Z} \tuple{X, Z} \in_2 \mathbf{X}\right) \rightarrow \opA{X} \opE{Z} \tuple{X, Z} \in_2 \mathbf{X}$.
\end{enumerate}
Let $\GWO_0$\inda{conceptual mathematics, $\CM$!with GWO0@with $\GWO_0$} denote the axioms $\text{\ref{itm:cm-gwo-ord-axiom}} + \text{\ref{itm:cm-gwo-res-arith}} + \text{\ref{itm:cm-gwo-init-seg-axiom}}$, where the transfinite induction scheme is replaced by the single axiom \ref{itm:cm-gwo-res-arith}; and $\GWO_\Sigma$\inda{conceptual mathematics, $\CM$!with GWO1@with $\GWO_\Sigma$} denote the axioms $\text{\ref{itm:cm-gwo-ord-axiom}} + \text{\ref{itm:cm-gwo-res-sigma}} + \text{\ref{itm:cm-gwo-init-seg-axiom}}$ likewise.

I believe that these fragments may be argued as `somewhat predicative' in Weaver's standards, since by forbidding transfinite induction over a universal second-order quantifier, one essentially prevents construction using information about the entire second-order domain at each step of the inductive process; whereas on the other hand, querying about the well-ordering in the constructed initial segment so far should not be as dangerous.

$\CM + \GWO_\Sigma$, particularly, has good prospects in applications as well. Looking back on our use of the transfinite induction scheme, one should observe that, as remarked in \autoref{ftn:zorn-lemm-proof-sigma}, our proof of Zorn's lemma as \autoref{prop:cm-gwo-zorn} only uses the axiom once on a decidable formula and again on some $\opE{Z} \theta\mleft(X, Z\mright)$ where $\theta\mleft(X, Z\mright)$ is decidable, i.e.\ a $\Sigma$-formula. Thus, \autoref{prop:cm-gwo-zorn} and all of its consequences should still hold in $\CM + \GWO_\Sigma$. The fragment $\CM + \GWO_0$ though, is weaker and cannot accommodate as much.

As for their proof-theoretic strength, we shall first look at exactly the same interpretation as in \autoref{sec:oa-cm-gwo} and adapt \autoref{prop:cm-gwo-ind-real} to realise \ref{itm:cm-gwo-res-sigma}. Notice that the key use of bar induction in the proof for this instance is then on a claim of the form
\[ha {\downarrow} \Vdash \mleft(\opE{Z} \tuple{X, Z} \in_2 \mathbf{X}\mright)\mleft[H^{-1} a/X, f/\mathbf{X}\mright]\]
for all $a$ in a given well-ordering and fixed $h \in \Hyp$, $f \in S_3$. By definition of the realisability conditions, this is equivalent\footnote{Pedantically, one reads $\tuple{X, Z} \in_2 \mathbf{X}$ as an abbreviation for $\opE{W \in_2 \mathbf{X}} W = \tuple{X, Z}$ where the inner formula is additionally the defining formula of second-order ordered pairs, and thus it should have a realiser with some internal structure. However, such a realiser is easily computable in $\Hyp$ given correct assignments to the variables $X, Z$ and $W$. Thus it is not hard to modify the proof of \autoref{prop:cm-gwo-ind-real} in this specific case so that we just work with a function $h$ where the second component of every $ha {\downarrow}$ is trivial, as if $\tuple{X, Z} \in_2 \mathbf{X}$ is an atomic statement.} to saying
\[ha {\downarrow} \land \bp_0 \mleft(ha\mright) \in S_2 \land f \mleft(g \mleft(H^{-1} a\mright) \mleft(\bp_0 \mleft(ha\mright)\mright)\mright) = 1,\]
where $g = \oplam{den} \bdN \mleft(\bp_0 n\mright) 0 \mleft(d \mleft(\bp_1 n\mright)\mright) \mleft(\bdN \mleft(\bp_0 n\mright) 1 \mleft(e \mleft(\bp_1 n\mright)\mright) 0\mright)$ such that $gde \in S_2$ computes our ordered pair encoding
\[\tuple{X, Y} = \left\{\tuple{0, n} : n \in_1 X\right\} \cup \left\{\tuple{1, m} : m \in_1 Y\right\}\]
for two second-order objects represented by $d, e \in S_2$ respectively. The conjuncts $ha {\downarrow}$ and $\bp_0 \mleft(ha\mright) \in S_2$ are both $\Delta^1_2$, so we only need bar induction for $\Delta^1_2$-formulae to realise \ref{itm:cm-gwo-res-sigma}.

As it turns out, one gets nice base theories to work with if the formula in the bar induction scheme is restricted to $\Pi^1_n$ for some $n \in \mathbb{N}$. Using Simpson's notations in \cite{simpson09-soa}, these fragments shall be denoted as $\Pi^1_n\TI_0$, whereas $\Pi^1_n\TI$ denote their respective extensions with unrestricted induction on the natural numbers. We can conclude that:
\begin{proposition}
  $\Pi^1_2\TI$ suffices to realise \ref{itm:cm-gwo-res-sigma}, and thus all of $\CM + \GWO_\Sigma$.
\end{proposition}
The theories $\Pi^1_n\TI$ are studied in \cite{rathjen-weiermann93-kruskal-theorem} where it is concluded that $\left|\Pi^1_n\TI\right| = \vartheta_{\Omega^{\varepsilon_0}\mleft(n - 1\mright)}\mleft(0\mright)$ for $n \geq 2$, where $\Omega^\alpha\mleft(k\mright)$ denotes the power tower
\[\Omega^\alpha\mleft(0\mright) = \alpha, \qquad \Omega^\alpha\mleft(k + 1\mright) = \Omega^{\Omega^\alpha\mleft(k\mright)}.\]
It follows that the proof-theoretic strength of $\CM + \GWO_\Sigma$ is bounded above by $\vartheta_{\Omega^{\varepsilon_0}}\mleft(0\mright)$.

On the other hand, $\CM + \GWO_0$, while less capable, will require a weaker base theory to realise. This requires a different argument than \autoref{prop:cm-gwo-ind-real}:
\begin{proposition}
  $\ATR$ suffices to realise \ref{itm:cm-gwo-res-arith}, and thus all of $\CM + \GWO_0$.
\end{proposition}

\begin{proof}
  Fix some $h \in S_3$ and suppose that one has a realiser for the antecedent $\opA{X} \left(\opA{Y \prec X} Y \in_2 h \rightarrow X \in_2 h\right)$. This implies
  \[\opA{x \in S_2} \left(\opA{y \in S_2} \left(Hy \prec Hx \rightarrow hy = 1\right) \rightarrow hx = 1\right),\]
  and it suffices to show that $\opA{x \in S_2} hx = 1$ as then any $d \in \Hyp$ that is total on the appropriate inputs, e.g.\ $d = \oplam{hf} 0$, would realise the axiom as desired.

  Suppose that this is not the case, then there exists some $x \in S_2$ such that $hx = 0$, and it follows that we have some $y \in S_2$ such that $Hy \prec Hx$, where
  $h \mleft(H^{-1} \mleft(Hy\mright)\mright) = hy = 0$
  because $h \in S_3$ respects extensional equivalences, i.e.
  \[W = \left\{a \in \mathcal{H} : a \prec Hx \land a = H \mleft(H^{-1} a\mright) \land h \mleft(H^{-1} a\mright) = 0\right\}\]
  is nonempty.

  Now, just as in the proof of \autoref{prop:cm-gwo-ind-real}, we argue using the totality of $C^* \mleft(Hx\mright)$ that $W$ is a $\Delta^1_1$-subclass of $\mathbb{N}$ and thus a set in $\ATR$ by $\Delta^1_1$-comprehension. Since $\prec$ is a well-ordering on $\mathcal{H}$, this set $W$ must contain a $\prec$-minimal element $b$. However, there must then be some $y \in S_2$ such that $hy = 0$ and
  \[Hy \prec H \mleft(H^{-1} b\mright) = b \prec Hx,\]
  contradicting the minimality of $b$. Thus, we must have $\opA{x \in S_2} hx = 1$ by contradiction.
\end{proof}

It follows that the proof-theoretic strength of $\CM + \GWO_0$ is bounded above by the ordinal $\Gamma_{\varepsilon_0}$.

\label{loc:lim-pred-discuss}It is worth remarking briefly that while these ordinals lie beyond the well-known traditional `limit of predicativity', $\Gamma_0$ as worked out by Feferman \cite{feferman64-predicative-analysis} and Sch\"utte \cite{schutte65-predicative-well-orderings}, Weaver's mathematical conceptualist would not have difficulties seeing them as predicative. This is because Weaver views any countable iterations as concrete --- which is a stronger philosophical assumption than one can infer from Feferman or Sch\"utte's methodology --- and thus, as he made explicit in \cite{weaver09-predicativity}*{section 1.4}, he agrees with the following criticism in Howard \cite{howard96-proof-theory}: suppose that $\left(\gamma_n\right)_{n \in \mathbb{N}}$ is a sequence of ordinals with limit $\Gamma_0$, and that a predicativist has a uniform way of justifying each $\gamma_n$ as well-ordered from the well-foundedness of preceding ordinals, then he should immediately see the entire, countable process as a predicatively valid way to justify that `for all $n \in \mathbb{N}$, $\gamma_n$ is well-ordered', and thus derive that $\Gamma_0$ is well-ordered as well.

Assuming that this argument is philosophically sound, a predicativist like Weaver should easily reach the conclusion that ordinals as large as $\Gamma_{\varepsilon_0}$ are predicatively well-ordered; and following his claims about iterated truth theories in \cite{weaver09-predicativity}*{sections 2.8--2.9}, he might be able to say the same for $\vartheta_{\Omega^{\varepsilon_0}}\mleft(0\mright)$ as well. His methods, however, are unrelated to the focus of this thesis and will not be elaborated here.

On the other hand, I do not presently have lower bound arguments for these theories, thus it remains open whether these upper bounds for the weak fragments of $\CM + \GWO$ are tight; it could well be the case that $\CM + \GWO_\Sigma$ is actually weaker than $\ATR_0$.

Indeed, one could observe that the way we utilise the global well-ordering in $\CM + \GWO_\Sigma$ is very similar to how $\Sigma_1$-set induction\inds{set induction axiom!$\Sigma_1$-}, i.e.\ the axiom
\[\opA{x} \left(\opA{y \in x} \varphi\mleft(y\mright) \rightarrow \varphi\mleft(x\mright)\right) \rightarrow \opA{x} \varphi\mleft(x\mright)\]
for any $\Sigma_1$-formula $\varphi\mleft(x\mright)$, is used in Kripke--Platek set theory $\KP$ and the very same proof of Zorn's lemma can be replicated in $\KP^- + \Sigma_1\text{-set induction} + V = L$, where $\KP^-$ denotes the base theory of $\KP$ (with infinity) without the axiom of foundation. Since other axioms of $\CM$ like dependent choice do not contribute much to the strength of an intuitionistic theory (because they are trivially realised in a realisability interpretation using the recursion theorem and sufficient induction, cf.\ \autoref{lem:cm-rec-real}), we conjecture that:

\begin{conjecture}
  $\CM + \GWO_\Sigma$ has a much lower proof-theoretic ordinal than $\vartheta_{\Omega^{\varepsilon_0}}\mleft(0\mright)$; it might be equiconsistent with
  \[\KP^- + \Sigma_1\text{-set induction} + \text{full induction on $\omega$}.\]
\end{conjecture}

It is maintained by Rathjen \cite{rathjen93-kp-fragments} that the latter theory has the proof-theoretic ordinal\footnote{The precise result credited to Cantini in \cite{rathjen93-kp-fragments} is that $\alpha = \varphi_{\varepsilon_0}\mleft(0\mright)$ is the least ordinal such that $L_\alpha$ satisfies all $\Sigma_1$-consequences of the given extension of $\KP^-$. It then follows from \cite{rathjen93-kp-fragments}*{Theorem 2.1} and a folklore result in \cite{rathjen93-kp-fragments}*{Remark 2.2} that the theory has proof-theoretic ordinal $\varphi_\alpha\mleft(0\mright) = \varphi_{\varphi_{\varepsilon_0}\mleft(0\mright)}\mleft(0\mright)$.}
\[\varphi_{\varphi_{\varepsilon_0}\mleft(0\mright)}\mleft(0\mright),\]
credited to an obscure paper of Cantini \cite{cantini83-admissible-sets}; this is far below even $\Gamma_0$.

However, when we interpret the second-order objects as $S_2$ in the realisability model presented in this thesis, we need $\ATR_0$ at least to attach a non-trivial well-ordering to this set, as we did in \autoref{sec:oa-cm-gwo}. Thus, I suspect that an entirely different interpretation might be needed to generate results toward this conjecture and decide that it shall not be resolved within the scope of this thesis.

\chapter{Preliminaries in intuitionistic set theory}
\label{chap:ist-intro}

For the remainder of this thesis, we will shift our focus to the realisability of first-order intuitionistic set theories, i.e.\ formal theories in the first-order language that consists of a sole binary relation symbol $\in$. In this chapter, we survey and prove some preliminary results crucial for the developments in all later chapters, including presentations and comparisons of the key fragments of set theory that will be investigated. We will reproduce a standard realisability argument for $\CZF$ in $\CZF$ itself in Rathjen \cite{rathjen06-czf-realizability}, which will serve as a template for the varied realisability models we shall construct later, and we will examine some basic properties of G\"odel's constructible universe $L$ inside this interpretation.

\section{Constructive Zermelo--Fraenkel set theory}
\label{sec:czf-intro}

The main results in the set-theoretic part of this thesis will be about a well-known system called \emph{constructive Zermelo--Fraenkel set theory}, $\CZF$, and its extension by adding the powerset axiom. This pair of first-order theories is sandwiched between \emph{intuitionistic Kripke--Platek set theory} $\IKP$, which is a proper subtheory of $\CZF$, and the much stronger \emph{intuitionistic Zermelo--Fraenkel set theory} $\IZF$. The book draft \cite{aczel-rathjen10-cst-book} by Aczel and Rathjen serves as a standard reference for these constructs.

We can begin defining them by letting $\IKP$ be the same recursive set of axioms as classical Kripke--Platek set theory $\KP$ over intuitionistic logic, with the condition that the set induction scheme must be used instead of the foundation scheme\footnote{The foundation scheme is an analogue of the axiom of foundation in $\ZF$ and written as
  \[\opE{x} \varphi\mleft(x\mright) \rightarrow \opE{x} \left(\varphi\mleft(x\mright) \land \opA{y \in x} \neg \varphi\mleft(y\mright)\right)\]
  for any first-order formula $\varphi$, such as in \cite{barwise75-admissible-sets}*{section I.2}. This is classically equivalent to set induction, but intuitionistically stronger, as it implies the law of excluded middle immediately when one applies it to
  \[\varphi\mleft(x\mright) \defeqv x = 1 \lor \left(x = 0 \land \psi\right)\]
  for an arbitrary formula $\psi$ (cf.\ \cite{aczel-rathjen10-cst-book}*{Proposition 10.4.1}). On the other hand, any realisability argument, such as the one we will soon see in \autoref{sec:czf-real-in-czf}, demonstrates that set induction cannot imply the law of excluded middle.}. In other words, we have the following axioms:
\begin{enumerate}[label=\uplabel{\arabic*},series=set-axioms]
  \item \label{itm:set-ext-axiom} the axiom of extensionality,
        \[\opA{a} \opA{b} \left(\opA{x} \left(x \in a \leftrightarrow x \in b\right) \leftrightarrow a = b\right);\]
  \item \label{itm:set-pair-axiom} the pairing axiom,
        \[\opA{a} \opA{b} \opE{y} \opA{x} \left(x \in y \leftrightarrow x = a \lor x = b\right);\]
  \item \label{itm:set-union-axiom} the union axiom,
        \[\opA{a} \opE{b} \opA{x} \left(x \in b \leftrightarrow \opE{u \in a} x \in u\right);\]
  \item \label{itm:set-inf-axiom} the axiom of infinity, $\opE{a} \Ind\mleft(a\mright)$, where
        \begin{align*}
          \Ind\mleft(a\mright) & \defeqv \opE{y \in a} \opA{z} \neg z \in y \land \opA{x \in a} \opE{y \in a} y = x^+, \\
          y = x^+              & \defeqv \opA{z} \left(z \in y \leftrightarrow z \in x \lor z = x\right);
        \end{align*}
  \item \label{itm:set-sep-axiom} the scheme of \emph{$\Delta_0$-separation},
        \[\opA{a} \opE{y} \opA{x} \left(x \in y \leftrightarrow x \in a \land \varphi\mleft(x\mright)\right),\]
        for any $\Delta_0$-formula $\varphi\mleft(x\mright)$, i.e.\ a formula that may contain bounded quantifiers $\opA{y \in z}$, $\opE{y \in z}$ but not unbounded quantifiers;
  \item \label{itm:set-coll-axiom} the scheme of \emph{$\Delta_0$-collection},
        \[\opA{x \in a} \opE{y} \theta\mleft(x, y\mright) \rightarrow \opE{b} \opA{x \in a} \opE{y \in b} \theta\mleft(x, y\mright),\]
        for any $\Delta_0$-formula $\theta\mleft(x, y\mright)$;
        \xdef\setcollaxiomlabel{\arabic{enumi}}
  \item \label{itm:set-ind-axiom} the \emph{set induction}\inds{set induction axiom} scheme,
        \[\opA{a} \left(\opA{x \in a} \varphi\mleft(x\mright) \rightarrow \varphi\mleft(a\mright)\right) \rightarrow \opA{a} \varphi\mleft(a\mright),\]
        for any first-order formula $\varphi$.
\end{enumerate}

For some caveats, we treat equality as a primitive logical symbol instead of an abbreviation for extensional equality. We also see bounded quantifiers $\opA{y \in z}$, $\opE{y \in z}$ as separate primitive constructs in the language instead of abbreviations for some corresponding formulae with unbounded quantifiers, since it is routine to provide different conditions for bounded and unbounded quantifiers in many realisability interpretations (cf.\ \autoref{sec:czf-real-in-czf}). This means that we will assume as axioms
\begin{align*}
  \opA{x} \left(x \in a \rightarrow \varphi\mleft(x\mright)\right) & \leftrightarrow \opA{x \in a} \varphi\mleft(x\mright), \\
  \opE{x} \left(x \in a \land \varphi\mleft(x\mright)\right)       & \leftrightarrow \opE{x \in a} \varphi\mleft(x\mright),
\end{align*}
for any formula $\varphi$ and need to check that they are realised by our realisability conditions. However, we still opt to not include them in axiom listings for particular theories like $\IKP$, since it remains ambiguous whether they are `logical' or `non-logical' axioms of the system.

We shall also remark here that, in the context of weak set theories like $\KP$ or $\IKP$, many authors prefer to use the axiom of strong infinity instead of \ref{itm:set-inf-axiom}, namely
\[\opE{a} \left(\Ind\mleft(a\mright) \land \opA{b} \left(\Ind\mleft(b\mright) \rightarrow \opA{x \in a} x \in b\right)\right)\]
(cf.\ \cite{aczel-rathjen10-cst-book}*{Chapter 3}), since in order to define the least inductive set $\omega$ as the intersection of all inductive sets, one seemingly needs to invoke an unavailable instance of $\Pi_1$-separation. However, this is in fact not necessary, as it follows from the following lemma that the usual axiom of infinity \ref{itm:set-inf-axiom} already implies the axiom of strong infinity:
\begin{lemma}[$\IKP$]
  \label{lem:ikp-omega-from-inf}
  There is a $\Delta_0$-formula $\varphi\mleft(x\mright)$ such that for any set $a$ satisfying $\opE{b \subseteq a} \Ind\mleft(b\mright)$, the set
  \[\omega = \left\{x \in a : \varphi\mleft(x\mright)\right\}\]
  is precisely the least inductive set in the universe, i.e.\ we have $\Ind\mleft(\omega\mright)$\ and $\omega \subseteq b$ for any $b$ satisfying $\Ind\mleft(b\mright)$. It follows immediately from the axiom of extensionality that such $\omega$ is unique and does not depend on the choice of $a$.
\end{lemma}

\begin{proof}
  We consider the formula\footnote{This construction is a simplified version of \cite{wang-sw26-plump}*{Lemma 3.7} in my arXiv manuscript.}
  \[\varphi\mleft(x\mright) \defeqv \left(x = \varnothing \lor \opE{y \in x} x = y^+\right) \land \opA{y \in x} \left(y = \varnothing \lor \opE{z \in x} y = z^+\right),\]
  where $x = \varnothing$ is the $\Delta_0$-formula $\opA{y \in x} \neg y \in x$ denoting that $x$ is empty.

  Assume that $\varphi\mleft(x\mright)$ holds for some set $x$. We first show by set induction that every $y \in x$ is transitive, i.e.\ $\opA{z \in y} \opA{w \in z} w \in y$. This is because each $y \in x$ must be either the empty set or a successor $z^+$ where $z \in x$. In the latter case, $z \in y$ so the inductive hypothesis implies that $z$ is transitive, which immediately entails that $y = z^+$ is also transitive. Since $x$ itself is either the empty set or a successor $y^+$ where $y \in x$, we know $x$ is transitive as well.

  For any $y \in x$ and $z \in y$, we then have $z \in x$ and thus $z = \varnothing \lor \opE{w \in x} z = w^+$ where the latter disjunct implies $\opE{w \in y} z = w^+$ by transitivity of $y$. It follows that
  \[\opA{x} \left(\varphi\mleft(x\mright) \rightarrow \opA{y \in x} \varphi\mleft(y\mright)\right).\]

  Fix any set $b$ satisfying $\Ind\mleft(b\mright)$, it then follows easily from set induction that $\opA{x} \left(\varphi\mleft(x\mright) \rightarrow x \in b\right)$, and thus $\omega = \left\{x \in a : \varphi\mleft(x\mright)\right\} \subseteq b$ as desired. It now remains to show that $\Ind\mleft(\omega\mright)$ holds.

  Fix some $b \subseteq a$ satisfying $\Ind\mleft(b\mright)$, and let $\varnothing$ denote the necessarily unique empty set in $b$, then clearly $\varnothing \in \omega$. For any $x \in \omega$, we have $\varphi\mleft(x\mright)$ and thus $x \in b$ as shown above. It follows that its successor $x^+ = x \cup \left\{x\right\}$ lies in $b$ and thus in $a$. Now, it is easy to check that $\varphi\mleft(x\mright)$ implies $\varphi\mleft(x^+\mright)$ and thus $x^+ \in \omega$ by definition.
\end{proof}

We then have the following new axiom schemes in $\CZF$:
\begin{enumerate}[label=\uplabel{\noexpand\enumdynlabel},resume=set-axioms]
  \def\enumdynlabel{\setcollaxiomlabel$'$}
  \addtocounter{enumi}{-1} % not a new axiom
  \item \label{itm:set-str-coll-axiom} the \emph{strong collection} scheme,
        \[\opA{x \in a} \opE{y} \theta\mleft(x, y\mright) \rightarrow \opE{b} \left(\opA{x \in a} \opE{y \in b} \theta\mleft(x, y\mright) \land \opA{y \in b} \opE{x \in a} \theta\mleft(x, y\mright)\right),\]
        for any first-order formula $\theta$;
        \def\enumdynlabel{\arabic{enumi}}
  \item \label{itm:set-sub-coll-axiom} the \emph{subset collection}\inds{subset collection axiom} scheme,
        \begin{align*}
          \opE{c} \opA{u} & \left(\opA{x \in a} \opE{y \in b} \psi\mleft(x, y, u\mright)\right.                                                                                                               \\
                          & \quad \left.{} \rightarrow \opE{d \in c} \left(\opA{x \in a} \opE{y \in d} \psi\mleft(x, y, u\mright) \land \opA{y \in d} \opE{x \in a} \psi\mleft(x, y, u\mright)\right)\right),
        \end{align*}
        for any first-order formula $\psi$;
        \xdef\setsubcollaxiomlabel{\arabic{enumi}}
\end{enumerate}
and we simply define $\CZF = \IKP + \text{\ref{itm:set-str-coll-axiom}} + \text{\ref{itm:set-sub-coll-axiom}}$\inds{constructive Zermelo--Fraenkel set theory, $\CZF$}. The strong collection scheme \ref{itm:set-str-coll-axiom} entails \ref{itm:set-coll-axiom} and is necessary here instead of just the full collection scheme, since in intuitionistic logic, we cannot derive separation instances from collection and thus lack the sufficient amount of separation in $\CZF$ to prove strong collection from collection alone.

Subset collection is a more technical assertion that is often mentioned in relation to two more commonly seen axioms:
\begin{enumerate}[label=\uplabel{\noexpand\enumdynlabel}]
  \def\enumdynlabel{\setsubcollaxiomlabel E}
  \item \label{itm:set-exp-axiom} the axiom of \emph{exponentiation}, $\Exp$\inds{axiom of exponentiation, $\Exp$},
        \[\opA{a} \opA{b} \opE{c} \opA{f} \left(f \in c \leftrightarrow \left(f : a \rightarrow b\right)\right),\]
        where $f : a \rightarrow b$ is the usual $\Delta_0$-formula expressing that $f$ is a set-theoretic function with domain $a$ and codomain $b$, essentially that $f$ consists of elements from the Cartesian product $a \times b$, such that every $x \in a$ is paired with exactly one $y \in b$;
        \def\enumdynlabel{\setsubcollaxiomlabel P}
  \item \label{itm:set-powerset-axiom} the \emph{powerset} axiom,
        \[\opA{a} \opE{y} \opA{x} \left(x \in y \leftrightarrow x \subseteq a\right).\]
\end{enumerate}

Following \cite{aczel-rathjen10-cst-book}*{Chapter 5}, we know that:
\begin{proposition}
  \label{prop:ikp-exp-sub-coll-powerset}
  \begin{enumerate}[label=\uplabel{\roman*}]
    \item \label{itm:sub-coll-imp-exp} $\IKP$ proves that \ref{itm:set-sub-coll-axiom} implies \ref{itm:set-exp-axiom};
    \item \label{itm:powerset-imp-sub-coll} $\IKP + \text{strong collection \ref{itm:set-str-coll-axiom}}$ proves that \ref{itm:set-powerset-axiom} implies \ref{itm:set-sub-coll-axiom}.
  \end{enumerate}
\end{proposition}

\begin{proof}
  For \ref{itm:sub-coll-imp-exp}, fix sets $a$ and $b$ and observe that the Cartesian product $a \times b$ exists by two applications of $\Delta_0$-collection. For any set-theoretic function $f : a \rightarrow b$, we must have
  \[\opA{x \in a} \opE{y \in a \times b} \left(y \in f \land \opE{z \in b} y = \tuple{x, z}\right),\]
  where $\tuple{x, z} = \left\{\left\{x\right\}, \left\{x, z\right\}\right\}$ denotes the usual Kuratowski ordered pair construction that populates our Cartesian product. By subset collection, there is some set $c$ (depending on $a$ and $b$) such that for any $f : a \rightarrow b$, there is $d \in c$ such that
  \[\opA{x \in a} \opE{y \in d} \left(y \in f \land \opE{z \in b} y = \tuple{x, z}\right) \land \opA{y \in d} y \in f,\]
  i.e.\ $d = f$, by virtue of $f$ being a function. It follows that
  \[\left\{f \in c : \left(f : a \rightarrow b\right)\right\},\]
  which is a set by $\Delta_0$-separation, consists precisely of all functions from $a$ to $b$.

  For \ref{itm:powerset-imp-sub-coll}, we simply take $c = \Pow\mleft(b\mright)$, the powerset of $b$, and consider any $u$ such that $\opA{x \in a} \opE{y \in b} \psi\mleft(x, y, u\mright)$ holds. By strong collection, there is some set $d$ such that
  \[\opA{x \in a} \opE{y \in d} \psi\mleft(x, y, u\mright) \land \opA{y \in d} \opE{x \in a} \left(y \in b \land \psi\mleft(x, y, u\mright)\right).\]
  Note that this ensures $d \subseteq b$, thus $d \in c = \Pow\mleft(b\mright)$ as desired.
\end{proof}

In classical set theory, we can trivially construct some powerset $\Pow\mleft(x\mright)$ from the set of functions $f : x \rightarrow \left\{0, 1\right\}$ as
\[\Pow\mleft(x\mright) = \left\{\left\{y \in x : \tuple{y, 1} \in f\right\} : \left(f : x \rightarrow \left\{0, 1\right\}\right)\right\},\]
by $\Delta_0$-collection, hence \ref{itm:set-exp-axiom} implies \ref{itm:set-powerset-axiom} and all three axioms are equivalent. However, this uses the instance of excluded middle $\opA{x} \opA{y} \left(y \in x \lor \neg y \in x\right)$ and does not hold in $\CZF$.

Notably, Peter Aczel in \cite{aczel78-type-theoretic-cst} interpreted $\CZF$ in a version of Martin-L\"of type theory with a single $W$-type, which can be easily implemented in any of the theories $\ID_1$, $\BI$ or $\KP$ through inductive definitions. More detail about this interpretation will be given in \autoref{sec:mltt-interp}, where we adapt the construction to build exotic intuitionistic ordinals within the model. For now, observe that $\IKP$ already supports $\Sigma$-inductive definitions (see \autoref{sec:czf-real-in-czf}) and, in a way analogous to \autoref{thm:cm-ind-def}, contains the intuitionistic $\ID_1^i$ as a subtheory. Thus, using again the result in \cite{buchholz-pohlers78-iterated-id} that $\ID_1^i \equiv_{\Con} \BI$, we know that $\CZF$ has exactly the same proof-theoretic strength as $\KP$ at the Bachmann--Howard ordinal.

On the other hand, we denote $\CZFP = \CZF + \text{\ref{itm:set-powerset-axiom}}$\inds{constructive Zermelo--Fraenkel set theory, $\CZF$!1@$\CZFP$} to be the result of replacing its subset collection scheme by the powerset axiom. Rathjen in \cite{rathjen12-czf-powerset} showed that $\CZFP$ is equiconsistent with the theories $\KPP$\inds{Kripke--Platek set theory with powerset, $\KPP$} and $\IKPP$, i.e. classical and intuitionistic Kripke--Platek set theory with a new primitive unary function symbol\footnote{\label{ftn:kpp-prim-pow} As stipulated in this manner, the primitive operation $\Pow$ will be allowed to occur in $\Delta_0$-formulae used in the separation and collection axioms, which makes these systems different from simply adjoining the powerset axiom \ref{itm:set-powerset-axiom} to $\KP$ or $\IKP$. The latter are weaker since powersets in general will not be definable inside a $\Delta_0$-formula; see \cite{rathjen12-czf-powerset}*{section 2.3} for more discussion.} $\Pow$ defined to be the powerset operation. The interpretation of $\CZFP$ in $\KPP$ involves a modification to the usual realisability argument known as \emph{weak realisability}, which will become a key ingredient in \autoref{sec:czf-weak-real} later.

We can immediately see that $\KPP$ proves the consistency of full second-order arithmetic $\Zs$, i.e.\ the base theory $\ACA_0$ together with unrestricted comprehension axiom for all second-order formulae. $\BI$ is a proper subtheory of $\Zs$ and thus correspondingly, $\CZF$ with only the subset collection axiom is much weaker than the variant $\CZFP$ with the powerset axiom.

Before concluding this section, we also note that:
\begin{theorem}[\cite{aczel-rathjen10-cst-book}*{Corollary 4.2.8}]
  \label{thm:czf-em-eq-zf}
  $\CZF$ and $\CZFP$ are both equivalent to $\ZF$ over classical logic.
\end{theorem}

\begin{proof}
  By \autoref{prop:ikp-exp-sub-coll-powerset}, $\CZF$ implies the exponentiation axiom $\Exp$. Observe that classically, strong collection implies separation, set induction implies foundation and exponentiation implies powerset. Thus we have all the axioms of $\ZF$.
\end{proof}

While this is the case, there exist stronger intuitionistic fragments of $\ZF$. We can let $\IZF$ denote $\CZFP + \text{full separation}$\inds{intuitionistic Zermelo--Fraenkel set theory, $\IZF$}, so that $\IZF$ consists of essentially the same axioms as $\ZF$ except the axiom of foundation is replaced by set induction.

By forcing over a complete Boolean algebra in $\IZF$, Robin Grayson in \cite{grayson79-heyting-valued-models}*{section 2.3} verified that $\IZF$ is equiconsistent with classical $\ZF$. The strength of $\ZF$, however, is enormous even in comparison to $\CZFP$ and $\KPP$. Rathjen showed as \cite{rathjen20-power-kp-choice}*{Theorem 4.6} that $\KPP + V = L$ suffices\footnote{This is a slightly surprising result, which invokes the recent framework of relativised ordinal analysis developed in \cite{rathjen14-power-kp}. While the proof is not directly related to the topics of this thesis, this will be revisited in \autoref{sec:overview-v-eq-l} in greater detail to provide background for our subsequent investigations.} to prove $\Con\mleft(\KPP\mright)$, and $\KPP + V = L$ is an obvious subtheory of $\ZF + V = L$, which is equiconsistent with $\ZF$ through the constructible universe construction.

In this thesis, we will mainly be concerned with $\CZF$ and $\CZFP$; we will not work with models of $\IZF$ as it is too strong for many of the questions we analyse here. For features of realisability for $\IZF$, we refer the readers to the PhD thesis of Charles McCarty \cite{mccarty84-realizability-recursive-maths}.

\section{A standard realisability model for \texorpdfstring{$\CZF$}{CZF}}
\label{sec:czf-real-in-czf}

We now outline the standard realisability argument for $\CZF$ inside $\CZF$ in \cite{rathjen06-czf-realizability}, i.e.\ we show in the given interpretation that for any axiom $\varphi$ of $\CZF$, $\CZF$ proves that $\varphi$ is realised.

Given any pca $\mathcal{A}$ over natural numbers --- such that both $\mathcal{A}$ and its application function are sets in our base theory --- we first define the domain $V\mleft(\mathcal{A}\mright)$ of the interpretation, which we shall borrow the tradition from forcing and call the class of \emph{names}. This will be an uninterpreted realisability model (cf.\ \autoref{sec:uninterp-real}), as with many other realisability interpretations of set theory, and the names will form a very large class, which we construct using $\Sigma$-inductive definitions in $\IKP$ (cf.\ \cite{aczel-rathjen10-cst-book}*{section 19.3}).

\subsection*{$\Sigma$-inductive definitions}

\begin{theorem}[$\IKP$, \cite{aczel-rathjen10-cst-book}*{Theorem 19.3.1}\inds{Sigma-inductive definition@$\Sigma$-inductive definition}]
  \label{thm:set-ind-def}
  Let $\Phi$ be a $\Sigma_1$-class of ordered pairs. Then there is a $\Sigma_1$-class $I\mleft(\Phi\mright)$ satisfying
  \begin{enumerate}[label=\uplabel{\roman*}]
    \item \label{itm:set-ind-def-closed} for any $\tuple{x, y} \in \Phi$, if $y \subseteq I\mleft(\Phi\mright)$, then $x \in I\mleft(\Phi\mright)$;
    \item \label{itm:set-ind-def-min} for any formula $\varphi\mleft(x\mright)$, if
          \[\opA{{\tuple{x, y} \in \Phi}} \left(\opA{z \in y} \varphi\mleft(z\mright) \rightarrow \varphi\mleft(x\mright)\right),\]
          then $\opA{x \in I\mleft(\Phi\mright)} \varphi\mleft(x\mright)$.
  \end{enumerate}
\end{theorem}

\begin{proof}
  We consider the formula
  \begin{align*}
    \psi\mleft(u, x, f\mright) \defeqv & \text{$f$ is a function on domain $u^+ = u \cup \left\{u\right\}$} \land x \in f\mleft(u\mright)                                                           \\
                                       & \quad {} \land \opA{v \in u^+} \opA{y \in f\mleft(v\mright)} \opE{z} \left(\tuple{y, z} \in \Phi \land \opE{w \in v} z \subseteq f\mleft(w\mright)\right),
  \end{align*}
  which is $\Sigma_1$ because we can move the bounded quantifiers in the last conjunct inside unbounded existential quantifiers by $\Delta_0$-collection. Let
  \[I\mleft(\Phi\mright) = \left\{x : \opE{u} \opE{f} \psi\mleft(u, x, f\mright)\right\}.\]

  To verify \ref{itm:set-ind-def-closed}, fix any $\tuple{x, y} \in \Phi$ such that $y \subseteq I\mleft(\Phi\mright)$. By strong $\Sigma_1$-collection (which is a theorem in $\IKP$, cf.\ \cite{aczel-rathjen10-cst-book}*{Theorem 19.1.5}), let $s$ be a set such that
  \[\opA{z \in y} \opE{{\tuple{u, f} \in s}} \psi\mleft(u, z, f\mright) \land \opA{{\tuple{u, f} \in s}} \opE{z \in y} \psi\mleft(u, z, f\mright).\]
  Take $q = \bigcup \left\{u : \tuple{u, f} \in s\right\}$ and consider
  \[g = \left\{\tuple{q^+, \left\{x\right\}}, \tuple{q, y}\right\} \cup \left\{\tuple{v, \bigcup_{\tuple{u, f} \in s, v \in u} f\mleft(v\mright)} : v \in q\right\}.\]
  Clearly, $g$ is a function since $q^+, q \not\in q$ by set induction. We also see that for each $\tuple{u, f} \in s$, $v \in u$, $f\mleft(v\mright) \subseteq g\mleft(v\mright)$, thus $\psi\mleft(q, z, g \upharpoonright q^+\mright)$ for each $z \in y$ and therefore $\psi\mleft(q^+, x, g\mright)$. It follows that $x \in I\mleft(\Phi\mright)$.

  For \ref{itm:set-ind-def-min}, fix any $x \in I\mleft(\Phi\mright)$ and corresponding $u, f$ such that $\psi\mleft(u, x, f\mright)$ holds. Using stipulations in $\psi$ we immediately show by set induction on $v$ that
  \[\opA{v \in u^+} \opA{y \in f\mleft(v\mright)} \varphi\mleft(y\mright),\]
  and it follows from $x \in f\mleft(u\mright)$ that $\varphi\mleft(x\mright)$ holds as well.
\end{proof}

The use of strong $\Sigma_1$-collection here is covered by the strong collection axiom when we work in $\CZF$, thus in $\CZF$ we can even generate inductive definitions from any definable class $\Phi$ (cf.\ \cite{aczel-rathjen10-cst-book}*{section 12.1}).

\begin{corollary}[$\IKP$]
  \label{cor:set-ind-def-exist-pre}
  Given the generated inductive definition $I\mleft(\Phi\mright)$, for any $x \in I\mleft(\Phi\mright)$, there is some $y \subseteq I\mleft(\Phi\mright)$ such that $\tuple{x, y} \in \Phi$.
\end{corollary}

\begin{proof}
  Simply take the formula
  \[\varphi\mleft(x\mright) \defeqv \opE{y \subseteq I\mleft(\Phi\mright)} \tuple{x, y} \in \Phi\]
  and observe that $\varphi\mleft(x\mright)$ implies $x \in I\mleft(\Phi\mright)$ by condition \ref{itm:set-ind-def-closed} in \autoref{thm:set-ind-def}. Thus for any $\tuple{x, y} \in \Phi$, if $\opA{z \in y} \varphi\mleft(z\mright)$, then $y \subseteq I\mleft(\Phi\mright)$, i.e.\ $\varphi\mleft(x\mright)$ holds. Therefore, the statement follows from condition \ref{itm:set-ind-def-min} in \autoref{thm:set-ind-def}.
\end{proof}

Now, we construct the class of names for our realisability model:
\begin{corollary}[$\IKP$]
  \label{cor:czf-in-czf-name-def}
  There is a $\Sigma_1$-class $V\mleft(\mathcal{A}\mright)$ satisfying
  \[\opA{x} \left(x \in V\mleft(\mathcal{A}\mright) \leftrightarrow x \subseteq \mathcal{A} \times V\mleft(\mathcal{A}\mright)\right).\]
\end{corollary}

\begin{proof}
  Simply let class $\Phi = \left\{\tuple{x, y} : x \subseteq \mathcal{A} \times y\right\}$ and set $V\mleft(\mathcal{A}\mright) = I\mleft(\Phi\mright)$. The forward direction is trivially ensured by \autoref{cor:set-ind-def-exist-pre} whereas for the backward direction, fix $x \subseteq \mathcal{A} \times V\mleft(\mathcal{A}\mright)$ and take $y = \left\{z : \tuple{a, z} \in x\right\}$ by $\Delta_0$-collection. Then $\tuple{x, y} \in \Phi$ and condition \ref{itm:set-ind-def-closed} in \autoref{thm:set-ind-def} ensures that $x \in V\mleft(\mathcal{A}\mright)$.
\end{proof}

\subsection*{Realisability for $\CZF$}

We proceed to define the atomic realisability conditions as class functions $F_\in, G_= : V\mleft(\mathcal{A}\mright)^2 \rightarrow \mathcal{P}\mleft(\mathcal{A}\mright)$ by simultaneous $\Sigma$-recursion\footnote{A pedantic formulation of this can be quite complicated. Observe that an `element name' of some $u \in V\mleft(\mathcal{A}\mright)$ is not a direct element of $u$, but rather some
  \[x \in \left\{f, x\right\} \in \tuple{f, x} \in u\]
  by the Kuratowski pair construction, thus the induction needs to be on the transitive closure $\TC\mleft(\left\{u\right\}\mright)$. Additionally, we perform the recursion on two parameters $u$ and $v$, while also swapping their positions from time to time (see the second conjunct in the definition of $G_=$), so the full formal presentation must be
  \begin{align*}
     & \opA{u} \opA{x \in \TC\mleft(\left\{u\right\}\mright) \cap V\mleft(\mathcal{A}\mright)} \opA{v} \opA{y \in \TC\mleft(\left\{v\right\}\mright) \cap V\mleft(\mathcal{A}\mright)} \\
     & \qquad \text{all four constructs $F_\in\mleft(x, y\mright)$, $F_\in\mleft(y, x\mright)$, $G_=\mleft(x, y\mright)$, $G_=\mleft(y, x\mright)$ are well-defined and unique},
  \end{align*}
  which is then proved by set induction on $u$ first and then on $v$ for each inductive step.} in $\IKP$ as follows:
\begin{align*}
  F_\in\mleft(u, v\mright) & = \left\{e \in \mathcal{A} : \opE{{\tuple{f, x} \in v}} \left(\bp_0 e {\downarrow} = f \land \bp_1 e {\downarrow} \in G_=\mleft(u, x\mright)\right)\right\}, \\
  G_=\mleft(u, v\mright)   & = \begin{aligned}[t]
                                  & {\left\{e \in \mathcal{A} : \opA{{\tuple{f, x} \in u}} \bp_0 e f {\downarrow} \in F_\in\mleft(x, v\mright) \right.}  \\
                                  & \qquad \qquad \left.{} \land \opA{{\tuple{f, x} \in v}} \bp_1 e f {\downarrow} \in F_\in\mleft(x, u\mright)\right\}.
                               \end{aligned}
\end{align*}

The atomic clauses of the realisability conditions are then
\begin{align*}
  d \Vdash u \in v & \defeqv d \in F_\in\mleft(u, v\mright), \\
  d \Vdash u = v   & \defeqv d \in G_=\mleft(u, v\mright).
\end{align*}

Additionally, as previously mentioned, we will specify realisability for bounded quantifiers as separate clauses given by
\begin{align*}
  d \Vdash \opA{x \in u} \varphi & \defeqv \opA{{\tuple{f, v} \in u}} df {\downarrow} \Vdash \varphi\mleft[v/x\mright],                                                  \\
  d \Vdash \opE{x \in u} \varphi & \defeqv \opE{{\tuple{f, v} \in u}} \left(\bp_0 d {\downarrow} = f \land \bp_1 d {\downarrow} \Vdash \varphi\mleft[v/x\mright]\right).
\end{align*}
Observe that the forms of realisers for bounded quantifiers are chosen to match the definitions of $F_\in$ and $G_=$ precisely, so that we immediately have
\begin{align*}
  \left(d \Vdash u \in v\right) & \leftrightarrow \left(d \Vdash \opE{x \in v} u = x\right),               \\
  \left(d \Vdash u = v\right)   & \leftrightarrow \left(d \Vdash u \subseteq v \land v \subseteq u\right),
\end{align*}
where $u \subseteq v$ is read as an abbreviation for $\opA{x \in u} x \in v$.

We will treat $V\mleft(\mathcal{A}\mright)$ under the realisability conditions as a `model' analogous to classical models in first-order model theory, and write $V\mleft(\mathcal{A}\mright) \vDash \varphi\mleft[\vec{u}/\vec{x}\mright]$ to abbreviate for
\[\opE{d \in \mathcal{A}} d \Vdash \varphi\mleft[\vec{u}/\vec{x}\mright]\]
where $\vec{u} \in V\mleft(\mathcal{A}\mright)^k$.

Now, we first check that the equality axioms for our primitive equality symbol (cf.\ \autoref{sec:real-con-int-logic}) are realised. Observe that there are precisely two cases for the atomic substitution scheme \ref{itm:eq-sub-axiom} (up to renaming of variables):
\begin{lemma}[\cite{rathjen06-czf-realizability}*{Lemma 4.2}]
  We have fixed $\ior, \ios, \iot, \ioso, \iost \in \mathcal{A}$ such that for any $x, y, z \in V\mleft(\mathcal{A}\mright)$,
  \begin{enumerate}[label=\uplabel{\roman*}]
    \item $\ior \Vdash x = x$;
    \item $\ios \Vdash x = y \rightarrow y = x$;
    \item $\iot \Vdash x = y \rightarrow \left(y = z \rightarrow x = z\right)$;
    \item $\ioso \Vdash x = y \rightarrow \left(x \in z \rightarrow y \in z\right)$;
    \item $\iost \Vdash x = y \rightarrow \left(z \in x \rightarrow z \in y\right)$.
  \end{enumerate}
  Thus the equality axioms \ref{itm:eq-refl-axiom}--\ref{itm:eq-sub-axiom} are realised.
\end{lemma}

\begin{proof}
  \begin{enumerate}[label=\uplabel{\roman*}]
    \item We let
          \[h = \br \mleft(\oplam{fg} \bp f \mleft(\bp gg\mright)\mright)\]
          by recursion, and simply define $\ior = \bp hh$. By set induction\footnote{Again since names in $V\mleft(\mathcal{A}\mright)$ are not direct elements of each other, this induction must be formalised on the transitive closure, i.e.\ proving the formula
            \[\opA{x} \opA{u \in \TC\mleft(\left\{x\right\}\mright) \cap V\mleft(\mathcal{A}\mright)} \left(h \Vdash u \subseteq u \land \ior \Vdash u = u\right).\]
            Same applies to subsequent induction arguments.} on $V\mleft(\mathcal{A}\mright)$, it is not hard to check that $h \Vdash x \subseteq x$ and $\ior \Vdash x = x$ for any $x \in V\mleft(\mathcal{A}\mright)$.
    \item $\ios = \oplam{d} \bp \mleft(\bp_1 d\mright) \mleft(\bp_0 d\mright)$ is trivial.
    \item We examine some functions
          \begin{align*}
            s_2 & = \oplam{gde} g \mleft(\bp_1 e\mright) \mleft(d \mleft(\bp_0 e\mright)\mright),                                                                    \\
            t   & = \oplam{gdef} s_2 ge \mleft(df\mright),                                                                                                           \\
            p   & = \oplam{gde} \bp \mleft(tg \mleft(\bp_0 d\mright) \mleft(\bp_0 e\mright)\mright) \mleft(tg \mleft(\bp_1 e\mright) \mleft(\bp_1 d\mright)\mright), \\
            s_1 & = \oplam{gde} \bp \mleft(\bp_0 e\mright) \mleft(pgd \mleft(\bp_1 e\mright)\mright),
          \end{align*}
          and observe that, fix\footnote{It is important that we focus on the variable $y$ through this induction, since when $p$ invokes $t$ above, it will swap the variables $x$ and $z$, but $y$ remains unchanged and provides us with the appropriate inductive hypothesis to work with.} $y \in V\mleft(\mathcal{A}\mright)$ and given some $g \in \mathcal{A}$ such that for any $\tuple{f, v} \in y$ and any $u, w \in V\mleft(\mathcal{A}\mright)$,
          \[g \Vdash u = v \rightarrow \left(v \in w \rightarrow u \in w\right),\]
          then we have, for any $x, z \in V\mleft(\mathcal{A}\mright)$,
          \begin{align*}
            s_2 g & \Vdash y \subseteq z \rightarrow \left(x \in y \rightarrow x \in z\right),             \\
            tg    & \Vdash x \subseteq y \rightarrow \left(y \subseteq z \rightarrow x \subseteq z\right), \\
            pg    & \Vdash x = y \rightarrow \left(y = z \rightarrow x = z\right),                         \\
            s_1 g & \Vdash x = y \rightarrow \left(y \in z \rightarrow x \in z\right).
          \end{align*}
          Thus, by recursion, if we take $g = \oplam{de} \br \mleft(\oplam{ah} s_1 h \mleft(\bp_0 a\mright) \mleft(\bp_1 a\mright)\mright) \mleft(\bp de\mright)$, then we can see that $g, s_2 g, tg, pg, s_1 g$ realise their corresponding formulae for all $x, y, z \in V\mleft(\mathcal{A}\mright)$ by set induction on $y$. It suffices now to take $\iot = pg$.
    \item Simply use $s_1$ from the induction above as $\ioso = s_1 g$.
    \item Simply use $s_2$ from the induction above as $\iost = \oplam{d} s_2 g \mleft(\bp_0 d\mright)$. \qedhere
  \end{enumerate}
\end{proof}

Through a standard (metatheoretic) induction on the complexity of formula $\varphi$, we know that the equality axioms prove
\[x = y \rightarrow \left(\varphi\mleft(x\mright) \rightarrow \varphi\mleft(y\mright)\right)\]
for any formula $\varphi$. By \autoref{prop:real-con-int-logic}, it follows that there is some $e \in \mathcal{A}$, depending on $\varphi$, such that
\[e \Vdash x = y \rightarrow \left(\varphi\mleft(x\mright) \rightarrow \varphi\mleft(y\mright)\right).\]
We use this to show that the axioms for the bounded quantifiers are realised:
\begin{lemma}
  \label{lem:set-bnd-quant-axioms-real}
  The following axioms are realised:
  \begin{align*}
    \opA{x} \left(x \in a \rightarrow \varphi\mleft(x\mright)\right) & \leftrightarrow \opA{x \in a} \varphi\mleft(x\mright), \\
    \opE{x} \left(x \in a \land \varphi\mleft(x\mright)\right)       & \leftrightarrow \opE{x \in a} \varphi\mleft(x\mright).
  \end{align*}
\end{lemma}

\begin{proof}
  First see that $\oplam{df} d \mleft(\bp f \ior\mright)$ realises $\opA{x} \left(x \in a \rightarrow \varphi\mleft(x\mright)\right) \rightarrow \opA{x \in a} \varphi\mleft(x\mright)$, whereas $\oplam{d} \bp \mleft(\bp \mleft(\bp_0 d\mright) \ior\mright) \mleft(\bp_1 d\mright)$ realises $\opE{x \in a} \varphi\mleft(x\mright) \rightarrow \opE{x} \left(x \in a \land \varphi\mleft(x\mright)\right)$.

  For the converses, fix $e \in \mathcal{A}$ such that $e \Vdash x = y \rightarrow \left(\varphi\mleft(x\mright) \rightarrow \varphi\mleft(y\mright)\right)$. Then $\oplam{df} e \mleft(\ios \mleft(\bp_1 f\mright)\mright) \mleft(d \mleft(\bp_0 f\mright)\mright)$ realises $\opA{x \in a} \varphi\mleft(x\mright) \rightarrow \opA{x} \left(x \in a \rightarrow \varphi\mleft(x\mright)\right)$, whereas $\oplam{d} \bp \mleft(\bp_{00}d\mright) \mleft(e \mleft(\bp_{01} d\mright) \mleft(\bp_1 d\mright)\mright)$ realises $\opE{x} \left(x \in a \land \varphi\mleft(x\mright)\right) \rightarrow \opE{x \in a} \varphi\mleft(x\mright)$.
\end{proof}

Moving on to the axioms of $\CZF$, where \cite{rathjen06-czf-realizability}*{section 5} serves as a primary reference, we have:
\begin{lemma}
  The axiom of extensionality \ref{itm:set-ext-axiom} is realised:
  \[\opA{a} \opA{b} \left(\opA{x} \left(x \in a \leftrightarrow x \in b\right) \leftrightarrow a = b\right).\]
\end{lemma}

\begin{proof}
  It suffices to realise $\opA{x} \left(x \in a \leftrightarrow x \in b\right) \rightarrow a = b$ since the converse is a consequence of the equality axioms. We use the fact that $\left(d \Vdash a = b\right) \leftrightarrow \left(d \Vdash a \subseteq b \land b \subseteq a\right)$, and it follows that
  \[\oplam{d} \bp \mleft(\oplam{f} \bp_0 d \mleft(\bp f \ior\mright)\mright) \mleft(\oplam{f} \bp_1 d \mleft(\bp f \ior\mright)\mright) \Vdash \opA{x} \left(x \in a \leftrightarrow x \in b\right) \rightarrow a = b. \qedhere\]
\end{proof}

For the pairing axiom \ref{itm:set-pair-axiom} and union axiom \ref{itm:set-union-axiom}, we define
\begin{align*}
  \left\{a, b\right\}^\bullet & = \left\{\tuple{0, a}\right\} \cup \left\{\tuple{1, b}\right\},            \\
  \bigcup\nolimits^\bullet a  & = \left\{\tuple{\bp fg, y}: \tuple{f, x} \in a, \tuple{g, y} \in x\right\}
\end{align*}
for any $a, b \in V\mleft(\mathcal{A}\mright)$. For pairing, observe that simply
\[\bp \mleft(\oplam{d} d\mright) \mleft(\oplam{d} d\mright) \Vdash x \in \left\{a, b\right\}^\bullet \leftrightarrow x = a \lor x = b,\]
whereas for union, note that our definition ensures
\[\left(\bp \mleft(\bp fg\mright) d \Vdash x \in \bigcup\nolimits^\bullet a\right) \leftrightarrow \left(\bp f \mleft(\bp gd\mright) \Vdash \opE{u \in a} x \in u\right),\]
and hence we have
\begin{align*}
  \oplam{d} \bp \mleft(\bp_{00} d\mright) \mleft(\bp \mleft(\bp_{01} d\mright) \mleft(\bp_1 d\mright)\mright) & \Vdash x \in \bigcup\nolimits^\bullet a \rightarrow \opE{u \in a} x \in u, \\
  \oplam{d} \bp \mleft(\bp \mleft(\bp_0 d\mright) \mleft(\bp_{10} d\mright)\mright) \mleft(\bp_{11} d\mright) & \Vdash \opE{u \in a} x \in u \rightarrow x \in \bigcup\nolimits^\bullet a.
\end{align*}

For the axiom of infinity \ref{itm:set-inf-axiom}, we recursively define
\begin{align*}
  \overline{n}      & = \left\{\tuple{m, \overline{m}} : m \in n\right\} \in V\mleft(\mathcal{A}\mright) \quad \text{for all $n \in \omega$}, \\
  \overline{\omega} & = \left\{\tuple{n, \overline{n}} : n \in \omega\right\} \in V\mleft(\mathcal{A}\mright).
\end{align*}

We see that $\overline{\omega}$ is realised as the least inductive set precisely:
\begin{lemma}
  \label{lem:czf-real-omega-least-ind}
  Both $\Ind\mleft(\overline{\omega}\mright)$ and $\Ind\mleft(a\mright) \rightarrow \overline{\omega} \subseteq a$ are realised. Note that the former suffices to show that $\opE{a} \Ind\mleft(a\mright)$ is realised.
\end{lemma}

\begin{proof}
  Let
  \[g = \oplam{n} \bp \mleft(\oplam{d} \bdN dn \mleft(\bp 1 \mleft(\bp_1 d\mright)\mright) \mleft(\bp 0 d\mright)\mright) \mleft(\oplam{d} \bdN \mleft(\bp_0 d\mright) 0d \mleft(\bp nd\mright)\mright)\]
  and verify that for any $n \in \omega$, $gn \Vdash \overline{n + 1} = \overline{n}^+$. It follows that
  \[\bp \mleft(\bp 00\mright) \mleft(\oplam{n} \bp \mleft(\bsN n\mright) \mleft(gn\mright)\mright) \Vdash \Ind\mleft(\overline{\omega}\mright).\]

  For the latter, we instead check that
  \[\opA{x \in \overline{\omega}} \left(\left(x = \varnothing \lor \opE{y \in x} x = y^+\right) \land \opA{y \in x} \left(y = \varnothing \lor \opE{z \in x} y = z^+\right)\right)\]
  is realised. This is clearly the case as\footnote{One caveat is that $x = \varnothing$ as defined in \autoref{lem:ikp-omega-from-inf} is the $\Delta_0$-formula $\opA{y \in x} \neg y \in x$ instead of $\opA{y} \neg y \in x$, so it demands a slightly more complicated realiser.} $\bdN^* n0 \mleft(\bp 0 \mleft(\oplam{a} 0\mright)\mright) \mleft(\bp 1 \mleft(\bp \mleft(\bpN n\mright) \mleft(g \mleft(\bpN n\mright)\mright)\mright)\mright)$ realises $\overline{n} = \varnothing \lor \opE{y \in \overline{m}} \overline{n} = y^+$ for any $m \geq n \in \omega$. Thus, once we show that all axioms of $\CZF$ are realised, it follows from \autoref{lem:ikp-omega-from-inf} inside the realisability model that $\overline{\omega}$ is precisely the least inductive set as we constructed in that lemma.
\end{proof}

We will leave the scheme of $\Delta_0$-separation \ref{itm:set-sep-axiom} to the end to apply a trick that depends on collection. So we first see that:
\begin{lemma}
  \label{lem:czf-real-set-ind}
  The set induction scheme \ref{itm:set-ind-axiom} is realised:
  \[\opA{a} \left(\opA{x \in a} \varphi\mleft(x\mright) \rightarrow \varphi\mleft(a\mright)\right) \rightarrow \opA{a} \varphi\mleft(a\mright).\]
\end{lemma}

\begin{proof}
  We claim that for any $e \in \mathcal{A}$, $\oplam{f} \br \mleft(\oplam{eg} fg\mright) e$ realises the axiom. Suppose that $f \Vdash \opA{a} \mleft(\opA{x \in a} \varphi\mleft(x\mright) \rightarrow \varphi\mleft(a\mright)\mright)$, then we invoke set induction on $x \in V\mleft(\mathcal{A}\mright)$ and, from the inductive hypothesis that for any $e \in \mathcal{A}$ and $\tuple{d, y} \in x$,
  \[\br \mleft(\oplam{eg} fg\mright) e \Vdash \varphi\mleft(y\mright),\]
  we derive that $f \mleft(\oplam{e} \br \mleft(\oplam{eg} fg\mright) e\mright) \Vdash \varphi\mleft(x\mright)$. Thus the claim follows.

  We can simply take $\oplam{f} \br \mleft(\oplam{eg} fg\mright) \bk$, for example, for a concrete realiser.
\end{proof}

\begin{lemma}
  \label{lem:czf-str-coll-real}
  The strong collection scheme \ref{itm:set-str-coll-axiom} is realised:
  \[\opA{x \in a} \opE{y} \theta\mleft(x, y\mright) \rightarrow \opE{b} \left(\opA{x \in a} \opE{y \in b} \theta\mleft(x, y\mright) \land \opA{y \in b} \opE{x \in a} \theta\mleft(x, y\mright)\right).\]
\end{lemma}

\begin{proof}
  For $a \in V\mleft(\mathcal{A}\mright)$, suppose that $f \Vdash \opA{x \in a} \opE{y} \theta\mleft(x, y\mright)$, i.e.\ for any $\tuple{e, x} \in a$, $fe {\downarrow} \Vdash \theta\mleft(x, y\mright)$ for some $y \in V\mleft(\mathcal{A}\mright)$. By strong collection in the base theory, there is some set $s$ such that
  \begin{enumerate}[label=\uplabel{\roman*}]
    \item for any $\tuple{e, x} \in a$, there is $\tuple{e, y} \in s$ such that $fe \Vdash \theta\mleft(x, y\mright)$, and
    \item \label{itm:czf-str-coll-real-con2} for any $u \in s$, we have $u = \tuple{e, y} \in \mathcal{A} \times V\mleft(\mathcal{A}\mright)$ and there is $\tuple{e, x} \in a$ such that $fe \Vdash \theta\mleft(x, y\mright)$.
  \end{enumerate}
  Note that condition \ref{itm:czf-str-coll-real-con2} ensures that $s \in V\mleft(\mathcal{A}\mright)$. We see that $\oplam{e} \bp e \mleft(fe\mright)$ realises both $\opA{x \in a} \opE{y \in s} \theta\mleft(x, y\mright)$ and $\opA{y \in s} \opE{x \in a} \theta\mleft(x, y\mright)$. It follows that $\oplam{f} \bp \mleft(\oplam{e} \bp e \mleft(fe\mright)\mright) \mleft(\oplam{e} \bp e \mleft(fe\mright)\mright)$ realises the axiom.
\end{proof}

\begin{lemma}
  \label{lem:czf-sub-coll-real}
  The subset collection scheme \ref{itm:set-sub-coll-axiom} is realised:
  \begin{align*}
    \opE{c} \opA{u} & \left(\opA{x \in a} \opE{y \in b} \psi\mleft(x, y, u\mright)\right.                                                                                                               \\
                    & \quad \left.{} \rightarrow \opE{d \in c} \left(\opA{x \in a} \opE{y \in d} \psi\mleft(x, y, u\mright) \land \opA{y \in d} \opE{x \in a} \psi\mleft(x, y, u\mright)\right)\right),
  \end{align*}
\end{lemma}

\begin{proof}
  Fix $a, b, u \in V\mleft(\mathcal{A}\mright)$ and suppose that $f \Vdash \opA{x \in a} \opE{y \in b} \psi\mleft(x, y, u\mright)$, i.e.\ for any $\tuple{e, x} \in a$, there exists $\tuple{d, y} \in b$ such that $\bp_0 \mleft(fe\mright) {\downarrow} = d$ and $\bp_1 \mleft(fe\mright) {\downarrow} \Vdash \psi\mleft(x, y, u\mright)$.

  We construct $b^* = \left\{\tuple{e, y} : e \in \mathcal{A} \land \opE{d \in \mathcal{A}} \tuple{d, y} \in b\right\}$. Now, for any $\tuple{e, x} \in a$, there exists $\tuple{e, y} \in b^*$ such that $\bp_1 \mleft(fe\mright) {\downarrow} \Vdash \psi\mleft(x, y, u\mright)$. By subset collection in the base theory, there is some set $c$, depending only on $a$ and $b^*$, such that for any $u \in V\mleft(\mathcal{A}\mright)$ and realiser $f$ as above, there is some $d \in c$ such that
  \begin{enumerate}[label=\uplabel{\roman*}]
    \item for any $\tuple{e, x} \in a$, there is $\tuple{e, y} \in d$ such that $\bp_1 \mleft(fe\mright) \Vdash \psi\mleft(x, y, u\mright)$, and
    \item \label{itm:czf-sub-coll-real-con2} for any $v \in d$, we have $v = \tuple{e, y} \in b^*$ and there is $\tuple{e, x} \in a$ such that $\bp_1 \mleft(fe\mright) \Vdash \psi\mleft(x, y, u\mright)$.
  \end{enumerate}
  Condition \ref{itm:czf-sub-coll-real-con2} ensures that such $d \subseteq b^*$ and thus we can safely restrict
  \[c^* = \left\{\tuple{0, d} : d \in c, d \subseteq b^*\right\} \in V\mleft(\mathcal{A}\mright)\]
  by $\Delta_0$-separation and then collection. Just as in the proof of \autoref{lem:czf-str-coll-real}, given any $u \in V\mleft(\mathcal{A}\mright)$, $\oplam{e} \bp e \mleft(\bp_1 \mleft(fe\mright)\mright)$ realises both $\opA{x \in a} \opE{y \in d} \psi\mleft(x, y, u\mright)$ and $\opA{y \in d} \opE{x \in a} \psi\mleft(x, y, u\mright)$ for the specific $d$ above, which satisfies $\tuple{0, d} \in c^*$. It follows that
  \[\oplam{f} \bp 0 \mleft(\bp \mleft(\oplam{e} \bp e \mleft(\bp_1 \mleft(fe\mright)\mright)\mright) \mleft(\oplam{e} \bp e \mleft(\bp_1 \mleft(fe\mright)\mright)\mright)\mright)\]
  realises the axiom.
\end{proof}

Finally, for $\Delta_0$-separation, we shall avoid the complexities of analysing the realisability condition for a $\Delta_0$-formula by invoking the following result:

Let the \emph{binary intersection}\inds{binary intersection axiom} axiom be the statement that
\[\opA{a} \opA{b} \opE{y} \opA{x} \left(x \in y \leftrightarrow x \in a \land x \in b\right),\]
which essentially states that the intersection $a \cap b$ is a set. Then we have:
\begin{theorem}[\cite{aczel-rathjen10-cst-book}*{Corollary 9.5.7}]
  The system $\IKP - \Delta_0\text{-separation} + \text{strong collection \ref{itm:set-str-coll-axiom}}$ suffices\footnote{To be exact, strong $\Delta_0$-collection already suffices for the constructions here.} to prove that $\Delta_0$-separation is equivalent to the binary intersection axiom.
\end{theorem}

\begin{proof}[Proof sketch]
  The proof is quite technical and we refer the readers to \cite{aczel-rathjen10-cst-book}*{section 9.5} for details. Roughly speaking, we observe that given a set $u = \left\{x : x = a \land \varphi\right\}$ where $x$ does not occur free in $\varphi$, strong collection implies that for any set $b$,
  \[\left\{x : x = b \land \varphi\right\} = \left\{b : x \in u\right\}\]
  exists. Thus, if we let $!\varphi$ denote the statement $\opA{a} \opE{b} b = \left\{x : x = a \land \varphi\right\}$, then we just need to find one such set in order to prove $!\varphi$.

  Now, notice that for any $a, b$,
  \[\left\{a\right\} \cap \left\{b\right\} = \left\{x : x = a \land a = b\right\}.\]
  Thus the binary intersection axiom implies that $!\mleft(a = b\mright)$. This is used as a starting point to derive $!\varphi$ for every $\Delta_0$-formula $\varphi$. It then follows that
  \[\left\{x \in a : \varphi\mleft(x\mright)\right\} = \bigcup \left\{\left\{y : y = x \land \varphi\mleft(x\mright)\right\} : x \in a\right\}\]
  by strong collection and union axioms.
\end{proof}

It follows that we just need to realise the binary intersection axiom:
\begin{lemma}
  \label{lem:czf-real-in-czf-bin-inter}
  The binary intersection axiom is realised:
  \[\opA{a} \opA{b} \opE{y} \opA{x} \left(x \in y \leftrightarrow x \in a \land x \in b\right).\]
\end{lemma}

\begin{proof}
  Given $a, b \in V\mleft(\mathcal{A}\mright)$, consider
  \[c = \left\{\tuple{\bp fe, x} : \tuple{f, x} \in a, e \Vdash x \in b\right\} \in V\mleft(\mathcal{A}\mright).\]
  Note that $c = \bigcup \left\{\left\{\tuple{\bp fe, x} : e \in F_\in\mleft(x, b\mright)\right\} : \tuple{f, x} \in a\right\}$ and hence a set. Observe that
  \[\left(\bp \mleft(\bp fg\mright) d \Vdash x \in c\right) \leftrightarrow \left(\bp fd \Vdash x \in a\right) \land \left(\bp \mleft(\bp_0 g\mright) \mleft(\iot d \mleft(\bp_1 g\mright)\mright) \Vdash x \in b\right),\]
  and thus
  \begin{align*}
    \oplam{d} \bp \mleft(\bp \mleft(\bp_{00} d\mright) \mleft(\bp_1 d\mright)\mright) \mleft(\bp \mleft(\bp_{010} d\mright) \mleft(\iot \mleft(\bp_1 d\mright) \mleft(\bp_{011} d\mright)\mright)\mright) & \Vdash x \in c \rightarrow x \in a \land x \in b, \\
    \oplam{d} \bp \mleft(\bp \mleft(\bp_{00} d\mright) \mleft(\ioso \mleft(\bp_{01} d\mright) \mleft(\bp_1 d\mright)\mright)\mright) \mleft(\bp_{01} d\mright)                                            & \Vdash x \in a \land x \in b \rightarrow x \in c
  \end{align*}
  as desired.
\end{proof}

Combining the results so far and \autoref{prop:real-con-int-logic}, we can conclude that:
\begin{theorem}[\cite{rathjen06-czf-realizability}*{Theorem 5.1}]
  Let $\varphi$ be any theorem of $\CZF$, then $\CZF$ proves $V\mleft(\mathcal{A}\mright) \vDash \varphi$, i.e.\ there exists $d \in \mathcal{A}$ such that $d \Vdash \varphi$.
\end{theorem}

\section{Intuitionistic \texorpdfstring{$L$}{L}}
\label{sec:intro-int-l}

Just as in classical set theory, we say that a set $x$ is \emph{transitive} if every element of $x$ is also a subset of $x$, i.e.
\[\opA{y \in x} \opA{z \in y} z \in x,\]
and we define the \emph{ordinals}, $\Ord$\inds{ordinals, $\Ord$}, as the class of transitive sets with transitive elements.

However, many important properties of ordinals in classical set theory rely on the law of excluded middle, and hence cannot be recovered in intuitionistic contexts. For example, we have:
\begin{lemma}
  \label{lem:ord-trichotomy-fail}
  The realisability model of $\CZF$ in $\CZF$ realises that
  \[\neg\!\opA{\alpha, \beta \in \Ord} \left(\alpha \subseteq \beta \rightarrow \alpha \in \beta \lor \alpha = \beta\right).\]
\end{lemma}

\begin{proof}
  It suffices to find a counterexample, i.e.\ $\alpha, \beta \in V\mleft(\mathcal{A}\mright)$ such that $\alpha \in \Ord$, $\beta \in \Ord$ and $\alpha \subseteq \beta$ are realised, but no $d \in \mathcal{A}$ realises $\alpha \in \beta$ or $\alpha = \beta$. We shall take
  \[\alpha = \overline{2} = \left\{\tuple{0, \overline{0}}, \tuple{1, \overline{1}}\right\}, \qquad \beta = \left\{\tuple{0, \overline{0}}, \tuple{0, \overline{1}}\right\}.\]

  It is easy to verify that $\opA{n \in \omega} n \in \Ord$ still holds in $\CZF$ (or even $\IKP$), hence $\alpha \in \Ord$ must be realised. In the case of $\beta$, simply observe that if $\tuple{a, x} \in \beta$ and $\tuple{b, y} \in x$, then we must have $x = \overline{1}$ and $y = \varnothing$, so $\bp 0 \mleft(\bp 00\mright) \Vdash y \in \beta$ and thus $\oplam{ab} \bp 0 \mleft(\bp 00\mright) \Vdash \text{$\beta$ is transitive}$. Since both $\overline{0}$ and $\overline{1}$ contain no elements other than empty sets (as seen in the realisability model), $\oplam{a} 0$ trivially realises that both are transitive. Putting them together,
  \[\bp \mleft(\oplam{ab} \bp 0 \mleft(\bp 00\mright)\mright) \mleft(\oplam{xa} 0\mright) \Vdash \beta \in \Ord.\]

  Now, it is easy to see that $\oplam{d} \bp 0 \ior \Vdash \alpha \subseteq \beta$. We also observe that no $d \in \mathcal{A}$ can realise $a = \overline{0}$ where $a \in V\mleft(\mathcal{A}\mright)$ is inhabited, thus no $d$ can realise $\alpha = \overline{0}$, nor $\alpha = \overline{1}$ (in the latter case, since $\tuple{1, \overline{1}} \in \alpha$ and the sole element of $\overline{1}$ is $\tuple{0, \overline{0}}$, thus $\bp_1 \mleft(\bp_0 d 1\mright)$ will then realise $\overline{1} = \overline{0}$). Therefore, no $d$ can realise $\alpha \in \beta$.

  Finally, suppose that some $f \Vdash \beta \subseteq \alpha$. Using the fact that no $d$ realises $\overline{0} = \overline{1}$, we know that $f$ must match up $\overline{0}$ with $\overline{0}$ and $\overline{1}$ with $\overline{1}$, i.e.\ $\bp_0 \mleft(f0\mright) = 0$ for $\tuple{0, \overline{0}} \in \beta$ and $\bp_0 \mleft(f0\mright) = 1$ for $\tuple{0, \overline{1}} \in \beta$. But then $0 = 1$ in the base theory, a contradiction. We conclude that $\alpha = \beta$ is not realised either.
\end{proof}

Intuitionistic set theory also does not prove the trichotomy that every ordinal is either zero, a successor or a limit --- simply see $\beta$ in the proof above and observe that it is not a successor.

Thus, the traditional definition for transfinite hierarchies like G\"odel's constructible universe $L$\inds{G\"odel's constructible universe, $L$} does not work, and we need to define them using a uniform clause instead. Namely, it is conventional to stipulate
\[L_\alpha = \bigcup_{\beta \in \alpha} \dd\mleft(L_\beta\mright)\]
for all $\alpha \in \Ord$ in intuitionistic settings, and then take the union $L = \bigcup_{\alpha \in \Ord} L_\alpha$ as usual. Here $\dd\mleft(x\mright)$ still intends to mean the definable subsets of the first-order structure $\tuple{x; \in}$ in the language of set theory.

Robert Lubarsky \cite{lubarsky93-intuitionistic-l} first studied intuitionistic $L$ and checked that constructions related to $\dd\mleft(x\mright)$ work as expected in intuitionistic set theory. Essentially, we consider the set of finite tuples
\[x^{< \omega} = \bigcup_{n \in \omega} x^n\]
and construct the truth predicates $\Tr_{\god{\varphi}}$, where $\god{\varphi} \in \omega$ is the G\"odel numbering of formula $\varphi$, internally as\footnote{We take a simplified approach with the guarding conditions on the length of $\tau$ here, which says precisely that when evaluating a formula $\varphi$ on a finite set of variable assignments $\tau$, any atomic or quantifier subformula containing a variable not occurring in the assignment is simply treated as $\bot$ --- even for bound occurrences. It is easy to see that this semantic simplification does not affect the expressive power of first-order logic.}
\begin{align*}
  \Tr_{\god{x_i \in x_j}}              & = \left\{\tau \in x^{< \omega} : \left|\tau\right| > i, j, \tau_i \in \tau_j\right\},                                         \\
  \Tr_{\god{x_i = x_j}}                & = \left\{\tau \in x^{< \omega} : \left|\tau\right| > i, j, \tau_i = \tau_j\right\},                                           \\
  \Tr_{\god{\varphi \land \psi}}       & = \Tr_{\god{\varphi}} \cap \Tr_{\god{\psi}},                                                                                  \\
  \Tr_{\god{\varphi \lor \psi}}        & = \Tr_{\god{\varphi}} \cup \Tr_{\god{\psi}},                                                                                  \\
  \Tr_{\god{\neg \varphi}}             & = x^{< \omega} \setminus \Tr_{\god{\varphi}},                                                                                 \\
  \Tr_{\god{\varphi \rightarrow \psi}} & = \left\{\tau \in x^{< \omega} : \tau \in \Tr_{\god{\varphi}} \rightarrow \tau \in \Tr_{\god{\psi}}\right\},                  \\
  \Tr_{\god{\opA{x_i} \varphi}}        & = \left\{\tau \in x^{< \omega} : \left|\tau\right| > i, \opA{y \in x} \tau\mleft[y/i\mright] \in \Tr_{\god{\varphi}}\right\}, \\
  \Tr_{\god{\opE{x_i} \varphi}}        & = \left\{\tau \in x^{< \omega} : \left|\tau\right| > i, \opE{y \in x} \tau\mleft[y/i\mright] \in \Tr_{\god{\varphi}}\right\}
\end{align*}
by $\Sigma$-recursion on natural numbers, where $\tau\mleft[y/i\mright]$ denotes the result of replacing the $i^\text{th}$ component of finite tuple $\tau$ by object $y$. Each $\Tr_{\god{\varphi}} \subseteq x^{< \omega}$ is a set by $\Delta_0$-separation. One can then define
\[\dd\mleft(x\mright) = \left\{\left\{y \in x : y^\frown \tau \in \Tr_{\god{\varphi}}\right\} : \text{$\varphi$ is a formula}, \tau \in x^{< \omega}\right\}\]
by $\Delta_0$-separation and collection.

$\IKP$ suffices for this construction, where we observe that the axiom of infinity is crucial for encoding the set of first-order formulae and thus the recursive process. Recently in \cite{matthews-rathjen24-constructible-universe}, Matthews and Rathjen studied the alternative formulation of $L$ using Barwise's fundamental operations in \cite{barwise75-admissible-sets}*{section II.6} in absence of the axiom of infinity, and verified that the two approaches still coincide in intuitionistic set theory. For the topics examined in this thesis, however, Lubarsky's construction suffices and is more convenient.

A first observation about intuitionistic $L$ is that the following classical proof makes an essential use of the property invalidated in \autoref{lem:ord-trichotomy-fail}:
\begin{lemma}[$\ZF$ or $\KP$]
  \label{lem:classical-ord-in-l}
  For any $\alpha \in \Ord$, $L_\alpha \cap \Ord = \alpha$.
\end{lemma}

\begin{proof}
  We prove this by set induction on $\alpha$. By the inductive hypothesis, $L_\beta \cap \Ord = \beta$ for any $\beta \in \alpha$. Since $\Ord$ is a $\Delta_0$-class, it follows that $\beta \in \dd\mleft(L_\beta\mright) \subseteq L_\alpha$ and hence $\alpha \subseteq L_\alpha \cap \Ord$.

  For the converse, suppose that $\gamma \in L_\alpha \cap \Ord$. There is some $\beta \in \alpha$ such that $\gamma \in \dd\mleft(L_\beta\mright)$, which means $\gamma \subseteq L_\beta \cap \Ord = \beta$. Using the classical fact that
  \[\opA{\alpha, \beta \in \Ord} \left(\alpha \subseteq \beta \rightarrow \alpha \in \beta \lor \alpha = \beta\right),\]
  it follows that $\gamma \in \beta \lor \gamma = \beta$, and thus $\gamma \in \alpha$.
\end{proof}

In \autoref{sec:ord-modal}, we will explain a formal proof of the fact that the opposite is indeed consistent with $\CZF$. This implies that there is no longer a canonical way to make every $\alpha \in \Ord$ a definable set --- thus an element --- in $L$. The following question, first raised by Lubarsky in \cite{lubarsky93-intuitionistic-l}, remains open:
\begin{problem}
\label{prob:ord-in-l}
Does any non-classical system of set theory prove $\Ord \subseteq L$?
\end{problem}

On the other hand, Lubarsky showed in \cite{lubarsky93-intuitionistic-l} that several other properties of classical $L$ can still be recovered in intuitionistic logic, even though the proof can be much more complicated than in the classical case. One key result is that:
\begin{theorem}
  \label{thm:l-omega-dec}
  Any $\Delta_0$-formula $\varphi$ is decidable on $L_\omega$, i.e.\ let $x_1, \ldots, x_k$ enumerate the free variables in $\varphi$, then $\CZF$ \uppar{or even $\IKP$} proves
  \[\opA{x_1, \ldots, x_k \in L_\omega} \left(\varphi \lor \neg \varphi\right).\]
\end{theorem}

\begin{proof}[Proof sketch]
  By simultaneous induction on $\omega$, one can still show the trichotomy for natural numbers
  \[\opA{n, m \in \omega} \left(n \in m \lor n = m \lor m \in n\right),\]
  and this decidability is preserved by logical connectives and natural-number-bounded quantifiers. It follows that any $\Delta_0$-formula evaluated in the structure $\tuple{\omega; \in}$, which can possibly contain $\Sigma_1$-definable functions such as arithmetic operations, satisfies excluded middle.

  Now, we build the (reverse) Ackermann encoding, a function $f$ on domain $\omega$, given as
  \[f\mleft(n\mright) = \left\{f\mleft(m\mright) : m \in n, n \mathbin{\mathrm{mod}} 2^{m + 1} \geq 2^m\right\}.\]
  One can verify by induction on $n \in \omega$ that every $L_n$ is precisely the image of $f$ on an appropriate initial segment of $\omega$, because given the inductive hypothesis, the map $f$ then allows us to translate any formula evaluated on $L_n$ into a $\Delta_0$-formula evaluated on $\omega$, which is decidable.

  Through this induction, we see that $f$ is a bijection between $\omega$ and $L_\omega$. Thus, every $\Delta_0$-formula on $L_\omega$ translates into a $\Delta_0$-formula with arithmetic operations on $\omega$, and is decidable as desired.

  We refer the readers to \cite{lubarsky93-intuitionistic-l} for more detail.
\end{proof}

\begin{corollary}
  \label{cor:l-omega-ord}
  For any $n \in \omega$, $L_n \cap \Ord = n$. It follows immediately that $L_\omega \cap \Ord = \omega$.
\end{corollary}

\begin{proof}
  Suppose that $\alpha \in L_\omega$ and $n \in \omega \cap L_\omega$ satisfies $\alpha \subseteq n$. We know from \autoref{thm:l-omega-dec} that
  \[\opE{k \in n} k \not\in \alpha \lor \opA{k \in n} k \in \alpha.\]
  The latter disjunct implies $\alpha = n$; in the former case, we show that $\alpha \subseteq k$. Suppose that this is not the case, i.e.\ there is some $\ell \in \alpha$ such that $\ell \not\in k$. However, both $k, \ell \in \omega$, thus by the trichotomy for natural numbers, we have either $k \in \ell$ or $k = \ell$, both contradicting $k \not\in \alpha$. Thus we must have $\alpha \subseteq k$ by decidability from \autoref{thm:l-omega-dec}.

  Hence, by induction on $n \in \omega$, we see that
  \[\opA{n \in \omega \cap L_\omega} \opA{\alpha \in L_\omega} \left(\alpha \subseteq n \rightarrow \alpha \in n \lor \alpha = n\right).\]
  The proof of \autoref{lem:classical-ord-in-l} now applies to give $\opA{n \in \omega} L_n \cap \Ord = n$.
\end{proof}

Another result of Lubarsky \cite{lubarsky93-intuitionistic-l} that is relevant is the following:
\begin{theorem}
  The relativised statement $\left(V = L\right)^L$ holds.
\end{theorem}

Note that this is proved without assuming that $\Ord \subseteq L$, unlike in classical set theory. For any $\alpha \in \Ord$, Lubarsky constructed an ordinal
\[\alpha^* = \left\{\beta \in \Ord : L_{\left(\alpha_\text{aug} +_\mathrm{H} n\right)^-} \vDash \opE{x} \left(x = L_\beta \land \dd\mleft(x\mright) \subseteq L_\alpha\right)\right\},\]
where $\alpha_\text{aug} = \bigcup \left\{\beta_\text{aug} : \beta \in \alpha\right\} \cup \left\{\alpha\right\} \cup \left\{\omega^+\right\}$, $\left(\alpha_\text{aug} +_\mathrm{H} n\right)^-$ is a variation of \emph{hereditary addition} (cf.\ \autoref{sec:incom-intro} for normal ordinal addition) given in terms of the standard variant as
\begin{align*}
  \alpha +_\mathrm{H} \gamma                & = \left(\alpha \cup \bigcup_{\beta \in \alpha} \left(\beta +_\mathrm{H} \gamma\right)\right) + \gamma, \\
  \left(\alpha +_\mathrm{H} \gamma\right)^- & = \alpha \cup \bigcup_{\beta \in \alpha} \left(\beta +_\mathrm{H} \gamma\right),
\end{align*}
and $n \in \omega$ is fixed large enough for technical existential witnesses. Since the only parameter $L_\alpha$ itself is a set in $L$, it follows that $\alpha^* \in L \cap \Ord$ and Lubarsky verified that $L_{\alpha^*} = L_\alpha$ to derive $L \vDash V = L$.

\section{Some exotic ordinals and a modality}
\label{sec:ord-modal}

Before ending this chapter, we will reason with $L$ inside the realisability model in \autoref{sec:czf-real-in-czf} to formally demonstrate that intuitionistic $L$ can indeed have wild behaviours beyond the decidable segment $L_\omega$. We shall show that $\CZF$ is consistent with the statement $\opE{\alpha \in \Ord} \neg L_\alpha \cap \Ord = \alpha$. We also see that the idea above of `constructing $\alpha^*$ so that $L_{\alpha^*} = L_\alpha$' is non-trivial, as it is consistent to have $\alpha \neq \beta$ yet $L_\alpha = L_\beta$.

To begin, we construct a family of names for `$\omega$-like' intuitionistic ordinals. For any $u \subseteq \omega$ in the base theory, we denote
\[\overline{\omega}_{(u)} = \left\{\tuple{n, \overline{n}} : n \in \omega\right\} \cup \left\{\tuple{0, \overline{n}} : n \in u\right\} \in V\mleft(\mathcal{A}\mright).\]
Notice that whether $0 \in u$ is irrelevant, as $\overline{\omega}_{(u)} = \overline{\omega}_{(u \cup \left\{0\right\})}$. Thus in most scenarios we will just work with sets $u \subseteq \omega$ satisfying $0 \in u$.

\begin{lemma}
  \label{lem:exo-omega-is-ord}
  For any $u \subseteq \omega$, $V\mleft(\mathcal{A}\mright) \vDash \overline{\omega}_{(u)} \in \Ord$.
\end{lemma}

\begin{proof}
  Simply observe that $\oplam{ab} \bp b \ior \Vdash \text{$\overline{\omega}_{(u)}$ is transitive}$, where the same realiser also realises $\overline{n}$ is transitive uniformly for any $n \in \omega$. Thus
  \[\bp \mleft(\oplam{ab} \bp b\mright) \mleft(\oplam{xab} \bp b \ior\mright) \Vdash \overline{\omega}_{(u)} \in \Ord. \qedhere\]
\end{proof}

\begin{lemma}
  \label{lem:exo-omega-subset}
  For any $u, v \subseteq \omega$ such that $0 \in u, v$, we have $V\mleft(\mathcal{A}\mright) \vDash \overline{\omega}_{(u)} \subseteq \overline{\omega}_{(v)}$ if and only if $u \subseteq v$.
\end{lemma}

\begin{proof}
  The backward direction is immediate, as $\oplam{n} \bp n \ior$ realises $\overline{\omega}_{(u)} \subseteq \overline{\omega}_{(v)}$ when $u \subseteq v$ holds.

  For the forward direction, observe by induction on $\omega$ that for any $n, m \in \omega$, $V\mleft(\mathcal{A}\mright) \vDash \overline{n} = \overline{m}$ if and only if $n = m$. Now, suppose that some $f \in \mathcal{A}$ realises $\overline{\omega}_{(u)} \subseteq \overline{\omega}_{(v)}$. The sole index $i \in \omega$ such that $\tuple{i, \overline{0}} \in \overline{\omega}_{(u)}$ is $i = 0$, and the same holds for $\overline{\omega}_{(v)}$. Thus, $f$ must match these indices up, i.e.\ $\bp_0 \mleft(f0\mright) = 0$. It follows that for any $n \in u$, i.e.\ $\tuple{0, \overline{n}} \in \overline{\omega}_{(u)}$, there is some $\tuple{j, \overline{m}} \in \overline{\omega}_{(v)}$ such that
  \[j = \bp_0 \mleft(f0\mright) = 0,\]
  and $\bp_1 \mleft(f0\mright) \Vdash \overline{n} = \overline{m}$, i.e.\ $m = n$. Hence $\tuple{0, \overline{n}} \in \overline{\omega}_{(v)}$.

  It follows that for any $n \in u$, if $n \neq 0$ then $n \in v$. Using the assumption that $0 \in u, v$, we have $u \subseteq v$.
\end{proof}

Most notably, $\overline{\omega}_{(\varnothing)} = \overline{\omega}_{(\left\{0\right\})} = \overline{\omega}$ is the least element in this family. In this section, we will focus on the behaviour of the opposite extreme, $\overline{\omega}_{(\omega)}$.

\subsection*{Realisability with a modal operator}

To completely characterise $\overline{\omega}_{(\omega)}$, we need to expand the formal language and introduce a modal operator $\Box$, so that $V\mleft(\mathcal{A}\mright) \vDash \phi \rightarrow \Box \psi$ means roughly `given a realiser $d \Vdash \phi$, I know that there will be a realiser for $\psi$, but I do not know how to compute it from $d$'. Strictly speaking, this is neither an instance of necessity nor possibility (cf.\ \autoref{lem:dn-modal-axioms} below), but more closely related to Kosta Do\v{s}en's system of modal logic for the intuitionistic double negation\indg{double negation as modality} in \cite{dosen84-intuitionistic-double-negation}, from which I took the inspiration.

Formally, we simply expand the first-order language by a new clause that, whenever $\varphi$ is a formula in the modal language, so is $\Box \varphi$. We extend the realisability conditions in \autoref{sec:czf-real-in-czf} by
\[d \Vdash \Box \varphi \defeqv \opE{e \in \mathcal{A}} e \Vdash \varphi.\]

\begin{lemma}
  \label{lem:dn-modal-axioms}
  The following modal postulates are realised for any formulae $\varphi$ and $\psi$:
  \begin{enumerate}[label=\uplabel{\arabic*}]
    \item $\Box \left(\varphi \rightarrow \psi\right) \rightarrow \left(\Box \varphi \rightarrow \Box \psi\right)$;
    \item $\varphi \rightarrow \Box \varphi$;
    \item $\Box \varphi \rightarrow \neg \neg \varphi$.
  \end{enumerate}
\end{lemma}

\begin{proof}
  Notice nothing needs to be computed to realise any assertion of the form $\Box \varphi$ that is `correct' (just like negated assertions). Thus:
  \begin{enumerate}[label=\uplabel{\arabic*}]
    \item $\oplam{de} 0$ realises $\Box \left(\varphi \rightarrow \psi\right) \rightarrow \left(\Box \varphi \rightarrow \Box \psi\right)$;
    \item $\oplam{d} 0$ realises $\varphi \rightarrow \Box \varphi$;
    \item $\oplam{d} 0$ realises $\Box \varphi \rightarrow \neg \neg \varphi$. \qedhere
  \end{enumerate}
\end{proof}

If the underlying base theory is classical, e.g.\ when the aforementioned realisability of $\CZF$ is performed in the base theory $\ZF$, we additionally realise $\Box \varphi \leftrightarrow \neg \neg \varphi$ trivially, since classically,
\[\left(\opE{e \in \mathcal{A}} e \Vdash \varphi\right) \leftrightarrow \left(\neg\!\opA{e \in \mathcal{A}} \neg e \Vdash \varphi\right).\]
Then we recover precisely the modality considered in Do\v{s}en \cite{dosen84-intuitionistic-double-negation}.

Meanwhile, returning to the base theory $\CZF$, we check that:
\begin{proposition}
  \label{prop:exo-omega-def-by-box}
  It is realised that
  \[x \in \overline{\omega}_{(\omega)} \leftrightarrow \left(x \subseteq \overline{\omega} \land \Box x \in \overline{\omega}\right) \land \opA{y \in \overline{\omega}} \left(\Box y \in x \rightarrow y \in x\right).\]
\end{proposition}

\begin{proof}
  For the forward direction, we just need to equivalently realise
  \[\opA{x \in \overline{\omega}_{(\omega)}} \left(x \subseteq \overline{\omega} \land \Box x \in \overline{\omega} \land \opA{y \in \overline{\omega}} \left(\Box y \in x \rightarrow y \in x\right)\right).\]
  Observe that $\oplam{n} \bp n \ior \Vdash \overline{m} \subseteq \overline{\omega}$ for any $m \in \omega$, thus $\oplam{xn} \bp n \ior \Vdash \opA{x \in \overline{\omega}_{(\omega)}} x \subseteq \overline{\omega}$. Next, $\oplam{x} 0$ trivially realises $\opA{x \in \overline{\omega}_{(\omega)}} \Box x \in \overline{\omega}$, since every element in $\overline{\omega}_{(\omega)}$ is of the form $\tuple{n, \overline{m}}$. Finally, for any $n, m \in \omega$, $\oplam{d} \bp n \ior \Vdash \Box \overline{n} \in \overline{m} \rightarrow \overline{n} \in \overline{m}$, thus
  \[\oplam{xnd} \bp n \ior \Vdash \opA{x \in \overline{\omega}_{(\omega)}} \opA{y \in \overline{\omega}} \left(\Box y \in x \rightarrow y \in x\right).\]

  Now, for the backward direction, fix $x \in V\mleft(\mathcal{A}\mright)$ and suppose that $f$ realises the entire right-hand side of the biconditional. It follows from the realisability of the modality that there exists some $d \in \mathcal{A}$ such that $d \Vdash x \in \overline{\omega}$, i.e.\ $\bp_1 d \Vdash x = \overline{\bp_0 d}$. We derive from $\bp_{00} f \Vdash x \subseteq \overline{\omega}$ that for any $\tuple{g, y} \in x$, $\bp_{00} fg \Vdash y \in \overline{\omega}$, i.e.
  \[\bp_1 \mleft(\bp_{00} fg\mright) \Vdash y = \overline{m}\]
  where $m = \bp_0 \mleft(\bp_{00} fg\mright)$. Using that $x = \overline{\bp_0 d}$ is realised with the fact that $\overline{m} \in \overline{n}$ is realised if and only if $m \in n$, we must have $m \in \bp_0 d$ and thus $\bp_{00} fg$ also realises $y \in \overline{\bp_0 d}$, i.e.\ $\bp_{00} f \Vdash x \subseteq \overline{\bp_0 d}$.

  On the other hand, since $\bp_1 d \Vdash x = \overline{\bp_0 d}$, for any $m \in \bp_0 d$, $\bp_{11} dm \Vdash \overline{m} \in x$. It follows that $0 \Vdash \Box \overline{m} \in x$ and thus $\bp_1 fm0 \Vdash \overline{m} \in x$. Therefore, $\oplam{m} \bp_1 fm0 \Vdash \overline{\bp_0 d} \subseteq x$, and we can conclude that
  \[\bp \mleft(\bp_{00} f\mright) \mleft(\oplam{m} \bp_1 fm0\mright) \Vdash x = \overline{\bp_0 d}\]
  where the realiser does not involve $d$. Since $\tuple{0, \overline{\bp_0 d}} \in \overline{\omega}_{(\omega)}$, we have
  \begin{align*}
     & \oplam{f} \bp 0 \mleft(\bp \mleft(\bp_{00} f\mright) \mleft(\oplam{m} \bp_1 fm0\mright)\mright)                                                                                                                                           \\
     & \qquad {} \Vdash \left(x \subseteq \overline{\omega} \land \Box x \in \overline{\omega}\right) \land \opA{y \in \overline{\omega}} \left(\Box y \in x \rightarrow y \in x\right) \rightarrow x \in \overline{\omega}_{(\omega)}. \qedhere
  \end{align*}
\end{proof}

This will be useful as the defining formula for $\overline{\omega}_{(\omega)}$ when we reason inside $V\mleft(\mathcal{A}\mright)$. However, we cannot use a formula containing modal operators to show definability inside $L$, thus we see that one of the $\Box$'s can be dropped:

\begin{lemma}
  \label{lem:exo-omega-box-dneg-equiv}
  $V\mleft(\mathcal{A}\mright) \vDash \opA{\alpha \in \overline{\omega}_{(\omega)}} \opA{n \in \overline{\omega}} \left(\neg \neg n \in \alpha \rightarrow \Box n \in \alpha\right)$.
\end{lemma}

\begin{proof}
  We use the fact that $\CZF$ proves the trichotomy for natural numbers, which implies
  \[\opA{m \in \omega} \opA{n \in \omega} \left(\neg \neg n \in m \rightarrow n \in m\right).\]
  Using the modal postulates in \autoref{lem:dn-modal-axioms} (and \autoref{lem:czf-real-omega-least-ind} that $\overline{\omega}$ is indeed the least inductive set), it follows that $V\mleft(\mathcal{A}\mright)$ realises
  \[\Box m \in \overline{\omega} \rightarrow \left(\Box n \in \overline{\omega} \rightarrow \left(\Box \neg \neg n \in m \rightarrow \Box n \in m\right)\right).\]

  Working inside $V\mleft(\mathcal{A}\mright)$, we consider any $\alpha \in \overline{\omega}_{(\omega)}$. It follows from \autoref{prop:exo-omega-def-by-box} that $\Box \alpha \in \overline{\omega}$.
  For any $n \in \overline{\omega}$, suppose $\neg \neg n \in \alpha$, then we have $\Box n \in \overline{\omega}$ and $\Box \neg \neg n \in \alpha$ by \autoref{lem:dn-modal-axioms}, and thus $\Box n \in \alpha$ as desired.
\end{proof}

\begin{corollary}
  \label{cor:exo-omega-def-by-dneg}
  It is realised that
  \[x \in \overline{\omega}_{(\omega)} \leftrightarrow \left(x \subseteq \overline{\omega} \land \Box x \in \overline{\omega}\right) \land \opA{y \in \overline{\omega}} \left(\neg \neg y \in x \rightarrow y \in x\right).\]
\end{corollary}

\begin{proof}
  Using \autoref{prop:exo-omega-def-by-box} as a basis, for the forward direction, it suffices to derive $\Box y \in x$ given $x \in \overline{\omega}_{(\omega)}$ and $y \in \overline{\omega}$ such that $\neg \neg y \in x$ holds. This follows from \autoref{lem:exo-omega-box-dneg-equiv}.

  For the backward direction, given $x \subseteq \overline{\omega}$ such that $\Box x \in \overline{\omega}$, it suffices to derive $\opA{y \in \overline{\omega}} \left(\Box y \in x \rightarrow y \in x\right)$ from $\opA{y \in \overline{\omega}} \left(\neg \neg y \in x \rightarrow y \in x\right)$. This follows from the modal postulate $\Box \varphi \rightarrow \neg \neg \varphi$ in \autoref{lem:dn-modal-axioms}.
\end{proof}

Before relating $\overline{\omega}_{(\omega)}$ to the constructible universe $L$, we need a technical lemma:

\begin{lemma}[$\IKP$]
  \label{lem:l-omega-can-compute-rank}
  $L$-ranks can be defined for elements in $L_\omega$, i.e.\ for any $x \in L_\omega$, there exists a natural number $k \in \omega$ such that
  \[\opA{\ell \in \omega} \left(x \in L_\ell \leftrightarrow k \in \ell\right).\]
\end{lemma}

\begin{proof}
  For any $x \in L_\omega$, there is $n \in \omega$ such that $x \in \dd\mleft(L_n\mright)$. We show
  \[\opA{x \in \dd\mleft(L_n\mright)} \opE{k \in \omega} \opA{\ell \in \omega} \left(x \in L_\ell \leftrightarrow k \in \ell\right)\]
  by induction on $n$.

  When $x \in \dd\mleft(L_0\mright)$, then $x \subseteq L_0 = \varnothing$ and thus $x = \varnothing$. Clearly we have $\opA{\ell \in \omega} \left(x \in L_\ell \leftrightarrow 0 \in \ell\right)$. Otherwise, $n = m^+$ is a successor, and for $x \in \dd\mleft(L_n\mright)$, by \autoref{thm:l-omega-dec} we know that either $x \in L_n = \dd\mleft(L_m\mright)$ or $x \in \dd\mleft(L_n\mright) \setminus L_n$. The former case is covered by the inductive hypothesis, whereas in the latter case, we use the trichotomy for natural numbers to see that
  \[\opA{\ell \in \omega} \left(x \in L_\ell \leftrightarrow n \in \ell\right),\]
  since $n \not\in \ell$ implies $\ell \subseteq n$ and thus $x \not\in L_\ell \subseteq L_n$.
\end{proof}

\begin{corollary}
  \label{cor:exo-omega-l-compute-witness}
  $V\mleft(\mathcal{A}\mright) \vDash \opA{\alpha \in \overline{\omega}_{(\omega)}} \opA{x \in L_{\overline{\omega}}} \left(\Box x \in L_\alpha \rightarrow x \in L_\alpha\right)$.
\end{corollary}

\begin{proof}
  Work inside $V\mleft(\mathcal{A}\mright)$ and consider any $\alpha \in \overline{\omega}_{(\omega)}$ and any $x \in L_{\overline{\omega}}$. By \autoref{lem:l-omega-can-compute-rank}, we can find $k \in \overline{\omega}$ such that $\opA{\ell \in \overline{\omega}} \left(x \in L_\ell \leftrightarrow k \in \ell\right)$, thus from the modal postulates,
  \[\Box \ell \in \overline{\omega} \rightarrow \left(\Box x \in L_\ell \leftrightarrow \Box k \in \ell\right).\]
  We have $\Box \alpha \in \overline{\omega}$ by \autoref{prop:exo-omega-def-by-box}, so assuming $\Box x \in L_\alpha$, we can derive $\Box k \in \alpha$. The last conjunct in \autoref{prop:exo-omega-def-by-box} then implies that $k \in \alpha$. It follows that $k^+ \subseteq \alpha$ and we have $x \in L_{k^+} \subseteq L_\alpha$.
\end{proof}

\subsection*{A family of coinciding $L_{\overline{\omega}_{(u)}}$'s}

Consider a $\Delta_0$-formula $\varphi\mleft(x, y\mright)$ in the language of set theory --- without the modal operator $\Box$ --- with no other free variables, such that
\[\CZF \vdash \opA{x \in \omega} \opE{!y \in L_\omega} \varphi\mleft(x, y\mright).\]
We say that a set $u \subseteq \omega$ \emph{encompasses $\omega$ through $\varphi$} if $0 \in u$ and, for any $n \in \omega$, there exists $m \in u$ such that
\[V\mleft(\mathcal{A}\mright) \vDash \left\{x \in L_{\overline{m}} \cap \Ord : \opE{y \in L_{\overline{m}}} \varphi\mleft(x, y\mright)\right\} = \overline{n}.\]

Informally speaking, this definition means that all elements of $\omega$ can be defined from elements of $u$ in a uniform way that is easy to implement in $L$. Clearly $\omega$ encompasses itself by taking
\[\varphi\mleft(x, y\mright) \defeqv x = x,\]
so that $V\mleft(\mathcal{A}\mright) \vDash \left\{x \in L_{\overline{n}} \cap \Ord : \opE{y \in L_{\overline{n}}} \varphi\mleft(x, y\mright)\right\} = \overline{n}$ for all $n \in \omega$. We first see that there are many other subsets of $\omega$ that encompass $\omega$:
\begin{lemma}
  Let $u \subseteq \omega$ be any subset containing all the even numbers. Any such $u$ encompasses $\omega$ through the same formula.
\end{lemma}

\begin{proof}
  Take $\varphi\mleft(x, y\mright)$ to express the following: $x$ is an ordinal and $y$ is a function on domain $x^+$, satisfying $y\mleft(\varnothing\mright) = x$ and $y\mleft(z^+\mright) = y\mleft(z\mright)^+$ for every $z \in x$. By induction on $\omega$, such $y$ exists and must be unique for any $x \in \omega$. Given the definition of Kuratowski pairs, we have
  \[x \cdot 2 \in \left\{x, x \cdot 2\right\} \in \tuple{x, x \cdot 2} \in y.\]
  It follows that
  \[\opA{n \in \omega} \left\{x \in L_{n \cdot 2 + 2} \cap \Ord : \opE{y \in L_{n \cdot 2 + 2}} \varphi\mleft(x, y\mright)\right\} = n\]
  is a theorem in $\CZF$. Since it is easy to check that $V\mleft(\mathcal{A}\mright)$ realises the correct arithmetic relations between names $\overline{n}$ and $\overline{n \cdot 2 + 2}$, we must have
  \[V\mleft(\mathcal{A}\mright) \vDash \left\{x \in L_{\overline{n \cdot 2 + 2}} \cap \Ord : \opE{y \in L_{\overline{n \cdot 2 + 2}}} \varphi\mleft(x, y\mright)\right\} = \overline{n}. \qedhere\]
\end{proof}

We now define a new formula
\[\Psi\mleft(x\mright) \defeqv x \in \Ord \land \opA{y \in \omega} \left(\neg \neg y \in x \rightarrow y \in x\right)\]
where, pedantically, we do not want to think of $\omega$ as a parameter to this formula, but rather read the quantifier $\opA{y \in \omega}$ as letting $y$ range over the $\Delta_0$-class for the least inductive set given in \autoref{lem:ikp-omega-from-inf}. We verify that:
\begin{theorem}
  \label{thm:exo-omega-l-con-max-omega}
  For any $u \subseteq \omega$ that encompasses $\omega$ through some $\Delta_0$-formula, we always realise
  \[\overline{\omega}_{(\omega)} = \left\{x \in L_{\overline{\omega}_{(u)}} : L_{\overline{\omega}_{(u)}} \vDash \Psi\mleft(x\mright)\right\} \in \dd\mleft(L_{\overline{\omega}_{(u)}}\mright).\]
\end{theorem}

\begin{proof}
  First note that $\overline{\omega} \subseteq \overline{\omega}_{(u)}$ by \autoref{lem:exo-omega-subset}, thus $\overline{\omega} \subseteq L_{\overline{\omega}_{(u)}}$ by \autoref{cor:l-omega-ord}, where $L_{\overline{\omega}_{(u)}}$ is a transitive set. Thus indeed we can define $\omega$ as a $\Delta_0$-definable subclass of $L_{\overline{\omega}_{(u)}}$ by \autoref{lem:ikp-omega-from-inf} and avoid using it as a parameter in $\Psi\mleft(x\mright)$.

  Now, for the forward direction, first observe that our assumption about $u$ encompassing $\omega$ implies that, for some $\Delta_0$-formula $\varphi\mleft(x, y\mright)$,
  \[\oplam{a} \bp 00 \Vdash \opA{\alpha \in \overline{\omega}_{(\omega)}} \opE{\beta \in \overline{\omega}_{(u)}} \Box \left\{x \in L_\beta \cap \Ord : \opE{y \in L_\beta} \varphi\mleft(x, y\mright)\right\} = \alpha.\]

  Work inside $V\mleft(\mathcal{A}\mright)$ and consider any $\alpha \in \overline{\omega}_{(\omega)}$. By \autoref{lem:exo-omega-is-ord} and \autoref{cor:exo-omega-def-by-dneg}, we know that $\Psi\mleft(\alpha\mright)$ holds, thus it suffices to show that $\alpha \in L_{\overline{\omega}_{(u)}}$. Our assumption for $u$ now means there is some $\beta \in \overline{\omega}_{(u)}$ such that $\Box \left\{x \in L_\beta \cap \Ord : \opE{y \in L_\beta} \varphi\mleft(x, y\mright)\right\} = \alpha$ holds and we know
  \[\left\{x \in L_\beta \cap \Ord : \opE{y \in L_\beta} \varphi\mleft(x, y\mright)\right\} \in \dd\mleft(L_\beta\mright) \subseteq L_{\overline{\omega}_{(u)}}.\]
  Thus it suffices to show that actually $\left\{x \in L_\beta \cap \Ord : \opE{y \in L_\beta} \varphi\mleft(x, y\mright)\right\} = \alpha$.

  We take any $x \in L_\beta \cap \Ord$ such that there exists $y \in L_\beta$ satisfying $\varphi\mleft(x, y\mright)$, then the modal postulates imply $\Box x \in \alpha$. Observe that $\beta \in \overline{\omega}_{(u)} \subseteq \overline{\omega}_{(\omega)}$ by \autoref{lem:exo-omega-subset} and thus $\beta \subseteq \omega$ by \autoref{prop:exo-omega-def-by-box}. Hence,
  \[x \in L_\beta \cap \Ord \subseteq L_\omega \cap \Ord = \omega\]
  by \autoref{cor:l-omega-ord}. Therefore, $\Box x \in \alpha$ implies $x \in \alpha$ by the last conjunct in \autoref{prop:exo-omega-def-by-box}.

  On the other hand, assume that $x \in \alpha$ holds. Since $\alpha \subseteq \omega$, we have $x \in \omega = L_\omega \cap \Ord$ by \autoref{cor:l-omega-ord} while the modal postulates give $\Box x \in L_\beta$. Here we know $\beta \in \overline{\omega}_{(u)} \subseteq \overline{\omega}_{(\omega)}$, and therefore \autoref{cor:exo-omega-l-compute-witness} would imply that $x \in L_\beta$ indeed.

  Likewise, our assumption on $\varphi$ implies that there is a unique $y \in L_\omega$ satisfying $\varphi\mleft(x, y\mright)$, while the modal postulates give
  \[\Box \left(\opE{y' \in L_\beta} \varphi\mleft(x, y'\mright)\right).\]
  Observe that $\opE{y' \in L_\beta} \varphi\mleft(x, y'\mright)$ would imply $y \in L_\beta$ by uniqueness of $y$, thus we instead have $\Box y \in L_\beta$ here. \autoref{cor:exo-omega-l-compute-witness} implies $y \in L_\beta$, so we have verified $x \in \left\{x \in L_\beta \cap \Ord : \opE{y \in L_\beta} \varphi\mleft(x, y\mright)\right\}$.

  For the backward direction, by \autoref{lem:exo-omega-subset}, $\overline{\omega}_{(u)} \subseteq \overline{\omega}_{(\omega)}$, thus $L_{\overline{\omega}_{(u)}} \subseteq L_{\overline{\omega}_{(\omega)}}$ and it suffices to show that
  \[V\mleft(\mathcal{A}\mright) \vDash \left\{x \in L_{\overline{\omega}_{(\omega)}} : L_{\overline{\omega}_{(\omega)}} \vDash \Psi\mleft(x\mright)\right\} \subseteq \overline{\omega}_{(\omega)}.\]

  Take any $x \in L_{\overline{\omega}_{(\omega)}}$ and assume that $\Psi\mleft(x\mright)$ holds. This means $x \in \dd\mleft(L_\alpha\mright)$ for some $\alpha \in \overline{\omega}_{(\omega)}$. It follows by \autoref{prop:exo-omega-def-by-box} that $\Box \alpha \in \omega$, and thus $\Box x \in L_\omega \cap \Ord$, i.e.\ $\Box x \in \omega$ by \autoref{cor:l-omega-ord}. Additionally, we know that $\alpha \subseteq \omega$, thus
  \[x \subseteq L_\alpha \cap \Ord \subseteq L_\omega \cap \Ord = \omega.\]
  We have shown that $x \subseteq \omega \land \Box x \in \omega$. Combined with the second conjunct of $\Psi\mleft(x\mright)$, we immediately know $x \in \overline{\omega}_{(\omega)}$ by \autoref{cor:exo-omega-def-by-dneg}.
\end{proof}

Take any $u \subsetneq \omega$ that encompasses $\omega$, then \autoref{lem:exo-omega-subset} implies that $V\mleft(\mathcal{A}\mright) \vDash \overline{\omega}_{(u)} \subsetneq \overline{\omega}_{(\omega)}$. By \autoref{thm:exo-omega-l-con-max-omega}, however, we have $\overline{\omega}_{(\omega)} \subseteq L_{\overline{\omega}_{(u)}}$. Therefore, it follows immediately that:
\begin{corollary}
  \label{cor:czf-real-alpha-l-alpha-mismatch}
  $V\mleft(\mathcal{A}\mright)$ realises $\opE{\alpha \in \Ord} \neg L_\alpha \cap \Ord = \alpha$.
\end{corollary}

We finally show that additionally, $L_{\overline{\omega}_{(u)}} = L_{\overline{\omega}_{(\omega)}}$. We need another technical lemma:
\begin{lemma}[$\IKP$]
  \label{lem:sub-omega-l-def}
  For any $\alpha, \beta \in \Ord$ such that $\alpha, \beta \subseteq \omega$, if $\alpha \in \dd\mleft(L_\beta\mright)$, then $L_\alpha \in \dd\mleft(L_\beta\mright)$.
\end{lemma}

\begin{proof}
  Given such $\alpha$ and $\beta$, first observe that using the process described in the proof sketch of \autoref{thm:l-omega-dec}, it is not hard to show that for any $n \in \omega$, $L_n$ is finite and thus $L_{n^+} = \dd\mleft(L_n\mright) = \Pow\mleft(L_n\mright) \cap L_\omega$, and this characterisation is downwards absolute to $L_\beta$.

  For any $n \in \alpha$, $n \in \dd\mleft(L_m\mright) \subseteq L_\beta$ for some $m \in \beta$. It follows that $n \subseteq L_m \cap \Ord = m$ by \autoref{cor:l-omega-ord}, thus $n \in m$ or $n = m$ by trichotomy. Either way, we have $L_n \in \dd\mleft(L_m\mright) \subseteq L_\beta$. Additionally, $\dd\mleft(L_n\mright) \subseteq \dd\mleft(L_m\mright) \subseteq L_\beta$, thus $\dd\mleft(L_n\mright) = \Pow\mleft(L_n\mright) \cap L_\beta$ through our observation above.

  Now, we just need to construct
  \[L_\alpha = \left\{x \in L_\beta : \opE{n \in L_\beta} \left(n \in \alpha \land x \subseteq L_n\right)\right\} \in \dd\mleft(L_\beta\mright),\]
  and it suffices to check that there is a first-order formula $\psi\mleft(n, x\mright)$ such that
  \[\left(L_\beta \vDash \psi\mleft(n, x\mright)\right) \leftrightarrow x = L_n\]
  for any $n \in L_\beta \cap \Ord \subseteq \omega$, $x \in L_\beta$.

  We can easily characterise $L_n$ by asserting that there exists a function $f$ on domain $n^+$ such that $f\mleft(\varnothing\mright) = \varnothing$ and $f\mleft(m^+\mright) = \Pow\mleft(f\mleft(m\mright)\mright) \cap L_\beta$ for every $m \in n$. However, we need to take into account that the depth of Kuratowski pairs only guarantees the existence of such $f$ corresponding, informally, to $n - 3$ inside $\dd\mleft(L_m\mright)$ for any $m \geq n$. Therefore, let $\psi'\mleft(n, x\mright)$ be the statement that such $f$ exists on domain $n^+$ with $x = f\mleft(n\mright)$, then we need to take $\psi\mleft(n, x\mright)$ as
  \begin{align*}
     & \left(n = 0 \land x = \varnothing\right) \lor \left(n = 1 \land x = \Pow\mleft(\varnothing\mright)\right) \lor \left(n = 2 \land x = \Pow\Pow\mleft(\varnothing\mright)\right) \\
     & \qquad \qquad {} \lor \opE{m \in n} \opE{y \in x} \left(\psi'\mleft(m, y\mright) \land n = m^{+++} \land x = \Pow\Pow\Pow\mleft(y\mright)\right). \qedhere
  \end{align*}
\end{proof}

\begin{corollary}
  \label{cor:exo-omega-same-l}
  For any $u \subseteq \omega$ that encompasses $\omega$ through some $\Delta_0$-formula, we have
  \[V\mleft(\mathcal{A}\mright) \vDash L_{\overline{\omega}_{(u)}} = L_{\overline{\omega}_{(\omega)}}.\]
  It follows that $V\mleft(\mathcal{A}\mright)$ realises $\opE{\alpha, \beta \in \Ord} \left(\alpha \neq \beta \land L_\alpha = L_\beta\right)$.
\end{corollary}

\begin{proof}
  Working inside $V\mleft(\mathcal{A}\mright)$, by \autoref{lem:exo-omega-subset}, we have $\overline{\omega}_{(u)} \subseteq \overline{\omega}_{(\omega)}$, thus $L_{\overline{\omega}_{(u)}} \subseteq L_{\overline{\omega}_{(\omega)}}$. So we just need to check the opposite direction.

  Take any $\alpha \in \overline{\omega}_{(\omega)}$, then by \autoref{thm:exo-omega-l-con-max-omega}, there must be some $\beta \in \overline{\omega}_{(u)}$ such that $\alpha \in \dd\mleft(L_\beta\mright)$. We also know that $\alpha, \beta \subseteq \omega$ by \autoref{prop:exo-omega-def-by-box}, thus we know from \autoref{lem:sub-omega-l-def} that $L_\alpha \in \dd\mleft(L_\beta\mright)$. It follows that every set $x \in \dd\mleft(L_\alpha\mright)$ will be of the form
  \begin{align*}
    x & = \left\{y \in L_\alpha : L_\alpha \vDash \psi\mleft(y\mright)\right\}                                                                                                    \\
      & = \left\{y \in L_\beta : L_\beta \vDash y \in L_\alpha \land \psi\mleft(y\mright)^{L_\alpha}\right\} \in \dd\mleft(L_\beta\mright) \subseteq L_{\overline{\omega}_{(u)}}.
  \end{align*}
  It follows immediately that
  \[L_{\overline{\omega}_{(\omega)}} = \bigcup_{\alpha \in \overline{\omega}_{(\omega)}} \dd\mleft(L_\alpha\mright) \subseteq L_{\overline{\omega}_{(u)}}. \qedhere\]
\end{proof}

\chapter{Incomparable codings}
\label{chap:incom-coding}

We will now elaborate on general properties and operations on incomparable ordinals in intuitionistic set theory, and present the concrete mechanism I call `incomparable codings' that encodes information from the internal structure of an arbitrary set into complex intuitionistic ordinals $\alpha$ and thus renders highly `non-constructible' sets to be definable from $\alpha$.

We shall say that two sets $x$ and $y$ are \emph{weakly incomparable}\inds{incomparable set}, denoted $x \perp y$, if
\[\neg x \in y^+ \land \neg y \in x^+,\]
or equivalently $x \not\in y \land x \neq y \land y \not\in x$. Classically, there may be incomparable sets but never incomparable ordinals. However, in intuitionistic models, for example, we see the ordinals $\alpha$ and $\beta$ used in the proof of \autoref{lem:ord-trichotomy-fail} are easily realised to be weakly incomparable, thus $\CZF$ is consistent with the existence of weakly incomparable ordinals.

We say that two sets $x$ and $y$ are \emph{strongly incomparable} if
\[\neg x \subseteq y \land \neg y \subseteq x.\]
Using \autoref{lem:exo-omega-subset}, we see that the realisability model in the previous chapter gives
\[V\mleft(\mathcal{A}\mright) \vDash \text{$\overline{\omega}_{(u)}$, $\overline{\omega}_{(v)}$ are strongly incomparable ordinals}\]
if and only if $u, v \subseteq \omega$ are strongly incomparable sets in the base theory. We have plenty of such sets, for example, $u = \omega \setminus \left\{1\right\}$ and $v = \omega \setminus \left\{2\right\}$, thus $\CZF$ is also consistent with the existence of strongly incomparable ordinals.

General theory on such ordinals will be developed in \autoref{sec:incom-intro}, which will be utilised to implement a specific realisability model of $\CZF + V = L$ through the rest of the sections.

\section{Incomparable ordinals and \texorpdfstring{$L$}{L}}
\label{sec:incom-intro}

The key focus of our endeavour will be to build and use large collections of pairwise incomparable ordinals, also known as antichains. Classically, given an arbitrary set $s$, an antichain of size $s$ could usually be given through a function $f$ on domain $s$ such that
\[\opA{x, y \in s} \left(\neg x = y \rightarrow f\mleft(x\mright) \perp f\mleft(y\mright)\right),\]
so that the range of $f$ would be an antichain of the appropriate size.

However, observe that this definition is intuitionistically weak, essentially because both sides of the implication are negated statements. Thus, from some assumptions such as $f\mleft(x\mright) = f\mleft(y\mright)$, we can derive $\neg \neg x = y$ by contraposition but, in intuitionistic logic which lacks double negation elimination, cannot eventually derive $x = y$. Thus, we must rephrase the definition of an incomparable sequence to be in positive terms whenever possible.

Formally, we shall thus say that a function $f$ on domain $s$ is \emph{pairwise weakly incomparable}\inds{pairwise incomparable function} if
\[\opA{x, y \in s} \left(\neg f\mleft(x\mright) \in f\mleft(y\mright) \land \left(f\mleft(x\mright) = f\mleft(y\mright) \rightarrow x = y\right)\right).\]
It is easy to see that this implies the classical definition above, whereas the converse implication would only hold in classical logic.

Likewise, we say that a function $f$ on domain $s$ is \emph{pairwise strongly incomparable} if
\[\opA{x, y \in s} \left(f\mleft(x\mright) \subseteq f\mleft(y\mright) \rightarrow x = y\right).\]

\subsection*{Theory of weak incomparables}

Important properties of weak incomparability are explored in my manuscript \cite{wang-sw25-antichain-pre}, which has been accepted for publication by the Journal of Symbolic Logic by the time of writing this thesis. There, an internal mechanism to systematically generate large amounts of weakly incomparable ordinals is demonstrated through ordinal arithmetic.

Firstly, we define \emph{ordinal addition} intuitionistically as
\[\alpha + \beta = \alpha \cup \bigcup_{\gamma \in \beta} \left(\alpha + \gamma\right)^+.\]
This is the convention used in, for example, \cite{aczel-rathjen10-cst-book}*{Chapter 9}. As intuitionistic logic does not prove the trichotomy that every ordinal is either zero, a successor or a limit, we cannot define addition for the three cases separately.

Observe that:
\begin{lemma}[\cite{wang-sw25-antichain-pre}*{Lemma 5}]
  \label{lem:set-ord-plus-inj}
  For any $\alpha, \beta, \gamma \in \Ord$, we have $\alpha + \beta \in \alpha + \gamma \rightarrow \beta \in \gamma$ and $\alpha + \beta = \alpha + \gamma \rightarrow \beta = \gamma$.
\end{lemma}

\begin{proof}
  We prove the conjunction of the two desired properties by simultaneous set induction on the parameters $\beta$ and $\gamma$. First suppose that
  \[\alpha + \beta \in \alpha + \gamma = \alpha \cup \bigcup_{\delta \in \gamma} \left(\alpha + \delta\right)^+.\]
  Observe that if $\alpha + \beta \in \alpha$, then also $\alpha + \beta \in \alpha + \beta$ which contradicts the set induction axiom. Thus we must have $\alpha + \beta \in \left(\alpha + \delta\right)^+$ for some $\delta \in \gamma$, i.e.\ $\alpha + \beta \in \alpha + \delta$ or $\alpha + \beta = \alpha + \delta$. By the inductive hypothesis, we have $\beta \in \delta$ or $\beta = \delta$ and hence $\beta \in \gamma$ either way.

  On the other hand, suppose that $\alpha + \beta = \alpha + \gamma$, then
  \[\alpha + \delta \in \left(\alpha + \delta\right)^+ \subseteq \alpha + \beta = \alpha + \gamma\]
  for any $\delta \in \beta$. By the inductive hypothesis, we have $\delta \in \gamma$ for any $\delta \in \beta$, i.e.\ $\beta \subseteq \gamma$. An entirely symmetrical argument shows that $\gamma \subseteq \beta$ as well, hence $\beta = \gamma$ as desired.
\end{proof}

\begin{lemma}[\cite{wang-sw25-antichain-pre}*{Lemma 6}]
  \label{lem:set-perp-plus-still-perp}
  Suppose $\alpha \perp \beta \in \Ord$, for any $\gamma \in \Ord$, also $\alpha \perp \beta^+ + \gamma$.
\end{lemma}

\begin{proof}
  Firstly, suppose that $\beta^+ + \gamma \in \alpha^+$. Then $\beta \in \beta^+ + \gamma \subseteq \alpha$, contradicting our assumption that $\alpha \perp \beta$. Observe that this already implies $\neg \alpha = \beta^+ + \gamma$ for any $\gamma \in \Ord$.

  It suffices to further show by set induction on $\gamma$ that $\neg \alpha \in \beta^+ + \gamma$. Suppose that $\alpha \in \beta^+ + \gamma$ yet $\neg \alpha \in \beta^+ + \delta$ for every $\delta \in \gamma$. Since also $\neg \alpha = \beta^+ + \delta$, i.e.\ we have $\neg \alpha \in \left(\beta^+ + \delta\right)^+$, thus we must have $\alpha \in \beta^+$. This again contradicts $\alpha \perp \beta$.
\end{proof}

\begin{corollary}[\cite{wang-sw25-antichain-pre}*{Corollary 7}]
  \label{cor:set-perp-ord-pair}
  Suppose $\alpha \perp \beta \in \Ord$, then for any $\gamma_1, \gamma_2, \delta_1, \delta_2 \in \Ord$, we have
  \[\left(\alpha^+ + \gamma_1\right) \cup \left(\beta^+ + \gamma_2\right) = \left(\alpha^+ + \delta_1\right) \cup \left(\beta^+ + \delta_2\right) \rightarrow \gamma_1 = \delta_1 \land \gamma_2 = \delta_2.\]
\end{corollary}

\begin{proof}
  Assume $\left(\alpha^+ + \gamma_1\right) \cup \left(\beta^+ + \gamma_2\right) = \left(\alpha^+ + \delta_1\right) \cup \left(\beta^+ + \delta_2\right)$. For any $\eta \in \gamma_1$, we have $\alpha^+ + \eta \in \alpha^+ + \gamma_1 \subseteq \left(\alpha^+ + \delta_1\right) \cup \left(\beta^+ + \delta_2\right)$. Suppose that $\alpha^+ + \eta \in \beta^+ + \delta_2$, then $\alpha \in \beta^+ + \delta_2$ by transitivity, and this contradicts \autoref{lem:set-perp-plus-still-perp}. This means we must have $\alpha^+ + \eta \in \alpha^+ + \delta_1$ and thus $\eta \in \delta_1$ by \autoref{lem:set-ord-plus-inj}, i.e.\ $\gamma_1 \subseteq \delta_1$. Entirely symmetrical arguments show that $\delta_1 \subseteq \gamma_1$, $\gamma_2 \subseteq \delta_2$ and $\delta_2 \subseteq \gamma_2$ as well, hence $\gamma_1 = \delta_1$ and $\gamma_2 = \delta_2$ as desired.
\end{proof}

\autoref{cor:set-perp-ord-pair} essentially claims that given two weakly incomparable ordinals, there is a method to represent any pair of two ordinals as a single ordinal uniformly. We proceed to see that this generates an encoding that is not only injective, but also (pairwise) weakly incomparable:
\begin{proposition}[\cite{wang-sw25-antichain-pre}*{Proposition 8}]
  \label{prop:set-perp-ineq-to-perp}
  Suppose $\alpha \perp \beta \in \Ord$, for any set $s \subseteq \Ord$ of ordinals, the function $f : s \rightarrow \Ord$ given by
  \[f\mleft(\gamma\mright) = \alpha^+ \cup \left(\beta^+ + \gamma\right)\]
  is pairwise weakly incomparable.
\end{proposition}

\begin{proof}
  We first show that $\neg f\mleft(\gamma\mright) \in f\mleft(\delta\mright)$ for any $\gamma, \delta \in s$. Suppose that $\alpha^+ \cup \left(\beta^+ + \gamma\right) \in \alpha^+ \cup \left(\beta^+ + \delta\right)$. Notice that $\alpha^+ \cup \left(\beta^+ + \gamma\right) \in \alpha^+$ implies $\alpha \in \alpha^+ \cup \left(\beta^+ + \gamma\right) \subseteq \alpha$, contradicting set induction, thus we must have $\alpha^+ \cup \left(\beta^+ + \gamma\right) \in \beta^+ + \delta$ instead. However, by \autoref{lem:set-perp-plus-still-perp}, we have $\neg \alpha \in \beta^+ + \delta$ and hence another contradiction. Thus we must have $\neg \alpha^+ \cup \left(\beta^+ + \gamma\right) \in \alpha^+ \cup \left(\beta^+ + \delta\right)$.

  Now, it remains to show that
  \[\alpha^+ \cup \left(\beta^+ + \gamma\right) = \alpha^+ \cup \left(\beta^+ + \delta\right) \rightarrow \gamma = \delta.\]
  This is simply an instance of \autoref{cor:set-perp-ord-pair} where $\gamma_1 = \delta_1 = \varnothing$.
\end{proof}

Given fixed $\alpha \perp \beta \in \Ord$ and any function $f : x \rightarrow \Ord$, we shall define $f^\perp : x \rightarrow \Ord$ by
\[f^\perp\mleft(y\mright) = \alpha ^+ \cup \left(\beta^+ + f\mleft(y\mright)\right).\]
By \autoref{prop:set-perp-ineq-to-perp}, $f^\perp$ is pairwise weakly incomparable if and only if $f$ is injective, i.e.\ $f$ satisfies $\opA{y, z \in x} \left(f\mleft(y\mright) = f\mleft(z\mright) \rightarrow y = z\right)$. In \cite{wang-sw25-antichain-pre}, I continued by iterating similar constructions and derive that:
\begin{theorem}[$\IKP$]
  \label{thm:set-perp-class-func}
  Suppose $\alpha \perp \beta \in \Ord$, then there is a $\Sigma_1$-class function $F$ such that for any set $x$, $F\mleft(x\mright) : x \rightarrow \Ord$ is pairwise weakly incomparable.
\end{theorem}

The precise constructions are quite technical and are omitted in this thesis.

Now, incomparable ordinal functions are important because they allow us to encode membership information for subsets of arbitrarily complex sets through ordinal constructions. This idea is briefly commented on by Matthews and Rathjen in the final section of \cite{matthews-rathjen24-constructible-universe}, and can be made precise in the following way:
\begin{proposition}
  \label{prop:set-perp-encode-subset}
  Suppose that $f : x \rightarrow \Ord$ is pairwise weakly incomparable and $f \in L$, then for any $u \subseteq x$, there exists a corresponding ordinal $\delta_u$ such that $\delta_u \in L$ implies $u \in L$.
\end{proposition}

\begin{proof}
  One simply takes $\delta_u = \bigcup_{y \in u} f\mleft(y\mright)^+$, which is an ordinal trivially. Let
  \[v = \left\{y \in x : f\mleft(y\mright) \in \delta_u\right\},\]
  then $u \subseteq v$ trivially. Observe that for any $y \in v$, we must have $z \in u$ such that $f\mleft(y\mright) \in f\mleft(z\mright)^+$ by definition of $\delta_u$. However, weak incomparability means that $\neg f\mleft(y\mright) \in f\mleft(z\mright)$ and that $f\mleft(y\mright) = f\mleft(z\mright)$ must imply $y = z$. Therefore, we have $y = z \in u$ and thus $v = u$ as a consequence.

  It suffices to observe that $f, x \in L$ (since $x$ is the domain of $f$). Thus if also $\delta_u \in L$, then $v$ is a definable set from parameters in $L$ and thus $u = v \in L$.
\end{proof}

Using this construction, we can derive the main result in \cite{wang-sw25-antichain-pre}:
\begin{theorem}[$\IKP$, \cite{wang-sw25-antichain-pre}*{Theorem 2}]
  \label{thm:set-perp-v-eq-l-equiv}
  Assume there exist $\alpha \perp \beta \in \Ord$, then the following are equivalent:
  \begin{enumerate}[label=\uplabel{\roman*}]
    \item $\Ord \subseteq L$,
    \item $V = L$.
  \end{enumerate}
\end{theorem}

\begin{proof}
  The backward direction is trivial. For the forward direction, we use set induction on $x$ to show that $\opA{x} x \in L$.

  From the inductive hypothesis, we will know that $x \subseteq L$, so by strong $\Sigma_1$-collection (which is a theorem in $\IKP$, cf.\ \cite{aczel-rathjen10-cst-book}*{Theorem 19.1.5}) we have a set $s \subseteq \Ord$ such that for all $y \in x$, there exists $\gamma \in s$ such that $y \in L_\gamma$. Take $\sigma = \bigcup_{\gamma \in s} \gamma$, then $x \subseteq L_\sigma$.

  Now, we use the crucial fact that the construction of the $\Sigma_1$-class function $F$ in \autoref{thm:set-perp-class-func} uses only the weak theory $\IKP$, and it is verified in \cite{matthews-rathjen24-constructible-universe}*{section 5} that for every theorem $\varphi$ of $\IKP$, $\IKP \vdash \varphi^L$. It follows that we can compute $F^L\mleft(L_\sigma\mright) \in L$ and, since $\Sigma_1$-statements are upwards absolute, we must have
  \[F\mleft(L_\sigma\mright) = F^L\mleft(L_\sigma\mright) \in L.\]
  It follows from \autoref{prop:set-perp-encode-subset} that there exists an ordinal $\delta_x$ such that $\delta_x \in L$ implies $x \in L$. Since we assumed $\Ord \subseteq L$, this must be the case.
\end{proof}

\begin{corollary}[\cite{wang-sw25-antichain-pre}*{Corollary 14}\footnote{This is a simplified variant of the original statement in the published manuscript. The simplification is proposed by Jem Lord through private communication.}]
  \label{cor:izf-2-not-pow-1-conseq}
  Let theory $T$ be any extension of $\IKP$ that proves the powerset $\Pow\mleft(1\mright)$ exists. If $T + V \neq L + 2 \neq \Pow\mleft(1\mright)$ is consistent, then $T \not\vdash \Ord \subseteq L$.
\end{corollary}

\begin{proof}
  Working inside the theory $T$, it is easy to verify that $2 = \left\{0, 1\right\}$ and $\Pow\mleft(1\mright)$ are both ordinals, and clearly $2 \not\in \Pow\mleft(1\mright)$ and $\Pow\mleft(1\mright) \not\in 2$. Therefore,
  \[T + 2 \neq \Pow\mleft(1\mright) \vdash 2 \perp \Pow\mleft(1\mright).\]
  Suppose that $T \vdash \Ord \subseteq L$, then by \autoref{thm:set-perp-v-eq-l-equiv}, we can derive $V = L$ inside $T + 2 \neq \Pow\mleft(1\mright)$. This contradicts the assumption that $T + V \neq L + 2 \neq \Pow\mleft(1\mright)$ is consistent.
\end{proof}

Observe here that, if we work with the strong theory $T = \IZF$, then by its full separation scheme, for any formula $\varphi$ which does not contain the variable $x$, we can consider some ordinal
\[\alpha_\varphi = \left\{x \in 1 : \varphi\right\} \in \Pow\mleft(1\mright),\]
thus $2 = \Pow\mleft(1\mright)$ would imply $\alpha_\varphi = 0 \lor \alpha_\varphi = 1$ and thus the law of excluded middle, $\varphi \lor \neg \varphi$. Therefore, $\IZF + 2 = \Pow\mleft(1\mright)$ is simply equivalent to classical $\ZF$ over intuitionistic logic.

Informally, we can call some statement $\psi$ an \emph{anticlassical axiom}\indg{anticlassical axiom}\footnote{There are many well-known anticlassical axioms --- for example, Church's thesis, the uniformity principle, etc. We will not present any details here, but refer the readers to Chen and Rathjen \cite{chen-rm-rathjen12-lifschitz-realizability} for elaboration, with explicit realisability constructions showing that they are consistent with $\IZF$.} if $\ZF$ (or any reasonable classical base theory) proves $\neg \psi$, yet $\psi$ is consistent with some intuitionistic theory such as $\IZF$. It follows from the previous remark that, for any anticlassical axiom $\psi$, $\IZF + 2 = \Pow\mleft(1\mright) \vdash \neg \psi$ over intuitionistic logic. By contraposition,
\[\IZF + \psi \vdash 2 \neq \Pow\mleft(1\mright).\]
Therefore, if for any such axiom $\psi$, $\IZF + V \neq L + \psi$ is consistent, then the presumptions in \autoref{cor:izf-2-not-pow-1-conseq} are obtained and we have $\IZF \not\vdash \Ord \subseteq L$.

Analogously, when we consider the $\Delta_0$-instances of separation in the weaker theory $\CZFP$, we see that if any model of $\CZFP + V \neq L$ also violates some $\Delta_0$-instance of excluded middle, for example, if
\[\CZFP + V \neq L + \neg\!\opA{x, y} \left(x \in y \lor x \not\in y\right)\]
is consistent, then $\CZFP \not\vdash \Ord \subseteq L$.

Unfortunately, however, I know of no such interpretations, and thus \autoref{prob:ord-in-l} remains open. The main difficulty is precisely as we show in this chapter, that the presence of specific coding objects in $L$ can cause it to be much larger than what we traditionally conceive in classical set theory, thus making the implementation or verification of $V \neq L$ very complicated in intuitionistic models.

\subsection*{Theory of strong incomparables}

Now, \autoref{thm:set-perp-v-eq-l-equiv} is a surprising result in itself, since it relates in intuitionistic set theory the statement $\Ord \subseteq L$, which is completely trivial in classical set theory, with $V = L$, which is easily demolished classically through forcing. However, working with weak incomparables can be unsatisfying since whether $\Ord \subseteq L$ remains an open question in intuitionistic set theory (cf.\ \autoref{sec:intro-int-l}) and thus, constructions in \autoref{prop:set-perp-encode-subset} do not produce ordinal codes that lie readily in $L$.

If the ordinals in the range of the incomparable function are particularly nice, then there are analogous constructions which guarantee membership in $L$. For example, suppose that $f : x \rightarrow \Ord$ has all of its values given by some $g : x \rightarrow \Ord$ as
\[\opA{y \in x} f\mleft(y\mright) = L_{g\mleft(y\mright)} \cap \Ord,\]
then for any $u \subseteq x$, we can construct
\[\delta^*_u = \bigcup_{y \in u} \dd\mleft(L_{g\mleft(y\mright)}\mright) \cap \Ord = L_{\bigcup_{y \in u} g\mleft(y\mright)^+} \cap \Ord \in L.\]

We can prove an analogue of \autoref{prop:set-perp-encode-subset} using the ordinals $\delta^*_u$. Observe carefully that now, it is necessary to assume that the function is pairwise strongly incomparable:
\begin{proposition}
  \label{prop:l-ord-strong-incom-code}
  Suppose that $f : x \rightarrow \Ord$ is pairwise strongly incomparable, $f \in L$, and $f$ is given by some $g : x \rightarrow \Ord$ as $f\mleft(y\mright) = L_{g\mleft(y\mright)} \cap \Ord$ for every $y \in x$. Then for any $u \subseteq x$, we have $u \in L$.
\end{proposition}

\begin{proof}
  We know that $f, x$ and $\delta^*_u$ as constructed above are in $L$. Let
  \[v  = \left\{y \in x : f\mleft(y\mright) \in \delta^*_u\right\},\]
  then $v \in L$. We likewise trivially have $u \subseteq v$. For any $y \in v$, we must have $z \in u$ such that $f\mleft(y\mright) \in \dd\mleft(L_{g\mleft(z\mright)}\mright) \cap \Ord$ by definition of $\delta^*_u$. This implies that $f\mleft(y\mright) \subseteq L_{g\mleft(z\mright)} \cap \Ord = f\mleft(z\mright)$, and strong incomparability implies that $y = z$. Therefore, we have $y = z \in u$ and thus $u = v \in L$ as desired.
\end{proof}

Specifically, recall that \autoref{cor:l-omega-ord} states $\opA{n \in \omega} L_n \cap \Ord = n$. Thus, for any $\alpha \in \Ord$ such that $\alpha \subseteq \omega$, we immediately have
\[\alpha = \bigcup_{n \in \alpha} n^+ = \bigcup_{n \in \alpha} \dd\mleft(L_n\mright) \cap \Ord = L_{\bigcup_{n \in \alpha} n^+} \cap \Ord = L_\alpha \cap \Ord.\]
It follows immediately that:
\begin{corollary}
  \label{cor:omega-sub-strong-incom-code}
  Suppose that $f : x \rightarrow \Pow\mleft(\omega\mright) \cap \Ord$ is pairwise strongly incomparable and $f \in L$, then for any $u \subseteq x$, we have $u \in L$.
\end{corollary}

However, given these additional requirements, it is no longer possible to automatically generate large collections of incomparables using methods similar to \autoref{prop:set-perp-ineq-to-perp}. Later in \autoref{sec:czf-v-eq-l}, we will make use of recursion-theoretic constructions to manually build a pairwise strongly incomparable function $\omega \rightarrow \Pow\mleft(\omega\mright) \cap \Ord$ in $L$; in the next chapter, we will then develop some semi-automatic mechanisms to construct larger sets of strongly incomparable ordinals given a base theory containing the powerset axiom.

\section{The case for \texorpdfstring{$V = L$}{V = L}}
\label{sec:overview-v-eq-l}

In the remainder of this chapter, we shall apply the method of incomparable codings to analyse the strength implications of the axiom $V = L$, following the constructions in my arXiv manuscript \cite{wang-sw26-czf-v-l-pre}. Thus, let us first understand why it is interesting --- and non-trivial --- to study $V = L$ for certain set theories. We reproduce the introduction to \cite{wang-sw26-czf-v-l-pre} here:

Traditionally, most set theories $T$ we work with admit the theorem
\[T \vdash T^L,\]
and it follows immediately that the `inner model' $L$ is a model of the extension $T + V = L$, thus $T$ and $T + V = L$ are equiconsistent. This is the case for well-known classical theories $\KP$, $\ZF$ and $\ZFC$; as we know from Lubarsky \cite{lubarsky93-intuitionistic-l} and Matthews and Rathjen \cite{matthews-rathjen24-constructible-universe}, this is also the case for intuitionistic theories $\IKP$ and $\IZF$.

However, there are known theories $T$ that do not satisfy $T \vdash T^L$ and violate the equiconsistency claim. For example, Rathjen proved that:
\begin{theorem}[\cite{rathjen20-power-kp-choice}*{Theorem 4.6}]
  \label{thm:kpp-v-eq-l}
  $\KPP + V = L$ is much stronger than $\KPP$, where $\KPP$ is classical Kripke--Platek set theory\inds{Kripke--Platek set theory with powerset, $\KPP$} with a new primitive unary function symbol $\Pow$ for powerset \uppar{see remarks in \autoref{sec:czf-intro} and \autoref{ftn:kpp-prim-pow}}.
\end{theorem}

The key step in this proof is the following lemma:
\begin{lemma}
  \label{lem:kpp-unc-ord-str}
  $\KPP + \text{$\omega_1$ exists}$ proves the consistency of $\KPP$.
\end{lemma}

\begin{proof}[Proof sketch]
  Using the technique in \cite{rathjen05-replacement-collection}*{section 4}, we can construct the Bachmann--Howard ordinal $\vartheta_{\varepsilon_{\Omega + 1}}\mleft(0\mright)$ as some class of ordinals $\mathrm{Acc}$ in $\KPP$. Assuming $\omega_1$ exists, $\mathrm{Acc}$ is a $\Sigma_1$-definable subclass of $\omega_1$ whose existential quantifier can be easily bounded to range over sets of functions on $\omega_1$, thus $\mathrm{Acc}$ is actually a subset of $\omega_1$ and $\vartheta_{\varepsilon_{\Omega + 1}}\mleft(0\mright)$ exists as a set-theoretic ordinal.

  One then applies the method of relativised ordinal analysis in \cite{rathjen14-power-kp}. By \cite{rathjen14-power-kp}*{Theorem 8.2}, if $\KPP$ proves a $\Sigma_1$-sentence $\varphi$ (possibly containing the powerset symbol), then
  \[V_{\vartheta_{\varepsilon_{\Omega + 1}}\mleft(0\mright)} \vDash \varphi.\]
  Therefore, $\Con\mleft(\KPP\mright)$ follows immediately.
\end{proof}

\begin{proof}[Proof of \autoref{thm:kpp-v-eq-l}]
  It is easy to verify in classical set theory that for any countable ordinal $\alpha \in \Ord$, $L_\alpha$ is also countable. Thus, assuming $V = L$, then $\Pow\mleft(\omega\mright) \in L$ is uncountable, and we must have $\Pow\mleft(\omega\mright) \in L_\alpha$ for some uncountable $\alpha \in \Ord$. It follows that $\omega_1$ exists, and by \autoref{lem:kpp-unc-ord-str}, $\KPP + V = L$ proves the consistency of $\KPP$.
\end{proof}

On the intuitionistic side, we similarly know that:
\begin{proposition}
  \label{prop:czf-l-not-exp-con}
  $\CZFP + \neg \Exp^L$ is consistent. Since $\CZF \vdash \Exp$ by \autoref{prop:ikp-exp-sub-coll-powerset}, $\CZFP \not\vdash \CZF^L$.
\end{proposition}

\begin{proof}
  Let $\REM$ denote the axiom scheme of \emph{restricted excluded middle}\inds{restricted excluded middle, $\REM$}, i.e.\ $\varphi \lor \neg \varphi$ for any $\Delta_0$-formula $\varphi$. Observe that $\CZF + \REM$ suffices for all uses of the classical logic in \autoref{thm:kpp-v-eq-l}, namely that the ordinals are linearly ordered and $L_\alpha$ is countable for any countable $\alpha \in \Ord$. Therefore, inside\footnote{$\REM$ also suffices for deriving the equivalence between the axioms of exponentiation, subset collection and powerset in $\CZF$ (cf.\ \autoref{prop:ikp-exp-sub-coll-powerset} or \cite{aczel-rathjen10-cst-book}*{section 10.1}), thus $\CZFP + \REM$ and $\CZF + \REM$ are actually the same theory. However, we will continue to write $\CZFP + \REM$ for clarity reasons.} $\CZFP + \REM + V = L$, we similarly consider $\Pow\mleft(\omega\mright) \in L_\alpha$ and $\alpha$ must then contain all countable ordinals; so $\CZFP + \REM + V = L$ also proves that the Bachmann--Howard ordinal $\vartheta_{\varepsilon_{\Omega + 1}}\mleft(0\mright)$ exists.

  We know that the relativised ordinal analysis can be performed in an intuitionistic base theory (cf.\ remarks in \cite{rathjen23-ep-intuitionistic-set-theory}*{section 3(c)}). On the other hand, $\CZFP + \REM$ is equiconsistent with $\KPP$, using the same weak realisability argument we mentioned in \autoref{sec:czf-intro} and will explain in more detail in \autoref{sec:czf-weak-real}. Therefore, it shall follow that
  \[\CZFP + \REM + V = L \vdash \Con\mleft(\CZFP + \REM\mright).\]

  Now, suppose that $\CZFP + \neg \Exp^L$ is not consistent, i.e.\ $\CZFP \vdash \neg \neg \Exp^L$. It follows that $\CZFP + \REM \vdash \neg \neg \Exp^L$. However, the axioms of exponentiation, subset collection and powerset are the sole obstacles preventing $L$ from being a model of $\CZFP$ here, as extending the proof of $L \vDash \IKP$ in \cite{matthews-rathjen24-constructible-universe}*{section 5}, it is easy to verify that $\CZFP + \REM$ proves
  \[L \vDash \CZF - \text{subset collection} + \REM + V = L.\]
  Therefore, combining with $\CZFP + \REM \vdash \neg \neg \Exp^L$ and the fact that $\REM$ and exponentiation together imply the powerset axiom, we have $\CZFP + \REM \vdash \neg \neg \varphi^L$ for every axiom $\varphi$ of $\CZFP + \REM + V = L$. It follows that, if there is a proof of $\CZFP + \REM + V = L \vdash \bot$, then there is also a proof of $\CZFP + \REM \vdash \neg \neg \bot^L$, which is simply $\CZFP + \REM \vdash \bot$ since intuitionistic logic admits triple negation elimination.

  In other words, $\Con\mleft(\CZFP + \REM\mright)$ implies $\Con\mleft(\CZFP + \REM + V = L\mright)$, and both are proved in $\CZFP + \REM + V = L$ by the argument above. This is a contradiction to G\"odel's second incompleteness theorem.
\end{proof}

This is a slightly stronger result than \cite{matthews-rathjen24-constructible-universe}*{Theorem 7.12}, where a variant of weak realisability is used to directly show $\CZFP \not\vdash \Exp^L$. Either way, we know that for any fragment $T \subseteq \CZFP$ that proves $\Exp$, $T$ cannot prove $T^L$. This leaves open the question whether we still have $T \equiv_{\Con} T + V = L$ for such theories, most notably when $T = \IKPP$, $\CZF$ and $\CZFP$. In what follows, we shall resolve the case for $\CZF$ and prove that, despite the proposition above, we still have
\[\CZF \equiv_{\Con} \CZF + V = L.\]
The case for $\IKPP$ and $\CZFP$ will then be treated in \autoref{chap:plump-power}.

\section{A universe over Martin-L\"of type theory}
\label{sec:mltt-interp}

We will construct a realisability interpretation of $\CZF + V = L$ following the structure of Aczel's interpretation of $\CZF$ in Martin-L\"of type theory with a single $W$-type (sometimes known as $\MLV$) in \cite{aczel78-type-theoretic-cst}. Roughly speaking, the set-theoretic universe is defined by the type-theoretic constructor
\[A : \text{Type}, f : A \rightarrow V \vdash \text{sup}\mleft(A, f\mright) : V,\]
where $A$ ranges over the `small' types in the theory.

As a basis, we use a presentation of the type theory given in \cite{rathjen14-czf-lpo}. There, it is implemented as a definable structure over the natural numbers using inductive definitions in the classical theory $\BI$ of second-order arithmetic. Since we have already presented $\Sigma$-inductive definitions for $\IKP$ in \autoref{thm:set-ind-def}, we would conveniently perform the same constructions inside classical Kripke--Platek set theory $\KP$ instead.

We fix the pca $\Kl$, Kleene's first algebra, which is clearly a $\Delta_0$-definable partial structure over $\omega$ in $\KP$. We shall implement the semantics of the type theory by interpreting the function types as collections of elements (with appropriate domains) in $\Kl$. In the type theory, we shall have the following types\footnote{We follow the notation used in \cite{wang-sw26-czf-v-l-pre}, which was taken from Rathjen \cite{rathjen14-czf-lpo}. The overlined type names here are unrelated to the similarly named objects in \autoref{sec:czf-real-in-czf}.}:
\begin{enumerate}[label=\uplabel{\arabic*}]
  \item for each $n \in \omega$, $\overline{n}$ names a finite type with $n$ elements;
  \item $\overline{\omega}$ names the type of natural numbers;
  \item $\overline{X}$ names a new, arbitrary primitive type that will later be specified through its interpretation in the base theory --- this is our addition to $\MLV$ in order to introduce incomparable codings into the interpretation of $\CZF$;
  \item the usual type formers $\sigma, \pi$ that build dependent $\Sigma$- and $\Pi$-types.
\end{enumerate}
All type names are encoded as natural numbers in the following way:
\begin{alignat*}{2}
  \overline{n}      & = \bp 0n \quad \mathrlap{\text{for each $n \in \omega$},}                                                             \\
  \overline{\omega} & = \bp 10,                                                 & \qquad \sigma & = \oplam{nm} \bp 2 \mleft(\bp nm\mright), \\
  \overline{X}      & = \bp 11,                                                 & \pi           & = \oplam{nm} \bp 3 \mleft(\bp nm\mright).
\end{alignat*}

Fix any set $A \subseteq \omega$, we first invoke $\Sigma$-inductive definitions to define the universe of small types, $\bU_A = I\mleft(\Phi\mright)$, where $\Phi$ is a $\Delta_0$-class given by the following clauses:
\begin{enumerate}[label=\uplabel{\arabic*}]
  \item $\tuple{\tuple{\overline{n}, n}, \varnothing} \in \Phi$ for each $n \in \omega$, which immediately implies $\tuple{\overline{n}, n} \in \bU_A$;
  \item $\tuple{\tuple{\overline{\omega}, \omega}, \varnothing} \in \Phi$, which implies $\tuple{\overline{\omega}, \omega} \in \bU_A$;
  \item $\tuple{\tuple{\overline{X}, A}, \varnothing} \in \Phi$, which implies $\tuple{\overline{X}, A} \in \bU_A$;
  \item for any $n, e \in \omega$, if $f : s \rightarrow \Pow\mleft(\omega\mright)$ is a set-theoretic function on some domain $s \subseteq \omega$, such that $n \in s$ and for any $k \in f\mleft(n\mright)$,
        \[ek {\downarrow} \in s,\]
        then $\tuple{\tuple{\sigma ne, p}, f} \in \Phi$ where
        \[p = \left\{\bp ku : k \in f\mleft(n\mright), u \in f\mleft(ek\mright)\right\},\]
        and also $\tuple{\tuple{\pi ne, q}, f} \in \Phi$ where
        \[q = \left\{d \in \omega : \opA{k \in f\mleft(n\mright)} dk {\downarrow} \in f\mleft(ek\mright)\right\}.\]
\end{enumerate}

The idea is that $\bU_A$ shall be a partial function on $\omega$, such that for any type name $n$, $\bU_A\mleft(n\mright)$ is the collection of all elements in that type, where elements in $\Pi$-types are represented by objects in the pca $\Kl$. We check that the inductive definition ensures that $\bU_A$ is a function:
\begin{lemma}[$\KP$]
  \label{lem:mltt-ind-def-func}
  For any $n \in \omega$, if $\tuple{n, p}, \tuple{n, q} \in \bU_A$, then $p = q$.
\end{lemma}

\begin{proof}
  By condition \ref{itm:set-ind-def-min} in \autoref{thm:set-ind-def}, it suffices to verify that
  \[\opA{{\tuple{\tuple{n, p}, y} \in \Phi}} \left(\opA{{\tuple{m, p'} \in y}} \opA{{\tuple{m, q} \in \bU_A}} p' = q \rightarrow \opA{{\tuple{n, q} \in \bU_A}} p = q\right).\]
  This is clear since by \autoref{cor:set-ind-def-exist-pre}, any $\tuple{n, q} \in \bU_A$ must correspond to some $\tuple{\tuple{n, q}, y} \in \Phi$ where $y \subseteq \bU_A$. Looking at the definition of $\Phi$, we see that the set $p$ corresponding to $n \in \omega$ is uniquely determined by relevant information in $y$, and therefore, a different $\tuple{n, q}$ cannot lie in $\bU_A$ if no alternative values for types in $y$ can be found.
\end{proof}

Now, $\bU_A$ is given as a $\Sigma_1$-class, so we can use it to again inductively define the $W$-type in the theory, implemented as a class $\bV_A = I\mleft(\Psi\mright) \subseteq \omega$ in $\KP$, using the $\Sigma_1$-class $\Psi \subseteq \omega \times \Pow\mleft(\omega\mright)$ given by
\[\tuple{n, y} \in \Psi \defeqv n \in \omega \land \opE{{\tuple{m, p} \in \bU_A}} \left(\bp_0 n {\downarrow} = m \land \opA{k \in p} \bp_1 nk {\downarrow} \in y\right).\]
This enforces that each element in $\bV_A$ should be of the form $\bp ne$ where $n$ lies in the domain of $\bU_A$ and $e \in \Kl$ maps each $k \in \bU_A\mleft(n\mright)$ to $ek {\downarrow} \in \bV_A$.

$\bV_A$ will serve as the domain of set names, on which we build an interpreted realisability model of $\CZF$. Just as in \autoref{sec:czf-real-in-czf}, we specify realisers for the atomic formulae by simultaneous recursion; however, since the names for sets are now given as objects in $\bV_A \subseteq \Kl$, it suffices to use \autoref{cor:recursion-theorem} for the pca to build $f_\in, g_= \in \Kl$ as
\begin{align*}
  f_\in & = \oplam{\alpha \beta} \sigma \mleft(\bp_0 \beta\mright) \mleft(\oplam{x} g_= \alpha \mleft(\bp_1 \beta x\mright)\mright),                                                                                                                                                  \\
  g_=   & = \oplam{\alpha \beta} \sigma \mleft(\pi \mleft(\bp_0 \alpha\mright) \mleft(\oplam{y} f_\in \mleft(\bp_1 \alpha y\mright) \beta\mright)\mright) \mleft(\oplam{x} \pi \mleft(\bp_0 \beta\mright) \mleft(\oplam{y} f_\in \mleft(\bp_1 \beta y\mright) \alpha\mright)\mright).
\end{align*}
Using induction on $\bV_A$ as given by condition \ref{itm:set-ind-def-min} in \autoref{thm:set-ind-def}, it is easy to verify that $f_\in \alpha \beta$ and $g_= \alpha \beta$ lie in the domain of $\bU_A$ for any $\alpha, \beta \in \bV_A$, i.e.\ they are proper type names. We can thus define the atomic clauses of the realisability conditions here as
\begin{align*}
  d \Vdash \alpha \in \beta & \defeqv d \in \bU_A\mleft(f_\in \alpha \beta\mright), \\
  d \Vdash \alpha = \beta   & \defeqv d \in \bU_A\mleft(g_= \alpha \beta\mright).
\end{align*}

The conditions for bounded quantifiers are now
\begin{align*}
  d \Vdash \opA{x \in \alpha} \varphi & \defeqv \opA{i \in \bU_A\mleft(\bp_0 \alpha\mright)} di {\downarrow} \Vdash \varphi\mleft[\bp_1 \alpha i/x\mright],                                 \\
  d \Vdash \opE{x \in \alpha} \varphi & \defeqv \bp_0 d {\downarrow} \in \bU_A\mleft(\bp_0 \alpha\mright) \land \bp_1 d \Vdash \varphi\mleft[\bp_1 \alpha \mleft(\bp_0 d\mright)/x\mright],
\end{align*}
so that it is again not hard to verify that
\begin{align*}
  \left(d \Vdash \alpha \in \beta\right) & \leftrightarrow \left(d \Vdash \opE{x \in \beta} \alpha = x\right),                        \\
  \left(d \Vdash \alpha = \beta\right)   & \leftrightarrow \left(d \Vdash \alpha \subseteq \beta \land \beta \subseteq \alpha\right).
\end{align*}

Just like in \autoref{sec:czf-real-in-czf}, we treat $\bV_A$ as the name of this realisability model, and write $\bV_A \vDash \varphi$ to abbreviate for
\[\opE{d \in \Kl} d \Vdash \varphi,\]
where $\varphi$ is a first-order formula possibly containing parameters from $\bV_A$.

Proceeding to the axioms, we see that the realisers for the equality axioms are entirely analogous to the constructions in \autoref{sec:czf-real-in-czf}. The realisers for the bounded quantifier axioms need to be adjusted slightly to accommodate for the interpreted realisability conditions for unbounded quantifiers, but still use the equality axioms in essentially the same way as \autoref{lem:set-bnd-quant-axioms-real}. The axioms of extensionality and set induction are also analogous.

For the pairing axiom, we will utilise the primitive type $\overline{2}$ in $\bU_A$ and construct a function in $\Kl$ that maps any $\alpha, \beta \in \bV_A$ to an object
\[\left\{\alpha, \beta\right\}^V = \bp \overline{2} \mleft(\oplam{x} \bdN x0 \alpha \beta\mright) \in \bV_A.\]
For the union axiom, we use a $\Sigma$-type and build
\[\bigcup\nolimits^V \alpha = \bp \mleft(\sigma \mleft(\bp_0 \alpha\mright) \mleft(\oplam{x} \bp_0 \mleft(\bp_1 \alpha x\mright)\mright)\mright) \mleft(\oplam{x} \bp_1 \mleft(\bp_1 \alpha \mleft(\bp_0 x\mright)\mright) \mleft(\bp_1 x\mright)\mright) \in \bV_A.\]
For infinity, we build a map $g \in \Kl$ from $\omega$ to $\bV_A$ as
\[g = \br \mleft(\oplam{af} \bp \overline{a} f\mright) = \br \mleft(\oplam{af} \bp \mleft(\bp 0a\mright) f\mright)\]
by recursion, so that $gn = \bp \overline{n} \mleft(\oplam{m} gm\mright)$. Then the names for natural numbers and $\omega$ can be given as
\begin{align*}
  n^V      & = gn \in \bV_A \quad \text{for all $n \in \omega$}, \\
  \omega^V & = \bp \overline{\omega} g \in \bV_A.
\end{align*}
Verifying that these constructions work is again analogous to \autoref{sec:czf-real-in-czf}.

The realiser for strong collection is slightly different, as we observe that we do not need assumptions from the base theory for the collection scheme in an \emph{interpreted} realisability model of set theory:
\begin{lemma}
  The strong collection scheme \ref{itm:set-str-coll-axiom} is realised in $\bV_A$:
  \[\opA{x \in a} \opE{y} \theta\mleft(x, y\mright) \rightarrow \opE{b} \left(\opA{x \in a} \opE{y \in b} \theta\mleft(x, y\mright) \land \opA{y \in b} \opE{x \in a} \theta\mleft(x, y\mright)\right).\]
\end{lemma}

\begin{proof}
  Let $\vec{z}$ enumerate the free variables in $\theta\mleft(x, y\mright)$ other than $x$, $y$ and $a$. For any $\alpha \in \bV_A$, $\vec{\beta} \in \bV_A^k$ and $f \Vdash \opA{x \in a} \opE{y} \theta\mleft(x, y\mright)\mleft[\alpha/a, \vec{\beta}/\vec{z}\mright]$, we can define
  \[g = \oplam{\alpha f} \bp \mleft(\bp_0 \alpha\mright) \mleft(\oplam{x} \bp_0 \mleft(fx\mright)\mright),\]
  then for any $k \in \bU_A\mleft(\bp_0 \alpha\mright)$, $\bp_1 \mleft(fk\mright) \Vdash \theta\mleft[\alpha/a, \bp_1 \alpha k/x, \bp_1 \mleft(g \alpha f\mright) k/y, \vec{\beta}/\vec{z}\mright]$. Thus
  \[\oplam{\alpha \vec{\beta} f} \bp \mleft(g \alpha f\mright) \mleft(\bp \mleft(\oplam{k} \bp k \mleft(\bp_1 \mleft(fk\mright)\mright)\mright) \mleft(\oplam{k} \bp k \mleft(\bp_1 \mleft(fk\mright)\mright)\mright)\mright)\]
  realises the axiom.
\end{proof}

For subset collection, we use a $\Pi$-type:
\begin{lemma}
  The subset collection scheme \ref{itm:set-sub-coll-axiom} is realised in $\bV_A$:
  \begin{align*}
    \opE{c} \opA{u} & \left(\opA{x \in a} \opE{y \in b} \psi\mleft(x, y, u\mright)\right.                                                                                                               \\
                    & \quad \left.{} \rightarrow \opE{d \in c} \left(\opA{x \in a} \opE{y \in d} \psi\mleft(x, y, u\mright) \land \opA{y \in d} \opE{x \in a} \psi\mleft(x, y, u\mright)\right)\right),
  \end{align*}
\end{lemma}

\begin{proof}
  Let $\vec{z}$ enumerate the free variables in $\psi\mleft(x, y, u\mright)$ other than $x$, $y$, $u$, $a$, $b$ and $c$. We define
  \[c^* = \oplam{\alpha \beta} \bp \mleft(\pi \alpha \mleft(\oplam{x} \beta\mright)\mright) \mleft(\oplam{h} \bp \mleft(\bp_0 \alpha\mright) \mleft(\oplam{x} \bp_1 \beta \mleft(hx\mright)\mright)\mright),\]
  such that for any $\alpha, \beta \in \bV_A$, $c^* \alpha \beta \in \bV_A$. Now, suppose that some
  \[f \Vdash \opA{x \in a} \opE{y \in b} \psi\mleft(x, y, u\mright)\mleft[\alpha/a, \beta/b, c^* \alpha \beta/c, \delta/u, \vec{\gamma}/\vec{z}\mright],\]
  and we will let
  \[d^* = \bp \mleft(\bp_0 \alpha\mright) \mleft(\oplam{x} \bp_1 \beta \mleft(\bp_0 \mleft(fx\mright)\mright)\mright) = \bp_1 \mleft(c^* \alpha \beta\mright) \mleft(\oplam{x} \bp_0 \mleft(fx\mright)\mright).\]
  Then $\bp \mleft(\oplam{x} \bp x \mleft(\bp_1 \mleft(fx\mright)\mright)\mright) \mleft(\oplam{x} \bp x \mleft(\bp_1 \mleft(fx\mright)\mright)\mright)$ realises the corresponding conjunction that $d^*$ needs to satisfy as $d$. Therefore,
  \[\oplam{\alpha \beta \vec{\gamma}} \bp \mleft(c^* \alpha \beta\mright) \mleft(\oplam{\delta f} \bp \mleft(\oplam{x} \bp_0 \mleft(fx\mright)\mright) \mleft(\bp \mleft(\oplam{x} \bp x \mleft(\bp_1 \mleft(fx\mright)\mright)\mright) \mleft(\oplam{x} \bp x \mleft(\bp_1 \mleft(fx\mright)\mright)\mright)\mright)\mright)\]
  realises the axiom.
\end{proof}

For $\Delta_0$-separation, the usual argument depends on the observation that the collection of realisers for each $\Delta_0$-formula forms a type in $\bU_A$. However, we use the same trick as in \autoref{sec:czf-real-in-czf} and just treat binary intersection instead, where we only need to use the type for realisers in the atomic case:
\begin{lemma}
  The binary intersection axiom is realised in $\bV_A$:
  \[\opA{a} \opA{b} \opE{y} \opA{x} \left(x \in y \leftrightarrow x \in a \land x \in b\right).\]
\end{lemma}

\begin{proof}
  Just as in \autoref{lem:czf-real-in-czf-bin-inter}, we consider
  \[\gamma = \oplam{\alpha \beta} \bp \mleft(\sigma \mleft(\bp_0 \alpha\mright) \mleft(\oplam{x} f_\in \mleft(\bp_1 \alpha x\mright) \beta\mright)\mright) \mleft(\oplam{x} \bp_1 \alpha \mleft(\bp_0 x\mright)\mright),\]
  and we will have
  \[\left(\bp \mleft(\bp fg\mright) d \Vdash \delta \in \gamma \alpha \beta\right) \leftrightarrow \left(\bp fd \Vdash \delta \in \alpha\right) \land \left(\bp \mleft(\bp_0 g\mright) \mleft(\iot d \mleft(\bp_1 g\mright)\mright) \Vdash \delta \in \beta\right).\]
  Thus it is completely analogous to verify that the biconditional is realised.
\end{proof}

Therefore, we can conclude that:
\begin{theorem}
  \label{thm:mltt-real-czf}
  Let $\varphi$ be any theorem of $\CZF$, then $\KP$ proves $\bV_A \vDash \varphi$ for any subset $A \subseteq \omega$, i.e.\ there exists $d \in \Kl$ such that $d \Vdash \varphi$.
\end{theorem}

Finally, the reason we need this new realisability interpretation is that the semantics we give $\MLV$ here ensures every type is a subset of $\omega$, so that the realisability model admits the following \emph{axiom of subcountability}\inds{axiom of subcountability}:
\begin{lemma}
  $\bV_A$ realises that
  \[\opA{x} \opE{u, f} \left(u \subseteq \omega \land \text{$f$ is a surjection from $u$ onto $x$}\right).\]
\end{lemma}

\begin{proof}
  For any $\alpha \in \bV_A$ that is assigned to $x$, we can compute $u$ as
  \[\bp \mleft(\bp_0 \alpha\mright) \mleft(\oplam{x} x^V\mright),\]
  and $f$ as $\bp \mleft(\bp_0 \alpha\mright) \mleft(\oplam{x} \tuple{x^V, \bp_1 \alpha x}^V\mright)$, where $\tuple{\beta, \gamma}^V = \left\{\left\{\beta, \beta\right\}^V, \left\{\beta, \gamma\right\}^V\right\}^V$. Observe that $\oplam{x} \bp x \ior \Vdash u \subseteq \omega^V$, and it is also easy to realise that $f$ is a surjection from $u$ onto $x$.
\end{proof}

In the next section, we shall use the axiom of subcountability together with the distinguished primitive type $\overline{X}$ to construct an incomparable coding for all sets in this realisability model.

\section{An equiconsistency proof for \texorpdfstring{$\CZF + V = L$}{CZF + V = L}}
\label{sec:czf-v-eq-l}

Following \cite{wang-sw26-czf-v-l-pre}, we now select a specific set $A \subseteq \omega$ and see what effect it has on the realisability interpretation. To begin, we fix an arbitrary, recursive way to encode finite $\omega$-tuples as a single natural number, such that the following notions are definable as applicative terms in the pca:
\begin{itemize}
  \item $\bl_s t = \left|t\right|$ is the length of the tuple,
  \item $\bt_s tn$ is the truncated tuple of first $n$ components in $t$ for any $0 \leq n \leq \bl_s t$,
\end{itemize}
whenever $t$ is a valid code for some tuple. For example, we can encode the tuple $\tuple{a_0, \ldots, a_{n - 1}}$ as
\[\bp n \mleft(\bp a_{n - 1} \mleft(\bp a_{n - 2} \mleft(\cdots \bp a_1 \mleft(\bp a_0 0\mright) \cdots \mright)\mright)\mright),\]
conveniently reversed for back truncations, so that $\bl_s = \bp_0$, and
\[\bt_s = \oplam{tn} \bp n \mleft(f \mleft(\bp_0 t - n\mright) \mleft(\bp_1 t\mright)\mright),\]
where $f = \br \mleft(\oplam{agt} \bdN^* a0t \mleft(g \mleft(\bpN a\mright) \mleft(\bp_1 t\mright)\mright)\mright)$ strips a number of leading components from the given tuple.

We also need an alternative, recursive pairing function $\bp'$, which has a specific property. We define it directly as follows\indg{pairing function}:
\[\bp' ab = \left\{\begin{aligned}
     & \max\mleft\{a, b\mright\} \cdot \left(\max\mleft\{a, b\mright\} + 1\right) - a + b &  & \text{if $\max\mleft\{a, b\mright\}$ is even}, \\
     & \max\mleft\{a, b\mright\} \cdot \left(\max\mleft\{a, b\mright\} + 1\right) + a - b &  & \text{otherwise}.
  \end{aligned}\right.\]

We first observe that:
\begin{lemma}[$\KP$]
  \label{lem:max-pairing-bij}
  For any $n \in \omega$, the map $\tuple{a, b} \mapsto \bp' ab$ is a bijection between the Cartesian product $n \times n$ and $n^2 \in \omega$.
\end{lemma}

\begin{proof}
  Just observe that for any $a, b \in \omega$,
  \[\max\mleft\{a, b\mright\}^2 \leq \bp' ab < \left(\max\mleft\{a, b\mright\} + 1\right)^2,\]
  so for each $n \in \omega$, $n > 0$, $\bp'$ bijects precisely the set of $\left(2 \cdot n - 1\right)$ pairs
  \[\left\{\tuple{n - 1, i} : 0 \leq i < n\right\} \cup \left\{\tuple{i, n - 1} : 0 \leq i < n - 1\right\}\]
  onto $\left\{\left(n - 1\right)^2, \ldots, n^2 - 1\right\}$.
\end{proof}

It follows that $\tuple{a, b} \mapsto \bp' ab$ is a bijection $\omega^2 \rightarrow \omega$. Since $\bp'$ is a total, recursive function, it is also a recursive process to search for the unique numbers $a, b \in \omega$ such that $\bp' ab = x$ for a given $x \in \omega$. Therefore, by \autoref{lem:rec-in-pca}, it follows that we have the corresponding unpairing functions $\bp'_0, \bp'_1 \in \Kl$ and we do not need to specify concrete mathematical expressions to compute them. Additionally, this specific $\bp'$ is chosen to satisfy the following property:
\begin{lemma}
  \label{lem:max-pairing-property}
  For any $i, j \in \omega$ where $i \neq j$, there exist arbitrarily large $n \in \omega$ such that $\bp' in < \bp' jn$.
\end{lemma}

\begin{proof}
  It is easy to check that if $i > j$, then for any even $n > \max\mleft\{i, j\mright\}$ we will have $\bp' in < \bp' jn$; otherwise, $i < j$ and for any odd $n > \max\mleft\{i, j\mright\}$ we will have $\bp' in < \bp' jn$.
\end{proof}

We now build the infinite set $A$ in the style of some recursion-theoretic priority argument, as the set of an increasing sequence of finite tuples
\[t_0 \sqsubseteq t_1 \sqsubseteq t_2 \sqsubseteq \cdots\]
where each $t_i \in \omega^{< \omega}$ and $t_i \sqsubseteq t_j$, or equivalently, $\bl_s t_i \leq \bl_s t_j$ and $\bt_s t_j \left(\bl_s t_i\right) = t_i$ for any $i \leq j$. We then propose a countable number of requirements on finite tuples: for any $i, j, f \in \omega$ where $i \neq j$, we say that $t \in \omega^{< \omega}$ satisfies the requirement with index $\tuple{i, j, f}$ if there exists some $n \in \omega$ such that $i, j < n$, $\bp' in, \bp' jn \leq \bl_s t$, and
\[f \mleft(\bt_s t \mleft(\bp' in\mright)\mright) {\downarrow} \rightarrow f\mleft(\bt_s t \mleft(\bp' in\mright)\mright) \neq \bt_s t \mleft(\bp' jn\mright),\]
i.e.\ that $f$, as a function in the pca $\Kl$, does not correctly compute the first $\bp' jn$ components of the sequence $t$ given its first $\bp' in$ components, for some $n \in \omega$ where $t$ has enough components. Clearly if some tuple $t$ satisfies a requirement here, then every extension of $t$ also satisfies it; we would like that for every such requirement, $A$ contains a tail of tuples satisfying it. Formally, we prove the following statement:
\begin{proposition}[$\KP$]
  \label{prop:mltt-real-non-comp-path}
  There is some $h : \omega \rightarrow \omega$ such that for any $i, j, f \in \omega$ where $i \neq j$, the requirement with index $\tuple{i, j, f}$ is satisfied by some finite initial segment of $h$.
\end{proposition}

\begin{proof}
  As we mentioned, the collection of finite tuples $t \in \omega^{< \omega}$ satisfying each single requirement is an upwards closed set, and \autoref{lem:max-pairing-property} ensures that this set will be dense: for any $i, j \in \omega$ where $i \neq j$ and any tuple $t$, take some $n > i, j, \bl_s t$ such that $\bp' in < \bp' jn$. It follows that
  \[\bp' jn > \bp' in \geq n^2 > \bl_s t\]
  by \autoref{lem:max-pairing-bij}. We extend $t$ to an arbitrary $t' \sqsupseteq t$ such that $\bp' in \leq \bl_s t' < \bp' jn$. Then for any further extension $t'' \sqsupseteq t'$, the value $f \mleft(\bt_s t'' \mleft(\bp' in\mright)\mright) = f \mleft(\bt_s t' \mleft(\bp' in\mright)\mright)$ (or the non-termination thereof) is already determined. Clearly by choosing a suitable extension $t''$ of length $\bp' jn$, different from the output of $f$, we can make sure that $f$ does not compute $\bt_s t'' \mleft(\bp' jn\mright) = t''$ correctly.

  Therefore, we can now construct an increasing sequence of tuples
  \[t_0 \sqsubseteq t_1 \sqsubseteq t_2 \sqsubseteq \cdots\]
  where $t_0$ is the empty tuple and each $t_{i + 1}$ is an extension of $t_i$ found using the procedure above to satisfy the requirement with index $\tuple{\bp_{00} i, \bp_{01} i, \bp_1 i}$ if $\bp_{00} i \neq \bp_{01} i$, and simply $t_{i + 1} = t_i^\frown 0$ if $\bp_{00} i = \bp_{01} i$. This will not involve choice principles since at each step, we can always take an appropriate $t_{i + 1}$ whose code is the least in $\omega$. By $\Delta_0$-separation, clearly there exists $h : \omega \rightarrow \omega$ such that $h\mleft(n\mright)$ is the $n^\text{th}$ component in sufficiently long tuples in this sequence.
\end{proof}

The distinguished type $\overline{X}$ shall then be given by the set
\[A = \left\{t \in \omega : t \sqsubseteq h\right\},\]
i.e.\ the (codes of) finite initial segments of $h$. We now examine
\[\alpha_0 = \bp \overline{X} \mleft(\oplam{t} \left(\bl_s t\right)^V\mright) \in \bV_A\]
in the realisability model, and see that it has the following basic properties:
\begin{lemma}
  $\bV_A$ realises the following:
  \begin{enumerate}[label=\uplabel{\roman*}]
    \item \label{itm:mltt-incom-ord-basic-ord} $\alpha_0 \in \Ord$;
    \item \label{itm:mltt-incom-ord-basic-sub} $\alpha_0 \subseteq \omega$.
  \end{enumerate}
\end{lemma}

\begin{proof}
  First check that
  \[\oplam{t} \bp \mleft(\bl_s t\mright) \ior \Vdash \opA{\alpha \in \alpha_0} \alpha \in \omega^V,\]
  thus \ref{itm:mltt-incom-ord-basic-sub} is realised.

  Now, for \ref{itm:mltt-incom-ord-basic-ord}, observe that for any $t \in A$ and $0 \leq n \leq \bl_s t$, $\bp \mleft(\bt_s tn\mright) \ior \Vdash n^V \in \alpha_0$. Therefore,
  \[\oplam{tx} \bp \mleft(\bt_s tx\mright) \ior \Vdash \opA{\alpha \in \alpha_0} \opA{\beta \in \alpha} \beta \in \alpha_0,\]
  realising that $\alpha_0$ is transitive. Now it follows immediately from $\alpha_0 \subseteq \omega$ within $\CZF$ that any element in $\alpha_0$ is transitive, thus $\alpha_0 \in \Ord$.
\end{proof}

Next, we use the fact that our pairing function $\bp'$ is recursive; this means we can write down the same definition inside the realisability model (or even inside $L$) to define a set-theoretic function from $\omega^2$ to $\omega$ that performs exactly the same computations. More explicitly, we construct the following name of a set-theoretic pairing function
\[\overline{p} = \bp \overline{\omega} \mleft(\oplam{x} \tuple{\tuple{\left(\bp'_0 x\right)^V, \left(\bp'_1 x\right)^V}^V, x^V}^V\mright) \in \bV_A,\]
then the realisability model will interpret $\overline{p}\mleft(i^V, j^V\mright) = \left(\bp' ij\right)^V$ for the names $i^V, j^V, \left(\bp' ij\right)^V \in \bV_A$, for any $i, j \in \omega$.

We use $\overline{p}$ inside the realisability model to extract the complexity coded into $\alpha_0$ in \autoref{prop:mltt-real-non-comp-path}, by defining a set-theoretic function $f_0 : \omega \rightarrow \Pow\mleft(\omega\mright)$ as
\[f_0\mleft(i\mright) = \left\{n \in \omega : \opA{k \in n} \overline{p}\mleft(i, k\mright) \in \alpha_0\right\}\]
within $\bV_A$. It is easy to check that this definition guarantees that $\opA{i \in \omega} f_0\mleft(i\mright) \in \Pow\mleft(\omega\mright) \cap \Ord$ is realised. We want to invoke \autoref{cor:omega-sub-strong-incom-code} on $f_0$, so we check that it is realised to be pairwise strongly incomparable:

\begin{proposition}
  \label{prop:f-subset-eq}
  For any $i, j \in \omega$, if some $d \Vdash f_0\mleft(i^V\mright) \subseteq f_0\mleft(j^V\mright)$, then $i = j$. Namely, it follows immediately that
  \[\oplam{ijx} \ior \Vdash \opA{i, j \in \omega} \left(f_0\mleft(i\mright) \subseteq f_0\mleft(j\mright) \rightarrow i = j\right).\]
\end{proposition}

\begin{proof}
  Assume otherwise, i.e.\ $i \neq j$. Reasoning inside $\bV_A$, we know that $f_0\mleft(i^V\mright) \subseteq f_0\mleft(j^V\mright)$, by definition of $f_0$, is equivalent to
  \[\opA{n \in \omega^V} \left(\opA{k \in n} \overline{p}\mleft(i^V, k\mright) \in \alpha_0 \rightarrow \opA{k \in n} \overline{p}\mleft(j^V, k\mright) \in \alpha_0\right),\]
  so suppose that this is realised by some $d \in \Kl$. Pedantically, we take the formula $\overline{p}\mleft(i^V, k^V\mright) \in \alpha_0$ as a shorthand for
  \[\opE{x \in \overline{p}} \opE{n \in \alpha_0} x = \tuple{\tuple{i^V, k^V}, n},\]
  thus if some $e \Vdash \overline{p}\mleft(i^V, k^V\mright) \in \alpha_0$, then $\bp_0 e = \bp' ik$, and $\bp_{10} e$ must be the unique finite sequence in $A$ of length precisely $\bp' ik$.

  Conversely, we can also find some function $q \in \Kl$ such that, suppose $t \in A$ is any finite sequence satisfying $\bl_s t \geq \bp' ik$, then
  \[\bp \mleft(\bp' ik\mright) \mleft(\bp \mleft(\bt_s t \mleft(\bp' ik\mright)\mright) \mleft(qik\mright)\mright) \Vdash \overline{p}\mleft(i^V, k^V\mright) \in \alpha_0.\]
  (Here we use $q$ to conceal the cumbersome process of writing out the internal structure of realisers for the defining formula of set-theoretic ordered pairs, but it is routine to check that this can be done.)

  Consequently, for any $n > \mathrm{max}\mleft\{i, j\mright\}$, we let $t \in A$ be any long enough finite sequence such that $\bp' in, \bp' jn \leq \bl_s t$. Observe that our stipulation of the pairing function $\bp'$ also ensures $\opA{k \leq n} \bp' ik \leq \bp' in$. Therefore,
  \[\oplam{k} \bp \mleft(\bp' ik\mright) \mleft(\bp \mleft(\bt_s t \mleft(\bp' ik\mright)\mright) \mleft(qik\mright)\mright) \Vdash \opA{k \in \left(\bsN n\right)^V} \overline{p}\mleft(i^V, k\mright) \in \alpha_0.\]
  In other words, if we pass this to the realiser $d$ at the beginning of this proof, i.e.\ we take $f = \oplam{nt} d\mleft(\bsN n\mright) \mleft(\oplam{k} \bp \mleft(\bp' ik\mright) \mleft(\bp \mleft(\bt_s t \mleft(\bp' ik\mright)\mright) \mleft(qik\mright)\mright)\mright) n$, then we have $fnt \Vdash \overline{p}\mleft(j^V, n^V\mright) \in \alpha_0$ and $\bp_{10} \mleft(fnt\mright) = \bt_s t \mleft(\bp' jn\mright)$.

  Observe that $\bl_s \mleft(\bt_s t \mleft(\bp' in\mright)\mright) = \bp' in$, i.e.\ $\bp'_1 \mleft(\bl_s \mleft(\bt_s t \mleft(\bp' in\mright)\mright)\mright) = n$. This means we can have a single-parameter function $g = \oplam{s} \bp_{10}\mleft(f \mleft(\bp'_1 \mleft(\bl_s s\mright)\mright) s\mright)$, such that
  \[g\mleft(\bt_s t \mleft(\bp' in\mright)\mright) {\downarrow} \land g\mleft(\bt_s t \mleft(\bp' in\mright)\mright) = \bt_s t \mleft(\bp' jn\mright).\]

  However, this recursive $g$ is now universal for any appropriate $n$ and long enough $t \in A$, whereas in our construction of $A$ through \autoref{prop:mltt-real-non-comp-path} we have required that there exist $n$ and $t \in A$ where $g$ does not compute $\bt_s t \mleft(\bp' jn\mright)$ from $\bt_s t \mleft(\bp' in\mright)$ correctly. Therefore, $i \neq j$ implies a contradiction and thus we must have $i = j$.
\end{proof}

Finally, it follows from the remark before \autoref{cor:omega-sub-strong-incom-code} that $\alpha_0$, as an element of $\Pow\mleft(\omega\mright) \cap \Ord$, satisfies
\[\bV_A \vDash \alpha_0 = L_{\alpha_0} \cap \Ord \in L.\]
Therefore, $f_0$, which is defined using parameters $\omega$ (which lies in $L$ by \autoref{cor:l-omega-ord}), $\overline{p}$ (which is recursively definable from $\omega$) and $\alpha_0 \in L$, also lies in $L$. By \autoref{cor:omega-sub-strong-incom-code}, we have $\bV_A \vDash \Pow\mleft(\omega\mright) \subseteq L$.

\begin{theorem}
  \label{thm:mltt-real-v-l}
  $\bV_A$ realises $V = L$.
\end{theorem}

\begin{proof}
  We work inside $\bV_A$, which is a model of $\CZF + \text{subcountability}$ by \autoref{sec:mltt-interp}.

  For any set $s$ consider the transitive closure $\TC\mleft(\left\{s\right\}\mright)$ and correspondingly a surjection $f : u \rightarrow \TC\mleft(\left\{s\right\}\mright)$ where $u \subseteq \omega$, by the axiom of subcountability. Without loss of generality, we can assume conveniently that $0 \in u$ and $f\mleft(0\mright) = s$. We define
  \[\sigma_s = \left\{n \in \omega : \opE{i, j \in u} \left(f\mleft(i\mright) \in f\mleft(j\mright) \land n = \overline{p}\mleft(i, j\mright)\right)\right\}.\]

  Given $\sigma_s$, we can recursively define some (possibly partial) function $g$ on domain $u$ as
  \[g\mleft(n\mright) = \left\{g\mleft(m\mright) : m \in u \land \overline{p}\mleft(m, n\mright) \in \sigma_s\right\}\]
  if $g\mleft(m\mright)$ exists for all such $m \in u$. The function $g$ is defined by a $\Sigma_1$-formula using parameters $u, \sigma_s$, where $u, \sigma_s \in L$ because $\bV_A$ realises $\Pow\mleft(\omega\mright) \subseteq L$. Using collection, we can easily prove by set induction on $x$ that, for any set $x$,
  \[\opA{n \in u} \left(f\mleft(n\mright) = x \rightarrow g\mleft(n\mright) = x \land L \vDash \text{$g\mleft(n\mright)$ exists}\right).\]
  Since $g$, as a $\Sigma_1$-definable function, must be upwards absolute between transitive classes, clearly $s = f\mleft(0\mright) = g\mleft(0\mright) \in L$.
\end{proof}

\begin{corollary}
  \label{cor:czf-equi-v-eq-l}
  The consistency of $\KP$ implies that of $\CZF + V = L$. It follows that
  \[\CZF + V = L \equiv_{\Con} \KP \equiv_{\Con} \CZF.\]
\end{corollary}

\begin{proof}
  This follows immediately from \autoref{thm:mltt-real-czf} and \autoref{thm:mltt-real-v-l}.
\end{proof}

\chapter{Plump ordinals and powersets}
\label{chap:plump-power}

From \autoref{chap:incom-coding}, we observe that the primary difficulty when using the method of incomparable codings to work with $L$ is that, given a collection of incomparable ordinals $s \subseteq L \cap \Ord$ and an arbitrary subset $u \subseteq s$, it can be hard to ensure the single ordinal that can be used to uniquely identify this set, $\bigcup_{\alpha \in u} \alpha^+$, is guaranteed to lie in $L$ without knowing that $u \in L$, essentially due to \autoref{prob:ord-in-l} that there is no known proof of $\Ord \subseteq L$ in intuitionistic logic. Consequently, we often need to resort to more complicated, non-systematic constructions such as previously done in \autoref{prop:l-ord-strong-incom-code} and \autoref{cor:omega-sub-strong-incom-code}.

In this chapter, we will formalise the notion of \emph{plump ordinals}, originally discussed in a category-theoretic context by Paul Taylor \cite{taylor96-intuitionistic-ordinals}. We show that assuming the powerset axiom, the unbounded subclass of plump ordinals will form a smaller `class of ordinals' that are always in $L$, and we will see how this simplifies arguments of incomparable codings. Using a weak realisability interpretation of $\CZFP$ inside classical $\KPP$, we will then extend the equiconsistency proof in \autoref{sec:czf-v-eq-l} to $\CZFP + V = L$.

\section{Theory of plump ordinals}
\label{sec:plump}

To begin, we present the general results on plump ordinals from my manuscript \cite{wang-sw26-plump}. Following Taylor \cite{taylor96-intuitionistic-ordinals}, we informally say that a class $\mathcal{O} \subseteq \Ord$ is \emph{the class of plump ordinals} if it contains precisely those ordinals $\alpha \in \Ord$ satisfying the following two conditions:
\begin{enumerate}[label=\uplabel{\arabic*}]
  \item \label{itm:pl-cls-con-trans} $\alpha \subseteq \mathcal{O}$,
  \item \label{itm:pl-cls-closed-under-pl-ord-subset} for any $\beta \in \alpha$, the class $\Pow\mleft(\beta\mright) \cap \mathcal{O} \subseteq \alpha$.
\end{enumerate}

Formally speaking, this is not a definition for the class $\mathcal{O}$, because it is self-referential and some form of recursion theorem is needed to justify that such a class $\mathcal{O}$ is indeed definable. But first observe that the class of plump ordinals, if it exists, is at least unique:
\begin{proposition}
  Suppose that $\mathcal{O}_1, \mathcal{O}_2$ are both `classes of plump ordinals' as just defined, then $\mathcal{O}_1 = \mathcal{O}_2$.
\end{proposition}

\begin{proof}
  It suffices to prove by set induction that for any $\alpha \in \mathcal{O}_1 \cup \mathcal{O}_2$, the classes
  \[\Pow\mleft(\alpha\mright) \cap \mathcal{O}_1 = \Pow\mleft(\alpha\mright) \cap \mathcal{O}_2.\]

  To begin, assume that $\alpha \in \mathcal{O}_1$, and we shall first show using the inductive hypothesis that $\Pow\mleft(\alpha\mright) \cap \mathcal{O}_1 \subseteq \Pow\mleft(\alpha\mright) \cap \mathcal{O}_2$. Consider any $\beta \subseteq \alpha$ such that $\beta \in \mathcal{O}_1$, then for any $\gamma \in \beta \subseteq \alpha$, we also have $\gamma \in \mathcal{O}_1$. By the inductive hypothesis,
  \[\gamma \in \Pow\mleft(\gamma\mright) \cap \mathcal{O}_1 = \Pow\mleft(\gamma\mright) \cap \mathcal{O}_2,\]
  i.e.\ $\gamma \in \mathcal{O}_2$ as well. Therefore, $\beta \subseteq \mathcal{O}_2$.

  Additionally, for any $\gamma \in \beta \subseteq \alpha$, we have $\Pow\mleft(\gamma\mright) \cap \mathcal{O}_1 \subseteq \beta$. Thus again by the inductive hypothesis,
  \[\Pow\mleft(\gamma\mright) \cap \mathcal{O}_2 = \Pow\mleft(\gamma\mright) \cap \mathcal{O}_1 \subseteq \beta.\]
  Therefore, both criteria are verified, and by the definition of $\mathcal{O}_2$, we must also have $\beta \in \mathcal{O}_2$.

  Now, by an entirely symmetrical argument, we can show that $\alpha \in \mathcal{O}_2$ implies $\Pow\mleft(\alpha\mright) \cap \mathcal{O}_2 \subseteq \Pow\mleft(\alpha\mright) \cap \mathcal{O}_1$. Notice that as $\alpha \in \Pow\mleft(\alpha\mright)$, this suffices for the conclusion that for any $\alpha \in \mathcal{O}_1 \cup \mathcal{O}_2$, we have
  \[\alpha \in \mathcal{O}_1 \leftrightarrow \alpha \in \mathcal{O}_2.\]
  In other words, $\alpha \in \mathcal{O}_1 \cup \mathcal{O}_2$ implies both $\alpha \in \mathcal{O}_1$ and $\alpha \in \mathcal{O}_2$, and thus $\Pow\mleft(\alpha\mright) \cap \mathcal{O}_1 = \Pow\mleft(\alpha\mright) \cap \mathcal{O}_2$ as desired.
\end{proof}

Inside the theory $\IKPP$ where $\Pow$ is a primitive function symbol that is allowed to occur in $\Sigma_1$-formulae, it is not hard to imitate \cite{taylor96-intuitionistic-ordinals}*{Proposition 4.3} and use $\Sigma$-recursion on $\alpha \in \Ord$ to define the family of sets $\theta_\alpha \subseteq \Pow\mleft(\alpha\mright)$ satisfying
\[\theta_\alpha = \left\{\beta \in \Pow\mleft(\alpha\mright) \cap \Ord : \opA{\gamma \in \beta} \left(\gamma \in \theta_\gamma \land \theta_\gamma \subseteq \beta\right)\right\}\]
for each $\alpha \in \Ord$. Then it is easy to verify that $\mathcal{O} = \left\{\alpha \in \Ord : \alpha \in \theta_\alpha\right\}$ will be the class of plump ordinals.

However, for our purpose of studying the interaction of the plump ordinals and G\"odel's constructible universe $L$, it can be difficult to adapt this construction since the powerset axiom is known to fail in $L$ in extensions of $\IKPP$ (cf.\ \cite{matthews-rathjen24-constructible-universe}*{section 7}, and previous discussions in \autoref{sec:overview-v-eq-l}). Thus, we will use the alternative presentation in \cite{wang-sw26-plump} where we show that the class of plump ordinals can be defined by a $\Pi_1$-formula --- in fact, a $\Delta_0$-formula if we consider $\Pow$ as a primitive function symbol --- in any extension of $\IKP$.

To this end, we denote
\[\relpl_\alpha\mleft(\gamma\mright) \defeqv \opA{\delta \in \gamma} \opA{\varepsilon \in \alpha} \left(\varepsilon \subseteq \delta \rightarrow \varepsilon \in \gamma\right)\]
and say that \emph{$\gamma$ is plump relative to $\alpha$} when this holds. We define $\PlOrd$\inds{plump ordinals, $\PlOrd$} as the $\Pi_1$-subclass of $\alpha \in \Ord$ satisfying
\[\opA{\beta \in \alpha} \opA{\gamma \subseteq \beta} \left(\relpl_\alpha\mleft(\gamma\mright) \rightarrow \gamma \in \alpha \land \opA{\delta \in \alpha} \left(\beta \in \delta \rightarrow \gamma \in \delta\right)\right).\]
We shall show that $\PlOrd$ is precisely the class of plump ordinals.

Let us first observe the following basic properties:
\begin{lemma}
  \label{lem:pl-ord-elem-rel-pl}
  For any $\alpha \in \PlOrd$, if $\beta \in \alpha$, then $\relpl_\alpha\mleft(\beta\mright)$ holds.
\end{lemma}

\begin{proof}
  We show by set induction on $\beta$ that if $\beta \in \alpha$, then
  \[\opA{\gamma \in \alpha} \left(\gamma \subseteq \beta \rightarrow \relpl_\alpha\mleft(\gamma\mright)\right).\]

  Fix $\gamma, \beta \in \alpha$ such that $\gamma \subseteq \beta$, and assume that $\delta \in \gamma$, $\varepsilon \in \alpha$ and $\varepsilon \subseteq \delta$. It follows that $\delta \in \beta$ (and $\delta \in \alpha$ by transitivity), thus by the inductive hypothesis $\relpl_\alpha\mleft(\varepsilon\mright)$ holds. We know that $\varepsilon \subseteq \delta$, $\gamma \in \alpha$ and $\delta \in \gamma$, then the assumption $\alpha \in \PlOrd$ must imply that $\varepsilon \in \gamma$, precisely what we need to conclude that $\relpl_\alpha\mleft(\gamma\mright)$ holds.
\end{proof}

\begin{lemma}
  \label{lem:pl-ord-subset-rel-pl-is-pl-ord}
  For any $\alpha \in \PlOrd$, if $\beta \subseteq \alpha$ and $\relpl_\alpha\mleft(\beta\mright)$ holds, then $\beta \in \PlOrd$.
\end{lemma}

\begin{proof}
  First notice that $\beta \subseteq \alpha$ and $\relpl_\alpha\mleft(\beta\mright)$ imply that $\beta$ is transitive: consider any $\gamma \in \beta$, $\delta \in \gamma$, we know that $\delta \in \alpha$ since $\alpha \in \Ord$ and also $\delta \subseteq \gamma$ since $\gamma \in \Ord$. It follows from $\relpl_\alpha\mleft(\beta\mright)$ that we must have $\delta \in \beta$.

  Now suppose that $\beta' \in \beta$, $\gamma \subseteq \beta'$ and $\relpl_\beta\mleft(\gamma\mright)$ holds. For any $\delta \in \gamma$ and $\varepsilon \in \alpha$ such that $\varepsilon \subseteq \delta$, we know by transitivity that $\delta \in \beta$, so it follows from $\relpl_\alpha\mleft(\beta\mright)$ that $\varepsilon \in \beta$, and it follows again from $\relpl_\beta\mleft(\gamma\mright)$ that $\varepsilon \in \gamma$. In other words, $\relpl_\alpha\mleft(\gamma\mright)$ holds.

  From the assumption that $\alpha \in \PlOrd$, we then know that $\gamma \in \alpha$. Since also $\beta' \in \beta$ and $\gamma \subseteq \beta'$, it follows from $\relpl_\alpha\mleft(\beta\mright)$ that $\gamma \in \beta$. For any $\delta \in \beta$, we now have $\beta' \in \alpha$, $\gamma \subseteq \beta'$ and $\delta \in \alpha$, so it follows from $\relpl_\alpha\mleft(\gamma\mright)$ that $\beta' \in \delta$ implies $\gamma \in \delta$. We have shown that $\beta \in \PlOrd$.
\end{proof}

\begin{proposition}
  \label{prop:pl-ord-class-trans}
  For any $\alpha \in \PlOrd$, if $\beta \in \alpha$, then $\beta \in \PlOrd$.
\end{proposition}

\begin{proof}
  We have $\relpl_\alpha\mleft(\beta\mright)$ by \autoref{lem:pl-ord-elem-rel-pl}, and we also have $\beta \subseteq \alpha$ by transitivity. Thus we have $\beta \in \PlOrd$ by \autoref{lem:pl-ord-subset-rel-pl-is-pl-ord}.
\end{proof}

\begin{lemma}
  \label{lem:subset-pl-ord-sat-relpl}
  For any $\alpha, \beta \in \PlOrd$, if $\beta \subseteq \alpha$, then $\relpl_\alpha\mleft(\beta\mright)$ holds.
\end{lemma}

\begin{proof}
  Suppose that $\delta \in \beta$, $\varepsilon \in \alpha$ and $\varepsilon \subseteq \delta$. By \autoref{lem:pl-ord-elem-rel-pl}, we know that $\relpl_\alpha\mleft(\varepsilon\mright)$ holds. Since $\beta \subseteq \alpha$, it is easy to check that $\relpl_\beta\mleft(\varepsilon\mright)$ also holds. Then the assumption that $\beta \in \PlOrd$ immediately implies $\varepsilon \in \beta$, precisely what we need to show.
\end{proof}

\begin{proposition}
  \label{prop:pl-ord-closed-under-pl-ord-subset}
  For any $\alpha, \gamma \in \PlOrd$, if $\gamma \subseteq \beta$ for some $\beta \in \alpha$, then $\gamma \in \alpha$.
\end{proposition}

\begin{proof}
  Since $\alpha \in \PlOrd$, it suffices to show that $\relpl_\alpha\mleft(\gamma\mright)$ holds. But we know that $\gamma \subseteq \beta \subseteq \alpha$ by transitivity, thus \autoref{lem:subset-pl-ord-sat-relpl} suffices for the proof.
\end{proof}

Here \autoref{prop:pl-ord-class-trans} states that any $\alpha \in \PlOrd$ satisfies the aforementioned condition \ref{itm:pl-cls-con-trans} for the `class of plump ordinals', whereas \autoref{prop:pl-ord-closed-under-pl-ord-subset} states that any $\alpha \in \PlOrd$ satisfies condition \ref{itm:pl-cls-closed-under-pl-ord-subset}. For the converse, we need to verify that:
\begin{proposition}
  \label{prop:pl-ord-crit}
  For any set $x$, we have $x \in \PlOrd$ if both of the following hold:
  \begin{enumerate}[label=\uplabel{\arabic*}]
    \item \label{itm:pl-ord-proof-con-trans} for any $\beta \in x$, $\beta \in \PlOrd$;
    \item \label{itm:pl-ord-proof-closed-under-pl-ord-subset} for any $\gamma \in \PlOrd$, if $\gamma \subseteq \beta$ for some $\beta \in x$, then $\gamma \in x$.
  \end{enumerate}
\end{proposition}

\begin{proof}
  Condition \ref{itm:pl-ord-proof-con-trans} already ensures that $x \subseteq \Ord$. For transitivity, note that for any $\beta \in x$, $\gamma \in \beta$, we know that $\gamma \subseteq \beta$ since $\beta \in \Ord$, and $\gamma \in \PlOrd$ by \autoref{prop:pl-ord-class-trans}, thus condition \ref{itm:pl-ord-proof-closed-under-pl-ord-subset} implies that $\gamma \in x$ as well.

  Now suppose that $\beta \in x$, $\gamma \subseteq \beta$ and $\relpl_x\mleft(\gamma\mright)$ holds. Since $\beta \subseteq x$ by transitivity, it follows immediately from definition that $\relpl_\beta\mleft(\gamma\mright)$ also holds. Here $\beta \in \PlOrd$, so by \autoref{lem:pl-ord-subset-rel-pl-is-pl-ord}, $\gamma \in \PlOrd$, thus condition \ref{itm:pl-ord-proof-closed-under-pl-ord-subset} ensures that $\gamma \in x$.

  Additionally, for any $\delta \in x$, suppose that $\beta \in \delta$. Observe that we also have $\relpl_\delta\mleft(\gamma\mright)$ and $\delta \in \PlOrd$, so it follows immediately that $\gamma \in \delta$ as well.
\end{proof}

This allows us to finally conclude that:
\begin{corollary}[$\IKP$]
  $\PlOrd$ is the class of plump ordinals.
\end{corollary}

Observe that combining \autoref{lem:pl-ord-subset-rel-pl-is-pl-ord} and \autoref{lem:subset-pl-ord-sat-relpl}, we know that for any $\alpha \in \PlOrd$ and any $\beta \subseteq \alpha$,
\[\beta \in \PlOrd \leftrightarrow \relpl_\alpha\mleft(\beta\mright),\]
where the right-hand side of the equivalence is a $\Delta_0$-formula. Therefore, for any $\alpha \in \PlOrd$, we can define its \emph{plump successor} as
\[\alpha^\pls = \left\{\beta \in \Pow\mleft(\alpha\mright) : \relpl_\alpha\mleft(\beta\mright)\right\} = \Pow\mleft(\alpha\mright) \cap \PlOrd\]
if this exists as a set; and we know in $\IKP$ that whenever the powerset $\Pow\mleft(\alpha\mright)$ exists, the plump successor $\alpha^\pls$ also exists by $\Delta_0$-separation. We verify that:
\begin{lemma}
  \label{lem:pl-successor-still-pl}
  For any $\alpha \in \PlOrd$, if $\alpha^\pls$ exists, then $\alpha^\pls \in \PlOrd$.
\end{lemma}

\begin{proof}
  We just need to check the two criteria in \autoref{prop:pl-ord-crit}: for any $\beta \in \alpha^\pls$, we have $\beta \in \PlOrd$ by definition; for any $\gamma \in \PlOrd$, if $\gamma \subseteq \beta$ for some $\beta \in \alpha^\pls$, then $\gamma \subseteq \alpha$ and $\gamma \in \alpha^\pls$ again by definition.
\end{proof}

\begin{lemma}
  \label{lem:pl-union-still-pl}
  For any set $s \subseteq \PlOrd$, $\bigcup s \in \PlOrd$.
\end{lemma}

\begin{proof}
  We again check the two criteria in \autoref{prop:pl-ord-crit}: for any $\beta \in \bigcup s$, we have $\beta \in \alpha$ for some $\alpha \in s \subseteq \PlOrd$. Then $\beta \in \PlOrd$ by \autoref{prop:pl-ord-class-trans}. For any $\gamma \in \PlOrd$, if $\gamma \subseteq \beta$ for some $\beta \in \bigcup s$, then again we have $\beta \in \alpha$ for some $\alpha \in s \subseteq \PlOrd$, and it follows that $\gamma \in \alpha \subseteq \bigcup s$ by \autoref{prop:pl-ord-closed-under-pl-ord-subset}.
\end{proof}

Note that $\IKP$ alone, however, does not suffice to prove that many plump ordinals exist. We can observe that:
\begin{lemma}
  Any set $\alpha \subseteq 1 = \left\{\varnothing\right\}$ is a plump ordinal.
\end{lemma}

\begin{proof}
  Check the two criteria in \autoref{prop:pl-ord-crit}: for any $\beta \in \alpha$, we have $\beta = \varnothing$ and clearly $\beta \in \PlOrd$ by definition; for any $\gamma \subseteq \beta$, we also have $\gamma = \beta = \varnothing$ and thus $\gamma \in \alpha$ as well.
\end{proof}

It follows that $1 \in \PlOrd$ and also $1^\pls = \Pow\mleft(1\mright) \in \PlOrd$ (if it exists) by \autoref{lem:pl-successor-still-pl}. So $\IKP$ cannot even prove that $1^\pls$ exists.

In contrast, in $\IKPP$ or any extension of it, we can define class functions by $\Sigma$-recursion on a defining formula that contains the primitive symbol $\Pow$, so we can define the \emph{plump operator} $\left(-\right)^\pl : \Ord \rightarrow \PlOrd$ recursively as
\[\alpha^\pl = \bigcup_{\beta \in \alpha} \left(\beta^\pl\right)^\pls.\]
It follows from \autoref{lem:pl-successor-still-pl} and \autoref{lem:pl-union-still-pl} that this is well-defined and $\alpha^\pl \in \PlOrd$ for every $\alpha \in \Ord$. It is easy to check that:
\begin{proposition}
  \label{prop:pl-op-surj}
  For any $\alpha \in \PlOrd$, $\alpha^\pl = \alpha$. Thus $\left(-\right)^\pl$ is a surjection onto $\PlOrd$.
\end{proposition}

\begin{proof}
  It follows immediately from \autoref{prop:pl-ord-closed-under-pl-ord-subset} that for any $\alpha, \beta \in \PlOrd$, if $\beta \in \alpha$ then also $\beta^\pls \subseteq \alpha$. Thus, by the inductive hypothesis, we have
  \[\alpha^\pl = \bigcup_{\beta \in \alpha} \left(\beta^\pl\right)^\pls = \bigcup_{\beta \in \alpha} \beta^\pls = \alpha\]
  for any $\alpha \in \PlOrd$.
\end{proof}

\subsection*{Plump ordinals and $L$}

We now proceed to demonstrate that, unlike the complicated scenario with arbitrary ordinals and $L$ in intuitionistic logic, as we presented in \autoref{sec:ord-modal}, the behaviours of plump ordinals are much simpler to characterise. We will always have that:
\begin{lemma}[$\IKP$, \cite{wang-sw26-plump}*{Lemma 3.5}]
  \label{lem:uniform-def-pl-ord-in-l}
  For any $\alpha \in \PlOrd$, we have
  \[\alpha = \left\{x \in L_\alpha : L_\alpha \vDash x \in \PlOrd\right\} \in \dd\mleft(L_\alpha\mright).\]
\end{lemma}

\begin{proof}
  We shall prove this by set induction on $\alpha$. Fix some $\alpha \in \PlOrd$, then the inductive hypothesis says that for any $\beta \in \alpha$,
  \[\beta = \left\{x \in L_\beta : L_\beta \vDash x \in \PlOrd\right\} \in \dd\mleft(L_\beta\mright),\]
  so immediately $\alpha \subseteq L_\alpha$. Notice additionally that $\PlOrd$ is a $\Pi_1$-class, which is downwards absolute between transitive sets and classes, thus for any $\beta \in \alpha$, since $\beta \in \PlOrd$, we must also have $L_\alpha \vDash \beta \in \PlOrd$.

  It remains to prove the converse: for any $x \in L_\alpha$ such that $L_\alpha \vDash x \in \PlOrd$, we want to show that $x \in \alpha$. There must exist some $\beta \in \alpha$ such that $x \in \dd\mleft(L_\beta\mright)$. Consider any $y \in x$, we observe that \autoref{lem:pl-ord-elem-rel-pl} and \autoref{lem:pl-ord-subset-rel-pl-is-pl-ord} use no other assumptions than the axiom of set induction, thus their relativisations to any transitive set still hold, and we can relativise \autoref{prop:pl-ord-class-trans}:
  \[L_\alpha \vDash \opA{x \in \PlOrd} \opA{y \in x} y \in \PlOrd.\]
  Here, $\alpha$ is transitive, thus $\beta \subseteq \alpha$ and $L_\beta \subseteq L_\alpha$. Again by downward absoluteness, we know that $y \in L_\beta$ satisfies $L_\beta \vDash y \in \PlOrd$. By the inductive hypothesis, this means $y \in \beta$, i.e.\ $x \subseteq \beta$.

  Finally, we observe that \autoref{lem:subset-pl-ord-sat-relpl} also relativises to the transitive set $L_\alpha$, thus we have $L_\alpha \vDash \relpl_\beta\mleft(x\mright)$. However, $\relpl_\beta\mleft(x\mright)$ is a $\Delta_0$-formula, so it also holds in the original universe $V$. Thus, $x \in \PlOrd$ in $V$ by \autoref{lem:pl-ord-subset-rel-pl-is-pl-ord}, and since $x \subseteq \beta$, $\alpha \in \PlOrd$ immediately implies that $x \in \alpha$.
\end{proof}

\begin{corollary}[$\IKP$]
  \label{cor:l-preserve-pl-ord}
  $\PlOrd = \PlOrd^L$.
\end{corollary}

\begin{proof}
  For any $\alpha \in \PlOrd$, by \autoref{lem:uniform-def-pl-ord-in-l}, $\alpha \in \dd\mleft(L_\alpha\mright)$, thus $\alpha \in L_{\alpha^+} \in L$. We also have $L \vDash \alpha \in \PlOrd$ by downward absoluteness, thus $\PlOrd \subseteq \PlOrd^L$.

  Conversely, we just show that $\alpha \in \PlOrd^L$ implies $\alpha \in \PlOrd$ by set induction on $\alpha$. Using the fact that $L \vDash \IKP$ (cf.\ \autoref{sec:overview-v-eq-l}), we know that our lemmata on plump ordinals all work inside $L$. Thus, assuming $L \vDash \alpha \in \PlOrd$, it follows that for any $\beta \in \alpha$, we have $L \vDash \beta \in \PlOrd$ by \autoref{prop:pl-ord-class-trans} and thus actually $\beta \in \PlOrd$ by the inductive hypothesis.

  Additionally, for any $\beta \in \alpha$ and $\gamma \in \PlOrd$ such that $\gamma \subseteq \beta$, we have $\gamma \in L$ by the forward direction and thus $\gamma \in \alpha$ since $L \vDash \alpha \in \PlOrd$. It follows from \autoref{prop:pl-ord-crit} that $\alpha \in \PlOrd$ indeed.
\end{proof}

In \cite{wang-sw26-plump}, we further considered the class\footnote{It remains open, however, whether $L \neq L_\pl$ is a consistent axiom in any intuitionistic set theory.}
\[L_\pl = \bigcup_{\alpha \in \PlOrd} L_\alpha \subseteq L\]
and verified that also $\PlOrd = \PlOrd^{L_\pl}$. It is argued that, when $T$ denotes $\IKP$ together with a technical axiom scheme of \emph{plump $\Sigma_1$-bounding}, $\PlBSig$,
\[\opA{x \in a} \opE{\alpha \in \PlOrd} \varphi\mleft(x, \alpha\mright) \rightarrow \opE{\alpha \in \PlOrd} \opA{x \in a} \opE{\beta \in \alpha} \varphi\mleft(x, \beta\mright)\]
for any $\Sigma_1$-formula $\varphi$, then $L_\pl$ can be seen as some kind of `inner model' for $T$, in the sense that $T$ verifies both of the following:
\begin{enumerate}[label=\uplabel{\arabic*}]
  \item $L_\pl$ satisfies $T = \IKP + \PlBSig$ and also $V = L_\pl$;
  \item $L_\pl$ contains all plump ordinals.
\end{enumerate}

It is worth noting that restricted excluded middle, $\REM$, implies that ordinals are linearly ordered, and thus every ordinal is plump. Therefore $\IKP + \REM$ implies that $L_\pl = L$ and thus, by an entirely analogous argument as \autoref{prop:czf-l-not-exp-con}, we know that $\CZFP + \neg \Exp^{L_\pl}$ is consistent. It follows that if $T$ is any fragment of $\CZFP$ that proves $\Exp$, then $L_\pl$ may no longer be an inner model in the sense above.

However, we will not go further into the details as the only key fact relevant to this thesis is \autoref{lem:uniform-def-pl-ord-in-l}, since it implies that every plump ordinal is automatically in $L$ (or even $L_\pl$), and thus if we work with an incomparable encoding where all coding objects are plump ordinals, then every coded set will be definable within $L$ as well.

\subsection*{Plump ordinal arithmetic}

It is apparent that the automatic process with the weak incomparables in \autoref{sec:incom-intro} cannot be replicated within the plump ordinals --- especially there is no analogue of \autoref{lem:set-perp-plus-still-perp} where every involved ordinal is plump. However, there are still some semi-automatic processes that can simplify the construction and treatment of incomparable plump ordinals.

To present these results, we need to define a new notion of ordinal arithmetic\inds{plump ordinals, $\PlOrd$!arithmetic} that only produces plump ordinals, by replacing the successor operation used in usual ordinal arithmetic with the plump successor $\left(-\right)^\pl$. Specifically, we denote
\begin{align*}
  \alpha +^\pl \beta     & = \alpha \cup \bigcup_{\gamma \in \beta} \left(\alpha +^\pl \gamma\right)^\pls, \\
  \alpha \cdot^\pl \beta & = \bigcup_{\gamma \in \beta} \left(\alpha \cdot^\pl \gamma +^\pl \alpha\right)
\end{align*}
(cf.\ the usual ordinal addition in intuitionistic logic presented in \autoref{sec:incom-intro}). These operations are definable in $\IKPP$ by $\Sigma$-recursion on formulae containing the primitive function symbol $\Pow$, or by recursion using unrestricted strong collection in $\CZF$. Do notice that $\beta \in \Ord$ suffices for the following:
\begin{lemma}
  For any $\alpha \in \PlOrd$, $\beta \in \Ord$, we have $\alpha +^\pl \beta, \alpha \cdot^\pl \beta \in \PlOrd$.
\end{lemma}

\begin{proof}
  This follows from easy set induction on $\beta$, using \autoref{lem:pl-successor-still-pl} and \autoref{lem:pl-union-still-pl}.
\end{proof}

\label{loc:l-pl-rec-impl-remark}
Before we proceed, I want to make a technical remark on the recursive process involved here: since its defining formula invokes plumpness and thus the powerset operation, such definitions cannot be replicated as given inside $L$ in $\IKPP$ or $\CZF$, because in $L$, $\Exp$ and thus the powerset axiom may fail. However, we should observe that, given fixed $\alpha$ and $\beta$, there will be a large enough plump ordinal $\sigma$ which contains all the plump ordinals used in the process to define $\gamma +^\pl \delta$ for all $\gamma \in \alpha$, $\delta \in \beta$ --- namely, $\sigma = \alpha +^\pl \beta$ will work. It then follows that in the defining formula of this recursive process, every occurrence of $\gamma \in \PlOrd$ can be replaced by $\gamma \in \sigma$ and likewise, for example, every occurrence of $\gamma = \delta^\pls$ can be replaced by
\[\gamma \in \sigma \land \delta \in \gamma \land \opA{\eta \in \gamma} \eta \subseteq \delta.\]
It follows that plump ordinal addition up to $\alpha$ and $\beta$ can be defined by a $\Sigma_1$-formula
\[\gamma +^\pl \delta = \eta \defeqv \varphi\mleft(\gamma, \delta, \eta; \sigma\mright),\]
not containing the powerset symbol at all, but parametrised on $\sigma$. Since the plump ordinal parameter $\sigma$ must lie in $L$, by $\Sigma$-recursion in $\IKP$ we know that this operation is well-defined with its existential witnesses all inside $L$. By \autoref{cor:l-preserve-pl-ord}, we can then formulate plump ordinal addition in general as
\[\alpha +^\pl \beta = \eta \defeqv \opE{\sigma \in \PlOrd} \varphi\mleft(\alpha, \beta, \eta; \sigma\mright),\]
which will define the same relation when interpreted in both $V$ and $L$. Similar constructions work for plump ordinal multiplication and also for the recursive procedure we later employ in \autoref{thm:l-pl-ord-surj}, and we see that the same process can also be carried out in the subclass $L_\pl = \bigcup_{\alpha \in \PlOrd} L_\alpha$ as well.

Now, we check that the following basic properties of arithmetic hold. Observe that many short proofs here work specifically for plump ordinals and are not valid for the usual ordinal arithmetic on $\Ord$ in an intuitionistic context:
\begin{lemma}
  \label{lem:pl-add-inj}
  For any $\alpha, \beta, \gamma \in \PlOrd$, $\alpha +^\pl \beta \subseteq \alpha +^\pl \gamma$ implies $\beta \subseteq \gamma$. In other words, the class function $F_\alpha : \PlOrd \rightarrow \PlOrd$ given by $\beta \mapsto \alpha +^\pl \beta$ is injective.
\end{lemma}

\begin{proof}
  We show this by set induction on $\gamma$. Consider any $\delta \in \beta$, then $\alpha +^\pl \delta \in \left(\alpha +^\pl \delta\right)^\pls \subseteq \alpha +^\pl \beta$. By definition, $\alpha \subseteq \alpha +^\pl \delta$, so $\alpha +^\pl \delta \in \alpha$ would imply $\alpha \in \alpha$ by plumpness, thus a contradiction. Therefore, we must have $\alpha +^\pl \delta \in \left(\alpha +^\pl \varepsilon\right)^\pls$ for some $\varepsilon \in \gamma$. This means $\alpha +^\pl \delta \subseteq \alpha +^\pl \varepsilon$, and thus $\delta \subseteq \varepsilon$ by the inductive hypothesis. Since $\gamma \in \PlOrd$, we must have $\delta \in \gamma$.
\end{proof}

\begin{lemma}
  \label{lem:pl-mul-inj}
  For any $\alpha, \beta, \beta', \gamma \in \PlOrd$, if $\gamma \in \alpha$, then $\alpha \cdot^\pl \beta \subseteq \alpha \cdot^\pl \beta' +^\pl \gamma$ implies $\beta \subseteq \beta'$.
\end{lemma}

\begin{proof}
  We show this by set induction on $\beta'$. Consider any $\delta \in \beta$, we know that
  \[\alpha \cdot^\pl \delta +^\pl \gamma \in \left(\alpha \cdot^\pl \delta +^\pl \gamma\right)^\pls \subseteq \alpha \cdot^\pl \delta +^\pl \alpha.\]
  Thus $\alpha \cdot^\pl \delta +^\pl \gamma \in \alpha \cdot^\pl \beta \subseteq \alpha \cdot^\pl \beta' +^\pl \gamma$.

  Now, it suffices to show, by another set induction on $\gamma$ here, that $\alpha \cdot^\pl \delta +^\pl \gamma \in \alpha \cdot^\pl \beta' +^\pl \gamma$ implies $\delta \in \beta'$ (for the fixed $\beta'$ above). By the definition of plump ordinal addition, we must have either $\alpha \cdot^\pl \delta +^\pl \gamma \in \alpha \cdot^\pl \beta'$, or $\alpha \cdot^\pl \delta +^\pl \gamma \in \left(\alpha \cdot^\pl \beta' +^\pl \gamma'\right)^\pls$ for some $\gamma' \in \gamma$.

  In the first case, there must exist $\varepsilon \in \beta'$ such that $\alpha \cdot^\pl \delta +^\pl \gamma \in \alpha \cdot^\pl \varepsilon +^\pl \alpha$, i.e.\ there exists $\gamma' \in \alpha$ such that $\alpha \cdot^\pl \delta +^\pl \gamma \in \left(\alpha \cdot^\pl \varepsilon +^\pl \gamma'\right)^\pls$ and hence $\alpha \cdot^\pl \delta \subseteq \alpha \cdot^\pl \delta +^\pl \gamma \subseteq \alpha \cdot^\pl \varepsilon +^\pl \gamma'$. By the inductive hypothesis of the outer induction, $\delta \subseteq \varepsilon$, thus $\delta \in \beta'$ by plumpness.

  In the second case, we must have
  \[\alpha \cdot^\pl \delta +^\pl \gamma' \in \left(\alpha \cdot^\pl \delta +^\pl \gamma'\right)^\pls \subseteq \alpha \cdot^\pl \delta +^\pl \gamma \subseteq \alpha \cdot^\pl \beta' +^\pl \gamma',\]
  and thus $\delta \in \beta'$ by the inductive hypothesis of the inner induction.
\end{proof}

\begin{corollary}
  \label{cor:pl-ord-code-pairs}
  For any $\alpha, \beta \in \PlOrd$, the function $f : \alpha \times \beta \rightarrow \alpha \cdot^\pl \beta$ given by $\tuple{\gamma, \delta} \mapsto \alpha \cdot^\pl \delta +^\pl \gamma$ is injective.
\end{corollary}

\begin{proof}
  We first check that $f$ is indeed well-defined on codomain $\alpha \cdot^\pl \beta$: for any $\gamma \in \alpha, \delta \in \beta$, we see that
  \[\alpha \cdot^\pl \delta +^\pl \gamma \in \left(\alpha \cdot^\pl \delta +^\pl \gamma\right)^\pls \subseteq \alpha \cdot^\pl \delta +^\pl \alpha \subseteq \alpha \cdot^\pl \beta.\]

  Now, suppose that $f\mleft(\gamma, \delta\mright) = f\mleft(\gamma', \delta'\mright)$. Observe that we have $\alpha \cdot^\pl \delta \subseteq f\mleft(\gamma, \delta\mright) = \alpha \cdot^\pl \delta' +^\pl \gamma'$ and thus $\delta \subseteq \delta'$ by \autoref{lem:pl-mul-inj}. Symmetrically, we have $\delta' \subseteq \delta$, i.e.\ $\delta = \delta'$. This implies $\alpha \cdot^\pl \delta = \alpha \cdot^\pl \delta'$ and hence furthermore $\gamma = \gamma'$ by \autoref{lem:pl-add-inj}.
\end{proof}

We will first use this construction to define an intuitionistic analogue of the classical result that every $L_\alpha$ is in ($\Sigma_1$-definable) bijection with an initial segment of ordinals:
\begin{theorem}[$\IKPP$, \cite{wang-sw26-plump}*{Theorem 4.5}]
  \label{thm:l-pl-ord-surj}
  There exist class functions $\Theta, I$ on domain $\PlOrd$, such that for any $\alpha \in \PlOrd$, we have $\Theta\mleft(\alpha\mright) \in \PlOrd$ and $I\mleft(\alpha\mright) \subseteq \Theta\mleft(\alpha\mright) \times L_\alpha$ is a surjection from a subset of $\Theta\mleft(\alpha\mright)$ onto $L_\alpha$.
\end{theorem}

\begin{proof}
  We define $\Theta$ and $I$ simultaneously by $\Sigma$-recursion, using the following process where $\Pow$ is treated as a primitive symbol (and thus $\PlOrd$ is $\Delta_0$). Fix some $\alpha \in \PlOrd$, denote $\gamma = \bigcup_{\beta \in \alpha} \Theta\mleft(\beta\mright)$ and let
  \[\Theta\mleft(\alpha\mright) = \alpha \cdot^\pl \omega^\pl \cdot^\pl \gamma^\omega,\]
  where $\omega^\pl$ equals the union of all finite plump successors of $\varnothing$ and $\gamma^\omega$ is the union of all finite plump products $\gamma \cdot^\pl \gamma \cdot^\pl \cdots$. $I\mleft(\alpha\mright)$ will precisely contain those $\tuple{\delta, x} \in \Theta\mleft(\alpha\mright) \times L_\alpha$ such that there exists $\beta \in \alpha$, formula $\varphi$ and $\varepsilon_1, \ldots, \varepsilon_n \in \mathrm{dom}\mleft(I\mleft(\beta\mright)\mright) \subseteq \Theta\mleft(\beta\mright)$ such that
  \begin{enumerate}[label=\uplabel{\roman*}]
    \item $\delta = \sum_{i = 1}^n \alpha \cdot^\pl \omega^\pl \cdot^\pl \gamma^{n - i} \cdot^\pl \varepsilon_i +^\pl \alpha \cdot^\pl q^\pl +^\pl \beta$, where $q \in \omega$ is the G\"odel numbering of $\varphi$;
    \item $x = \left\{z \in L_\beta : L_\beta \vDash \varphi\mleft(z, I\mleft(\beta\mright)\mleft(\varepsilon_1\mright), \ldots, I\mleft(\beta\mright)\mleft(\varepsilon_n\mright)\mright)\right\} \in \dd\mleft(L_\beta\mright)$.
  \end{enumerate}
  This is expressible internally in $\IKPP$ using the standard G\"odel numbering and the internal truth predicates for set models presented in \autoref{sec:intro-int-l}.

  Now, repeated applications of \autoref{cor:pl-ord-code-pairs} shows immediately that if $\tuple{\delta, x}, \tuple{\delta', x'} \in I\mleft(\alpha\mright)$ with $\delta = \delta'$, then $x$ and $x'$ are defined through exactly the same formula and parameters, thus $I\mleft(\alpha\mright)$ is indeed a well-defined function. The other desired properties are all straightforward to verify.
\end{proof}

It is not hard to see that by the remark on page \pageref{loc:l-pl-rec-impl-remark}, the class functions $\Theta$ and $I$ will be definable within $L$ (or $L_\pl$) as well, through a formula which contains an existential quantifier that ranges over the plump ordinals.

On the other hand, we observe that \autoref{cor:pl-ord-code-pairs} additionally allows us to code a pairwise strongly incomparable function between plump ordinals by a single plump ordinal. More specifically:
\begin{proposition}
  \label{prop:pl-str-incom-func-as-pl-ord}
  There is a formula $\varphi\mleft(x, y; \alpha, \beta, \sigma\mright)$ such that, for any $\alpha, \beta \in \PlOrd$, if $f : s \rightarrow \beta$ is pairwise strongly incomparable on domain $s \subseteq \alpha$, then there exists some $\sigma \in \PlOrd$ such that
  \[\left(\gamma \in s \land f\mleft(\gamma\mright) = \delta\right) \leftrightarrow \varphi\mleft(\gamma, \delta; \alpha, \beta, \sigma\mright).\]
\end{proposition}

\begin{proof}
  We take
  \[\sigma = \bigcup_{\gamma \in s} \left(\alpha \cdot^\pl f\mleft(\gamma\mright) +^\pl \gamma\right)^\pls,\]
  and define $\varphi\mleft(x, y; \alpha, \beta, \sigma\mright)$ as the formula
  \[x \in \alpha \land y \in \beta \land \opE{\tau \in \sigma} \left(\tau = \alpha \cdot^\pl y +^\pl x \land \opA{\tau' \in \sigma} \left(\tau \subseteq \tau' \rightarrow \tau = \tau'\right)\right).\]

  For the forward direction, we see for any $\gamma \in s$ that indeed $\alpha \cdot^\pl f\mleft(\gamma\mright) +^\pl \gamma \in \sigma$. Additionally, for any $\tau \in \sigma$, there must exist $\gamma' \in s$ such that $\tau \subseteq \alpha \cdot^\pl f\mleft(\gamma'\mright) +^\pl \gamma'$. Now, if $\alpha \cdot^\pl f\mleft(\gamma\mright) +^\pl \gamma \subseteq \tau$, i.e.
  \[\alpha \cdot^\pl f\mleft(\gamma\mright) \subseteq \alpha \cdot^\pl f\mleft(\gamma\mright) +^\pl \gamma \subseteq \tau \subseteq \alpha \cdot^\pl f\mleft(\gamma'\mright) +^\pl \gamma',\]
  then by \autoref{lem:pl-mul-inj}, we must also have $f\mleft(\gamma\mright) \subseteq f\mleft(\gamma'\mright)$, which by pairwise strong incomparability implies $\gamma = \gamma'$. It then follows trivially that $\tau = \alpha \cdot^\pl f\mleft(\gamma\mright) +^\pl \gamma$, i.e.\ $\varphi\mleft(\gamma, f\mleft(\gamma\mright); \alpha, \beta, \sigma\mright)$ holds.

  For the backward direction, assume that $\varphi\mleft(\gamma, \delta; \alpha, \beta, \sigma\mright)$ holds. Then we must have $\alpha \cdot^\pl \delta +^\pl \gamma \in \sigma$, i.e.\ $\alpha \cdot^\pl \delta +^\pl \gamma \subseteq \alpha \cdot^\pl f\mleft(\gamma'\mright) +^\pl \gamma'$ for some $\gamma' \in s$, which also implies $\alpha \cdot^\pl \delta +^\pl \gamma = \alpha \cdot^\pl f\mleft(\gamma'\mright) +^\pl \gamma'$ by stipulation in $\varphi$. Since also $\gamma, \delta \in \PlOrd$ and $\gamma \in \alpha$, it follows from \autoref{cor:pl-ord-code-pairs} that $\delta = f\mleft(\gamma'\mright)$ and $\gamma = \gamma'$. Thus we have $\gamma \in s \land \delta = f\mleft(\gamma\mright)$ as desired.
\end{proof}

Since plump ordinals are always in $L$ by \autoref{lem:uniform-def-pl-ord-in-l}, and plump ordinal arithmetic is definable in $L$ by the remark on page \pageref{loc:l-pl-rec-impl-remark}, these coded functions must also be in $L$. We thus have the following:
\begin{theorem}[$\IKPP$, \cite{wang-sw26-plump}*{Theorem 4.7}]
  \label{thm:pl-str-incom-code-power}
  Suppose that $s \subseteq \PlOrd$ and $f : s \rightarrow \PlOrd$ is pairwise strongly incomparable, then $\Pow\mleft(s\mright) \subseteq L$. More specifically, there exists a plump ordinal $\alpha \in \PlOrd$ such that $\Pow\mleft(s\mright) \subseteq L_\alpha$.
\end{theorem}

\begin{proof}
  Let $\alpha = \bigcup_{\beta \in s} \beta^\pls$ so that $s \subseteq \alpha$ and let $x \subseteq s$ be any subset. We consider
  \[\sigma_x = \bigcup_{\beta \in x} f\mleft(\beta\mright)^\pls,\]
  then $\sigma_x$ is a plump ordinal, thus $\sigma_x \in L$ by \autoref{lem:uniform-def-pl-ord-in-l}. We can then define
  \[s_x = \left\{\beta \in \alpha : \beta \in s \land f\mleft(\beta\mright) \in \sigma_x\right\}.\]
  By \autoref{prop:pl-str-incom-func-as-pl-ord}, the condition $\beta \in s \land f\mleft(\beta\mright) \in \sigma_x$ can be equivalently written as $\opE{\gamma \in \sigma_x} \varphi\mleft(\beta, \gamma\mright)$ where $\varphi$ is a $\Delta_0$-formula except it contains plump ordinal additions and multiplications. Since these operations are definable inside $L$, we know that $s_x \in L$.

  Now obviously, for any $\beta \in x$, we have $f\mleft(\beta\mright) \in \sigma_x$, thus $\beta \in s_x$. On the other hand, for any $\beta \in s_x$ we have $\beta \in s$ and some $\gamma \in x$ such that $f\mleft(\beta\mright) \in f\mleft(\gamma\mright)^\pls$, i.e.\ $f\mleft(\beta\mright) \subseteq f\mleft(\gamma\mright)$, thus $\beta = \gamma \in x$ by incomparability. Therefore, we indeed have $x = s_x \in L$.

  Now, observe that every ordinal parameter we have used in the process is a plump ordinal, i.e.\ some $\delta \in \dd\mleft(L_\delta\mright)$ by \autoref{lem:uniform-def-pl-ord-in-l}. Thus, by taking plump successors, we can ensure that each $s_x$ lies in some $L_{\delta_x}$ where $\delta_x \in \PlOrd$. It follows that $\Pow\mleft(s\mright) \subseteq L_{\bigcup_{x \in \Pow\mleft(s\mright)} \delta_x}$, or simply $\Pow\mleft(s\mright) \subseteq L_\pl$ using the aforementioned class notation $L_\pl = \bigcup_{\alpha \in \PlOrd} L_\alpha$.
\end{proof}

\begin{corollary}[$\IKPP$]
  \label{cor:many-pl-incom-v-eq-l}
  Suppose that for any $\alpha \in \PlOrd$, there exists a pairwise strongly incomparable function $f : \alpha \rightarrow \PlOrd$, then we have $V = L$.
\end{corollary}

\begin{proof}
  We prove by set induction that every set $x$ lies in $L_\pl$ (hence in $L$), i.e.\ $x \in L_\beta$ for some plump ordinal $\beta$. By $\Sigma_1$-collection, it suffices to show that for any set $x \subseteq L_\alpha$ where $\alpha \in \PlOrd$, we also have $x \in L_\beta$ for some $\beta \in \PlOrd$. Consider the plump ordinal $\Theta\mleft(\alpha\mright)$ given in \autoref{thm:l-pl-ord-surj}, then \autoref{thm:pl-str-incom-code-power} implies that $\Pow\mleft(\Theta\mleft(\alpha\mright)\mright) \subseteq L_\beta$ for some plump ordinal $\beta$. Specifically, the following set
  \[\sigma = \left\{\gamma \in \Theta\mleft(\alpha\mright) : \opE{y \in x} \tuple{\gamma, y} \in I\mleft(\alpha\mright)\right\} \in L_\beta.\]

  However, as we remarked just after \autoref{thm:l-pl-ord-surj}, the class functions $\Theta$ and $I$ are definable in $L_\pl$ as well. It follows that
  \[s = \left\{I\mleft(\alpha\mright)\mleft(\gamma\mright) : \gamma \in \sigma\right\} \in L_\pl.\]
  For any $y \in x$, there exists $\gamma \in \Theta\mleft(\alpha\mright)$ such that $I\mleft(\alpha\mright)\mleft(\gamma\mright) = y$ by surjectivity and $\gamma \in \sigma$ by its definition, thus $x \subseteq s$. On the other hand, since $I\mleft(\alpha\mright)$ is a function, for any $\gamma \in \sigma$, $I\mleft(\alpha\mright)\mleft(\gamma\mright) \in x$ by definition of $\sigma$. Therefore, indeed $x = s \in L_\pl$.
\end{proof}

\section{Weak realisability with powerset}
\label{sec:czf-weak-real}

In the remainder of this chapter, we shall present a realisability model of $\CZFP$ where we indeed have as many incomparable objects as needed in \autoref{cor:many-pl-incom-v-eq-l}. We will work in a classical base theory $\KPP$ where it is more convenient to manipulate large sets. We consider a nested realisability implementation, where we first construct a variant of weak realisability in \cite{rathjen12-czf-powerset} to interpret $\CZFP$, and then build within this model the realisability interpretation described in \autoref{sec:czf-real-in-czf} to introduce the incomparable ordinals.

The weak realisability conditions shall use a class-sized pca $E_\wp\mleft(X\mright)$ of \emph{set recursive functions}\inds{set recursive function} extended by the powerset operation and an additional (partial) $\Sigma_1$-class function $X$, where the latter will be used later to supply some incomparable sets from the classical base theory to the intuitionistic models.

Since we will simultaneously interact with two notions of realisability using different pcas, we will distinguish their application terms by underlining the terms in $E_\wp\mleft(X\mright)$ (in this chapter only), i.e.\ we write $\underline{st}$ for the application of function $s$ on $t$ in $E_\wp\mleft(X\mright)$; we will write terms without underlines for the pca of the inner realisability model --- which will be Kleene's first algebra $\Kl$.

Now, $E_\wp\mleft(X\mright)$ will have the entire universe $V$ as its domain and the partial application function will be a class function given by $\Sigma$-inductive definition. We pick arbitrary distinct $\bk, \bs, \bp, \bp_0, \bp_1, \bsN, \bpN, \bdN, \bm{\gamma}, \bm{\rho}, \bm{\nu}, \bm{\pi}, \bi_1, \bi_2, \bi_3, \bar{\wp}, \mathbf{x} \in \omega$, and first informally specify the following application behaviours in $E_\wp\mleft(X\mright)$:
\begin{align*}
  \underline{\bk xy} {\downarrow}         & = x,                                                                                                       \\
  \underline{\bs xyz} {\downarrow}        & = \underline{xz\mleft(yz\mright)} \quad \text{when the right-hand side converges},                         \\
  \underline{\bp xy} {\downarrow}         & = \tuple{x, y},                                                                                            \\
  \underline{\bp_0 x} {\downarrow}        & = y \quad \text{when $x = \tuple{y, z}$},                                                                  \\
  \underline{\bp_1 x} {\downarrow}        & = z \quad \text{when $x = \tuple{y, z}$},                                                                  \\
  \underline{\bsN n} {\downarrow}         & = n + 1 \quad \text{when $n \in \omega$},                                                                  \\
  \underline{\bpN n} {\downarrow}         & = m \quad \text{when $n = m + 1$, $m \in \omega$},                                                         \\
  \underline{\bdN nmxy} {\downarrow}      & = \left\{\begin{aligned}
                                                        & x &  & \text{when $n = m \in \omega$},            \\
                                                        & y &  & \text{when $n, m \in \omega$, $n \neq m$},
                                                     \end{aligned}\right.                                              \\
  \underline{\bm{\gamma} xy} {\downarrow} & = x \cap \left(\bigcap y\right),                                                                           \\
  \underline{\bm{\rho} xy} {\downarrow}   & = \left\{\underline{xu} : u \in y\right\} \ \ \text{when $\underline{xu} {\downarrow}$ for all $u \in y$}, \\
  \underline{\bm{\nu} x} {\downarrow}     & = \bigcup x,                                                                                               \\
  \underline{\bm{\pi} xy} {\downarrow}    & = \left\{x, y\right\},                                                                                     \\
  \underline{\bi_1 xyz} {\downarrow}      & = \left\{\begin{aligned}
                                                        & x           &  & \text{when $y \in z$},     \\
                                                        & \varnothing &  & \text{when $y \not\in z$}, \\
                                                     \end{aligned} \right.                                                    \\
  \underline{\bi_2 xyz} {\downarrow}      & = \left\{u \in x : u \in y \rightarrow u \in z\right\},                                                    \\
  \underline{\bi_3 xyz} {\downarrow}      & = \left\{u \in x : u \in y \rightarrow z \in u\right\},                                                    \\
  \underline{\bar{\wp} x} {\downarrow}    & = \Pow\mleft(x\mright),                                                                                    \\
  \underline{\bx x} {\downarrow}          & = y \quad \text{when $y$ is some unique set such that $\tuple{x, y} \in X$},
\end{align*}
for any set $x, y, z$. Any application term not listed here is expected to diverge.

Formally speaking, the partial application function that satisfies these stipulations will be some inductive definition $I\mleft(\Phi\mright)$ where $\Phi$ is a $\Sigma_1$-class given as follows, using the same technique as \cite{rathjen12-czf-powerset}:

For any clause $\underline{e x_1 \cdots x_k} {\downarrow} = z$ appearing in the listing above, we will have
\begin{enumerate}[label=\uplabel{\arabic*}]
  \item if $k > 1$, then $\tuple{\tuple{\tuple{e, x}, \tuple{e, x}}, \varnothing} \in \Phi$ for any set $x$, i.e.\ we stipulate that $\underline{ex} {\downarrow} = \tuple{e, x}$;
  \item for any $1 < i < k$, we have
        \[\tuple{\tuple{\tuple{\tuple{e, x_1, \ldots, x_{i - 1}}, x_i}, \tuple{e, x_1, \ldots, x_i}}, \varnothing} \in \Phi\]
        for any set $x_1, \ldots, x_i$, i.e.\ we stipulate that
        \[\underline{\tuple{e, x_1, \ldots, x_{i - 1}} x_i} {\downarrow} = \tuple{e, x_1, \ldots, x_i};\]
  \item finally, we have $\tuple{\tuple{\tuple{e, x_1, \ldots, x_{k - 1}, x_k}, z}, y} \in \Phi$ for any set $x_1, \ldots, x_k$, or simply $\tuple{\tuple{\tuple{e, x_1}, z}, y}$ when $k = 1$, where the computation should converge and $y$ is the set of partial applications that are needed to produce output $z$.
\end{enumerate}
For example, for $\bs$, we need the following entries
\begin{align*}
  \tuple{\tuple{\tuple{\bs, x}, \tuple{\bs, x}}, \varnothing}                                                                               & \in \Phi, \\
  \tuple{\tuple{\tuple{\tuple{\bs, x}, y}, \tuple{\bs, x, y}}, \varnothing}                                                                 & \in \Phi, \\
  \tuple{\tuple{\tuple{\tuple{\bs, x, y}, z}, w}, \left\{\tuple{\tuple{x, z}, u}, \tuple{\tuple{y, z}, v}, \tuple{\tuple{u, v}, w}\right\}} & \in \Phi
\end{align*}
for any $x, y, z, u, v, w$. Since every distinguished operation $e$ is a natural number, every nonempty $e$ must contain $\varnothing$ as an element; on the other hand, ordered pairs and tuples implemented in the Kuratowski style are nonempty sets that do not contain $\varnothing$ as an element, thus every distinguished operation $e$ is distinct from every ordered pair, and it is easy to verify that $I\mleft(\Phi\mright)$ is indeed a partial function $V^2 \rightharpoonup V$ (cf.\ the proof that the inductive definition $\bU_A$ is a function in \autoref{lem:mltt-ind-def-func}).

Now, it is easy to see that $E_\wp\mleft(X\mright)$ is a pca over natural numbers by explicit stipulation. We can then define an interpreted realisability interpretation on it where the variables also range over the domain of all sets $V$. However, the pca cannot compute choice functions and thus cannot `pick an arbitrary element' from a nonempty set\footnote{Admittedly, our construction contains a custom $\Sigma_1$-class function, which could encode a choice function when such a function is definable in the base theory, e.g.\ when $V = L$ holds. However, since weak realisability itself is an interesting notion, we will present it here precisely as given in \cite{rathjen12-czf-powerset}.}, which will cause problems (for example, in \autoref{lem:ep-real-str-coll} later).

Therefore, we follow \cite{rathjen12-czf-powerset} and consider \emph{weak realisability}\indg{weak realisability} conditions, usually denoted as $\Vdash_\mathfrak{w}$, which use modified conditions for the disjunction and existential clauses that essentially allow one to specify a nonempty set of possible realisers without picking a specific one.

To begin, the atomic clauses are trivial:
\begin{align*}
  d \Vdash_\mathfrak{w} x \in y & \defeqv x \in y, \\
  d \Vdash_\mathfrak{w} x = y   & \defeqv x = y.
\end{align*}
The conjunction, negation, implication and universal clauses still follow \autoref{sec:real-con-int-logic}, whereas the disjunction and existential clauses are as follows:
\begin{align*}
  d \Vdash_\mathfrak{w} \varphi \lor \psi & \defeqv d \neq \varnothing \land \opA{e \in d} \begin{aligned}[t]
                                                                                              & {\left(\left(\underline{\bp_0 e} = 0 \land \underline{\bp_1 e} {\downarrow} \Vdash_\mathfrak{w} \varphi\right)\right.}     \\
                                                                                              & \left.{} \lor \left(\underline{\bp_0 e} = 1 \land \underline{\bp_1 e} {\downarrow} \Vdash_\mathfrak{w} \psi\right)\right),
                                                                                           \end{aligned}                                                                                                     \\
  d \Vdash_\mathfrak{w} \opE{x} \varphi   & \defeqv d \neq \varnothing \land \opA{e \in d} \left(\underline{\bp_0 e} {\downarrow} \land \underline{\bp_1 e} {\downarrow} \Vdash_\mathfrak{w} \varphi\mleft[\underline{\bp_0 e}/x\mright]\right).
\end{align*}
Finally, the conditions for bounded quantifiers are
\begin{align*}
  d \Vdash_\mathfrak{w} \opA{x \in b} \varphi & \defeqv \opA{c \in b} \underline{dc} {\downarrow} \Vdash_\mathfrak{w} \varphi\mleft[c/x\mright],                                                                                                           \\
  d \Vdash_\mathfrak{w} \opE{x \in b} \varphi & \defeqv d \neq \varnothing \land \opA{e \in d} \left(\underline{\bp_0 e} {\downarrow} \in b \land \underline{\bp_1 e} {\downarrow} \Vdash_\mathfrak{w} \varphi\mleft[\underline{\bp_0 e}/x\mright]\right).
\end{align*}

Here, the realisability model shares the same domain as the base theory, thus we will give it a new name $V_{E_\wp\mleft(X\mright)}$. We write $V_{E_\wp\mleft(X\mright)} \vDash \varphi$ to abbreviate for $\opE{d} d \Vdash_\mathfrak{w} \varphi$, where $\varphi$ is a first-order formula possibly containing any set as parameters.

We now verify that $V_{E_\wp\mleft(X\mright)}$ is a realisability model for the theory $\CZFP$. While some realisability conditions have changed, we observe that most logical axioms and rules can still be verified in the same way as in \autoref{prop:real-con-int-logic}, as long as we observe that singleton sets can be constructed trivially as an unordered pair
\[\mathfrak{sg} = \underline{\oplam{x} \bm{\pi} xx},\]
so that $\underline{\mathfrak{sg} x} {\downarrow} = \left\{x\right\}$. However, the axiom \ref{itm:int-logic-disj-elim} of disjunction elimination and the inference rule of $\exists$-elimination will require some additional treatment, using the following technical lemma:
\begin{lemma}[\cite{rathjen12-czf-powerset}*{Lemma 3.2}]
  For any formula $\varphi\mleft(\vec{x}\mright)$ in the language of set theory with all free variables indicated, there exists a term $\chi_\varphi$ in the pca $E_\wp\mleft(X\mright)$ such that $\KPP$ proves
  \[\opA{\vec{x}} \opA{c} \left(c \neq \varnothing \land \left(\opA{d \in c} d \Vdash_\mathfrak{w} \varphi\mleft(\vec{x}\mright)\right) \rightarrow \underline{\chi_\varphi \vec{x} c} {\downarrow} \Vdash_\mathfrak{w} \varphi\mleft(\vec{x}\mright)\right).\]
\end{lemma}

\begin{proof}
  We construct $\chi_\varphi$ by induction on the complexity of $\varphi$: for the atomic and negation cases, the actual realiser is irrelevant, so we can simply take $\chi_\varphi = \underline{\oplam{\vec{x} c} \varnothing}$.

  For the conjunction, implication and universal cases, we just defer the construction to some $\chi_\psi$ of a subformula $\psi$. Explicitly,
  \begin{align*}
    \chi_{\varphi \land \psi}                                                             & = \underline{\oplam{\vec{x} c} \bp \mleft(\chi_\varphi \vec{x} \mleft(\bm{\rho} \bp_0 c\mright)\mright) \mleft(\chi_\psi \vec{x} \mleft(\bm{\rho} \bp_1 c\mright)\mright)}, \\
    \chi_{\varphi \rightarrow \psi}                                                       & = \underline{\oplam{\vec{x} cd} \chi_\psi \vec{x} \mleft(\bm{\rho} \mleft(\oplam{f} fd\mright) c\mright)},                                                                  \\
    \chi_{\opA{y} \varphi\mleft(y\mright)} = \chi_{\opA{y \in t} \varphi\mleft(y\mright)} & = \underline{\oplam{\vec{x} cy} \chi_\varphi \vec{x} y \mleft(\bm{\rho} \mleft(\oplam{f} fy\mright) c\mright)},
  \end{align*}
  where in the bounded universal quantifier case, $t$ denotes any variable from $\vec{x}$.

  Finally, for the disjunction and existential cases, observe that a realiser is simply defined to be a `set of realisers' by stipulation, and thus it suffices to simply take $\chi_\varphi = \underline{\oplam{\vec{x} c} \bm{\nu} c}$.
\end{proof}

\begin{corollary}
  Any logical axioms or inference rules are realised in $V_{E_\wp\mleft(X\mright)}$.
\end{corollary}

\begin{proof}
  The only non-trivial axiom is disjunction elimination
  \[\left(\varphi \rightarrow \sigma\right) \rightarrow \left(\left(\psi \rightarrow \sigma\right) \rightarrow \left(\varphi \lor \psi \rightarrow \sigma\right)\right),\]
  realised by $\underline{\oplam{\vec{x} fgc} \chi_\sigma \vec{x} \mleft(\bm{\rho} \mleft(\oplam{d} \bdN \mleft(\bp_0 d\mright) 0 fg \mleft(\bp_1 d\mright)\mright) c\mright)}$ where we use $\chi_\sigma$ to deal with the set of realisers for $\varphi \lor \psi$.

  Likewise, for the $\exists$-elimination rule, assume that $f \Vdash_\mathfrak{w} \varphi \rightarrow \psi$ where the free variables are $x, \vec{y}$ and $x$ does not occur free in $\psi$, then $\opE{x} \varphi \rightarrow \psi$ is realised by $\underline{\oplam{\vec{y} c} \chi_\psi \vec{y} \mleft(\bm{\rho} \mleft(\oplam{d} f \mleft(\bp_0 d\mright) \vec{y} \mleft(\bp_1 d\mright)\mright) c\mright)}$.
\end{proof}

Now, equality is realised trivially here, so the equality axioms are already covered by \autoref{lem:eq-axiom-trivial}. For the axioms for bounded quantifiers:
\begin{lemma}
  The bounded quantifier axioms are realised in $V_{E_\wp\mleft(X\mright)}$:
  \begin{align*}
    \opA{x} \left(x \in a \rightarrow \varphi\mleft(x\mright)\right) & \leftrightarrow \opA{x \in a} \varphi\mleft(x\mright), \\
    \opE{x} \left(x \in a \land \varphi\mleft(x\mright)\right)       & \leftrightarrow \opE{x \in a} \varphi\mleft(x\mright).
  \end{align*}
\end{lemma}

\begin{proof}
  For any formula $\varphi\mleft(x\mright)$ with no other free variables (but possibly set parameters), we simply observe that
  \begin{align*}
    \left(d \Vdash_\mathfrak{w} \opA{x \in b} \varphi\mleft(x\mright)\right) & \leftrightarrow \left(\underline{\oplam{xe} dx} \Vdash_\mathfrak{w} \opA{x} \left(x \in b \rightarrow \varphi\mleft(x\mright)\right)\right),                          \\
    \left(d \Vdash_\mathfrak{w} \opE{x \in b} \varphi\mleft(x\mright)\right) & \leftrightarrow \left\{\tuple{x, \tuple{\varnothing, e}} : \tuple{x, e} \in d\right\} \Vdash_\mathfrak{w} \opE{x} \left(x \in b \land \varphi\mleft(x\mright)\right).
  \end{align*}

  For the former, $\underline{\oplam{dxe} dx}$ simply computes a realiser for the right-hand side from the left-hand side, whereas $\underline{\oplam{dx} dx \varnothing}$ computes the converse.

  For the latter, $\underline{\bm{\rho} \mleft(\oplam{e} \bp \mleft(\bp_0 e\mright) \mleft(\bp \varnothing \mleft(\bp_1 e\mright)\mright)\mright)}$ will then compute a realiser for the right-hand side from the left-hand side, whereas for the converse computation, $\underline{\bm{\rho} \mleft(\oplam{e} \bp \mleft(\bp_0 e\mright) \mleft(\bp_{11} e\mright)\mright)}$ will suffice.
\end{proof}

For the non-logical axioms, we shall begin by proving a variant of \cite{rathjen12-czf-powerset}*{Proposition 3.7}. This is a technical result that will be quite useful when reasoning with the axioms:
\begin{proposition}
  \label{prop:ep-real-delta0-correctness}
  For any $\Delta_0$-formula $\varphi\mleft(\vec{x}\mright)$ with all free variables indicated, possibly containing the primitive symbol $\Pow$, there exist terms $d_\varphi, k_\varphi$ in the pca $E_\wp\mleft(X\mright)$ such that $\KPP$ proves, for every tuple $\vec{x} \in V^k$ \uppar{of appropriate, finite length},
  \begin{enumerate}[label=\uplabel{\roman*}]
    \item $\underline{d_\varphi \vec{x}} {\downarrow} \in \left\{0, 1\right\}$ and $\underline{d_\varphi \vec{x}} = 1$ if and only if $\varphi\mleft(\vec{x}\mright)$ holds;
    \item if $\underline{d_\varphi \vec{x}} = 1$, then $\underline{k_\varphi \vec{x}} {\downarrow} \Vdash_\mathfrak{w} \varphi\mleft(\vec{x}\mright)$; if $\underline{d_\varphi \vec{x}} = 0$, then $\neg\!\opE{e} e \Vdash_\mathfrak{w} \varphi\mleft(\vec{x}\mright)$.
  \end{enumerate}

  It follows immediately that $\KPP$ proves that $\varphi\mleft(\vec{x}\mright)$ holds if and only if $\opE{e} e \Vdash_\mathfrak{w} \varphi\mleft(\vec{x}\mright)$.
\end{proposition}

\begin{proof}
  We construct $d_\varphi, k_\varphi$ by induction on the complexity of $\varphi$:
  \begin{enumerate}[label=\uplabel{\arabic*}]
    \item When $\varphi$ is the atomic formula $s \in t$, where $s$ and $t$ are either a variable or a finite number of applications of the symbol $\Pow$ to a variable, it suffices to take $d_\varphi = \underline{\oplam{\vec{x}} \bi_1 1 s' t'}$, where $\vec{x}$ enumerates the free variables of $\varphi$ and $s', t'$ are corresponding (formal) terms in $E_\wp\mleft(X\mright)$ by replacing occurrences of $\Pow$ in $s$ and $t$ with $\bar{\wp}$ instead; when $\varphi$ is the formula $s = t$, we observe that it suffices to simply take
          \[d_\varphi = d_{\opA{z \in s} z \in t \land \opA{z \in t} z \in s}\]
          by extensionality in the base theory, which can be handled later in the bounded quantifier case. In either case, we just need $k_\varphi = \underline{\oplam{\vec{x}} \varnothing}$.

    \item When $\varphi$ is of the form $\psi \land \chi$, it suffices to take\footnote{It is trivial to resolve the differences between how free variables are enumerated for $\psi$ and $\chi$. So we can assume without loss of generality that $d_\psi, k_\psi$, etc.\ all take the same inputs as $d_\varphi, k_\varphi$ that we are constructing here. Same applies to other cases below.}
          \[d_\varphi = \underline{\oplam{\vec{x}} \bm{\gamma} 1 \mleft(\bm{\pi} \mleft(d_\psi \vec{x}\mright) \mleft(d_\chi \vec{x}\mright)\mright)}, \]
          which enforces $\underline{d_\varphi \vec{x}} = \underline{d_\psi \vec{x}} \cap \underline{d_\chi \vec{x}}$, and simply $k_\varphi = \underline{\oplam{\vec{x}} \bp \mleft(k_\psi \vec{x}\mright) \mleft(k_\chi \vec{x}\mright)}$.

    \item When $\varphi$ is of the form $\psi \lor \chi$, it suffices to take
          \[d_\varphi = \underline{\oplam{\vec{x}} \bm{\nu} \mleft(\bm{\pi} \mleft(d_\psi \vec{x}\mright) \mleft(d_\chi \vec{x}\mright)\mright)},\]
          i.e.\ $\underline{d_\varphi \vec{x}} = \underline{d_\psi \vec{x}} \cup \underline{d_\chi \vec{x}}$, and
          \[k_\varphi = \underline{\oplam{\vec{x}} \bm{\nu} \mleft(\bm{\pi} \mleft(\bdN^* \mleft(d_\psi \vec{x}\mright) 1 \mleft(\mathfrak{sg} \mleft(\bp 0 \mleft(k_\psi \vec{x}\mright)\mright)\mright) \varnothing\mright) \mleft(\bdN^* \mleft(d_\chi \vec{x}\mright) 1 \mleft(\mathfrak{sg} \mleft(\bp 1 \mleft(k_\chi \vec{x}\mright)\mright)\mright) \varnothing\mright) \mright)},\]
          where $\mathfrak{sg} = \underline{\oplam{x} \bm{\pi} xx}$ is the singleton constructor, and $\underline{k_\varphi \vec{x}}$ evaluates to
          \[\left\{\tuple{0, \underline{k_\psi \vec{x}}} : \underline{d_\psi \vec{x}} = 1\right\} \cup \left\{\tuple{1, \underline{k_\chi \vec{x}}} : \underline{d_\chi \vec{x}} = 1\right\}.\]

    \item When $\varphi$ is of the form $\neg \psi$, we can simply handle it in the implication case below by taking $d_\varphi = d_{\psi \rightarrow \bot}$, where $\bot$ is a new propositional symbol with no possible realisers and $d_\bot = k_\bot = \varnothing$; clearly $k_\varphi = \underline{\oplam{\vec{x}} \varnothing}$ works.

    \item When $\varphi$ is of the form $\psi \rightarrow \chi$, observe that we just have to take $\underline{d_\varphi \vec{x}} = \left\{\varnothing : \varnothing \in \underline{d_\psi \vec{x}} \rightarrow \varnothing \in \underline{d_\chi \vec{x}}\right\}$, i.e.
          \[d_\varphi = \underline{\oplam{\vec{x}} \bi_2 1 \mleft(d_\psi \vec{x}\mright) \mleft(d_\chi \vec{x}\mright)}.\]
          By the inductive hypothesis, suppose that $\psi\mleft(\vec{x}\mright) \rightarrow \chi\mleft(\vec{x}\mright)$ is true and we have any realiser $e \Vdash_\mathfrak{w} \psi\mleft(\vec{x}\mright)$, then $\underline{k_\chi \vec{x}} \Vdash_\mathfrak{w} \chi\mleft(\vec{x}\mright)$; so for $k_\varphi$ all we need is
          \[k_\varphi = \underline{\oplam{\vec{x} e} k_\chi \vec{x}}.\]

    \item When $\varphi$ is of the form $\opA{y \in t} \psi\mleft(y\mright)$ where $t$ is either a variable or a finite number of applications of the symbol $\Pow$ to a variable, it suffices to take
          \begin{align*}
            d_\varphi & = \underline{\oplam{\vec{x}} \bm{\gamma} 1 \mleft(\bm{\rho} \mleft(\oplam{y} d_\psi \vec{x} y\mright) t'\mright)}, \\
            k_\varphi & = \underline{\oplam{\vec{x} y} k_\psi \vec{x} y},
          \end{align*}
          where $t'$ is the (formal) term in $E_\wp\mleft(X\mright)$ by replacing occurrences of $\Pow$ in $t$ with $\bar{\wp}$ instead.

    \item \label{itm:ep-real-delta0-correctness-exist} When $\varphi$ is of the form $\opE{y \in t} \psi\mleft(y\mright)$ where $t$ is either a variable or a finite number of applications of the symbol $\Pow$ to a variable, we take
          \begin{align*}
            d_\varphi & = \underline{\oplam{\vec{x}} \bm{\nu} \mleft(\bm{\rho} \mleft(\oplam{y} d_\psi \vec{x} y\mright) t'\mright)},                                                                                                                       \\
            k_\varphi & = \underline{\oplam{\vec{x}} \bm{\nu} \mleft(\bm{\rho} \mleft(\oplam{y} \bdN^* \mleft(d_\psi \vec{x} y\mright) 1 \mleft(\mathfrak{sg} \mleft(\bp y \mleft(k_\psi \vec{x} y\mright)\mright)\mright) \varnothing\mright) t'\mright)},
          \end{align*}
          where $\underline{k_\varphi \vec{x}}$ evaluates to
          \[\bigcup_{y \in t} \left\{\tuple{y, \underline{k_\psi \vec{x} y}} : \underline{d_\psi \vec{x} y} = 1\right\}. \qedhere\]
  \end{enumerate}
\end{proof}

It follows that any $\Pi_1$-formula has a uniform realiser for all variable assignments on which it is true in the base theory. Therefore, for example, we immediately know that the $\Pi_1$-formula stating the axiom of extensionality is realised simply by virtue of it being true in the base theory. Likewise, we know that the unordered pairs and unions are computable by $\bm{\pi}, \bm{\nu}$ in the pca, and it follows that we can uniformly realise the $\Delta_0$-statements `$\underline{\bm{\pi} xy}$ is the unordered pair of $x$ and $y$' or `$\underline{\bm{\nu} x}$ is the union of $x$' in $V_{E_\wp\mleft(X\mright)}$. The axiom of infinity (read as the $\Delta_0$-claim that $\omega$ is an inductive set), the binary intersection axiom and the powerset axiom are likewise trivial.

The case for the set induction scheme is completely analogous to \autoref{lem:czf-real-set-ind}, so all that remains is the strong collection axiom:
\begin{lemma}
  \label{lem:ep-real-str-coll}
  The strong collection scheme is realised in $V_{E_\wp\mleft(X\mright)}$:
  \[\opA{x \in a} \opE{y} \theta\mleft(x, y\mright) \rightarrow \opE{b} \left(\opA{x \in a} \opE{y \in b} \theta\mleft(x, y\mright) \land \opA{y \in b} \opE{x \in a} \theta\mleft(x, y\mright)\right).\]
\end{lemma}

\begin{proof}
  Let $\vec{z}$ enumerate the free variables in $\theta\mleft(x, y\mright)$ other than $x$, $y$ and $a$. For any set $a, \vec{c}$ and $f \Vdash_\mathfrak{w} \opA{x \in a} \opE{y} \theta\mleft(x, y\mright)\mleft[\vec{c}/\vec{z}\mright]$, we can define
  \[g = \underline{\oplam{af} \bm{\nu} \mleft(\bm{\rho} \mleft(\oplam{x} \bm{\rho} \bp_0 \mleft(fx\mright)\mright) a\mright)},\]
  then for any $x \in a$, $\underline{fx} \Vdash_\mathfrak{w} \opE{y \in b} \theta\mleft(x, y\mright)\mleft[\underline{gaf}/b, \vec{c}/\vec{z}\mright]$. Conversely, for any $y \in \underline{gaf}$, we compute
  \[h = \underline{\oplam{afy} \bm{\nu} \mleft(\bm{\rho} \mleft(\oplam{x} \bm{\nu} \mleft(\bm{\rho} \mleft(\oplam{d} \bdN \mleft(d_{x = y} y \mleft(\bp_0 d\mright)\mright) 1 \mleft(\mathfrak{sg} \mleft(\bp x \mleft(\bp_1 d\mright)\mright)\mright) \varnothing\mright) \mleft(fx\mright)\mright)\mright) a\mright)},\]
  where we use $d_{x = y}$ from \autoref{prop:ep-real-delta0-correctness} to select those entries in $\underline{fx}$ where the first component is $y$, i.e.\ we compute
  \[\underline{hafy} = \bigcup_{x \in a} \left\{\tuple{x, \underline{\bp_1 d}} : d \in \underline{fx}, \underline{\bp_0 d} = y\right\},\]
  so that for any $y \in \underline{gaf}$, $\underline{hafy} \Vdash_\mathfrak{w} \opE{x \in a} \theta\mleft(x, y\mright)\mleft[\vec{c}/\vec{z}\mright]$. It follows that
  \[\underline{\oplam{a \vec{b} f} \mathfrak{sg} \mleft(\bp \mleft(gaf\mright) \mleft(\bp f \mleft(haf\mright)\mright)\mright)}\]
  realises the axiom.
\end{proof}

By \autoref{prop:ikp-exp-sub-coll-powerset}, subset collection is implied by powerset and hence does not need to be realised separately. Thus, we have so far verified that:
\begin{theorem}
  \label{thm:czfp-weak-real}
  Let $\varphi$ be any theorem of $\CZFP$, then for any $\Sigma_1$-class function $X$, $\KPP$ proves $V_{E_\wp\mleft(X\mright)} \vDash \varphi$, i.e.\ there exists $d \in E_\wp\mleft(X\mright)$ such that $d \Vdash_\mathfrak{w} \varphi$.
\end{theorem}

Now, we will remark that one may consider a weaker pca $E\mleft(X\mright)$ already inside $\KP$ (or even $\IKP$), which omits the powerset operation $\bar{\wp}$. Let $\CZF^-$\inds{constructive Zermelo--Fraenkel set theory, $\CZF$!0@$\CZF^-$} denote the subtheory $\CZF - \text{subset collection}$ (which is sometimes also denoted $\CZF^0$, e.g.\ in \cite{rathjen05-replacement-collection}), then $\KP$ or $\IKP$ would suffice to verify that all axioms of $\CZF^-$ are realised (cf.\ Rathjen \cite{rathjen12-weak-strong-existence-property}). However, we cannot prove that subset collection is realised there, because even for the axiom of exponentiation, $\Exp$, the statement claiming `$f$ is the set of all functions between $x$ and $y$' is $\Pi_1$. Therefore, being able to compute a realiser for the existence of such $f$ from arbitrary $x, y$ would imply, by \autoref{prop:ep-real-delta0-correctness}, that $f$ is indeed the set of all functions between $x$ and $y$ in the base theory, which $\KP$ does not suffice to prove.

Additionally, observe that \autoref{prop:ep-real-delta0-correctness} implies the following:
\begin{lemma}
  $V_{E_\wp\mleft(X\mright)}$ realises the axiom scheme of restricted excluded middle, $\REM$\inds{restricted excluded middle, $\REM$},
  \[\varphi \lor \neg \varphi\]
  for any $\Delta_0$-formula $\varphi$.
\end{lemma}

\begin{proof}
  Let $\vec{x}$ enumerate the free variables in $\varphi$, then
  \[\oplam{\vec{x}} \mathfrak{sg} \mleft(\bp \mleft(1 - d_\varphi \vec{x}\mright) \mleft(\bdN \mleft(d_\varphi \vec{x}\mright) 1 k_\varphi k_{\neg \varphi} \vec{x}\mright)\mright)\]
  simply realises the axiom.
\end{proof}

Since $\REM$ implies the linearity of ordinals, etc., $V_{E_\wp\mleft(X\mright)}$ is not a `very intuitionistic model', and it will not contain any incomparable ordinals. Thus, in the next section, we will need to consider another realisability interpretation nested in this model of $\CZFP + \REM$ to continue our construction.

We will additionally remark here that it is not hard to make $V_{E_\wp\mleft(X\mright)}$ itself realise $V = L$. Simply observe that the statement $V = L$ is a $\Pi_2$-formula, and we can write it in some form
\[\opA{x} \opE{y} \psi\mleft(x, y\mright)\]
where $\psi$ is $\Delta_0$ (such that this formula and the original formulation of $V = L$ are provably equivalent even in $\IKP$) and it is provable in $\IKP$ that for every $x$, the corresponding $y$ satisfying $\psi\mleft(x, y\mright)$, if it exists, is unique. Now assume that the base theory satisfies $\KPP + V = L$, then it follows that $x \mapsto y$ given by this formula is a total $\Delta_0$-class function, so we can build $V_{E_\wp\mleft(X\mright)}$ where $X$ is precisely this function. Now, by \autoref{prop:ep-real-delta0-correctness}, we then have
\[\underline{k_\psi x \mleft(\bx x\mright)} \Vdash_\mathfrak{w} \psi\mleft(x, \underline{\bx x}\mright),\]
and we have that $\underline{\oplam{x} \mathfrak{sg} \mleft(\bp \mleft(\bx x\mright) \mleft(k_\psi x \mleft(\bx x\mright)\mright)\mright)}$ realises $V = L$ for $V_{E_\wp\mleft(X\mright)}$.

However, this is an interpretation of $\CZFP + \REM + V = L$ inside $\KPP + V = L$, where both theories are much stronger than $\KPP$ or $\CZFP$, as we have discussed in \autoref{sec:overview-v-eq-l}. We will instead present a construction in a weaker base theory in the next section.

\section{An equiconsistency proof for \texorpdfstring{$\CZFP + V = L$}{CZF(P) + V = L}}
\label{sec:czfp-v-eq-l}

We now consider the realisability model $V\mleft(\mathcal{A}\mright)$ defined in \autoref{sec:czf-real-in-czf}. We first additionally see that:
\begin{lemma}
  \label{lem:czfp-real-powerset}
  $\CZFP$ proves that the powerset axiom is realised in $V\mleft(\mathcal{A}\mright)$:
  \[\opA{a} \opE{y} \opA{x} \left(x \in y \leftrightarrow x \subseteq a\right).\]
\end{lemma}

\begin{proof}
  Given any $a \in V\mleft(\mathcal{A}\mright)$, we consider
  \[\tilde{a} = \left\{\tuple{e, z} : \tuple{d, z} \in a, e \in \mathcal{A}\right\}\]
  and let $y = \left\{\tuple{f, u} : u \in \Pow\mleft(\tilde{a}\mright), f \in \mathcal{A}, f \Vdash u \subseteq a\right\} \in V\mleft(\mathcal{A}\mright)$. Note that for fixed $u, a$, the formula $f \Vdash u \subseteq a$ is essentially
  \[\opA{{\tuple{d, x} \in u}} fd {\downarrow} \in F_\in\mleft(x, a\mright)\]
  using $F_\in$ from \autoref{sec:czf-real-in-czf}, and thus a $\Delta_0$-formula. Therefore, $y$ is indeed a set by $\Delta_0$-separation and collection.

  It now suffices to come up with realisers $t_1, t_2 \in \mathcal{A}$ such that
  \[t_1 \Vdash \opA{x \in y} x \subseteq a, \qquad t_2 \Vdash \opA{x} \left(x \subseteq a \rightarrow x \in y\right).\]

  The first direction is trivial, as $t_1 = \oplam{f} f$ simply realises $\opA{x \in y} x \subseteq a$ by the definition of $y$.

  To compute $t_2$, assume some $f \Vdash x \subseteq a$. By strong collection, let $z$ be a set such that
  \begin{align*}
     & \opA{{\tuple{d, u} \in x}} \opE{{\tuple{d, v} \in z}} \left(\tuple{\bp_0 \mleft(fd\mright), v} \in a \land \bp_1 \mleft(fd\mright) \Vdash u = v\right)               \\
     & \qquad \land \opA{{\tuple{d, v} \in z}} \opE{{\tuple{d, u} \in x}} \left(\tuple{\bp_0 \mleft(fd\mright), v} \in a \land \bp_1 \mleft(fd\mright) \Vdash u = v\right).
  \end{align*}
  Clearly $g_1 = \bp \mleft(\oplam{x} \bp x \mleft(\bp_1 \mleft(fx\mright)\mright)\mright) \mleft(\oplam{x} \bp x \mleft(\iot \mleft(\bp_1 \mleft(fx\mright)\mright)\mright) \mright) \Vdash x = z$. We also have $g_2 = \oplam{x} \bp \mleft(\bp_0 \mleft(fx\mright)\mright) \ior \Vdash z \subseteq a$, thus $\tuple{g_2, z} \in y$ and it follows immediately that $\bp g_2 g_1 \Vdash x \in y$.
\end{proof}

Specifically, in this section we will focus only on the pca $\mathcal{A} = \Kl$, Kleene's first algebra. Observe that its partial application function is definable by an arithmetic formula, i.e.\ a $\Delta_0$-formula in the language of set theory. Thus by \autoref{prop:ep-real-delta0-correctness}, we know that both the classical base universe $V$ and the weak realisability interpretation $V_{E_\wp\mleft(\Kl\mright)}$ will agree on which set defines this pca. Likewise, we can define the collections of names
\[V\mleft(\Kl\mright)_\alpha = \bigcup_{\beta \in \alpha} \Pow\mleft(\omega \times V\mleft(\Kl\mright)_\beta\mright),\]
which is now a hierarchy of sets and can be constructed using $\Sigma$-recursion instead of inductive definitions, since we have the powerset axiom. We take $V\mleft(\Kl\mright) = \bigcup_{\alpha \in \Ord} V\mleft(\Kl\mright)_\alpha$ and it is trivial to check (in any extension of $\IKPP$) that this coincides with the construction of $V\mleft(\mathcal{A}\mright) = I\mleft(\Phi\mright)$ in \autoref{sec:czf-real-in-czf}, by verifying that the class $V\mleft(\Kl\mright)$ satisfies both requirements in \autoref{thm:set-ind-def} using set induction on $\Ord$. We will additionally see that:
\begin{lemma}
  \label{lem:ep-real-v-kl-agree}
  There is a term $t \in E_\wp\mleft(X\mright)$ such that $x \in V\mleft(\Kl\mright)$ if and only if $\underline{tx} {\downarrow} \Vdash_\mathfrak{w} x \in V\mleft(\Kl\mright)$.
\end{lemma}

\begin{proof}
  Here $V\mleft(\Kl\mright)$ is a class, so the claim is not trivial. However, since we phrased the construction of $V\mleft(\Kl\mright)$ through $\Sigma$-recursion, we see that the statement `$z = V\mleft(\Kl\mright)_\alpha$' can be explicitly expressed as `there exists $f : \alpha^+ \rightarrow V$ such that each $f\mleft(\beta\mright) = \bigcup_{\gamma \in \beta} \Pow\mleft(\omega \times f\mleft(\gamma\mright)\mright)$ and $f\mleft(\alpha\mright) = z$', where the part inside the existential quantifier is a $\Delta_0$-formula containing the powerset symbol $\Pow$.

  Now, we can compute the function $\alpha \mapsto V\mleft(\Kl\mright)_\alpha$ in $E_\wp\mleft(X\mright)$ by recursion as
  \[\underline{\br \mleft(\oplam{\alpha f} \bm{\nu} \mleft(\bm{\rho} \mleft(\oplam{\beta} \bar{\wp} \mleft(\bm{\nu} \mleft(\bm{\rho} \mleft(\oplam{n} \bm{\rho} \mleft(\bp n\mright) \mleft(f \beta\mright)\mright) \omega\mright)\mright)\mright) \alpha\mright)\mright)},\]
  so it is also easy to compute the (set-theoretic) witness function $f$ above. By \autoref{prop:ep-real-delta0-correctness}, $V$ and $V_{E_\wp\mleft(X\mright)}$ will agree on which $f$ is the witness to $z = V\mleft(\Kl\mright)_\alpha$, and thus they will agree on which set is $V\mleft(\Kl\mright)_\alpha$ as well for every ordinal $\alpha$.

  Therefore, it suffices to show that, given any $x \in V\mleft(\Kl\mright)$, $E_\wp\mleft(X\mright)$ can compute some $\alpha$ such that $x \in V\mleft(\Kl\mright)_\alpha$. However, we just need to compute its rank function,
  \[\mathrm{rk}\mleft(x\mright) = \sup \left\{\mathrm{rk}\mleft(y\mright)^+ : \tuple{n, y} \in x\right\},\]
  which is simply
  \[\underline{\br \mleft(\oplam{xf} \bm{\nu} \mleft(\bm{\rho} \mleft(\oplam{y} \bm{\nu} \mleft(\bm{\pi} \mleft(f \mleft(\bp_1 y\mright)\mright) \mleft(\mathfrak{sg} \mleft(f \mleft(\bp_1 y\mright)\mright)\mright)\mright)\mright) x\mright)\mright)}. \qedhere\]
\end{proof}

From now onwards, we will read $V\mleft(\Kl\mright) \vDash \varphi$ as the claim in the base theory that
\[\opE{e} e \Vdash_\mathfrak{w} \opE{n \in \Kl} n \Vdash \varphi.\]
The formula $\varphi$ can contain parameters from $V\mleft(\Kl\mright)$, since by \autoref{lem:ep-real-v-kl-agree} the base theory knows correctly which sets are names inside $V_{E_\wp\mleft(X\mright)}$.

\subsection*{A strongly incomparable $\omega$-sequence}

The first task towards $V = L$ now is to recover the pairwise strongly incomparable function on $\omega$ that we constructed in \autoref{sec:czf-v-eq-l}, which will be implemented here as a sequence of names $S_0 : \omega \rightarrow V\mleft(\Kl\mright)$ inside the base theory $\KPP$. As enabled by the theory of plump ordinals, which we will use later, we do not care about a uniform presentation $\alpha_0 \in V\mleft(\Kl\mright)$ and hence do not need the specific pairing function. However, we will still work with finite $\omega$-tuples and the operations $\bl_s, \bt_s \in \Kl$. For a function $f : \omega \rightarrow \omega$, we shall write $f \upharpoonright n \in \omega$ to denote (the code for) the initial segment $\tuple{f\mleft(0\mright), \ldots, f\mleft(n - 1\mright)}$.

\begin{proposition}
  \label{prop:kpp-non-com-omega-seq}
  There exists a family of functions $\tuple{f_n : \omega \rightarrow \omega}_{n \in \omega}$ such that for any $i, j, e \in \omega$, if $i \neq j$ then there exists some $m \in \omega$ satisfying
  \[e \mleft(f_i \upharpoonright m\mright) {\downarrow} \rightarrow e \mleft(f_i \upharpoonright m\mright) \neq f_j \upharpoonright m.\]
\end{proposition}

\begin{proof}
  Similar to \autoref{prop:mltt-real-non-comp-path}, we construct the family $f_n$ as an increasing sequence of finite objects $t^{(0)}, t^{(1)}, t^{(2)}, \ldots$ coded as natural numbers. Here we will pick each $t^{(i)}$ as a finite tuple of finite tuples of natural numbers, i.e.\ $t^{(i)} = \tuple{t^{(i)}_0, \ldots, t^{(i)}_{k - 1}} \in \left(\omega^{< \omega}\right)^{< \omega}$, and require that for $i < j$, we have $\bl_s t^{(i)} \leq \bl_s t^{(j)}$ and $t^{(i)}_k \sqsubseteq t^{(j)}_k$ for each $0 \leq k < \bl_s t^{(i)}$.

  Specifically, we will let $t^{(0)}$ be the empty tuple, and for each $i$ such that $\bp_{00} i \neq \bp_{01} i$, we will first extend $t^{(i)}_{\bp_{00} i}$ to an extension $s$ strictly longer than $\bl_s t^{(i)}_{\bp_{01} i}$ (adding empty tuples to $t^{(i)}$ if it does not yet contain a $\left(\bp_{00} i\right)^\text{th}$ or $\left(\bp_{01} i\right)^\text{th}$ component), and then extend $t^{(i)}_{\bp_{01} i}$ to an $s'$ such that $\bl_s s' = \bl_s s$ and $\bp_1 i s$ does not compute $s'$ correctly. We should build $t^{(i + 1)}$ by replacing the $\left(\bp_{00} i\right)^\text{th}$ and $\left(\bp_{01} i\right)^\text{th}$ components in $t^{(i)}$ with $s$ and $s'$ respectively.

  In order to ensure that the limit is a family of functions on domain $\omega$, we additionally require that for each $i$, $\bl_s t^{(i)} \geq i$ and also $\bl_s t^{(i)}_k \geq i$ for each $0 \leq k < i$, by filling in extra $0$'s whenever necessary. Since each $t^{(i)}$ is a finite object, we can again choose the one satisfying the requirements above with the least code at each step. By $\Delta_0$-separation, we can then define each function $f_n$ such that $f_n\mleft(k\mright)$ is the $k^\text{th}$ component of $t^{(m)}_n$ where $m = \max\mleft\{n, k\mright\} + 1$.
\end{proof}

Now, in the base theory $\KPP$, we shall say that a function $f : x \rightarrow V\mleft(\Kl\mright)$ is a \emph{sequence of pairwise strongly incomparable names} if for any $u, v \in x$,
\[\left(V\mleft(\Kl\mright) \vDash f\mleft(u\mright) \subseteq f\mleft(v\mright)\right) \rightarrow u = v.\]
We see it follows from the construction above that:
\begin{corollary}
  \label{cor:kpp-incom-omega-seq}
  There is a sequence of pairwise strongly incomparable names $S_0 : \omega \rightarrow V\mleft(\Kl\mright)$ such that for any $n \in \omega$, $V\mleft(\Kl\mright) \vDash S_0\mleft(n\mright) \in \Pow\mleft(\omega\mright) \cap \Ord$.
\end{corollary}

\begin{proof}
  Take the family of functions $\tuple{f_n}_{n \in \omega}$ as in \autoref{prop:kpp-non-com-omega-seq}. For any $n \in \omega$, we let
  \[S_0\mleft(n\mright) = \left\{\tuple{f_n \upharpoonright m, \overline{m}} : m \in \omega\right\} \in V\mleft(\Kl\mright)\]
  where $\overline{m} = \left\{\tuple{k, \overline{k}} : k \in m\right\} \in V\mleft(\Kl\mright)$ are the same recursively defined standard names for natural numbers introduced in \autoref{sec:intro-int-l}, and likewise
  \[\overline{\omega} = \left\{\tuple{m, \overline{m}} : m \in \omega\right\}\]
  is a standard name for $\omega$ in $V\mleft(\Kl\mright)$.

  Now, Turing computability, the properties we asserted for $f_n$ in \autoref{prop:kpp-non-com-omega-seq} and the definition of $S_0$ are all expressible as $\Delta_0$-formulae, so by \autoref{prop:ep-real-delta0-correctness}, $V_{E_\wp\mleft(X\mright)}$ knows that they satisfy the same properties. We can thus work with these sets inside $V_{E_\wp\mleft(X\mright)}$, a model of $\CZFP + \REM$.

  Clearly we see $\oplam{t} \bp \mleft(\bl_s t\mright) \ior \Vdash S_0\mleft(n\mright) \subseteq \overline{\omega}$, and likewise, $\oplam{tm} \bp \mleft(\bt_s tm\mright) \ior \Vdash \text{$S_0\mleft(n\mright)$ is transitive}$. Reasoning inside $V\mleft(\Kl\mright)$, it follows immediately that also
  \[V\mleft(\Kl\mright) \vDash S_0\mleft(n\mright) \in \Ord.\]

  For strong incomparability, since the base theory $\KPP$ is classical, it suffices to assume $i \neq j \in \omega$ and show that $V_{E_\wp\mleft(X\mright)}$ realises $\neg\!\opE{e \in \Kl} e \Vdash S_0\mleft(i\mright) \subseteq S_0\mleft(j\mright)$. Working inside $V_{E_\wp\mleft(X\mright)}$ which knows $i \neq j$, suppose that some $e \Vdash S_0\mleft(i\mright) \subseteq S_0\mleft(j\mright)$. Now, by induction on $\omega$, we know that for $m, n \in \omega$, if any $d \Vdash \overline{m} = \overline{n}$, then $m = n$ must be true. It follows that for any $m \in \omega$, we must have $\bp_0 \mleft(e \mleft(f_i \upharpoonright m\mright)\mright) {\downarrow} = f_j \upharpoonright m$. However, this contradicts the requirement in \autoref{prop:kpp-non-com-omega-seq} that any function in $\Kl$, specifically $\oplam{t} \bp_0 \mleft(et\mright)$ here, cannot compute $f_j \upharpoonright m$ correctly from $f_i \upharpoonright m$ for all $m \in \omega$.
\end{proof}

Afterwards, we will describe an iterated process to build the specific partial $\Sigma_1$-class function $X$ we need in the base theory $\KPP$, which starts from $S_0$ and combines plump mechanisms inside $V\mleft(\Kl\mright)$ and techniques from classical set theory in $\KPP$. Thus we shall first say a bit about how one can conveniently extract complicated constructions inside $V\mleft(\Kl\mright)$ to play with in the base theory.

Since our realisability structure is nested inside the model of weak realisability $V_{E_\wp\mleft(X\mright)}$, we will need to add some choice principles to the base theory to extract a single, canonical name from the existential realisers. We will consider the extension of $\KPP$ by the \emph{axiom of global choice}, $\ACg$\inds{axiom of global choice} in \cite{rathjen20-power-kp-choice}, which adds a new binary relation symbol $R$ and the new axioms
\begin{align*}
  \opA{x} \opA{y} \opA{z} & \left(R\mleft(x, y\mright) \land R\mleft(x, z\mright) \rightarrow y = z\right), \\
  \opA{x}                 & \left(x \neq \varnothing \rightarrow \opE{y \in x} R\mleft(x, y\mright)\right),
\end{align*}
specifying that $R$ is a choice class function. The symbol $R$ is also allowed to appear in $\Delta_0$-formulae in axiom schemes such as $\Delta_0$-separation or $\Delta_0$-collection. It follows that $\left\{\tuple{x, y} : R\mleft(x, y\mright)\right\}$ is a $\Delta_0$-class function (containing the primitive symbol $R$) that computes an element in $x$ for every nonempty set $x$.

We observe that, whenever $\CZFP$ proves a formula $\varphi$ of any complexity, then $\CZFP$ also proves that this formula is realised by some $f \in \Kl$, and it follows from \autoref{thm:czfp-weak-real} that there is a realiser $h \in E_\wp\mleft(X\mright)$ for
\[\opE{f \in \Kl} f \Vdash \varphi\]
or any provably equivalent forms. Since $E_\wp\mleft(X\mright)$ is inductively defined, $h$ and derived application terms can be expressed as partial $\Sigma_1$-class functions, possibly involving the primitive powerset symbol $\Pow$.

Additionally, since the realisers for axioms of $\CZFP$ never use the specific operation $\bx$ in $E_\wp\mleft(X\mright)$, it is easy to verify that the values of $h$ remain the same regardless of which class function we use as $X$ --- most notably, we can define the same $h$ in $E_\wp\mleft(\varnothing\mright)$, and it will be unproblematic to, conversely, construct the class $X$ using $\Sigma_1$-class functions like $h$.

For example, in \autoref{sec:plump} we defined the plump operator $\left(-\right)^\pl : \Ord \rightarrow \PlOrd$ in $\IKPP$ and observe that, using \autoref{lem:pl-successor-still-pl} and \autoref{lem:pl-union-still-pl}, it is very easy to generalise it to any set $x$ such that
\[x^\pl = \bigcup_{y \in x} \left(y^\pl\right)^\pls\]
and verify that still $x^\pl \in \PlOrd$. In other words, we have shown
\[\CZFP \vdash \opA{x} \opE{y} \left(y = x^\pl \land y \in \PlOrd\right),\]
and it follows that, by expanding the definition of realisability conditions for $V_{E_\wp\mleft(X\mright)}$, there is some $h \in E_\wp\mleft(\varnothing\mright)$ that realises
\[\opE{d \in \omega} \opA{x} \left(x \in V\mleft(\Kl\mright) \rightarrow \opE{y} \left(y \in V\mleft(\Kl\mright) \land d \Vdash y = x^\pl \land y \in \PlOrd\right)\right).\]
Let $h'$ be any element in $h$. For any $x \in V\mleft(\Kl\mright)$ in the base theory $\KPP$, we have the term $t$ in \autoref{lem:ep-real-v-kl-agree} which computes a realiser for $x \in V\mleft(\Kl\mright)$. Using global choice, we can define $e_x$ as some unique element (in the realiser for the inner existential subformula) such that
\[R\mleft(\underline{\mleft(\bp_1 h'\mright) x \mleft(tx\mright)}, e_x\mright)\]
holds, then we will have $\underline{\bp_0 e_x} \in V\mleft(\Kl\mright)$ and $V\mleft(\Kl\mright) \vDash \underline{\bp_0 e_x} \in \PlOrd$. We will abuse notation and denote this $\Sigma_1$-class function (containing the primitive symbol $R$) on $V\mleft(\Kl\mright)$ given by $x \mapsto \underline{\bp_0 e_x}$ also as $\left(-\right)^\pl$. Notice that $e_x$ additionally computes realisers for $\underline{\bp_0 e_x} = x^\pl$ being plump, or more precisely:
\begin{lemma}[$\KPP + \ACg$]
  \label{lem:pl-in-pl-ord-unif-real}
  There is some $d_0 \in \omega$ such that for any function $f : x \rightarrow V\mleft(\Kl\mright)$, let $g : x \rightarrow V\mleft(\Kl\mright)$ be given by $g\mleft(z\mright) = f\mleft(z\mright)^\pl$, then there is $\mathfrak{p} \in E_\wp\mleft(X\mright)$ such that
  \[\mathfrak{p} \Vdash_\mathfrak{w} \opA{z \in x} d_0 \Vdash g\mleft(z\mright) \in \PlOrd.\]
\end{lemma}

\begin{proof}
  Consider the realiser $h$ for
  \[\opE{d \in \omega} \opA{x} \left(x \in V\mleft(\Kl\mright) \rightarrow \opE{y} \left(y \in V\mleft(\Kl\mright) \land d \Vdash y = x^\pl \land y \in \PlOrd\right)\right).\]
  and fixed $h' \in h$ as on the previous page. Observe that the desired $d_0 \in \omega$ is then simply $\bp_1 \mleft(\underline{\bp_0 h'}\mright)$ --- note that the inner unpairing function is in $E_\wp\mleft(X\mright)$, whereas the outer one is in $\Kl$.

  On the other hand, note that $e : x \rightarrow E_\wp\mleft(\varnothing\mright)$ given by $z \mapsto \underline{\bp_1 e_z}$, where $e_z$ satisfies $R\mleft(\underline{\mleft(\bp_1 h'\mright) z \mleft(tz\mright)}, e_z\mright)$, is a set-theoretic function by strong $\Sigma_1$-collection. Here $\underline{\bp_1 e_z}$ realises the statement
  \[\underline{\bp_0 h'} \Vdash g\mleft(z\mright) = f\mleft(z\mright)^\pl \land g\mleft(z\mright) \in \PlOrd\]
  in $V_{E_\wp\mleft(X\mright)}$. Because $E_\wp\mleft(X\mright)$ can evaluate set-theoretic functions\footnote{\label{ftn:ep-rec-set-func}Specifically, we have $e_z = e\mleft(z\mright) = \underline{\bp_1 \mleft(\bm{\nu} \mleft(\bm{\gamma} \mleft(\bm{\rho} \mleft(\oplam{t} \bp z \mleft(\bp_1 t\mright)\mright) e\mright) \mleft(\bm{\pi} ee\mright)\mright)\mright)}$.}, $\mathfrak{p}$ just need to take $z \in x$ and extract the needed information from the set parameter $e$.
\end{proof}

\begin{corollary}
  \label{cor:kpp-incom-pl-omega-seq}
  There is a sequence of pairwise strongly incomparable names $T_0 : \omega \rightarrow V\mleft(\Kl\mright)$ such that $V_{E_\wp\mleft(X\mright)} \vDash \opA{n \in \omega} d_0 \Vdash T_0\mleft(n\mright) \in \PlOrd$.
\end{corollary}

The natural choice is to let $T_0\mleft(n\mright) = S_0\mleft(n\mright)^\pl$ using $S_0$ from \autoref{cor:kpp-incom-omega-seq}. We need a lemma:
\begin{lemma}[$\IKPP$]
  \label{lem:lin-ord-pl-pres-incom}
  If $\alpha$ is a linear ordinal, i.e.\ it satisfies the trichotomy
  \[\opA{\beta, \gamma \in \alpha} \left(\beta \in \gamma \lor \beta = \gamma \lor \gamma \in \beta\right),\]
  such as $\omega$, then for any $\beta, \gamma \in \Pow\mleft(\alpha\mright) \cap \Ord$, we have
  \[\beta^\pl \subseteq \gamma^\pl \rightarrow \beta \subseteq \gamma.\]
\end{lemma}

\begin{proof}
  Assume $\beta^\pl \subseteq \gamma^\pl$ and consider any $\delta \in \beta$. By definition, we have $\delta^\pl \in \gamma^\pl$, i.e.\ $\delta^\pl \in \left(\varepsilon^\pl\right)^\pls$ for some $\varepsilon \in \gamma$. Observe that $\varepsilon \in \delta$ implies $\varepsilon^\pl \in \delta^\pl$, which implies $\varepsilon^\pl \in \delta^\pl \subseteq \varepsilon^\pl$ and thus a contradiction. By trichotomy, we must then have $\delta \in \varepsilon$ or $\delta = \varepsilon$. In either case, we have $\delta \in \gamma$ as desired.
\end{proof}

\begin{proof}[Proof of \autoref{cor:kpp-incom-pl-omega-seq}]
  Take $S_0$ in \autoref{cor:kpp-incom-omega-seq} and let
  \[T_0\mleft(n\mright) = S_0\mleft(n\mright)^\pl.\]
  Observe that \autoref{lem:pl-in-pl-ord-unif-real} gives the realisers for plumpness, so it suffices to check strong incomparability.

  For any $i, j \in \omega$, suppose that $T_0\mleft(i\mright) \subseteq T_0\mleft(j\mright)$ is realised. We know that $V\mleft(\Kl\mright) \vDash T_0\mleft(i\mright) = S_0\mleft(i\mright)^\pl \land T_0\mleft(j\mright) = S_0\mleft(j\mright)^\pl$, and it recognises additionally that $S_0\mleft(i\mright), S_0\mleft(j\mright)$ are subsets of $\omega$, where $\omega$ is a linear ordinal by induction on it. It follows from \autoref{lem:lin-ord-pl-pres-incom} that
  \[V\mleft(\Kl\mright) \vDash T_0\mleft(i\mright) \subseteq T_0\mleft(j\mright) \rightarrow S_0\mleft(i\mright) \subseteq S_0\mleft(j\mright),\]
  thus $V\mleft(\Kl\mright) \vDash S_0\mleft(i\mright) \subseteq S_0\mleft(j\mright)$ and $i = j$ by assumptions in \autoref{cor:kpp-incom-omega-seq}.
\end{proof}

\subsection*{Classical incomparables}

The iterated process we shall now describe will use $T_0$ as a starting point to build sequences of pairwise strongly incomparable names on arbitrarily large domains $x$.

For this process to work over the entire universe, first note that in \cite{rathjen20-power-kp-choice}, Rathjen modified the relativised ordinal analysis of $\KPP$ in \cite{rathjen14-power-kp} to prove that $\KPP + \ACg \equiv_{\Con} \KPP$. The precise proof-theoretic treatment will not be relevant to this thesis, but we will observe that, by \autoref{lem:kpp-unc-ord-str}, the theory
\[\KPP + \ACg + \text{$\omega_1$ exists}\]
proves $\Con\mleft(\KPP\mright)$ and thus $\Con\mleft(\KPP + \ACg\mright)$ as well. It follows that, assuming $\KPP + \ACg$ is consistent, then it cannot prove $\omega_1$ exists and thus
\[\KPP + \ACg + \text{every $\alpha \in \Ord$ is countable}\]
is consistent. In other words, this extension is equiconsistent with $\KPP$ itself. We will now work in this extended base theory for the remainder of the section, and look at the following constructions:
\begin{lemma}
  \label{lem:kpp-union-real-pl}
  There exists some $f \in E_\wp\mleft(\varnothing\mright)$ such that, if $T : x \rightarrow V\mleft(\Kl\mright)$ is a function in the base theory and $V_{E_\wp\mleft(X\mright)} \vDash \opA{y \in x} d_0 \Vdash T\mleft(y\mright) \in \PlOrd$, then $\underline{fT} \in V\mleft(\Kl\mright)$ is the `union' of elements in the image of $T$, in the sense that $V\mleft(\Kl\mright) \vDash \underline{fT} \in \PlOrd$ and, for any $z \in V\mleft(\Kl\mright)$, $V\mleft(\Kl\mright) \vDash z \in \underline{fT}$ if and only if $V\mleft(\Kl\mright) \vDash z \in T\mleft(y\mright)$ for some $y \in x$.
\end{lemma}

\begin{proof}
  What we need is simply $\underline{fT} = \bigcup_{y \in x} T\mleft(y\mright)$. Reasoning inside $V_{E_\wp\mleft(X\mright)}$, it is easy to derive that for all $d \in \Kl$, $z \in V\mleft(\Kl\mright)$,
  \[d \Vdash z \in \underline{fT} \leftrightarrow \opE{y \in x} d \Vdash z \in T\mleft(y\mright).\]
  Now, observe that this additionally means
  \[\bp \mleft(\bp 0\mright) \bp_1 \Vdash \opA{z} \left(z \in \underline{fT} \leftrightarrow \opE{y \in T^*} z \in y\right),\]
  where $T^*$ is the name $\left\{\tuple{0, T\mleft(y\mright)} : y \in x\right\} \in V\mleft(\Kl\mright)$, i.e.\ $V\mleft(\Kl\mright) \vDash \underline{fT} = \bigcup T^*$. Since by assumptions, we know $V_{E_\wp\mleft(X\mright)} \vDash \oplam{y} d_0 \Vdash \opA{y \in T^*} y \in \PlOrd$, thus we can apply \autoref{lem:pl-union-still-pl} inside $V\mleft(\Kl\mright)$ to derive that $V\mleft(\Kl\mright) \vDash \underline{fT} \in \PlOrd$.
\end{proof}

\begin{lemma}
  \label{lem:kpp-incom-seq-ord-union}
  Suppose that $a, T, e$ are set-theoretic functions on domain $\alpha \in \Ord$ and for each $\beta \in \alpha$, $T\mleft(\beta\mright) : a\mleft(\beta\mright) \rightarrow V\mleft(\Kl\mright)$ is a sequence of pairwise strongly incomparable names, $e\mleft(\beta\mright) \in E_\wp\mleft(\varnothing\mright)$ and $e\mleft(\beta\mright) \Vdash_\mathfrak{w} \opA{x \in a\mleft(\beta\mright)} d_0 \Vdash T\mleft(\beta\mright)\mleft(x\mright) \in \PlOrd$, then there is some $T' : \bigcup_{\beta \in \alpha} a\mleft(\beta\mright) \rightarrow V\mleft(\Kl\mright)$ that is also pairwise strongly incomparable, and
  \[V_{E_\wp\mleft(X\mright)} \vDash \opA{x \in \bigcup_{\beta \in \alpha} a\mleft(\beta\mright)} d_0 \Vdash T'\mleft(x\mright) \in \PlOrd.\]
  Additionally, this $T'$ \uppar{and the relevant realiser, which must lie in $E_\wp\mleft(\varnothing\mright)$} can be uniquely determined from $a, T, e$ using a $\Sigma_1$-formula containing the primitive symbols $\Pow$ and $R$.
\end{lemma}

\begin{proof}
  Our extended base theory claims that $\alpha \in \Ord$ must be countable, so we can uniquely find some injection $\rho : \alpha \rightarrow \omega$ such that $R\mleft(u, \rho\mright)$ holds where $u$ is the set of all injections from $\alpha$ to $\omega$.

  Now, first observe that $e$ is a set-theoretic function and, as $E_\wp\mleft(X\mright)$ can evaluate set-theoretic functions (cf.\ \autoref{ftn:ep-rec-set-func}), $V_{E_\wp\mleft(X\mright)} \vDash \opA{\beta \in \alpha} \opA{x \in a\mleft(\beta\mright)} d_0 \Vdash T\mleft(\beta\mright)\mleft(x\mright) \in \PlOrd$. Similar to how we previously extracted the plump operator $\left(-\right)^\pl$ into the base theory, we know that the plump successor $\left(-\right)^\pls$, addition $+^\pl$ and multiplication $\cdot^\pl$ are also computed by some terms in $E_\wp\mleft(X\mright)$, taking as input elements in $V\mleft(\Kl\mright)$ and the realisers that they are plump (consisting of entries from both $E_\wp\mleft(X\mright)$ and $\Kl$). Consequently, we can compute, for example,
  \[\eta = \bigcup_{\beta \in \alpha, x \in a\mleft(\beta\mright)} T\mleft(\beta\mright)\mleft(x\mright)^\pls,\]
  and by \autoref{lem:kpp-union-real-pl}, $V\mleft(\Kl\mright) \vDash \eta \in \PlOrd$. We can then define $S' : \bigcup_{\beta \in \alpha} a\mleft(\beta\mright) \rightarrow V\mleft(\Kl\mright)$ as
  \[S'\mleft(x\mright) = \eta \cdot^\pl T_0\mleft(\rho\mleft(\beta_x\mright)\mright) +^\pl T\mleft(\beta_x\mright)\mleft(x\mright),\]
  where $\beta_x = \min \left\{\beta \in \alpha : x \in a\mleft(\beta\mright)\right\}$ and $T_0$ is as given in \autoref{cor:kpp-incom-pl-omega-seq}.

  It is easy to verify that, for any $x \in \bigcup_{\beta \in \alpha} a\mleft(\beta\mright)$, $V\mleft(\Kl\mright) \vDash S'\mleft(x\mright) \in \PlOrd$ by applying results in \autoref{sec:plump} inside $V\mleft(\Kl\mright)$. To see that $S'$ is pairwise strongly incomparable, suppose that $V\mleft(\Kl\mright) \vDash S'\mleft(x\mright) \subseteq S'\mleft(y\mright)$. Inside $V\mleft(\Kl\mright)$, we know that $T\mleft(\beta_x\mright)\mleft(x\mright), T\mleft(\beta_y\mright)\mleft(y\mright) \in \eta$, thus by \autoref{lem:pl-mul-inj}, we must have $T_0\mleft(\rho\mleft(\beta_x\mright)\mright) \subseteq T_0\mleft(\rho\mleft(\beta_y\mright)\mright)$. This implies $\rho\mleft(\beta_x\mright) = \rho\mleft(\beta_y\mright)$ in the base theory and thus $\beta_x = \beta_y$ by injectivity. It then follows from \autoref{lem:pl-add-inj} inside $V\mleft(\Kl\mright)$ that $T\mleft(\beta_x\mright)\mleft(x\mright) \subseteq T\mleft(\beta_y\mright)\mleft(y\mright)$, and thus $x = y$.

  Now, observe that the plump ordinals generated by ordinal arithmetic do not guarantee that their plumpness can be realised by the uniform realiser $d_0$. Thus, we additionally set $T'\mleft(x\mright) = S'\mleft(x\mright)^\pl$, so that they have the correct realisers by \autoref{lem:pl-in-pl-ord-unif-real}. Notice that for each $x \in \bigcup_{\beta \in \alpha} a\mleft(\beta\mright)$, $V\mleft(\Kl\mright)$ does know $S'\mleft(x\mright) \in \PlOrd$ (although by a non-uniform realiser), so $V\mleft(\Kl\mright) \vDash T'\mleft(x\mright) = S'\mleft(x\mright)$ by \autoref{prop:pl-op-surj}. It is thus easy to verify that also $V\mleft(\Kl\mright) \vDash T'\mleft(x\mright) \subseteq T'\mleft(y\mright)$ implies $x = y$.
\end{proof}

\begin{corollary}
  \label{cor:kpp-incom-seq-times-2}
  Suppose that $T : x \rightarrow V\mleft(\Kl\mright)$ is a sequence of pairwise strongly incomparable names and $V_{E_\wp\mleft(X\mright)} \vDash \opA{y \in x} d_0 \Vdash T\mleft(y\mright) \in \PlOrd$, then there is some $T' : x \times 2 \rightarrow V\mleft(\Kl\mright)$ \uppar{which can be determined using a $\Sigma_1$-formula}, that is also pairwise strongly incomparable and realised to be a collection of plump ordinals.
\end{corollary}

\begin{proof}
  Simply use \autoref{lem:kpp-incom-seq-ord-union} where $\alpha = 2$, $a\mleft(i\mright) = \left\{\tuple{y, i} : y \in x\right\}$, on some function
  \[S\mleft(i\mright)\mleft(\tuple{y, i}\mright) = T\mleft(y\mright). \qedhere\]
\end{proof}

\begin{lemma}
  \label{lem:kpp-incom-seq-pow}
  Suppose that $T : x \rightarrow V\mleft(\Kl\mright)$ is a sequence of pairwise strongly incomparable names and $V_{E_\wp\mleft(X\mright)} \vDash \opA{y \in x} d_0 \Vdash T\mleft(y\mright) \in \PlOrd$, then there is some $T' : \Pow\mleft(x\mright) \rightarrow V\mleft(\Kl\mright)$ \uppar{which can be determined using a $\Sigma_1$-formula}, that is also pairwise strongly incomparable and uniformly realised to be a collection of plump ordinals \uppar{by $d_0 \in \omega$}.
\end{lemma}

\begin{proof}
  Since we are working in a classical base theory, we have a function $\chi : \Pow\mleft(x\mright) \rightarrow \Pow\mleft(x \times 2\mright)$ given by
  \[\chi\mleft(s\mright) = \left\{\tuple{y, 0} : y \in s\right\} \cup \left\{\tuple{y, 1} : y \in x \setminus s\right\}.\]
  It is easy to check that $\chi$ is a `classical strongly incomparable map', in the sense that $\chi\mleft(s\mright) \subseteq \chi\mleft(t\mright) \rightarrow s = t$, simply because each $\chi\mleft(s\mright)$ must be a function on domain $x$.

  Now, we let $S$ be the sequence of pairwise strongly incomparable names on $x \times 2$ given by \autoref{cor:kpp-incom-seq-times-2}, and define some $S' : \Pow\mleft(x\mright) \rightarrow V\mleft(\Kl\mright)$ as
  \[S'\mleft(s\mright) = \bigcup_{y \in \chi\mleft(s\mright)} S\mleft(y\mright)^\pls,\]
  such that $V\mleft(\Kl\mright) \vDash S'\mleft(s\mright) \in \PlOrd$ for all $s \subseteq x$ by \autoref{lem:kpp-union-real-pl}. For strong incomparability, suppose that $V\mleft(\Kl\mright) \vDash S'\mleft(s\mright) \subseteq S'\mleft(t\mright)$. For any $y \in \chi\mleft(s\mright)$, we have $V\mleft(\Kl\mright) \vDash S\mleft(y\mright) \in S'\mleft(s\mright)$ and thus there must exist some $z \in \chi\mleft(t\mright)$ such that
  \[V\mleft(\Kl\mright) \vDash S\mleft(y\mright) \in S\mleft(z\mright)^\pls,\]
  i.e.\ $V\mleft(\Kl\mright) \vDash S\mleft(y\mright) \subseteq S\mleft(z\mright)$. It follows from the strong incomparability of $S$ that $y = z$ and thus $\chi\mleft(s\mright) \subseteq \chi\mleft(t\mright)$, which implies $s = t$ immediately.

  Finally, using the same trick as in the proof of \autoref{lem:kpp-incom-seq-ord-union} and taking $T'\mleft(s\mright) = S'\mleft(s\mright)^\pl$ will ensure that the names have uniform realisers for plumpness.
\end{proof}

Now, we consider the powerset hierarchy
\begin{align*}
  V_0          & = \varnothing,                                                           \\
  V_{\alpha^+} & = \Pow\mleft(V_\alpha\mright),                                           \\
  V_\lambda    & = \bigcup_{\alpha \in \lambda} V_\alpha \quad \text{for limit $\lambda$}
\end{align*}
inside the classical base theory, and we define two class functions $e, T$ on $\Ord$ by $\Sigma$-recursion (using formulae containing symbols $\Pow$ and $R$):
\begin{enumerate}[label=(\arabic*)]
  \item Let $k : \omega \rightarrow V_\omega$ be an arbitrary bijection --- for example, the Ackermann encoding. Then $T\mleft(\omega\mright) : V_\omega \rightarrow V\mleft(\Kl\mright)$ is given by
        \[T\mleft(\omega\mright)\mleft(x\mright) = T_0\mleft(k^{-1}\mleft(x\mright)\mright)\]
        with $T_0$ as in \autoref{cor:kpp-incom-pl-omega-seq} and $e\mleft(\omega\mright) \Vdash_\mathfrak{w} \opA{x \in V_\omega} d_0 \Vdash T\mleft(\omega\mright)\mleft(x\mright) \in \PlOrd$. Additionally, $T\mleft(n\mright) = \left.T\mleft(\omega\mright)\right|_{V_n}$ and $e\mleft(n\mright) = e\mleft(\omega\mright)$ for $n \in \omega$.
  \item When $\alpha > \omega$ and $\alpha = \beta^+$ is a successor, $T\mleft(\alpha\mright) : \Pow\mleft(V_\beta\mright) \rightarrow V\mleft(\Kl\mright)$ and the corresponding $e\mleft(\alpha\mright) \Vdash_\mathfrak{w} \opA{x \in \Pow\mleft(V_\beta\mright)} d_0 \Vdash T\mleft(\alpha\mright)\mleft(x\mright) \in \PlOrd$ are uniquely determined from $T\mleft(\beta\mright)$ and $e\mleft(\beta\mright)$ using \autoref{lem:kpp-incom-seq-pow}.
  \item When $\alpha > \omega$ is a limit, then $V_\alpha = \bigcup_{\beta \in \alpha} V_\beta$ and $T\mleft(\alpha\mright) : V_\alpha \rightarrow V\mleft(\Kl\mright)$ and the corresponding $e\mleft(\alpha\mright) \Vdash_\mathfrak{w} \opA{x \in V_\alpha} d_0 \Vdash T\mleft(\alpha\mright)\mleft(x\mright) \in \PlOrd$ are uniquely determined from $T\mleft(\beta\mright)$'s and $e\mleft(\beta\mright)$'s using \autoref{lem:kpp-incom-seq-ord-union} and the choice function $R$.
\end{enumerate}

From here onwards, we fix $X$ as the specific $\Sigma_1$-class function on $V\mleft(\Kl\mright)$, such that for any $x \in V\mleft(\Kl\mright)$, we fix $\alpha = \min \left\{\beta \in \Ord : x \subseteq V_\beta\right\}$, and define $X\mleft(x\mright) = \tuple{f_x, e\mleft(\alpha\mright)}$ where
\[f_x = \left\{\tuple{n, \tuple{y, T\mleft(\alpha\mright)\mleft(\tuple{n, y}\mright)}^\bullet} : \tuple{n, y} \in x\right\} \in V\mleft(\Kl\mright),\]
where $\tuple{a, b}^\bullet = \left\{\left\{a\right\}^\bullet, \left\{a, b\right\}^\bullet\right\}^\bullet$ and $\left\{a, b\right\}^\bullet = \left\{\tuple{0, a}, \tuple{1, b}\right\}$ are the standard names for ordered and unordered pairs in $V\mleft(\Kl\mright)$. We check that:
\begin{lemma}
  \label{lem:kpp-incomp-relation-in-vkl}
  $V\mleft(\Kl\mright)$ realises that for every $x$, there exists some $u$ such that:
  \begin{enumerate}[label=\uplabel{\roman*}]
    \item \label{itm:kpp-real-internal-incom-prop-pl} $\opA{{\tuple{y, \alpha} \in u}} \left(y \in x \land \alpha \in \PlOrd\right)$;
    \item \label{itm:kpp-real-internal-incom-prop-surj} $\opA{y \in x} \opE{\alpha} \tuple{y, \alpha} \in u$;
    \item \label{itm:kpp-real-internal-incom-prop-incom} $\opA{{\tuple{y, \alpha}, \tuple{z, \beta} \in u}} \left(\alpha \subseteq \beta \rightarrow y = z \land \alpha = \beta\right)$.
  \end{enumerate}
\end{lemma}

\begin{proof}
  The realisers for $V\mleft(\Kl\mright)$ just have to realise the three statements separately, but the realisers for $V_{E_\wp\mleft(X\mright)}$ must output $u$ and appropriate additional information given input $x \in V\mleft(\Kl\mright)$, which is why we need the specific class function $X$. Fix any input $x \in V\mleft(\Kl\mright)$, we will simply assign to $u$ the value $f_x = \underline{\bp_0 \mleft(\bx x\mright)}$.

  Now, the realiser for \ref{itm:kpp-real-internal-incom-prop-pl} in $V\mleft(\Kl\mright)$ is simply\footnote{Here the routine component of a realiser for the defining formula of an ordered tuple is omitted.} $g = \oplam{n} \bp \mleft(\bp n \ior\mright) d_0$ using the constant $d_0$ from \autoref{lem:pl-in-pl-ord-unif-real}. To see that
  \[V_{E_\wp\mleft(X\mright)} \vDash g \Vdash \opA{{\tuple{y, \alpha} \in u}} \left(y \in x \land \alpha \in \PlOrd\right),\]
  we observe that $V_{E_\wp\mleft(X\mright)}$ knows about the definition of $f_x$, and for any $\tuple{n, y} \in x$, the fact that $\bp n \ior \Vdash y \in x$ is a theorem in $\CZFP$ and thus must be realised. To realise the second conjunct, we check that $E_\wp\mleft(X\mright)$ is capable of recovering each $\tuple{n, y} \in x$ from the corresponding element in $f_x$ since each $\left\{a, b\right\}^\bullet$ is essentially a function on $\left\{0, 1\right\}$ in the base theory. For any $t = \tuple{n, \tuple{y, \alpha}^\bullet} \in f_x$, we just need to apply the $E_\wp\mleft(X\mright)$-function $e = \underline{\bp_1 \mleft(\bx x\mright)}$ to
  \[\tuple{n, y} = \tuple{n, \underline{\bp_1 t}\mleft(0\mright)\mleft(0\mright)},\]
  which can be easily evaluated in $E_{\wp}\mleft(X\mright)$, to obtain a realiser of $d_0 \Vdash \alpha \in \PlOrd$.

  The realiser for \ref{itm:kpp-real-internal-incom-prop-surj} is trivial, both for the part inside $V\mleft(\Kl\mright)$ and the part inside $V_{E_\wp\mleft(X\mright)}$, as it is equivalently a simple $\Delta_0$-claim. For \ref{itm:kpp-real-internal-incom-prop-incom}, we claim that the realiser in $V\mleft(\Kl\mright)$ is simply $\oplam{nmd} \bp \ior \ior$. Suppose that given some $\tuple{n, \tuple{y, \alpha}^\bullet}, \tuple{m, \tuple{z, \beta}^\bullet} \in f_x$, $V\mleft(\Kl\mright)$ sees a realiser for $\alpha \subseteq \beta$, then this is a realiser for $T\mleft(\gamma\mright)\mleft(\tuple{n, y}\mright) \subseteq T\mleft(\gamma\mright)\mleft(\tuple{m, z}\mright)$ for some classical $\gamma \in \Ord$, and by our specific, recursive construction of $T$, this implies $\tuple{n, y} = \tuple{m, z}$ and thus also $y = z \land \alpha = \beta$ in the classical base theory. When these atomic statements are true, they are also trivially realised in $V_{E_\wp\mleft(X\mright)}$, thus $\underline{\oplam{tsd} \bp 00}$ realises\footnote{Again, we omitted the process of treating the tuple notations as abbreviations and adjusting the realiser accordingly here.}
  \[\opA{{\tuple{n, \tuple{y, \alpha}^\bullet}, \tuple{m, \tuple{z, \beta}^\bullet} \in f_x}} \left(\left(\opE{d \in \Kl} d \Vdash \alpha \subseteq \beta\right) \rightarrow y = z \land \alpha = \beta\right),\]
  and working inside $V_{E_\wp\mleft(X\mright)}$, it follows immediately that $\bp \ior \ior$ realises the two equalities in $V\mleft(\Kl\mright)$.
\end{proof}

We can now work entirely within $V\mleft(\Kl\mright)$ to show that a system satisfying $\CZFP$ and the conditions in \autoref{lem:kpp-incomp-relation-in-vkl} must satisfy $V = L$. This will involve very similar constructions as done in the proofs of \autoref{prop:pl-str-incom-func-as-pl-ord} and \autoref{thm:pl-str-incom-code-power}, but we are unfortunately unable to use the results directly here, since in \autoref{lem:kpp-incomp-relation-in-vkl} we cannot construct these objects $u$ as functions on $x$, but only as some kinds of `multi-valued functions'.

\begin{theorem}
  \label{thm:kpp-real-v-l}
  $V\mleft(\Kl\mright) \vDash V = L$.
\end{theorem}

\begin{proof}
  Looking at the proof of \autoref{cor:many-pl-incom-v-eq-l}, we see that all we need to do is to replace the use of \autoref{thm:pl-str-incom-code-power}, i.e.\ verify that for any $\alpha \in \PlOrd$, also $\Pow\mleft(\alpha\mright) \subseteq L_\pl$.

  Now, fix any $\alpha \in \PlOrd$. Let $u$ be the corresponding subset of $\alpha \times \PlOrd$ as given by \autoref{lem:kpp-incomp-relation-in-vkl}. We first check that $u \in L_\pl$ by imitating \autoref{prop:pl-str-incom-func-as-pl-ord}. We define
  \[\sigma = \bigcup_{\tuple{\beta, \gamma} \in u} \left(\alpha \cdot^\pl \gamma +^\pl \beta\right)^\pls, \qquad \tau = \bigcup_{\tuple{\beta, \gamma} \in u} \gamma^\pls.\]
  Clearly $\sigma, \tau \in \PlOrd$, thus $\sigma, \tau \in L_\pl$. We now extract $u$ from $\alpha, \sigma, \tau$ by defining
  \begin{align*}
    v = & \left\{\tuple{\beta, \gamma} : \beta \in \alpha, \gamma \in \tau, \right.                                                                                                                                 \\
        & \qquad \left.\opE{\delta \in \sigma} \left(\delta = \alpha \cdot^\pl \gamma +^\pl \beta \land \opA{\delta' \in \sigma} \left(\delta \subseteq \delta' \rightarrow \delta = \delta'\right)\right)\right\}.
  \end{align*}
  For any $\tuple{\beta, \gamma} \in u$, simply let $\delta = \alpha \cdot^\pl \gamma +^\pl \beta$. Observe that if $\delta \subseteq \delta' \in \sigma$ then there is $\tuple{\varepsilon, \zeta} \in u$ such that $\delta' \subseteq \alpha \cdot^\pl \zeta +^\pl \varepsilon$. By \autoref{lem:pl-mul-inj}, this implies $\gamma \subseteq \zeta$, which then implies $\beta = \varepsilon$ and $\gamma = \zeta$ by \autoref{lem:kpp-incomp-relation-in-vkl}\ref{itm:kpp-real-internal-incom-prop-incom}.

  Conversely, for any $\tuple{\beta, \gamma} \in v$, there must exist $\tuple{\varepsilon, \zeta} \in u$ such that $\alpha \cdot^\pl \gamma +^\pl \beta \subseteq \alpha \cdot^\pl \zeta +^\pl \varepsilon$. Then $\alpha \cdot^\pl \zeta +^\pl \varepsilon \in \sigma$ and stipulations in the definition of $v$ imply that
  \[\alpha \cdot^\pl \gamma +^\pl \beta = \alpha \cdot^\pl \zeta +^\pl \varepsilon.\]
  It follows that $\beta = \varepsilon$ and $\gamma = \zeta$ by \autoref{lem:pl-mul-inj} and \autoref{lem:pl-add-inj}. Therefore, we have $u = v \in L_\pl$ as desired.

  Now, consider any subset $x \subseteq \alpha$, and we define $\rho = \bigcup_{\tuple{\beta, \gamma} \in u, \beta \in x} \gamma^\pls$. Clearly $\rho \in \PlOrd$ and thus $\rho \in L_\pl$. Since $u \in L_\pl$, we also have
  \[y = \left\{\beta \in \alpha : \opE{\gamma \in \rho} \tuple{\beta, \gamma} \in u\right\} \in L_\pl.\]
  Obviously $x \subseteq y$ by \autoref{lem:kpp-incomp-relation-in-vkl}\ref{itm:kpp-real-internal-incom-prop-surj}. On the other hand, for any $\beta \in y$, there is $\gamma \in \rho$, i.e.\ $\tuple{\varepsilon, \zeta} \in u$ satisfying $\varepsilon \in x$, $\gamma \subseteq \zeta$, and also $\tuple{\beta, \gamma} \in u$. By \autoref{lem:kpp-incomp-relation-in-vkl}\ref{itm:kpp-real-internal-incom-prop-incom} we must have $\beta = \varepsilon$ and $\gamma = \zeta$, thus $\beta \in x$. We have now shown that $x = y \in L_\pl$, i.e.\ $\Pow\mleft(\alpha\mright) \subseteq L_\pl$.
\end{proof}

\begin{corollary}
  The consistency of $\KPP$ implies that of $\CZFP + V = L$. It follows that
  \[\CZFP + V = L \equiv_{\Con} \KPP \equiv_{\Con} \CZFP.\]
  Since $\IKPP$ is a subtheory of $\CZFP$ of the same strength, the same equiconsistency also extends easily to $\IKPP$ and $\IKPP + V = L$.
\end{corollary}

\begin{proof}
  This follows immediately from \autoref{thm:kpp-real-v-l}.
\end{proof}

\chapter{Realisability with large pcas}
\label{chap:large-pcas}

Observe that the key constructions in \autoref{chap:incom-coding} and \autoref{chap:plump-power} both use the specific pca $\Kl$, which means that we have to start from a countable collection of incomparable objects, and even with the iterated construction in \autoref{sec:czfp-v-eq-l}, we cannot obtain arbitrarily large incomparable collections without some constraints on the height of the universe. In contrast, in \cite{wang-sw26-plump}*{section 5}, I built large collections of strongly incomparable ordinals directly by forcing on a large Heyting algebra and proved that:
\begin{theorem*}[$\ZFC$, \cite{wang-sw26-plump}*{Corollary 5.9}]
  Fix any set $x$, then there exists a Heyting algebra $H$ such that the corresponding Heyting-valued model $V^H$ satisfies
  \[V^H \vDash \Pow\mleft(\check{x}\mright) \in L.\]
\end{theorem*}

In this final chapter, we will discuss how the use of larger pcas can allow us to develop similar constructions in realisability interpretations.

\section{Classical embeddings}

We would like to start with a classical system, just as in the forcing scenario, and embed fragments of the base theory into our realisability model. However, this is often inappropriate when studying weak set theories because following \autoref{thm:czf-em-eq-zf}, a classical theory with an unrestricted strong collection principle like in $\CZF$ may be much stronger than the intuitionistic counterpart which we would like to model. On the other hand, we may not want to perform interpretations in the form of \autoref{sec:czf-real-in-czf} directly in weaker classical theories like $\KP$, since we do need principles like unrestricted strong collection to hold in the base theory in order for them to be realised (cf.\ the proof of \autoref{lem:czf-str-coll-real}).

Instead, in this chapter we will work inside the base theory $\CZF^- + \REM$ where $\CZF^- = \CZF - \text{subset collection}$\inds{constructive Zermelo--Fraenkel set theory, $\CZF$!0@$\CZF^-$}, and $\REM$ denotes restricted excluded middle\inds{restricted excluded middle, $\REM$}, i.e.\ $\varphi \lor \neg \varphi$ for any $\Delta_0$-formula $\varphi$. As we have remarked in \autoref{sec:czf-weak-real}, one can always build a model of $\CZF^- + \REM$ from the universe of a model of $\KP$, by interpreting truth in the model through weak realisability conditions.

The construction of the inductive $\Sigma_1$-class $V\mleft(\mathcal{A}\mright)$ and the corresponding realisability conditions in \autoref{sec:czf-real-in-czf} only requires $\IKP$, so they will work inside $\CZF^-$ already. Observe that none of the proofs in that section require subset collection --- except verifying the realisability of subset collection itself --- thus we immediately have that:
\begin{theorem}
  Let $\varphi$ be any theorem of $\CZF^-$, then $\CZF^-$ proves $V\mleft(\mathcal{A}\mright) \vDash \varphi$, i.e.\ there exists $d \in \mathcal{A}$ such that $d \Vdash \varphi$.
\end{theorem}

Now, internally in any universe of intuitionistic set theory, we shall say that a set model $\tuple{x; \in}$, where $x$ is any transitive set, is \emph{$\Delta_0$-decidable} if it satisfies the restricted law of excluded middle, i.e.
\[\opA{y_1, \ldots, y_n \in x} \left(\varphi \lor \neg \varphi\right)\]
holds for every $\Delta_0$-formula $\varphi$ where $y_1, \ldots, y_n$ enumerate all free variables in $\varphi$. For any ordinal $\alpha \in \Ord$, we shall say that $\alpha$ \emph{embeds decidably} into the universe if $L_\alpha$ is a $\Delta_0$-decidable set model.

\begin{lemma}[$\IKP$]
  \label{lem:dec-emb-def-in-l}
  If $\alpha$ embeds decidably into the universe, then for any $\beta \in \alpha^+$, $\beta$ is a linear ordinal and we have
  \[\opA{\gamma \in \beta} \gamma^+ = \dd\mleft(L_\gamma\mright) \cap \Ord,\]
  i.e.\ $\beta = L_\beta \cap \Ord$ immediately.
\end{lemma}

\begin{proof}
  We prove this by set induction on $\beta$. Fix $\beta \in \alpha^+$, then by the inductive hypothesis, every $\gamma \in \beta$ is linear, and
  \[\gamma = L_\gamma \cap \Ord \in L_\alpha.\]

  We first check that, for any $\delta \in L_\alpha \cap \Ord$, if $\delta \subseteq \gamma$ then either $\delta = \gamma$ or $\delta \in \gamma$. This is completely analogous to the proof of \autoref{cor:l-omega-ord}: by $L_\alpha$ being $\Delta_0$-decidable, we have
  \[\opE{\varepsilon \in \gamma} \varepsilon \not\in \delta \lor \opA{\varepsilon \in \gamma} \varepsilon \in \delta.\]
  The latter disjunct implies $\delta = \gamma$; in the former case, we show that $\delta \subseteq \varepsilon$. Suppose that this is not the case, i.e.\ there is some $\eta \in \delta$ such that $\eta \not\in \varepsilon$. However, both $\varepsilon, \eta \in \gamma$, thus by the linearity of $\gamma$, we have either $\varepsilon \in \eta$ or $\varepsilon = \eta$, both contradicting $\varepsilon \not\in \delta$. Thus we must have $\delta \subseteq \varepsilon$, since we can prove $\Delta_0$-formulae by contradiction inside $L_\alpha$.

  Hence, $\delta \subseteq \gamma$ implies either $\delta = \gamma$ or $\delta \subseteq \varepsilon$ for some $\varepsilon \in \gamma$. By set induction on $\gamma \in \beta$, it follows that
  \[\opA{\delta \in \dd\mleft(L_\alpha\mright) \cap \Ord} \opA{\gamma \in \beta} \left(\delta \subseteq \gamma \rightarrow \delta \in \gamma \lor \delta = \gamma\right).\]
  Thus $\dd\mleft(L_\gamma\mright) \cap \Ord = \gamma^+$ for every $\gamma \in \beta$, and $\beta = L_\beta \cap \Ord$.

  Finally, for linearity, we observe that the argument above can also be applied to certain $\delta \in \dd\mleft(L_\alpha\mright) \cap \Ord$, where we think of $\delta$ as a $\Delta_0$-definable `proper class' in the $\Delta_0$-decidable set model $L_\alpha$. Specifically, we take any $\gamma_1, \gamma_2 \in \beta$ and consider\footnote{\label{ftn:int-l-delta-0-sep}Before we prove the linearity of $\alpha$, we cannot take two arbitrary parameters $x, y \in L_\alpha$ and expect that an instance of $\Delta_0$-separation with both parameters still hold in $L_\alpha$, this is because they could be from $\dd\mleft(L_{\beta_1}\mright), \dd\mleft(L_{\beta_2}\mright)$ respectively, with $\beta_1, \beta_2 \in \alpha$ strongly incomparable, so that we cannot find $\sup\mleft\{\beta_1, \beta_2\mright\} \in \alpha$ and the classical argument for $\Delta_0$-separation will not apply. See \autoref{lem:zf-l-delta-0-separation} for more detail.} $\gamma_1 \cap \gamma_2 \in \dd\mleft(L_\alpha\mright)$. The decidability of $L_\alpha$ then implies
  \[\opE{\varepsilon \in \gamma'} \varepsilon \not\in \gamma_1 \cap \gamma_2 \lor \opA{\varepsilon \in \gamma'} \varepsilon \in \gamma_1 \cap \gamma_2,\]
  for any $\gamma' \in \beta$, where $\varepsilon \in \gamma_1 \cap \gamma_2$ can be read as $\varepsilon \in \gamma_1 \land \varepsilon \in \gamma_2$ to not involve the set $\gamma_1 \cap \gamma_2$ explicitly. By set induction on $\gamma'$, the same argument as above then implies both $\gamma_1 \cap \gamma_2 \in \gamma_1 \lor \gamma_1 \cap \gamma_2 = \gamma_1$ and $\gamma_1 \cap \gamma_2 \in \gamma_2 \lor \gamma_1 \cap \gamma_2 = \gamma_2$. Since having $\gamma_1 \cap \gamma_2 \in \gamma_1$ and $\gamma_1 \cap \gamma_2 \in \gamma_2$ simultaneously would imply
  \[\gamma_1 \cap \gamma_2 \in \gamma_1 \cap \gamma_2,\]
  a contradiction, we must have exactly one of the three other cases, i.e.
  \[\gamma_1 \in \gamma_2 \lor \gamma_1 = \gamma_2 \lor \gamma_2 \in \gamma_1. \qedhere\]
\end{proof}

This essentially tells us that when $\alpha$ embeds decidably into an intuitionistic universe, then the initial segment of ordinals before $\alpha$ behaves just like its classical counterpart, even with respect to constructions like G\"odel's $L$. In particular, we immediately have the following analogue of \autoref{cor:omega-sub-strong-incom-code}:
\begin{corollary}
  \label{cor:dec-emb-sub-strong-incom-code}
  If $\alpha \in \Ord$ embeds decidably into the universe, $f : x \rightarrow \Pow\mleft(\alpha\mright) \cap \Ord$ is pairwise strongly incomparable and $f \in L$, then for any $u \subseteq x$, we have $u \in L$.
\end{corollary}

\begin{proof}
  For any $\beta \in \Pow\mleft(\alpha\mright) \cap \Ord$, we have
  \[\beta = \bigcup_{\gamma \in \beta} \gamma^+ = \bigcup_{\gamma \in \beta} \dd\mleft(L_\gamma\mright) \cap \Ord = L_{\bigcup_{\gamma \in \beta} \gamma^+} \cap \Ord = L_\beta \cap \Ord\]
  by \autoref{lem:dec-emb-def-in-l}. The claim then follows from \autoref{prop:l-ord-strong-incom-code}.
\end{proof}

Observe that the relativisation $\varphi^x$ of a formula $\varphi$ to any set model $x$ is a $\Delta_0$-formula, thus $\CZF^- + \REM$ proves that every set model satisfies full classical logic\footnote{Pedantically, one might need to distinguish the instances of excluded middle in the axiom scheme for formula $\varphi$ in the \emph{metatheory} and quantifications over all formulae (or all G\"odel numberings) within the set theory itself. However, through the construction of truth predicates (cf.\ \autoref{sec:intro-int-l} and Lubarsky \cite{lubarsky93-intuitionistic-l}), one can show that for any (G\"odel numbering of) formula $\varphi$, the collection of instances $\left\{\vec{y} \in x^k : \tuple{x; \in} \vDash \varphi\mleft(\vec{y}\mright)\right\}$ forms a set. Thus, all we need here is the single instance of excluded middle
  \[\opA{\vec{y} \in x^k} \left(\vec{y} \in \left\{\vec{y} \in x^k : \tuple{x; \in} \vDash \varphi\mleft(\vec{y}\mright)\right\} \lor \vec{y} \not\in \left\{\vec{y} \in x^k : \tuple{x; \in} \vDash \varphi\mleft(\vec{y}\mright)\right\}\right)\]
  for an atomic formula in $\REM$.} and every ordinal embeds decidably into the universe.

However, we see that this is not the case for the realisability model:
\begin{proposition}
  \label{prop:vkl-omega-suc-not-cl}
  For any pca \uppar{over natural numbers} $\mathcal{A}$, $\omega$ embeds decidably into $V\mleft(\mathcal{A}\mright)$. However, $\omega^+$ does not embed decidably into $V\mleft(\Kl\mright)$.
\end{proposition}

\begin{proof}
  The first statement follows simply from \autoref{thm:l-omega-dec}. For the second part, take any $\Pi_1$-definable subset $s \subseteq \omega$ in the language of $\PA$ that is not recursive, i.e.\ there is no $e \in \Kl$ such that for any $n \in \omega$, $en {\downarrow} \in \left\{0, 1\right\}$ and $en = 1$ if and only if $n \in s$. For example, simply take the Turing jump of the empty set $\TJ\mleft(\varnothing\mright)$ as given explicitly by a universal $\Pi_1$-formula in \autoref{sec:oa-cm-gwo}. Now, let $\varphi$ be the defining formula for $s$, then we have
  \[s^* = \left\{x \in L_\omega : L_\omega \vDash x \in \Ord \land \varphi^\Ord\mleft(x\mright)\right\} \in \dd\mleft(L_\omega\mright),\]
  where $\varphi^\Ord$ is the set-theoretic formula obtained from $\varphi$ by replacing unbounded first-order quantifiers by quantifiers over $\Ord$ and arithmetic operations by their $\Sigma_1$-definition as a subset of $L_\omega$.

  Now, suppose that $\omega^+$ embeds decidably into $V\mleft(\Kl\mright)$, then we must have $V\mleft(\Kl\mright) \vDash \opA{n \in \omega} \left(n \in s^* \lor n \not\in s^*\right)$, i.e.
  \[V\mleft(\Kl\mright) \vDash \opA{n \in \omega} \left(\varphi^\omega\mleft(n\mright) \lor \neg \varphi^\omega\mleft(n\mright)\right).\]
  Specifically, it is easy to produce canonical names for both $\omega$ (i.e.\ $\overline{\omega}$ as defined in \autoref{sec:czf-real-in-czf}) and arithmetic operations on them:
  \begin{align*}
    \oplus & = \left\{\tuple{\bp nm, \tuple{\tuple{\overline{n}, \overline{m}}^\bullet, \overline{n + m}}^\bullet} : n, m \in \omega\right\},     \\
    \odot  & = \left\{\tuple{\bp nm, \tuple{\tuple{\overline{n}, \overline{m}}^\bullet, \overline{n \cdot m}}^\bullet} : n, m \in \omega\right\},
  \end{align*}
  so we can equivalently use explicit names $\overline{\omega}, \oplus, \odot$ in $\varphi^\omega$, such that it is formally the relativisation of a $\Pi_1$-formula in the language of set theory to the domain $\omega$ with additional parameters $\oplus, \odot$.

  Now, we know that for any $n, m \in \omega$, $\overline{n} \in \overline{m}$ is realised if and only if $n \in m$, in which case the realiser is always trivially $\bp n \ior$; similarly, correct instances of $\overline{n} = \overline{m}$, $\tuple{\tuple{\overline{n}, \overline{m}}, \overline{k}} \in \oplus$ and $\tuple{\tuple{\overline{n}, \overline{m}}, \overline{k}} \in \odot$ are all realised trivially. It follows from easy induction on the complexity of $\psi$ that for any $\Delta_0$-formula $\psi$ in the language of $\PA$, we have\footnote{Observe that this claim is set up in the same style as \autoref{prop:ep-real-delta0-correctness}, and also proved by a similar induction; the construction of the realisers here, however, will be much simpler. Also cf.\ \autoref{lem:alpha-rec-real-delta0-correctness} for a similar process.} $d_\psi, k_\psi \in \Kl$ such that for any $n_1, \ldots, n_k \in \omega$,
  \begin{enumerate}[label=\uplabel{\roman*}]
    \item $d_\psi n_1 \cdots n_k {\downarrow} \in \left\{0, 1\right\}$ and $d_\psi n_1 \cdots n_k = 1$ if and only if $\psi\mleft(n_1, \ldots, n_k\mright)$ is true on $\omega$ in the base theory;
    \item if $d_\psi n_1 \cdots n_k = 1$, then $k_\psi n_1 \cdots n_k {\downarrow} \Vdash \psi^\omega\mleft(\overline{n_1}, \ldots, \overline{n_k}\mright)$; if $d_\psi n_1 \cdots n_k = 0$, then $\neg\!\opE{e \in \Kl} e \Vdash \psi^\omega\mleft(\overline{n_1}, \ldots, \overline{n_k}\mright)$.
  \end{enumerate}

  It follows that also for the $\Pi_1$-formula $\varphi\mleft(n\mright) = \opA{m} \psi\mleft(n, m\mright)$, $\opE{e \in \Kl} e \Vdash \varphi^\omega\mleft(\overline{n}\mright)$ if and only if $\oplam{m} k_\psi nm \Vdash \varphi^\omega\mleft(\overline{n}\mright)$, if and only if $\varphi\mleft(n\mright)$ is true on $\omega$ in the base theory, if and only if $n \in s$. Thus, if some $d \in \Kl$ satisfies
  \[d \Vdash \opA{n \in \omega} \left(\varphi^\omega\mleft(n\mright) \lor \neg \varphi^\omega\mleft(n\mright)\right),\]
  then we must have $\bp_0 \mleft(dn\mright) {\downarrow} \in \left\{0, 1\right\}$ and $\bp_0 \mleft(dn\mright) = 0$ if and only if $n \in s$ for all $n \in \omega$. However, this is impossible as $s$ is not recursive.
\end{proof}

\autoref{prop:vkl-omega-suc-not-cl} additionally shows that being a linear ordinal is not sufficient for the decidable embedding as we defined here. In what follows, we will be concerned with the process of embedding larger ordinals decidably into a realisability model.

\subsection*{Linearity and decidability}

We shall point out here that the linearity of an ordinal still simplifies the process of verifying that it embeds decidably into the universe. Examine the following proof of the classical result:
\begin{lemma}[$\ZF$ or $\KP$]
  \label{lem:zf-l-delta-0-separation}
  For any $\alpha \in \Ord$, $\tuple{L_\alpha; \in}$ satisfies the axiom scheme of $\Delta_0$-separation.
\end{lemma}

\begin{proof}
  Let $\varphi$ be any $\Delta_0$-formula and $y_1, \ldots, y_n \in L_\alpha$, we want to show that
  \[x = \left\{z \in y_1 : \varphi\mleft(z, y_1, y_2, \ldots, y_n\mright)\right\} \in L_\alpha.\]
  Take $\beta_1, \ldots, \beta_n \in \alpha$ such that each $y_i \in \dd\mleft(L_{\beta_i}\mright)$, and let $\beta = \max \left\{\beta_1, \ldots, \beta_n\right\}$, then we must have $y_1, \ldots, y_n \in \dd\mleft(L_\beta\mright)$, i.e.\ there will be formulae $\psi_1, \ldots, \psi_n$ and parameters $\vec{u}_1, \ldots, \vec{u}_n \in L_\beta^{< \omega}$ such that for each $y_i$,
  \[\opA{x \in L_\beta} \left(x \in y_i \leftrightarrow \psi_i^{L_\beta}\mleft(x, \vec{u}_i\mright)\right).\]
  We can now transform $\varphi$ to an equivalent formula $\tilde{\varphi}$ by the following steps:
  \begin{enumerate}[label=\uplabel{\arabic*}]
    \item replace every instance of $y_i \in v$ by $\opE{z \in v} z = y_i$, where $v$ denotes any variable;
    \item replace every instance of $v = y_i$ or $y_i = v$ by $\opA{z \in v} z \in y_i \land \opA{z \in y_i} z \in v$;
    \item replace every instance of $\opA{z \in y_i} \chi\mleft(z\mright)$ by $\opA{z} \left(z \in y_i \rightarrow \chi\mleft(z\mright)\right)$, and every instance of $\opE{z \in y_i} \chi\mleft(z\mright)$ by $\opE{z} \left(z \in y_i \land \chi\mleft(z\mright)\right)$;
    \item replace every instance of $v \in y_i$ by $\psi_i\mleft(v, \vec{u}_i\mright)$.
  \end{enumerate}
  Then $\tilde{\varphi}$ will contain no instances of $y_i$ but the parameters $\vec{u}_i$ instead, and it is easy to verify that
  \[x = \left\{z \in L_\beta : \psi_1^{L_\beta}\mleft(z, \vec{u}_1\mright) \land \tilde{\varphi}^{L_\beta}\mleft(z, \vec{u}_1, \ldots, \vec{u}_n\mright)\right\} \in \dd\mleft(L_\beta\mright) \subseteq L_\alpha. \qedhere\]
\end{proof}

As previously mentioned in \autoref{ftn:int-l-delta-0-sep}, the proof involves finding a maximal $\beta \in \alpha$ which requires classical logic. However, one can clearly see that if we demand $\alpha$ to be linear beforehand, then $\IKP$ suffices for the result. Notably, we can have the following:
\begin{theorem}[$\IKP$]
  \label{thm:ord-cla-emb-equiv-form}
  For any ordinal $\alpha$, the following are equivalent:
  \begin{enumerate}[label=\uplabel{\arabic*}]
    \item \label{itm:ord-cla-emb-equiv-form-l-emb} $\alpha$ embeds decidably into the universe;
    \item \label{itm:ord-cla-emb-equiv-form-ord-emb} $L_\alpha \cap \Ord$ is $\Delta_0$-decidable;
    \item \label{itm:ord-cla-emb-equiv-form-atom-dec} $\alpha$ is linear and for any $\beta, \gamma \in L_\alpha \cap \Ord$, we have
          \[\beta \in \gamma \lor \beta = \gamma \lor \gamma \in \beta,\]
          i.e.\ $L_\alpha \cap \Ord$ is also linear.
  \end{enumerate}
\end{theorem}

\begin{proof}
  Firstly, it is trivial that \ref{itm:ord-cla-emb-equiv-form-l-emb} implies \ref{itm:ord-cla-emb-equiv-form-ord-emb} because $L_\alpha \cap \Ord$ is a $\Delta_0$-definable subset of $L_\alpha$. For \ref{itm:ord-cla-emb-equiv-form-ord-emb} implying \ref{itm:ord-cla-emb-equiv-form-atom-dec}, we just need to make the observation that in the proof of \autoref{lem:dec-emb-def-in-l}, every instance of excluded middle we use only concerns $\Delta_0$-formulae on ordinals and thus \ref{itm:ord-cla-emb-equiv-form-ord-emb} suffices for the argument.

  Finally, assuming \ref{itm:ord-cla-emb-equiv-form-atom-dec}, take any $\Delta_0$-formula $\varphi$ and a corresponding list of sets $y_1, \ldots, y_n \in L_\alpha$. We use the linearity of $\alpha$ to find $\beta \in \alpha$ such that $y_1, \ldots, y_n \in \dd\mleft(L_\beta\mright)$, apply the constructions in the classical \autoref{lem:zf-l-delta-0-separation} and derive that the set
  \[\delta_{\vec{y}} = \left\{x \in \left\{\varnothing\right\} : \varphi\mleft(y_1, \ldots, y_n\mright)\right\} \in \dd\mleft(L_\beta\mright) \subseteq L_\alpha.\]
  Note that since $\varnothing \in \dd\mleft(L_\beta\mright) \subseteq L_\alpha$ trivially, by the linearity assumption we know that $L_\beta \cap \Ord = \varnothing \lor \varnothing \in L_\beta$. The first case is trivial since $\beta$ is linear, so by $\Delta_0$-separation from \autoref{lem:zf-l-delta-0-separation} we know that any $x \in L_\beta$ implies $\varnothing \in L_\beta \cap \Ord$, thus also $L_\beta = \varnothing$. It follows that $y_1 = \cdots = y_n = \varnothing$ and the $\Delta_0$-formula $\varphi\mleft(y_1, \ldots, y_n\mright)$ must satisfy excluded middle. Therefore, we can assume the second disjunct $\varnothing \in L_\beta$ without loss of generality and thus $\left\{\varnothing\right\} \in \dd\mleft(L_\beta\mright)$ here.

  In this case $\delta_{\vec{y}}$ is well-defined, clearly $\delta_{\vec{y}} \subseteq 1$ and thus $\delta_{\vec{y}} \in L_\alpha \cap \Ord$. Our assumption then implies
  \[\delta_{\vec{y}} = \varnothing \lor \varnothing \in \delta_{\vec{y}},\]
  which immediately gives $\varphi\mleft(y_1, \ldots, y_n\mright) \lor \neg \varphi\mleft(y_1, \ldots, y_n\mright)$ as desired.
\end{proof}

\section{The uninterpreted set-recursive model}
\label{sec:e-rec-unint}

In this section, let $E\mleft(X\mright)$ denote the pca of set recursive functions\inds{set recursive function} with an additional $\Sigma_1$-class function $X$, i.e.\ the class-sized pca given by $\Sigma$-inductive definition in \autoref{sec:czf-weak-real}, except that we omit the clauses for powerset $\bar{\wp}$.

Here, any set in the base universe $V$ can be a `function in the pca $E\mleft(X\mright)$', even though most of them --- the ones that do not contain a distinguished operation --- will not converge on any input. We can still construct the class of all names, $V\mleft(E\mleft(X\mright)\mright)$, as an inductive definition $I\mleft(\Phi\mright)$ where $\Phi$ is the $\Delta_0$-class
\[\left\{\tuple{x, y} : x \subseteq V \times y\right\}.\]
Similar to \autoref{cor:czf-in-czf-name-def}, it follows immediately that $x \in V\mleft(E\mleft(X\mright)\mright) = I\mleft(\Phi\mright)$ if and only if $x \subseteq V \times V\mleft(E\mleft(X\mright)\mright)$.

However, since realisers in the pca now range over the entire universe, we observe that $F_\in\mleft(u, v\mright), G_=\mleft(u, v\mright)$ as defined in \autoref{sec:czf-real-in-czf} are no longer sets\footnote{To be exact, the problem is caused by the definition of $G_=$, which expects realisers of $u = v$ to be pairs of functions $f, g$ in the pca. Since the interpretation only cares about values of these functions on $f$ such that $\tuple{d, x} \in u$ or $\tuple{d, x} \in v$, if we take some $d$ such that $\neg\!\opE{x} \tuple{d, x} \in u \cup v$ holds, then for any set $z$, $f' = \oplam{a} \bdN \mleft(\bi_1 1a \mleft(\bm{\pi} dd\mright)\mright) 1z \mleft(fa\mright)$ realises $u = v$ if and only if $f$ does, where $\bi_1 1a \mleft(\bm{\pi} dd\mright)$ evaluates to $1$ if $a \in \left\{d\right\}$, i.e.\ if $a = d$. It follows that $G_=\mleft(u, v\mright)$ cannot be a set as long as it is nonempty.}. Thus, they need to be defined through $\Sigma$-inductive definitions as well.

We shall let $\Psi$ be the $\Sigma_1$-class of tuples $\tuple{\tuple{u, v, d}, y}$ such that $u, v \in V\mleft(E\mleft(X\mright)\mright)$ and there exists $\tuple{f, x} \in v$ such that
\begin{enumerate}[label=\uplabel{\roman*}]
  \item $\bp_0 d {\downarrow} = f$,
  \item for any $\tuple{g, z} \in u$, $\tuple{z, x, \bp_{10} dg {\downarrow}} \in y$,
  \item for any $\tuple{g, z} \in x$, $\tuple{z, u, \bp_{11} dg {\downarrow}} \in y$.
\end{enumerate}
We shall then set $F_\in = I\mleft(\Psi\mright)$ and
\[G_= = \begin{aligned}[t]
     & {\left\{\tuple{u, v, d} : u, v \in V\mleft(E\mleft(X\mright)\mright), \opA{{\tuple{f, x} \in u}} \tuple{x, v, \bp_0 df {\downarrow}} \in F_\in \right.}, \\
     & \qquad \qquad \qquad \left.{} \opA{{\tuple{f, x} \in v}} \tuple{x, u, \bp_1 df {\downarrow}} \in F_\in\right\}.
  \end{aligned}\]
It follows from \autoref{cor:set-ind-def-exist-pre} that $F_\in, G_= \subseteq V\mleft(E\mleft(X\mright)\mright)^2 \times E\mleft(X\mright)$, with also
\[F_\in = \left\{\tuple{u, v, d} : u, v \in V\mleft(E\mleft(X\mright)\mright), \opE{{\tuple{f, x} \in v}} \left(\bp_0 d {\downarrow} = f \land \tuple{u, x, \bp_1 d {\downarrow}} \in G_=\right)\right\}.\]

We define the atomic realisability conditions as
\begin{align*}
  d \Vdash u \in v & \defeqv \tuple{u, v, d} \in F_\in, \\
  d \Vdash u = v   & \defeqv \tuple{u, v, d} \in G_=,
\end{align*}
and observe that the result of \autoref{sec:czf-real-in-czf} will follow without any change, again except \autoref{lem:czf-sub-coll-real} for the subset collection scheme. In other words, we have that:
\begin{theorem}
  Let $\varphi$ be any theorem of $\CZF^-$, then $\CZF^-$ proves that $V\mleft(E\mleft(X\mright)\mright) \vDash \varphi$, i.e.\ there exists $d \in E\mleft(X\mright)$ such that $d \Vdash \varphi$.
\end{theorem}

Here $V\mleft(E\mleft(X\mright)\mright)$ still contains incomparable ordinals that we constructed in \autoref{sec:ord-modal}, so it is not a model of $\REM$. However, we can construct a class function of `standard names' from $V$ to $V\mleft(E\mleft(X\mright)\mright)$ recursively as
\[\check{x} = \left\{\tuple{y, \check{y}} : y \in x\right\}.\]

We first observe that:
\begin{lemma}
  \label{lem:e-rec-real-check-name-true}
  For any set $x, y$, if some $d \Vdash \check{x} \in \check{y}$, then $x \in y$; if some $d \Vdash \check{x} = \check{y}$, then $x = y$.
\end{lemma}

\begin{proof}
  We prove the following three statements
  \begin{align*}
    \left(\opE{d} d \Vdash \check{x} \in \check{y}\right) & \rightarrow x \in y, \\
    \left(\opE{d} d \Vdash \check{y} \in \check{x}\right) & \rightarrow y \in x, \\
    \left(\opE{d} d \Vdash \check{x} = \check{y}\right)   & \rightarrow x = y
  \end{align*}
  together by simultaneous set induction on $x$ and $y$.

  Suppose that $d \Vdash \check{x} \in \check{y}$, then there is $z = \bp_0 d {\downarrow} \in y$ such that $\bp_1 d {\downarrow} \Vdash \check{x} = \check{z}$. By the inductive hypothesis, we have $x = z \in y$ as desired. The case where $d \Vdash \check{y} \in \check{x}$ is symmetrical.

  Now, suppose that $d \Vdash \check{x} = \check{y}$. For any $z \in x$, we have $\tuple{z, \check{z}} \in \check{x}$, so $\mleft(\bp_0 d\mright) z \Vdash \check{z} \in \check{y}$ and by the first clause of inductive hypothesis, we have $z \in y$. Similarly, for any $z \in y$, we have $z \in x$ by the second clause of inductive hypothesis. Thus, $x = y$.
\end{proof}

\begin{corollary}
  \label{cor:e-rec-real-check-ord-true}
  For any set $x$, $x \in \Ord$ if and only if $V\mleft(E\mleft(X\mright)\mright) \vDash \check{x} \in \Ord$.
\end{corollary}

\begin{proof}
  First notice for any transitive set $x$, $\oplam{yz} \bp z \ior \Vdash \text{$\check{x}$ is transitive}$. Conversely, if $V\mleft(E\mleft(X\mright)\mright) \vDash \text{$\check{x}$ is transitive}$, then for any $y \in x$, $z \in y$, we know that $\check{z} \in \check{x}$ is realised and thus $x$ is transitive by \autoref{lem:e-rec-real-check-name-true}.

  It follows then from the definition of $\check{x}$ that $x$ is a set of transitive sets if and only if $V\mleft(E\mleft(X\mright)\mright)$ realises that $\check{x}$ is a set of transitive sets.
\end{proof}

\begin{lemma}[$\CZF^- + \REM$]
  \label{lem:e-rec-real-check-ord-lin}
  There is a uniform realiser $d \in E\mleft(X\mright)$ such that for every $\alpha \in \Ord$,
  \[d \Vdash \opA{\beta, \gamma \in \check{\alpha}} \left(\beta \in \gamma \lor \left(\beta = \gamma \lor \gamma \in \beta\right)\right),\]
  i.e.\ every $\check{\alpha}$ is linear in $V\mleft(E\mleft(X\mright)\mright)$.
\end{lemma}

\begin{proof}
  For any $\beta, \gamma \in \alpha$, observe simply from the definition that if $\beta \in \gamma$, then $\bp \beta \ior \Vdash \check{\beta} \in \check{\gamma}$; if $\beta = \gamma$, then $\ior \Vdash \check{\beta} = \check{\gamma}$. Thus, we just need to take $\beta, \gamma \in \alpha$ and compute whether one is an element of the other or that they are equal, and
  \[d = \oplam{\beta \gamma} \bdN \mleft(\bi_1 1 \beta \gamma\mright) 1 \mleft(\bp 0 \mleft(\bp \beta \ior\mright)\mright) \mleft(\bdN \mleft(\bi_1 1 \gamma \beta\mright) 0 \mleft(\bp 1 \mleft(\bp 0 \ior\mright)\mright) \mleft(\bp 1 \mleft(\bp 1 \mleft(\bp \gamma \ior\mright)\mright)\mright)\mright)\]
  realises the desired statement.
\end{proof}

We would like to show that every $\check{\alpha}$ embeds decidably into $V\mleft(E\mleft(X\mright)\mright)$. For this purpose, we will introduce a new primitive class symbol $\St$ in the language of $V\mleft(E\mleft(X\mright)\mright)$, with the following realisability condition:
\[d \Vdash x \in \St \defeqv \opE{v} \left(\bp_0 d {\downarrow} = v \land \bp_1 d {\downarrow} \Vdash x = \check{v}\right).\]
Conveniently, we can also introduce the quantifiers
\begin{align*}
  d \Vdash \opA{x \in \St} \varphi & \defeqv \opA{v} dv {\downarrow} \Vdash \varphi\mleft[\check{v}/x\mright],                                                  \\
  d \Vdash \opE{x \in \St} \varphi & \defeqv \opE{v} \left(\bp_0 d {\downarrow} = v \land \bp_1 d {\downarrow} \Vdash \varphi\mleft[\check{v}/x\mright]\right).
\end{align*}
and check that they have the same relation to the unbounded quantifiers:
\begin{lemma}
  The following are realised in $V\mleft(E\mleft(X\mright)\mright)$:
  \begin{align*}
    \opA{x} \left(x \in \St \rightarrow \varphi\mleft(x\mright)\right) & \leftrightarrow \opA{x \in \St} \varphi\mleft(x\mright), \\
    \opE{x} \left(x \in \St \land \varphi\mleft(x\mright)\right)       & \leftrightarrow \opE{x \in \St} \varphi\mleft(x\mright).
  \end{align*}
\end{lemma}

\begin{proof}
  This is completely analogous to \autoref{lem:set-bnd-quant-axioms-real}.
\end{proof}

We want $V\mleft(E\mleft(X\mright)\mright)$ to think that $\St$ is a transitive model of $\CZF^-$. For this purpose, we will need that our pca $E\mleft(X\mright)$ computes some choice principles. We observe that $\CZF^-$ is equiconsistent with $\CZF^- + V = L$ since its $L$ does satisfy all axioms of $\CZF^-$. Thus we can conveniently move to the base theory $\CZF^- + \REM + V = L$ and observe that $\CZF^- + \REM$ suffices to prove the classical theorem that there exists a $\Sigma_1$-class bijection between $\Ord$ and $L$ (cf.\ Devlin \cite{devlin84-constructibility}*{section II.2}), so that each $L_\alpha$ is mapped to an initial segment of $\Ord$. We shall take $X : V \rightarrow \Ord$ as this bijection.

The transitivity of $\St$ is trivial:
\begin{lemma}
  \label{lem:e-rec-st-transitive}
  $V\mleft(E\mleft(X\mright)\mright)$ realises $\opA{x \in \St} \opA{y \in x} y \in \St$.
\end{lemma}

\begin{proof}
  This is simply realised by $\oplam{xy} \bp y \ior$.
\end{proof}

Now, both extensionality and set induction in $\St$ are immediate consequences of \autoref{lem:e-rec-st-transitive}. For the pairing axiom:
\begin{lemma}
  The pairing axiom in $\St$ is realised in $V\mleft(E\mleft(X\mright)\mright)$:
  \[\opA{a, b \in \St} \opE{y \in \St} \opA{x \in \St} \left(x \in y \leftrightarrow x = a \lor x = b\right).\]
\end{lemma}

\begin{proof}
  For any set $a, b$, observe that if $y = \left\{a, b\right\}$, then
  \[\check{y} = \left\{\tuple{a, \check{a}}, \tuple{b, \check{b}}\right\}.\]
  For any $x$, $\oplam{d} \bp \mleft(\bdN \mleft(\bp_0 d\mright) 0ab\mright) \ior$ realises
  \[\check{x} = \check{a} \lor \check{x} = \check{b} \rightarrow \check{x} \in \check{y},\]
  since if $\check{x} = \check{a} \lor \check{x} = \check{b}$ is realised, then $x = a$ or $b$ by \autoref{lem:e-rec-real-check-name-true}; whereas $\oplam{d} \bp \mleft(\bi_1 1x \mleft(\bm{\pi} bb\mright)\mright) \ior$ realises
  \[\check{x} \in \check{y} \rightarrow \left(\check{x} = \check{a} \lor \check{x} = \check{b}\right),\]
  where we check whether $x$ is $a$ or $b$ by testing $x \in \left\{b\right\}$.
\end{proof}

For the infinity axiom, it suffices to observe that $V\mleft(E\mleft(X\mright)\mright) \vDash \omega \in \St$, which is the case because the name $\overline{\omega}$ we constructed in \autoref{sec:czf-real-in-czf} is precisely $\check{\omega}$.
$\Ind\mleft(\omega\mright)$ is then a $\Delta_0$-statement and holds in the transitive class $\St$.

For union and strong collection, we would need the following choice function construction:
\begin{lemma}[$\CZF^- + \REM$]
  \label{lem:e-rec-choice-func}
  When $X : V \rightarrow \Ord$ is a $\Sigma_1$-class bijection, then there is a term $\delta \in E\mleft(X\mright)$ such that for any nonempty set $s$, we have
  \[\delta s {\downarrow} \in s.\]
\end{lemma}

\begin{proof}
  We know that there must exist a single $x \in s$ such that $X\mleft(x\mright)$ is the least element inside $\left\{X\mleft(y\mright) : y \in s\right\}$. We can find this $x$ by first letting $\tau = \oplam{sx} \bdN \mleft(\bm{\nu} \mleft(\bm{\rho} \mleft(\oplam{y} \bi_1 1 \mleft(\bx y\mright) \mleft(\bx x\mright)\mright) s\mright)\mright) 1 \varnothing \mleft(\bm{\pi} xx\mright)$, so that
  \[\tau sx {\downarrow} = \left\{\begin{aligned}
       & \varnothing      &  & \text{if there is $y \in s$ such that $X\mleft(y\mright) \in X\mleft(x\mright)$}, \\
       & \left\{x\right\} &  & \text{otherwise}
    \end{aligned}\right.\]
  for any $x \in s$. Then $\delta = \oplam{s} \bm{\nu} \mleft(\bm{\rho} \mleft(\oplam{x} \tau sx\mright) s\mright)$ is what we need.
\end{proof}

We then have:
\begin{lemma}
  The union axiom in $\St$ is realised in $V\mleft(E\mleft(X\mright)\mright)$:
  \[\opA{a \in \St} \opE{b \in \St} \opA{x \in \St} \left(x \in b \leftrightarrow \opE{u \in a} x \in u\right).\]
\end{lemma}

\begin{proof}
  For any set $a$, let $b = \bigcup a$. We observe that for any $x$, $\oplam{d} \bp x \ior$ realises
  \[\opE{u \in \check{a}} \check{x} \in u \rightarrow \check{x} \in \check{b}\]
  trivially by \autoref{lem:e-rec-real-check-name-true}. On the other hand, observe that
  \[\bi_3 aax = \left\{u \in a : x \in u\right\}\]
  Using the choice function $\delta$, $\oplam{d} \bp \mleft(\delta \mleft(\bi_3 aax\mright)\mright) \mleft(\bp x \ior\mright)$ realises
  \[\check{x} \in \check{b} \rightarrow \opE{u \in \check{a}} \check{x} \in u. \qedhere\]
\end{proof}

\begin{lemma}
  \label{lem:e-rec-st-str-coll}
  The strong collection scheme in $\St$ is realised in $V\mleft(E\mleft(X\mright)\mright)$: for any $a \in \St$,
  \begin{align*}
     & \opA{x \in a} \opE{y \in \St} \theta^\St\mleft(x, y\mright)                                                                                                                  \\
     & \quad {} \rightarrow \opE{b \in \St} \left(\opA{x \in a} \opE{y \in b} \theta^\St\mleft(x, y\mright) \land \opA{y \in b} \opE{x \in a} \theta^\St\mleft(x, y\mright)\right).
  \end{align*}
\end{lemma}

\begin{proof}
  For any set $a$, suppose that $f \Vdash \opA{x \in \check{a}} \opE{y \in \St} \theta^\St\mleft(x, y\mright)$, and let $b = \bm{\rho} \mleft(\oplam{x} \bp_0 \mleft(fx\mright)\mright) a$. It is easy to check that
  \[f \Vdash \opA{x \in \check{a}} \opE{y \in \check{b}} \theta^\St\mleft(x, y\mright).\]
  On the other hand, let $g = \oplam{ay} \bm{\nu} \mleft(\bm{\rho} \mleft(\oplam{x} \bi_1 \mleft(\bm{\pi} xx\mright) \mleft(\bp_0 \mleft(fx\mright)\mright) \mleft(\bm{\pi} yy\mright)\mright) a\mright)$, so that
  \[gay {\downarrow} = \left\{x \in a : \bp_0 \mleft(fx\mright) = y\right\},\]
  then $\oplam{y} \bp \mleft(\delta \mleft(gay\mright)\mright) \mleft(\bp_1 \mleft(f \mleft(\delta \mleft(gay\mright)\mright)\mright)\mright)$ realises
  \[\opA{y \in \check{b}} \opE{x \in \check{a}} \theta^\St\mleft(x, y\mright). \qedhere\]
\end{proof}

Finally, the binary intersection axiom is trivial:
\begin{lemma}
  The binary intersection axiom in $\St$ is realised in $V\mleft(E\mleft(X\mright)\mright)$:
  \[\opA{a, b \in \St} \opE{y \in \St} \opA{x \in \St} \left(x \in y \leftrightarrow x \in a \land x \in b\right).\]
\end{lemma}

\begin{proof}
  For any set $a, b$, take $y = a \cap b$. For any $x$, simply by \autoref{lem:e-rec-real-check-name-true}, $\oplam{d} \bp \mleft(\bp x \ior\mright) \mleft(\bp x \ior\mright)$ realises
  \[\check{x} \in \check{y} \rightarrow \check{x} \in \check{a} \land \check{x} \in \check{b},\]
  whereas $\oplam{d} \bp x \ior$ realises
  \[\check{x} \in \check{a} \land \check{x} \in \check{b} \rightarrow \check{x} \in \check{y}. \qedhere\]
\end{proof}

Therefore, we have verified that:
\begin{theorem}
  \label{thm:e-rec-st-model-czf}
  Let $\varphi$ be any theorem of $\CZF^-$, then $\CZF^- + \REM + V = L$ proves $V\mleft(E\mleft(X\mright)\mright) \vDash \varphi^\St$, where $X : V \rightarrow \Ord$ is a $\Sigma_1$-class bijection.
\end{theorem}

\begin{corollary}[$\CZF^- + \REM + V = L$]
  For any $\alpha \in \Ord$, $\check{\alpha}$ embeds decidably into $V\mleft(E\mleft(X\mright)\mright)$.
\end{corollary}

\begin{proof}
  For every $\alpha \in \Ord$, $V\mleft(E\mleft(X\mright)\mright)$ realises that $\check{\alpha}$ is an ordinal by \autoref{cor:e-rec-real-check-ord-true}, and realises that $\check{\alpha}$ is linear by \autoref{lem:e-rec-real-check-ord-lin}. Thus by \autoref{thm:ord-cla-emb-equiv-form}, it suffices to work inside $V\mleft(E\mleft(X\mright)\mright)$ and show that $L_{\check{\alpha}} \cap \Ord$ is also linear.

  Now, we use the fact that $V\mleft(E\mleft(X\mright)\mright) \vDash \left(\CZF^-\right)^\St$, thus there is a well-defined set $L_{\check{\alpha}}^\St$ in the class $\St$. Since G\"odel's $L$ is a $\Sigma_1$-definable hierarchy and thus upwards absolute across transitive classes, it follows that $L_{\check{\alpha}} = L_{\check{\alpha}}^\St$. Then
  \[L_{\check{\alpha}} \cap \Ord = L_{\check{\alpha}}^\St \cap \Ord \in \St\]
  is an ordinal in $\St$, and by \autoref{cor:e-rec-real-check-ord-true}, $L_{\check{\alpha}} \cap \Ord = \check{\beta}$ for some $\beta \in \Ord$ in the base theory. Therefore, by \autoref{lem:e-rec-real-check-ord-lin} again we know that $L_{\check{\alpha}} \cap \Ord$ is linear.
\end{proof}

However, the system $V\mleft(E\mleft(X\mright)\mright)$ also has its problems. Firstly, while we have \autoref{cor:dec-emb-sub-strong-incom-code} for arbitrarily large ordinals $\check{\alpha}$, one still cannot construct incomparable coding arguments in the form of \autoref{sec:czf-v-eq-l} here. This is because every set-theoretic function, when supplied as a parameter, is computable in the pca $E\mleft(X\mright)$. Consequently:
\begin{lemma}[$\CZF^- + \REM + V = L$]
  \label{lem:e-rec-real-inv-name-eq}
  For any set $x$ and name $y \in V\mleft(E\mleft(X\mright)\mright)$, suppose that for any $z \in x$ there exists $f$ such that $\tuple{f, \check{z}} \in y$, then $V\mleft(E\mleft(X\mright)\mright) \vDash \check{x} \subseteq y$.

  It follows that if the name $y$ can be seen as a surjective set-theoretic function onto $\left\{\check{z} : z \in x\right\}$, then $V\mleft(E\mleft(X\mright)\mright) \vDash \check{x} = y$.
\end{lemma}

\begin{proof}
  By $\REM$ and $V = L$ together, we have a global axiom of choice and can construct a set-theoretic function $h$ on domain $x$ such that for every $z \in x$, $\tuple{h\mleft(z\mright), \check{z}} \in y$. Since $E\mleft(X\mright)$ can evaluate set-theoretic functions (cf.\ \autoref{ftn:ep-rec-set-func} on page \pageref{ftn:ep-rec-set-func}), for every $z \in x$ it suffices to see that $\bp g \ior \Vdash \check{z} \in y$ where $g = h\mleft(z\mright)$. Then $V\mleft(E\mleft(X\mright)\mright) \vDash \check{x} \subseteq y$.

  Conversely, if $y$ is a surjective function onto $\left\{\check{z} : z \in x\right\}$, then we can easily build
  \[y' = \left\{\tuple{f, z} : \tuple{f, \check{z}} \in y\right\},\]
  by observing that $z$ can be obtained back as $z = \left\{y : \tuple{y, \check{y}} \in \check{z}\right\}$. For every $\tuple{f, \check{z}} \in y$, then $\bp z \ior \Vdash \check{z} \in \check{x}$, where $z = y'\mleft(f\mright)$ can be computed from the index $f$. It follows that $V\mleft(E\mleft(X\mright)\mright) \vDash y \subseteq \check{x}$ and thus $\check{x} = y$.
\end{proof}

By \autoref{lem:e-rec-real-inv-name-eq}, imitating the constructions of names in $\Pow\mleft(\omega\mright) \cap \Ord$ in \autoref{sec:czf-v-eq-l} would only lead to names in $V\mleft(E\mleft(X\mright)\mright)$ that are (realised to be) equal to $\check{\omega}$. Thus we cannot construct ordinal names $\beta \in V\mleft(E\mleft(X\mright)\mright)$ that are realised as incomparable, proper subsets of $\check{\alpha}$ easily as we did in \autoref{sec:czf-v-eq-l}.

Secondly, we also point out $V\mleft(E\mleft(X\mright)\mright)$ is one construction that cannot accommodate the powerset axiom --- even when the base theory does admit the powerset axiom --- and this limits its adaptability to the extensions of $\CZF^-$. Specifically, we can prove the following:
\begin{theorem}[$\CZF^-$]
  $V\mleft(E\mleft(X\mright)\mright)$ realises the negation of the powerset axiom. Namely, it realises that
  \[\neg\!\opE{y} \opA{x} \left(x \in y \leftrightarrow x \subseteq 1\right).\]
\end{theorem}

\begin{proof}
  It suffices to derive a contradiction from the assumption that there is $y \in V\mleft(E\mleft(X\mright)\mright)$ with some $f \in E\mleft(X\mright)$ realising
  \[x \in y \leftrightarrow x \subseteq \overline{1}\]
  for all $x \in V\mleft(E\mleft(X\mright)\mright)$, where $\overline{1} = \left\{\tuple{\varnothing, \varnothing}\right\}$ is as given in \autoref{sec:czf-real-in-czf}.

  Let $g = \oplam{e} \bp \varnothing \ior$ such that for any set $a$,
  \[g \Vdash \left\{\tuple{a, \varnothing}\right\} \subseteq \overline{1}.\]
  It follows that there must exist some $\tuple{d, z} \in y$ such that
  \[\bp_{11} \mleft(\bp_1 fg\mright) \Vdash z \subseteq \left\{\tuple{a, \varnothing}\right\},\]
  where $z$ is nonempty because we also realise $\left\{\tuple{a, \varnothing}\right\} \subseteq z$. We now consider
  \[s = \left\{\bp_0 \mleft(\bp_{11} \mleft(\bp_1 fg\mright) e\mright) : \tuple{d, z} \in y, \tuple{e, u} \in z\right\},\]
  where we know that the computation must converge and thus $s$ exists by strong collection. Then also
  \[\bp_{11} \mleft(\bp_1 fg\mright) \Vdash z \subseteq \left\{\tuple{s, \varnothing}\right\},\]
  for some $\tuple{d, z} \in y$ where $z$ is nonempty, i.e.\ for every $\tuple{e, u} \in z$, we have $\bp_0 \mleft(\bp_{11} \mleft(\bp_1 fg\mright) e\mright) {\downarrow} = s$. However, by our definition of $s$ above, we also have
  \[\bp_0 \mleft(\bp_{11} \mleft(\bp_1 fg\mright) e\mright) \in s,\]
  a contradiction.
\end{proof}

\section{Realisability over \texorpdfstring{$\alpha$}{\textalpha}-recursion}

We will suggest that the problems of $V\mleft(E\mleft(X\mright)\mright)$ can be resolved by considering instead a large set-sized pca, for example the (weak) pca of $\alpha$-recursive functions for an admissible ordinal $\alpha$.

Following Sacks \cite{sacks90-higher-recursion}*{Chapter VII}, we say that an ordinal $\alpha$ is \emph{$\Sigma_1$-admissible} (or just \emph{admissible}) if $L_\alpha$ satisfies the $\Delta_0$-collection axiom. Working in the base theory $\CZF^- + \REM$, we observe that $\REM$ suffices to ensure that $\tuple{L_\alpha; \in}$ satisfies full excluded middle, whereas \autoref{lem:zf-l-delta-0-separation} applies to justify that $L_\alpha$ also satisfies $\Delta_0$-separation. Thus, it is easy to see that any $\alpha \in \Ord$ is $\Sigma_1$-admissible if and only if $L_\alpha$ is a model of classical $\KP - \text{Infinity}$.

An $\alpha$-recursive function is essentially a partial function on $L_\alpha$ that is $\Sigma_1$-definable in the structure $\tuple{L_\alpha; \in}$. Formally, we want to fix a universal $\Sigma_1$-formula $\pi\mleft(e, m, a, b\mright)$ in the language of set theory, such that for any $\Sigma_1$-formula $\varphi\mleft(m, a, b\mright)$ with all free variables indicated, $\KP$ (even without the axiom of infinity) proves that there exists a set $e$ such that
\[\opA{m} \opA{a} \opA{b} \left(\varphi\mleft(m, a, b\mright) \leftrightarrow \pi\mleft(e, m, a, b\mright)\right).\]
One could follow the construction in \cite{barwise75-admissible-sets}*{section V.1} or simply construct a truth predicate for $\Delta_0$-formulae by $\Sigma$-recursion, using an approach similar to the predicate $\Tr_{\left\ulcorner \varphi \right\urcorner}$ in \autoref{sec:intro-int-l}.

We then perform the same stipulations as we did in \autoref{sec:oa-cm} for $\Sigma^1_1$-definable functions in second-order arithmetic. We let $\Sigma L_\alpha$ denote the domain $L_\alpha$ with a partial application function given as $ea {\downarrow} = b$ if and only if $b \in L_\alpha$ is the \emph{unique} element such that
\[\pi^{L_\alpha}\mleft(e_0, e_1, a, b\mright)\]
holds, where $e = \tuple{e_0, e_1}$ is the set-theoretic, Kuratowski ordered pair. We observe that:
\begin{lemma}
  \label{lem:sigma-l1-smn}
  Given any $\Sigma_1$-formula $\varphi\mleft(m, a, b\mright)$ with all free variables indicated, then there exists $e \in L_\alpha$ such that for any $m, a \in L_\alpha$, $em {\downarrow}$ always holds, and $ema {\downarrow} = b$ if and only if $b \in L_\alpha$ is the unique element such that $\varphi^{L_\alpha}\mleft(m, a, b\mright)$ holds.
\end{lemma}

\begin{proof}
  Just like \autoref{lem:sigma11-smn}, pick $f, n \in L_\alpha$ by universality such that
  \begin{gather*}
    \opA{m} \opA{a} \opA{b} \left(b = \tuple{m, a} \leftrightarrow \pi\mleft(f, m, a, b\mright)\right), \\
    \opA{m} \opA{a} \opA{b} \left(\varphi\mleft(m, a, b\mright) \leftrightarrow \pi\mleft(n, m, a, b\mright)\right),
  \end{gather*}
  where the set-theoretic ordered pair is expressed using its $\Delta_0$ defining formula. Let $e = \tuple{f, n}$, then $em {\downarrow} = \tuple{n, m}$ for any $m \in L_\alpha$ and thus we have $ema {\downarrow} = b$ precisely when $b \in L_\alpha$ is the unique element such that $\varphi^{L_\alpha}\mleft(m, a, b\mright)$ holds.
\end{proof}

\begin{corollary}
  \label{cor:sigma-l1-pca}
  For any admissible ordinal $\alpha$, $\Sigma L_\alpha$ is a weak pca over natural numbers.
\end{corollary}

\begin{proof}
  First observe that any $\alpha \in \omega$ is not admissible. Thus by $\REM$, we know that $\omega \subseteq \alpha$ and hence $\omega \subseteq L_\alpha$ as well.

  We now use \autoref{lem:sigma-l1-smn} to identify the distinguished elements in the algebra. For operations that involve one parameter only, we simply ignore the parameter $m$ in the proposition:
  \begin{enumerate}[label=\uplabel{\roman*}]
    \item let $s \in L_\alpha$ be the index such that for any $m, a \in L_\alpha$, $sma {\downarrow} = b$ if and only if $b = a^+ = a \cup \left\{a\right\}$, and set $\bsN = s \varnothing$;
    \item let $p \in L_\alpha$ be the index such that for any $m, a \in L_\alpha$, $pma {\downarrow} = b$ if and only if $a = b^+$, and set $\bpN = p \varnothing$;
  \end{enumerate}

  For $\bk$ which involves two parameters, we use precisely the parameters in \autoref{lem:sigma-l1-smn}:
  \begin{enumerate}[label=\uplabel{\roman*},resume]
    \item let $\bk \in L_\alpha$ simply be the index such that for any $a, b \in L_\alpha$, $\bk ab {\downarrow} = c$ if and only if $c = a$.
  \end{enumerate}

  For $\bs$ which involves three parameters, we invoke the proposition twice:
  \begin{enumerate}[label=\uplabel{\roman*},resume]
    \item let $s' \in L_\alpha$ be the index such that for any $a, b \in L_\alpha$, $s' ab {\downarrow} = c$ if and only if $c \in L_\alpha$ is the unique element such that there exist $d, e \in L_\alpha$ and it is possible to decompose
          \[a = \tuple{\tuple{p_0, p_1}, \tuple{q_0, q_1}}, \qquad d = \tuple{d_0, d_1},\]
          such that the following holds\footnote{Just as in \autoref{sec:oa-cm}, the existence of a unique $c$ here does not also guarantee the uniqueness of the intermediate values $d$ and $e$. Thus the existence of a unique $c$ does not imply $a_0 b \mleft(a_1 b\mright) {\downarrow}$ and thus $\Sigma L_\alpha$ as constructed here is only a weak pca.}:
          \[\pi^{L_\alpha}\mleft(a_{00}, a_{01}, b, d\mright) \land \pi^{L_\alpha}\mleft(a_{10}, a_{11}, b, e\mright) \land \pi^{L_\alpha}\mleft(d_0, d_1, e, c\mright).\]
          In other words, if $\tuple{p_0, p_1} b {\downarrow} = d$, $\tuple{q_0, q_1} b {\downarrow} = e$ and $de {\downarrow} = c$, then also $s' ab {\downarrow} = c$. Now just let $\bs \in L_\alpha$ be such that $\bs ab {\downarrow} = s' \tuple{a, b}$.
  \end{enumerate}

  Finally, similarly for $\bdN$ which involves four parameters:
  \begin{enumerate}[label=\uplabel{\roman*},resume]
    \item let $d' \in L_\alpha$ be the index such that for any $a, b \in L_\alpha$, if we can decompose $a = \tuple{\tuple{p, q}, r}$, then
          \[d' ab {\downarrow} = \left\{\begin{aligned}
               & r &  & \text{if $p = q$},    \\
               & b &  & \text{if $p \neq q$}.
            \end{aligned}\right.\]
          Then let $d'' \in L_\alpha$ be such that $d'' ab {\downarrow} = d' \tuple{a, b}$ and finally $\bdN \in L_\alpha$ be such that $\bdN ab {\downarrow} = d'' \tuple{a, b}$. \qedhere
  \end{enumerate}
\end{proof}

We will also stipulate that the pairing function $\bp \in \Sigma L_\alpha$ computes precisely the Kuratowski ordered pair, i.e.\ $\bp ab {\downarrow} = \tuple{a, b}$.

We can work with precisely the realisability model $V\mleft(\Sigma L_\alpha\mright)$ as constructed in \autoref{sec:czf-real-in-czf}. Just as in \autoref{sec:e-rec-unint}, we define standard names $\check{x}$ for any $x \in L_\alpha$ recursively as
\[\check{x} = \left\{\tuple{y, \check{y}} : y \in x\right\}.\]
But now we have sets of `all standard names' or `all standard ordinals' as well, given as
\begin{align*}
  \check{L}_\alpha & = \left\{\tuple{x, \check{x}} : x \in L_\alpha\right\} \in V\mleft(\Sigma L_\alpha\mright),           \\
  \check{\alpha}   & = \left\{\tuple{\beta, \check{\beta}} : \beta \in \alpha\right\} \in V\mleft(\Sigma L_\alpha\mright).
\end{align*}

When $\alpha > \omega$ is admissible, i.e.\ $L_\alpha$ is a model of $\KP$, we want to repeat the process in \autoref{sec:e-rec-unint} and check that $\check{\alpha}$ embeds decidably into $V\mleft(\Sigma L_\alpha\mright)$. We shall first note that $\KP$ suffices to define all the set-recursive functions, i.e.\ $E\mleft(\varnothing\mright)$, as a $\Sigma_1$-class. It follows from \autoref{lem:sigma-l1-smn} that every specific function computable by $E\mleft(\varnothing\mright)$ inside $L_\alpha$ is also computable by some $\alpha$-recursive $e \in \Sigma L_\alpha$. Thus we can have most results from \autoref{sec:e-rec-unint} by essentially the same proof. We have that:
\begin{lemma}
  \label{lem:alpha-rec-real-check-name-true}
  For any $x, y \in L_\alpha$, if some $d \Vdash \check{x} \in \check{y}$, then $x \in y$; if some $d \Vdash \check{x} = \check{y}$, then $x = y$.
\end{lemma}

\begin{corollary}
  \label{cor:alpha-rec-real-check-ord-true}
  For any $x \in L_\alpha$, $x \in \Ord$ if and only if $V\mleft(\Sigma L_\alpha\mright) \vDash \check{x} \in \Ord$.
\end{corollary}

\begin{lemma}
  \label{lem:alpha-rec-real-check-ord-lin}
  $V\mleft(\Sigma L_\alpha\mright)$ realises that $\check{\alpha}$ is linear.
\end{lemma}

\begin{proof}
  Just take the construction from \autoref{lem:e-rec-real-check-ord-lin}. Note the uniformity in \autoref{lem:e-rec-real-check-ord-lin} corresponds precisely to the fact here that the realiser does not need to take $\alpha$ itself as a parameter (which it cannot, since $\alpha \not\in L_\alpha$).
\end{proof}

Now, the analogue of $\St$ here is simply the set $\check{L}_\alpha$, and $V\mleft(\Sigma L_\alpha\mright)$ should realise that $\tuple{\check{L}_\alpha; \in}$ is a model of $\CZF^-$. This would again require a choice principle, but we should observe that the following is provable in $\CZF^- + \REM$:
\begin{proposition}[\cite{sacks90-higher-recursion}*{Theorem VII.1.8}]
  \label{prop:alpha-rec-bij-l-alpha}
  When $\alpha$ is any admissible ordinal, there exists an $\alpha$-recursive bijection from $\alpha$ onto $L_\alpha$.
\end{proposition}

\begin{proof}[Proof sketch]
  Using the fact that $L_\alpha$ is a model of $\KP - \text{Infinity}$, then the claim is essentially the existence of a $\Sigma_1$-class bijection between $\Ord$ and $L$ in $\KP - \text{Infinity}$, which we already discussed in \autoref{sec:e-rec-unint}.
\end{proof}

\begin{corollary}
  There is a term $\delta \in \Sigma L_\alpha$ such that for any nonempty $s \in L_\alpha$, we have
  \[\delta s {\downarrow} \in s.\]
\end{corollary}

\begin{proof}
  This is completely analogous to \autoref{lem:e-rec-choice-func}.
\end{proof}

We can then verify that all axioms of $\CZF^-$ hold for $\check{L}_\alpha$ in the same way as \autoref{sec:e-rec-unint}, where we should specifically observe that the proof for the strong collection scheme in $\St$ (correspondingly $\check{L}_\alpha$ here), \autoref{lem:e-rec-st-str-coll}, does not require\footnote{To be more precise, the fact that we can compute
  \[\bm{\rho} xy = \left\{xu : u \in y\right\}\]
  invokes an instance of $\Sigma_1$-collection, which is provided in any model of $\IKP$ or $\KP$, such as $L_\alpha$ here.} collection from the base theory, even though the fact that the entire uninterpreted model ($V\mleft(E\mleft(X\mright)\mright)$ or $V\mleft(\Sigma L_\alpha\mright)$) satisfies a collection scheme would involve strong collection in the base theory. Therefore, the constructions in \autoref{sec:e-rec-unint} can be replicated under just the assumption that $L_\alpha$ satisfies $\KP$ and we have:
\begin{theorem}
  \label{thm:alpha-rec-real-l-alpha-czf}
  For admissible ordinal $\alpha > \omega$, $V\mleft(\Sigma L_\alpha\mright)$ sees that $\check{L}_\alpha$ is a transitive set; and let $\varphi$ be any theorem of $\CZF^-$, then $\CZF^- + \REM$ proves $V\mleft(\Sigma L_\alpha\mright) \vDash \varphi^{\check{L}_\alpha}$.
\end{theorem}

\begin{corollary}[$\CZF^- + \REM$]
  For any admissible ordinal $\alpha$, $\check{\alpha}$ embeds decidably into $V\mleft(\Sigma L_\alpha\mright)$.
\end{corollary}

\begin{proof}
  The case $\alpha = \omega$ is already covered by \autoref{prop:vkl-omega-suc-not-cl}. For $\alpha > \omega$, $V\mleft(\Sigma L_\alpha\mright)$ realises that $\check{\alpha}$ is linear by \autoref{lem:alpha-rec-real-check-ord-lin}. By \autoref{thm:ord-cla-emb-equiv-form}, it suffices to work inside $V\mleft(\Sigma L_\alpha\mright)$ and show that $L_{\check{\alpha}} \cap \Ord$ is also linear.

  We take any ordinal $\gamma \in L_{\check{\alpha}}$, then there exists $\beta \in \check{\alpha}$ such that $\gamma \in \dd\mleft(L_\beta\mright)$. However, using the fact that $V\mleft(\Sigma L_\alpha\mright) \vDash \left(\CZF^-\right)^{\check{L}_\alpha}$ (and that trivially $\check{\alpha} \subseteq \check{L}_\alpha$), we know there exists some well-defined $L_{\beta^+}^{\check{L}_\alpha} \in \check{L}_\alpha$, and
  \[\dd\mleft(L_\beta\mright) = L_{\beta^+} = L_{\beta^+}^{\check{L}_\alpha}\]
  by upward absoluteness. It follows that $\gamma \in \check{L}_\alpha$ and thus $V\mleft(\Sigma L_\alpha\mright)$ realises $\gamma = \check{\delta}$ for some $\delta \in \alpha$ in the base theory, by \autoref{cor:alpha-rec-real-check-ord-true}. In other words, $V\mleft(\Sigma L_\alpha\mright)$ realises that $L_{\check{\alpha}} \cap \Ord \subseteq \check{\alpha}$ and its linearity follows immediately.
\end{proof}

\subsection*{Incomparable ordinals in $V\mleft(\Sigma L_\alpha\mright)$}

Fix some admissible ordinal $\alpha > \omega$. We are now ready to build a pairwise strongly incomparable function $f$ to apply \autoref{cor:dec-emb-sub-strong-incom-code}. Just like in \autoref{sec:czf-v-eq-l}, we need to build some non-$\alpha$-recursive objects. First observe that:
\begin{lemma}
  \label{lem:alpha-rec-real-delta0-correctness}
  $V\mleft(\Sigma L_\alpha\mright)$ realises $\check{L}_\alpha = L_{\check{\alpha}}$. Additionally, for any $\Delta_0$-formula $\varphi\mleft(x_1, \ldots, x_n\mright)$ with all free variables indicated, there exist terms $d_\varphi, k_\varphi \in \Sigma L_\alpha$ such that for any $a_1, \ldots, a_n \in L_\alpha$,
  \begin{enumerate}[label=\uplabel{\roman*}]
    \item $d_\varphi a_1 \cdots a_n {\downarrow} \in \left\{0, 1\right\}$ and $d_\varphi a_1 \cdots a_n = 1$ if and only if $\varphi\mleft(a_1, \ldots, a_n\mright)$ holds;
    \item if $d_\varphi a_1 \cdots a_n = 1$, then $k_\varphi a_1 \cdots a_n {\downarrow} \Vdash \varphi\mleft(\check{a}_1, \ldots, \check{a}_n\mright)$; if $d_\varphi a_1 \cdots a_n = 0$, then $\neg\!\opE{e} e \Vdash \varphi\mleft(\check{a}_1, \ldots, \check{a}_n\mright)$.
  \end{enumerate}

  It follows immediately that $\varphi\mleft(a_1, \ldots, a_n\mright)$ if and only if $\opE{e} e \Vdash \varphi\mleft(\check{a}_1, \ldots, \check{a}_n\mright)$.
\end{lemma}

\begin{proof}
  The existence of $d_\varphi, k_\varphi$ is proved in a similar style to \autoref{prop:ep-real-delta0-correctness}, although the case here is much simpler since we are not dealing with weak realisability. The atomic case is simply due to \autoref{lem:alpha-rec-real-check-name-true} and the trivial realisers of $\check{x} \in \check{y}$ and $\check{x} = \check{y}$. In the binary connective cases, it suffices for $d_\varphi$ to compute binary operations since $L_\alpha$ is a model of classical logic. The bounded quantifier cases are due to the fact that we can search through a set $a \in L_\alpha$ (which corresponds to the `finite sets' in usual Turing machine computations) in $\Sigma L_\alpha$ using the operation $\bm{\rho}$ (cf.\ the corresponding constructions in \autoref{prop:ep-real-delta0-correctness}).

  Now, it suffices to prove by set induction on $\beta$ that for any $\beta \in \alpha^+$, $V\mleft(\Sigma L_\alpha\mright) \vDash \check{L}_\beta = L_{\check{\beta}}$ where
  \[\check{L}_\beta = \left\{\tuple{x, \check{x}} : x \in L_\beta\right\}.\]
  The only non-trivial case is the successor case, where we need to verify that if $y = \dd\mleft(x\mright)$ where $x \in L_\alpha$, then $V\mleft(\Sigma L_\alpha\mright) \vDash \check{y} = \dd\mleft(\check{x}\mright)$. This is precisely because, for any $z \in y = \dd\mleft(x\mright)$, we have that
  \[z = \left\{u \in x : \varphi^x\mleft(u, v_1, \ldots, v_n\mright)\right\}\]
  for $v_1, \ldots, v_n \in x$. Now, the relativisation of $\varphi$ (to a set variable) will be a $\Delta_0$-formula, thus
  \[\oplam{u} \bp \mleft(\oplam{d} k_{\varphi^x} xu v_1 \cdots v_n\mright) \mleft(\oplam{d} \bp u \ior\mright)\]
  will realise $\opA{u \in \check{x}} \left(u \in \check{z} \leftrightarrow \varphi^{\check{x}}\mleft(u, \check{v}_1, \ldots, \check{v}_n\mright)\right)$, and thus $V\mleft(\Sigma L_\alpha\mright) \vDash \check{z} \in \dd\mleft(\check{x}\mright)$. Conversely, suppose that $v \in V\mleft(\Sigma L_\alpha\mright)$ is a name such that\footnote{Formally, the formula $\varphi$ here is also represented by its G\"odel numbering as a name in $V\mleft(\Sigma L_\alpha\mright)$. However, since $V\mleft(\Sigma L_\alpha\mright)$ has the same $\omega$ as the base theory, it is not hard to ensure that the numbering $\check{n} \in \check{\omega}$ for $\varphi$ in $V\mleft(\Sigma L_\alpha\mright)$ and the corresponding $n$ in the base theory will represent precisely the same formula.}
  \[V\mleft(\Sigma L_\alpha\mright) \vDash v = \left\{u \in \check{x} : \varphi^{\check{x}}\mleft(u, \check{v}_1, \ldots, \check{v}_n\mright)\right\},\]
  for some $v_1, \ldots, v_n \in x$. We use \autoref{thm:alpha-rec-real-l-alpha-czf} that $\check{L}_\alpha$ is a model of $\CZF^-$ to derive that, given $\check{x} \in \check{L}_\alpha$, it must follow that $v \in \dd\mleft(\check{x}\mright) \in \check{L}_\alpha$, i.e.\ there is $z \in x$ such that $v = \check{z}$, as realised in $V\mleft(\Sigma L_\alpha\mright)$. It follows again from $\varphi^x$ being $\Delta_0$ that we will have
  \[\opA{u \in x} \left(u \in z \leftrightarrow \varphi^x\mleft(u, v_1, \ldots, v_n\mright)\right),\]
  i.e.\ $z \in y = \dd\mleft(x\mright)$ and thus $v = \check{z} \in \check{y}$ is realised.
\end{proof}

\begin{corollary}
  \label{cor:alpha-rec-real-sigma1-rel-correctness}
  Consider a relation $R \subseteq L_\alpha^n$ that is $\Sigma_1$-definable in $L_\alpha$. Then $V\mleft(\Sigma L_\alpha\mright)$ realises that the corresponding name
  \[\overline{R} = \left\{\tuple{\tuple{x_1, \ldots, x_n}, \tuple{\check{x}_1, \ldots, \check{x}_n}^\bullet} : \tuple{x_1, \ldots, x_n} \in R\right\} \in V\mleft(\Sigma L_\alpha\mright)\]
  lies in $L$.
\end{corollary}

\begin{proof}
  Let $R$ be defined by
  \[\tuple{x_1, \ldots, x_n} \in R \leftrightarrow \opE{y \in L_\alpha} \varphi\mleft(x_1, \ldots, x_n, y\mright),\]
  where $\varphi$ is $\Delta_0$. Let $h : \alpha \rightarrow L_\alpha$ be the $\alpha$-recursive bijection from $\alpha$ onto $L_\alpha$ given by \autoref{prop:alpha-rec-bij-l-alpha}, then for any $\tuple{x_1, \ldots, x_n} \in R$, there exists a unique $y \in L_\alpha$ such that $\varphi\mleft(x_1, \ldots, x_n, y\mright)$ holds and $y$ is $h$-least, i.e.\ for any $\beta \in h^{-1}\mleft(y\mright)$, we have $\neg \varphi\mleft(x_1, \ldots, x_n, h\mleft(\beta\mright)\mright)$. Observe that here the unique $y$ is defined by a $\Sigma_1$-formula relativised to $L_\alpha$, thus \autoref{lem:sigma-l1-smn} applies and there exists $g \in \Sigma L_\alpha$ such that for any $\tuple{x_1, \ldots, x_n} \in R$,
  \[\varphi\mleft(x_1, \ldots, x_n, g x_1 \cdots x_n {\downarrow}\mright).\]
  By \autoref{lem:alpha-rec-real-delta0-correctness}, then
  \[\bp \mleft(g x_1 \cdots x_n\mright) \mleft(k_\varphi x_1 \cdots x_n \mleft(g x_1 \cdots x_n\mright)\mright)\Vdash \opE{y \in \check{L}_\alpha} \varphi\mleft(\check{x}_1, \ldots, \check{x}_n, y\mright).\]
  It is easy to verify that conversely, for any $x_1, \ldots, x_n \in L_\alpha$, if $V\mleft(\Sigma L_\alpha\mright)$ realises $\opE{y \in \check{L}_\alpha} \varphi\mleft(\check{x}_1, \ldots, \check{x}_n, y\mright)$, then $\opE{y \in L_\alpha} \varphi\mleft(x_1, \ldots, x_n, y\mright)$ is true in the base theory by \autoref{lem:alpha-rec-real-delta0-correctness} again. Thus $V\mleft(\Sigma L_\alpha\mright)$ realises
  \[\opA{x_1, \ldots, x_n \in \check{L}_\alpha} \left(\tuple{x_1, \ldots, x_n} \in \overline{R} \leftrightarrow \opE{y \in \check{L}_\alpha} \varphi\mleft(\check{x}_1, \ldots, \check{x}_n, y\mright)\right),\]
  i.e.\ $\overline{R} \in \dd\mleft(\check{L}_\alpha\mright)$. By \autoref{lem:alpha-rec-real-delta0-correctness}, $\check{L}_\alpha = L_{\check{\alpha}}$ and thus $\overline{R} \in \dd\mleft(L_{\check{\alpha}}\mright) \subseteq L$.
\end{proof}

We can now proceed with the constructions. To begin, we work inside $L_\alpha$ as a model of $\KP$, fix a $\Sigma_1$-class bijection $H$ from $\Ord^{L_\alpha} = \alpha$ onto $V^{L_\alpha} = L_\alpha$ as mentioned in \autoref{prop:alpha-rec-bij-l-alpha}, and construct some $g_\beta \in L_\alpha$ for each $\beta \in \alpha$ through $\Sigma$-recursion as follows:

For each $\beta$, first let $g^*_\beta = \bigcup_{\gamma \in \beta} g_\gamma$ and $\beta^* = \left\{\gamma \in \alpha : \opE{a \in L_\alpha} \tuple{\gamma, a} \in g^*_\beta\right\}$ be the domain of $g^*_\beta$. Consider the following cases:
\begin{enumerate}[label=\uplabel{\roman*}]
  \item \label{itm:alpha-rec-pairing-case-wit} If $H\mleft(\beta\mright) = \tuple{e, \tuple{\gamma, \delta}}$ for $e \in L_\alpha$ and $\gamma \neq \delta \in \alpha$, then let
        \[\eta = \bigcup \left\{\varepsilon^+ : \opE{\zeta \in \alpha} \left(\tuple{\zeta, \tuple{\gamma, \varepsilon}} \in g^*_\beta \lor \tuple{\zeta, \tuple{\delta, \varepsilon}} \in g^*_\beta\right)\right\}.\]
        We must have $\eta \in \alpha$ since $g^*_\beta \in L_\alpha$ and $L_\alpha$ is a model of $\KP$. Take the unique enumeration (i.e.\ an order-preserving bijection)
        \[j_\beta : \mu_\beta \rightarrow \left\{\varepsilon \in \eta : \neg\!\opE{\zeta \in \alpha} \tuple{\zeta, \tuple{\gamma, \varepsilon}} \in g^*_\beta\right\}\]
        in the usual ordering, where $\mu_\beta \in \alpha$ is a proper initial segment, and let
        \begin{align*}
          g_\beta = g^*_\beta & \cup \left\{\tuple{\beta^* + \varepsilon, \tuple{\gamma, j_\beta\mleft(\varepsilon\mright)}} : \varepsilon \in \mu_\beta\right\}         \\
                              & \cup \left\{\tuple{\beta^* + \mu_\beta, \tuple{\gamma, \eta}}, \tuple{\left(\beta^* + \mu_\beta\right)^+, \tuple{\delta, \eta}}\right\}.
        \end{align*}
  \item \label{itm:alpha-rec-pairing-case-pair} If $H\mleft(\beta\mright) = \tuple{\gamma, \delta}$ for $\gamma, \delta \in \alpha$, and there does not exist $\varepsilon \in \alpha$ such that $\tuple{\varepsilon, \tuple{\gamma, \delta}} \in g^*_\beta$, then let
        \[g_\beta = g^*_\beta \cup \left\{\tuple{\beta^*, \tuple{\gamma, \delta}}\right\}.\]
  \item Otherwise, simply let $g_\beta = g^*_\beta$.
\end{enumerate}
Observe that we can distinguish cases \ref{itm:alpha-rec-pairing-case-wit} and \ref{itm:alpha-rec-pairing-case-pair} because in classical logic, any ordinal $\gamma$ must either be empty or contain the empty set $\varnothing$, whereas Kuratowski ordered pairs are always nonempty sets that do not contain $\varnothing$ as an element. Thus it cannot be the case that some $\tuple{e, \tuple{\gamma, \delta}} = \tuple{\gamma', \delta'}$.

We first check that:
\begin{lemma}[$\CZF^- + \REM$]
  \label{lem:alpha-rec-pairing}
  The $\Sigma_1$-definable $g = \bigcup_{\beta \in \alpha} g_\beta$ is a bijection between $\alpha$ and $\alpha \times \alpha$.
\end{lemma}

\begin{proof}
  It is easy to check by set induction that each $g^*_\beta$ and $g_\beta$ are functions on proper initial segments of $\alpha$, with $\gamma \in \beta \in \alpha$ implying $g_\gamma \subseteq g_\beta$. Thus, $g$ must also be a function on an initial segment of $\alpha$.

  Now, clause \ref{itm:alpha-rec-pairing-case-pair} in the definition of $g_\beta$ ensures that $g$ is a surjection onto $\alpha \times \alpha$. However, $g$ is a $\Sigma_1$-class in the set model $L_\alpha$ of $\KP$, and by $\Delta_0$-collection, there cannot be a surjection from an ordinal (in $\alpha$) onto $\left(\Ord \times \Ord\right)^{L_\alpha}$. Thus we know that the domain of $g$ must be $\alpha$ itself.

  Finally, for the injectivity of $g$, by the linearity of ordinals from $\REM$, it suffices to check that every $g_\beta$ is injective. In clause \ref{itm:alpha-rec-pairing-case-wit}, this is ensured by the choice of $\eta$, since for any $\tuple{\zeta, \tuple{\gamma, \varepsilon}}$ or $\tuple{\zeta, \tuple{\delta, \varepsilon}} \in g^*_\beta$, we must have $\varepsilon \in \eta$ and thus $\tuple{\gamma, \varepsilon} \neq \tuple{\gamma, \eta}$, $\tuple{\delta, \varepsilon} \neq \tuple{\delta, \eta}$. In clause \ref{itm:alpha-rec-pairing-case-pair}, this is ensured simply by stipulation.
\end{proof}

By \autoref{lem:sigma-l1-smn}, both this function $g$ and its inverse are then also computed by some terms in $\Sigma L_\alpha$. We can denote $G, G^{-1} \in \Sigma L_\alpha$ such that
\begin{align*}
  G \beta {\downarrow}             & = g\mleft(\beta\mright),             \\
  G^{-1} \beta \gamma {\downarrow} & = g^{-1}\mleft(\beta, \gamma\mright)
\end{align*}
for any $\beta, \gamma \in \alpha$. On the other hand, by \autoref{cor:alpha-rec-real-sigma1-rel-correctness}, there will also be a name $\overline{G} \in V\mleft(\Sigma L_\alpha\mright)$ that performs the same computations on $\check{\alpha}$, such that $V\mleft(\Sigma L_\alpha\mright) \vDash \overline{G} \in L$.

Next, in the base theory $\CZF^- + \REM$, outside of $L_\alpha$, we will recursively construct a corresponding sequence of objects $h_\beta$ --- which should be functions on the same domain as $g_\beta$ --- for each $\beta \in \alpha$ as follows:

For each $\beta$, first let $h^*_\beta = \bigcup_{\gamma \in \beta} h_\gamma$. Consider the following cases:
\begin{enumerate}[label=\uplabel{\roman*}]
  \item \label{itm:alpha-rec-incom-ind-case-wit} If $H\mleft(\beta\mright) = \tuple{e, \tuple{\gamma, \delta}}$ for $e \in L_\alpha$ and $\gamma \neq \delta \in \alpha$, we take $j_\beta, \mu_\beta$ from the corresponding clause in the construction of $g_\beta$ and let
        \[h'_\beta = h^*_\beta \cup \left\{\tuple{\beta^* + \varepsilon, \varnothing} : \varepsilon \in \mu_\beta\right\}.\]
        Evaluate $e h'_\beta$ in $\Sigma L_\alpha$. If $e h'_\beta {\downarrow} = h'_\beta \cup \left\{\tuple{\beta^* + \mu_\beta, f}\right\}$ for some $f \in L_\alpha$, then let
        \[h_\beta = h'_\beta \cup \left\{\tuple{\beta^* + \mu_\beta, f^+}, \tuple{\left(\beta^* + \mu_\beta\right)^+, \varnothing}\right\},\]
        where $f^+ = f \cup \left\{f\right\}$. If it does not converge or if its output does not have the correct form --- which we can test by $\REM$ --- simply let
        \[h_\beta = h'_\beta \cup \left\{\tuple{\beta^* + \mu_\beta, \varnothing}, \tuple{\left(\beta^* + \mu_\beta\right)^+, \varnothing}\right\}.\]
  \item If $H\mleft(\beta\mright) = \tuple{\gamma, \delta}$ for $\gamma, \delta \in \alpha$, and there does not exist $\varepsilon \in \alpha$ such that $\tuple{\varepsilon, \tuple{\gamma, \delta}} \in h^*_\beta$, let $h_\beta = h^*_\beta \cup \left\{\tuple{\beta^*, \varnothing}\right\}$.
  \item Otherwise, simply let $h_\beta = h^*_\beta$.
\end{enumerate}

\begin{lemma}
  \label{lem:alpha-rec-non-comp-path}
  The function $h = \bigcup_{\beta \in \alpha} h_\beta$ is a function from $\alpha$ to $L_\alpha$. Additionally, for any $e \in \Sigma L_\alpha$, and $\gamma, \delta \in \alpha$ where $\gamma \neq \delta$, there exists $\eta \in \alpha$ such that
  \[e \mleft(h \upharpoonright G^{-1} \gamma \eta\mright) {\downarrow} \rightarrow e \mleft(h \upharpoonright G^{-1} \gamma \eta\mright) \neq h \upharpoonright G^{-1} \delta \eta,\]
  where we also know that for any $\varepsilon \in \eta$, $G^{-1} \gamma \varepsilon \in G^{-1} \gamma \eta$.
\end{lemma}

\begin{proof}
  It is easy to check by set induction that each $h^*_\beta$ and $h_\beta$ are functions on the same domain as $g^*_\beta$ and $g_\beta$ respectively, with $\gamma \in \beta \in \alpha$ implying $h_\gamma \subseteq h_\beta$. Thus, $h$ is a function with the same domain as $g$ in \autoref{lem:alpha-rec-pairing}, i.e.\ $\alpha$.

  Now, given any $e \in \Sigma L_\alpha$ and $\gamma, \delta \in \alpha$ where $\gamma \neq \delta$, we consider the ordinal $\beta = H^{-1}\mleft(\tuple{e, \tuple{\gamma, \delta}}\mright) \in \alpha$. Then by the construction of $g_\beta$, we know that for some $\eta \in \alpha$, $G \mleft(\beta^* + \mu_\beta\mright) = \tuple{\gamma, \eta}$ and $G {\left(\beta^* + \mu_\beta\right)^+} = \tuple{\delta, \eta}$, whereas
  \[e \mleft(h \upharpoonright \beta^* + \mu_\beta\mright) {\downarrow} \rightarrow e \mleft(h \upharpoonright \beta^* + \mu_\beta\mright) \neq h \upharpoonright \left(\beta^* + \mu_\beta\right)^+\]
  holds, precisely because of our stipulation in clause \ref{itm:alpha-rec-incom-ind-case-wit} in the construction of $h_\beta$ above. The inclusion of the sequences $j_\beta$ in $g_\beta$ and $h_\beta$ ensures that for any $\varepsilon \in \eta$, $G^{-1} \gamma \varepsilon \in G^{-1} \gamma \eta$.
\end{proof}

We are now in a very familiar position to repeat the work in \autoref{sec:czf-v-eq-l}. Fix
\[\bl_\alpha = \oplam{f} \bm{\rho} \bp_0 f \in \Sigma L_\alpha\]
which computes the domain of the input $f$, and
\[\bt_\alpha = \oplam{fx} \bm{\nu} \mleft(\bm{\rho} \mleft(\oplam{t} \bi_1 \mleft(\bm{\pi} tt\mright) \mleft(\bp_0 t\mright) x\mright) f\mright) \in \Sigma L_\alpha\]
which computes the restriction $f \upharpoonright x$. Let
\[\alpha_0 = \left\{\tuple{h \upharpoonright \beta, \check{\beta}} : \beta \in \alpha\right\} \in V\mleft(\Sigma L_\alpha\mright)\]
using $h$ from \autoref{lem:alpha-rec-non-comp-path}, and we check that:
\begin{lemma}
  $V\mleft(\Sigma L_\alpha\mright)$ realises the following:
  \begin{enumerate}[label=\uplabel{\roman*}]
    \item \label{itm:alpha-rec-incom-ord-basic-ord} $\alpha_0 \in \Ord$;
    \item $\alpha_0 \subseteq \check{\alpha}$.
  \end{enumerate}
\end{lemma}

\begin{proof}
  Clearly
  \[\oplam{f} \bp \mleft(\bl_\alpha f\mright) \ior \Vdash \opA{\beta \in \alpha_0} \beta \in \check{\alpha}.\]

  For \ref{itm:alpha-rec-incom-ord-basic-ord}, we observe that for any $\tuple{f, \check{\beta}} \in \alpha_0$ and any $\gamma \in \beta$, $\bp \mleft(\bt_\alpha f \gamma\mright) \ior \Vdash \check{\gamma} \in \alpha_0$. Therefore,
  \[\oplam{f \gamma} \bp \mleft(\bt_\alpha f \gamma\mright) \ior \Vdash \opA{\beta \in \alpha_0} \opA{\gamma \in \beta} \gamma \in \alpha_0.\]
  It follows from $\alpha_0 \subseteq \check{\alpha}$ that every element in $\alpha_0$ is also transitive, thus $\alpha_0 \in \Ord$.
\end{proof}

Working inside $V\mleft(\Sigma L_\alpha\mright)$, we can then construct $f_0 : \check{\alpha} \rightarrow \Pow\mleft(\check{\alpha}\mright)$ as
\[f_0\mleft(\beta\mright) = \left\{\gamma \in \check{\alpha} : \opA{\delta \in \gamma} \overline{G}^{-1}\mleft(\beta, \delta\mright) \in \alpha_0\right\}.\]
It is easy to check that $V\mleft(\Sigma L_\alpha\mright)$ realises $\opA{\beta \in \check{\alpha}} f_0\mleft(\beta\mright) \in \Pow\mleft(\check{\alpha}\mright) \cap \Ord$, and \autoref{lem:alpha-rec-non-comp-path} ensures that $f_0$ is pairwise strongly incomparable:
\begin{proposition}
  For any $\gamma, \delta \in \alpha$, if some $d \Vdash f_0\mleft(\check{\gamma}\mright) \subseteq f_0\mleft(\check{\delta}\mright)$, then $\gamma = \delta$. Namely, it follows immediately that
  \[\oplam{\gamma \delta x} \ior \Vdash \opA{\gamma, \delta \in \check{\alpha}} \left(f_0\mleft(\gamma\mright) \subseteq f_0\mleft(\delta\mright) \rightarrow \gamma = \delta\right).\]
\end{proposition}

\begin{proof}
  This is entirely analogous to \autoref{prop:f-subset-eq}, but we present it here for completeness: assume otherwise, i.e.\ $\gamma \neq \delta$. Reasoning inside $V\mleft(\Sigma L_\alpha\mright)$, we know that $f_0\mleft(\check{\gamma}\mright) \subseteq f_0\mleft(\check{\delta}\mright)$, by definition of $f_0$, is equivalent to
  \[\opA{\beta \in \check{\alpha}} \left(\opA{\varepsilon \in \beta} \overline{G}^{-1}\mleft(\check{\gamma}, \varepsilon\mright) \in \alpha_0 \rightarrow \opA{\varepsilon \in \beta} \overline{G}^{-1}\mleft(\check{\delta}, \varepsilon\mright) \in \alpha_0\right),\]
  so suppose that this is realised by some $d \in \Sigma L_\alpha$. Pedantically, we take the formula $\overline{G}^{-1}\mleft(\check{\gamma}, \check{\varepsilon}\mright) \in \alpha_0$ as a shorthand for
  \[\opE{t \in \overline{G}} \opE{\eta \in \alpha_0} t = \tuple{\eta, \tuple{\check{\gamma}, \check{\varepsilon}}},\]
  thus if some $e \Vdash \overline{G}^{-1}\mleft(\check{\gamma}, \check{\varepsilon}\mright) \in \alpha_0$, then $\bp_{10} e = h \upharpoonright G^{-1} \gamma \varepsilon$ using the fact that $\overline{G}$ in $V\mleft(\Sigma L_\alpha\mright)$ computes precisely the same map as $G$ in the base theory.

  Conversely, we can also find some function $q \in \Sigma L_\alpha$ such that, for any $\tuple{f, \check{\beta}} \in \alpha_0$ where $G^{-1} \gamma \varepsilon \subseteq \bl_\alpha f$,
  \[q \gamma \varepsilon \mleft(\bt_\alpha f \mleft(G^{-1} \gamma \varepsilon\mright)\mright) \Vdash \overline{G}^{-1}\mleft(\check{\gamma}, \check{\varepsilon}\mright) \in \alpha_0,\]
  where the precise computation performed by $q$ is cumbersome to write but entirely routine.

  Consequently, for any $e \in \Sigma L_\alpha$, take $\eta \in \alpha$ as given by \autoref{lem:alpha-rec-non-comp-path} and $f = h \upharpoonright G^{-1} \gamma \eta$, then
  \[\oplam{\varepsilon} q \gamma \varepsilon \mleft(\bt_\alpha f \mleft(G^{-1} \gamma \varepsilon\mright)\mright) \Vdash \opA{\varepsilon \in \check{\chi}} \overline{G}^{-1}\mleft(\check{\gamma}, \varepsilon\mright) \in \alpha_0\]
  where $\chi = \eta^+ = \bm{\nu} \mleft(\bm{\pi} \eta \mleft(\bm{\pi} \eta \eta\mright)\mright)$. If we pass this to the realiser $d$ at the beginning of this proof, i.e.\ we take $p = \oplam{\eta f} d \mleft(\bm{\nu} \mleft(\bm{\pi} \eta \mleft(\bm{\pi} \eta \eta\mright)\mright)\mright) \mleft(\oplam{\varepsilon} q \gamma \varepsilon \mleft(\bt_\alpha f \mleft(G^{-1} \gamma \varepsilon\mright)\mright)\mright) \eta$, then we have $p \eta f \Vdash \overline{G}^{-1}\mleft(\check{\delta}, \check{\eta}\mright) \in \alpha_0$ and $\bp_{10} \mleft(p \eta f\mright) = h \upharpoonright G^{-1} \delta \eta$.

  Observe that $\eta$ can be recovered as $\bp_1 \mleft(G \mleft(\bl_\alpha f\mright)\mright)$, thus we have a single-parameter function $e^* = \oplam{f} \bp_{10}\mleft(p \mleft(\bp_1 \mleft(G \mleft(\bl_\alpha f\mright)\mright)\mright) f\mright)$, such that
  \[e^* \mleft(h \upharpoonright G^{-1} \gamma \eta\mright) {\downarrow} \land e^* \mleft(h \upharpoonright G^{-1} \gamma \eta\mright) = h \upharpoonright G^{-1} \delta \eta,\]
  which works for any $e \in \Sigma L_\alpha$ and its corresponding $\eta$. However, simply take $e = e^*$ and we have a contradiction to \autoref{lem:alpha-rec-non-comp-path}. Therefore, we must have $\gamma = \delta$ instead.
\end{proof}

It follows immediately that:
\begin{theorem}[$\CZF^- + \REM$]
  \label{thm:alpha-rec-pow-in-l}
  For any ordinal $\beta$, if there exists an admissible ordinal $\alpha$ such that $\beta \subseteq \alpha$, then there exists a \uppar{weak} pca $\mathcal{A}$ such that $\beta \subseteq \mathcal{A}$ and
  \[V\mleft(\mathcal{A}\mright) \vDash \Pow\mleft(\check{\beta}\mright) \subseteq L.\]
\end{theorem}

\begin{proof}
  The case where $\beta \subseteq \omega$ can be easily solved in $V\mleft(\Kl\mright)$ by adapting the similar proof in \autoref{sec:czf-v-eq-l}. Here, we will only consider $\beta > \omega$ and thus only admissible $\alpha > \omega$.

  Take $\mathcal{A} = \Sigma L_\alpha$ and work inside $V\mleft(\mathcal{A}\mright)$. Note that the argument in the proof of \autoref{cor:dec-emb-sub-strong-incom-code} would imply that $\alpha_0 \in L$, and since $\overline{G} \in L$ by \autoref{cor:alpha-rec-real-sigma1-rel-correctness}, we have $f_0 \in L$. Since $f_0$ is pairwise strongly incomparable, by \autoref{cor:dec-emb-sub-strong-incom-code}, we have $\Pow\mleft(\check{\alpha}\mright) \subseteq L$, and $\Pow\mleft(\check{\beta}\mright) \subseteq L$ follows from the trivial realisability of $\check{\beta} \subseteq \check{\alpha}$.
\end{proof}

Additionally, if the base theory satisfies exponentiation, subset collection or powerset (which have no difference under the assumption $\REM$), then $V\mleft(\mathcal{A}\mright)$ by \autoref{lem:czfp-real-powerset} also satisfies powerset, and we have $\Pow\mleft(\check{\beta}\mright) \in L$. So here the obstacles with set recursive functions are no longer present.

\backmatter

\printindex

\bibmain

\end{document}
